\ifx\danslelivre\undefined
\RequirePackage{luatex85}
\documentclass[10pt]{../fichiers-configuration-latex/smfart}
\usepackage{polyglossia}
%\usepackage[francais]{babel}
\usepackage{fontspec}
\usepackage[math-style=ISO,bold-style=upright,sans-style=upright]{unicode-math}
%\usepackage[mathletters]{ucs}
%\usepackage[utf8]{inputenc}
%\usepackage[T1]{fontenc}
\usepackage[unicode,colorlinks=true,linkcolor=blue,citecolor=Sepia,urlcolor=magenta]{hyperref}
\setmainlanguage{french}
\setmainfont[Ligatures=TeX]{Linux Libertine O}
% variante possible : TeXGyrePagella ?
\setmathfont{XITS Math}

% maths
\usepackage{amsmath}
%\usepackage{amssymb}
%\usepackage{mathrsfs}
%\usepackage{stmaryrd}
%\usepackage{euscript}

% couleurs
\usepackage[usenames,dvipsnames]{xcolor}

% graphiques
\usepackage{graphicx}
\usepackage[all]{xy}
% Xy-pic 3.8.9's PDF driver assumes that \fontdimen8\textfont3 is a
% legacy TeX rule width.  Under LuaTeX with unicode-math/XITS Math it is
% instead 6.62pt, which turns every arrow shaft into a thick black bar.
% Override only xypdf's private cached width; do not alter the math font.
\makeatletter
\AtBeginDocument{%
  \@ifundefined{xP@preclw}{}{%
    \edef\xP@preclw{0.4pt}%
    \edef\xP@lw{\xP@dim\xP@preclw}%
    \PackageInfo{sga4-xypdf-compat}{Xy-pic PDF line width set to \xP@preclw}%
  }%
}
\makeatother
%\usepackage{tikz}
%\usetikzlibrary{matrix,arrows,calc}
%\tikzset{column sep=3em,row sep=2em,text height=1.5ex, text depth=0.25ex,>=angle 90}
%\tikzset{description/.style={fill=white,inner sep=1pt}}

% hyperliens, 'etiquettes
\usepackage{xr-hyper}
\usepackage{../fichiers-configuration-latex/smfthm}
%\ifx\cacheretiquettes\undefined
%\usepackage[inline]{showlabels} \else \fi % pour rendre visibles les 'etiquettes
%\usepackage{../fichiers-configuration-latex/makeidx} % pour faire un index
%\usepackage[pagebackref,unicode]{hyperref}  % pour voir o`u sont cit'ees les r'ef'erences
\hypersetup{unicode,colorlinks=true,linkcolor=blue,citecolor=Sepia,urlcolor=magenta,
  hypertexnames=false}
%\hypersetup{breaklinks=true} %permet le retour `a la ligne dans les liens trop longs. Incompatible avec pdflatex
%\renewcommand*{\backref}[1]{{\small [cit\'e page(s)~#1.]}}

\usepackage{xspace}
\usepackage{pifont} % pour expose XVII

\renewcommand{\thefootnote}{\roman{footnote}} % pour mettre les notes en bas de page sous la forme (ii) et non (2), qui ressemble trop (?) a un exposant.

\newcounter{pageoriginal}
\newcommand{\pageoriginale}{\refstepcounter{pageoriginal}\marginpar{\footnotesize\textbf{\thepageoriginal}}\xspace}
\let\pageoriginaled\pageoriginale

\usepackage{ifthen}
\newboolean{original}
\def\sisi{\ifthenelse{\boolean{original}}}
\setboolean{original}{false}

\newcounter{cnote}
\def\nde#1{\ifthenelse{\boolean{original}}{}{\stepcounter{cnote}\def\thefootnote{\arabic{cnote}}\protect\footnote[\value{cnote}]{N.D.E.: #1}}}
\def\ndemark{\ifthenelse{\boolean{original}}{}{\stepcounter{cnote}\def\thefootnote{\arabic{cnote}}\protect\footnotemark[\value{cnote}]}}
\def\ndetext#1{\ifthenelse{\boolean{original}}{}{\def\thefootnote{\arabic{cnote}}\protect\footnotetext[\value{cnote}]{N.D.E.: #1}}}

\setmainlanguage{french}
\setotherlanguage{english}

% unicode-math deliberately disables the legacy \Arrowvert glyph.  The
% source uses it once as a rotated equality marker; \Vert has the intended
% appearance and keeps the historical diagram buildable with LuaLaTeX.
\renewcommand{\Arrowvert}{\Vert}


%%
%% Pour afficher le numéro de version Git (exige que LuaTeX soit appelé avec --shell-escape)
%%

\catcode`\%=12
\directlua{
function command_output(cmd)
  local output_pipe = io.popen(cmd)
  local line_iter = output_pipe and output_pipe:lines()
  local line = line_iter and line_iter()
  if output_pipe then output_pipe:close() end
  return line
end
local git_command1 = [[git log --pretty=format:"%h %ai" -1]]
local git_command2 = [[if git diff --quiet HEAD ; then echo clean ; else echo dirty ; fi]]
function git_status()
  local status = command_output(git_command1)
  if status then
    tex.tprint({luatexbase.catcodetables.CatcodeTableLaTeX, string.char(0x5c).."texttt{"}, {-2,status}, {luatexbase.catcodetables.CatcodeTableLaTeX, "}"})
    if not (command_output(git_command2) == "clean") then tex.sprint(-2, " (modifications non validées)") end
  else
    tex.sprint(-2, "?")
  end
end
}
\catcode`\%=14
\def\showgitstatus{\directlua{git_status()}}


% macros pour LaTeX
\renewcommand{\mathsf}{}

% Police différente pour U+2113 SCRIPT SMALL L
\setmathfont[range="2113]{latinmodern-math.otf}

% Give every explicit SGA label its own destination and preserve the printed
% number supplied as the optional argument.  The former definition inherited
% whichever counter anchor happened to precede it and ignored #1 entirely.
\makeatletter
\def\etiquette[#1]#2{%
  \phantomsection
  \protected@edef\@currentlabel{#1}%
  \label{#2}%
}
\makeatother
\def\eqhyperref[#1]#2{(\hyperref[#1]{#2})}
\def\pagehyperref[#1]#2{(page \hyperref[#1]{#2})}
\def\hyperrefext[#1][#2]#3{\hyperref[#2]{#1 #3}}
\def\eqhyperrefext[#1][#2]#3{\hyperref[#2]{#1 #3}}     %cf. exposes 17
\def\neweqhyperrefext[#1][#2]#3{\hyperref[#2]{#1 #3}}  %cf. exposes 17

\def\shat#1{\protect\widehat{#1}}
\def\stilde#1{\protect\widetilde{#1}}
%%%%%%%%%%%%%%%%%%


\def\etcn{etc.\xspace}
\def\etco{etc...\xspace}
\def\etc{\sisi{\etco}{\etcn}}
\def\ldotsn{\sisi{...}{\etcn}}

\def\cf{{cf.}\kern.3em}
\def\Cf{{Cf.}\kern.3em}
\def\ie{{i.e.}\kern.3em}
\def\iev{{i.e.},\xspace}
\def\loccito{loc.\kern2pt cit.\xspace}
\def\loccitn{\textit{loc.\kern2pt cit.}\xspace}
\def\loccit{\sisi{\loccito}{\loccitn}}
\def\EGA{\textup{EGA}\kern.2em}
\def\SGA{\textup{SGA}\kern.3em}

\def\quoi{}

\def\quav{\quad,\quad}
\def\quaet{\quad \text{et}\quad}
\def\fibreproduct#1#2#3{#1 ×_{#3} #2}
\let\fprod\fibreproduct
\def\ccheck#1#2{\mbox{${\check{\check{\rm #1}}}^{\lower2pt\hbox{$\scriptstyle #2$}}$}}
\def\sheaf#1{\protect\mathcal{#1}}
\def\vieuxsheaf#1{{#1}} % cf. problème avec sheaf{a\Phi_i} par exemple...
\DeclareMathOperator{\Hom}{Hom}
\DeclareMathOperator{\can}{can}
\DeclareMathOperator{\Ind}{Ind}
\DeclareMathOperator{\ind}{ind}
%\DeclareMathOperator{\Hhomtop}{\mathcal{H}\!\mathit{omtop}}
\DeclareMathOperator{\Hhomtop}{\mathbf{Homtop}}
\DeclareMathOperator{\Point}{Point}
\DeclareMathOperator{\Pro}{Pro}
\DeclareMathOperator{\Top}{Top}
\DeclareMathOperator{\Mod}{Mod}
\DeclareMathOperator{\Cart}{Cart}
%\DeclareMathOperator{\Mmorsite}{\mathcal{M}\!\mathit{orsite}}
\DeclareMathOperator{\Mmorsite}{\mathbf{Morsite}}
\DeclareMathOperator{\loc}{loc}
\DeclareMathOperator{\Loc}{Loc}
\DeclareMathOperator{\Pic}{Pic}
\DeclareMathOperator{\Et}{Et}
\DeclareMathOperator{\card}{card}
\DeclareMathOperator{\Ens}{Ens}
\DeclareMathOperator{\h}{\underline{H}}
\DeclareMathOperator{\Fonct}{Fonct}
\DeclareMathOperator{\Ker}{Ker}
\DeclareMathOperator{\but}{but}
\DeclareMathOperator{\Ob}{Ob}
\DeclareMathOperator{\ob}{ob}
\DeclareMathOperator{\Lib}{Lib}
\DeclareMathOperator{\esj}{esj}
\DeclareMathOperator{\id}{id}
\DeclareMathOperator{\Id}{Id}
\DeclareMathOperator{\ad}{ad}
\DeclareMathOperator{\Iim}{Im}
\DeclareMathOperator{\pr}{pr}
\DeclareMathOperator{\Cat}{Cat}
\DeclareMathOperator{\car}{car}
\DeclareMathOperator{\SysInd}{SysInd}
\DeclareMathOperator{\Sup}{Sup}
\DeclareMathOperator{\Coker}{Coker}
\DeclareMathOperator{\Filt}{Filt}
\DeclareMathOperator{\acc}{acc}
\DeclareMathOperator{\sur}{sur}
\DeclareMathOperator{\cat}{cat}
\DeclareMathOperator{\Ord}{Ord}
\DeclareMathOperator{\Gr}{Gr}
\DeclareMathOperator{\Cov}{Cov}
\DeclareMathOperator{\Caf}{Caf}
\DeclareMathOperator{\coker}{coker}
\DeclareMathOperator{\coim}{coim}
\DeclareMathOperator{\im}{im}
\DeclareMathOperator{\Ouv}{Ouv}
\DeclareMathOperator{\diag}{diag}
\DeclareMathOperator{\Esp}{Esp}
\DeclareMathOperator{\TOP}{TOP}
\DeclareMathOperator{\Unif}{Unif}
\DeclareMathOperator{\Iind}{\mathbf{Ind}}
\DeclareMathOperator{\Hhomcart}{\mathbf{Homcart}}
\DeclareMathOperator{\Hhomtopess}{\mathbf{Homtopess}}
\DeclareMathOperator{\Ppoints}{\mathbf{Points}}
\DeclareMathOperator{\Ppoint}{\mathbf{Point}}
\DeclareMathOperator{\Ttors}{\mathbf{Tors}}
\DeclareMathOperator{\Ssex}{\mathbf{Sex}}
\DeclareMathOperator{\Ffib}{\mathbf{Fib}}
\DeclareMathOperator{\Vvois}{\mathbf{Vois}}
\DeclareMathOperator{\Ppointess}{\mathbf{Pointess}}
\DeclareMathOperator{\scrHom}{\mathbf{Hom}}
\DeclareMathOperator{\sheafPic}{\mathbf{Pic}}
\DeclareMathOperator{\sheafAut}{\mathbf{Aut}}
\DeclareMathOperator{\Hhomtoplocan}{\mathbf{Homtoplocan}}
\DeclareMathOperator{\Hhomtopan}{\mathbf{Homtopan}}
\DeclareMathOperator{\Pptgeom}{\mathbf{Ptgéom}}
\DeclareMathOperator{\Morsite}{Morsite}
\DeclareMathOperator{\Fib}{Fib}
\DeclareMathOperator{\Fais}{Fais}
\DeclareMathOperator{\kar}{kar}
\DeclareMathOperator{\Groupes}{Groupes}
\DeclareMathOperator{\Restr}{Restr}
\DeclareMathOperator{\Isom}{Isom}
\DeclareMathOperator{\Prol}{Prol}
\DeclareMathOperator{\Res}{Res}
\DeclareMathOperator{\Tors}{Tors}
\DeclareMathOperator{\esp}{esp}
\DeclareMathOperator{\es}{es}
\DeclareMathOperator{\od}{od}
\DeclareMathOperator{\gr}{gr}
\DeclareMathOperator{\sob}{sob}
\DeclareMathOperator{\dfn}{dfn}
\DeclareMathOperator{\Sch}{Sch}
\DeclareMathOperator{\ouv}{ouv}
\DeclareMathOperator{\cons}{cons}
\DeclareMathOperator{\Nerf}{Nerf}
\DeclareMathOperator{\Points}{Points}
\DeclareMathOperator{\Espsob}{Espsob}
\DeclareMathOperator{\Homtop}{Homtop}
\DeclareMathOperator{\Pt}{Pt}
\DeclareMathOperator{\Quot}{Quot}
\DeclareMathOperator{\glob}{glob}
\DeclareMathOperator{\Homfib}{Homfib}
\DeclareMathOperator{\Qouv}{Qouv}
\DeclareMathOperator{\Quotouv}{Quotouv}
\DeclareMathOperator{\ens}{ens}
\DeclareMathOperator{\ab}{ab}
\DeclareMathOperator{\an}{an}
\DeclareMathOperator{\fs}{fs}
\DeclareMathOperator{\red}{red}
\DeclareMathOperator{\Ab}{Ab}
\DeclareMathOperator{\cosq}{cosq}
\DeclareMathOperator{\Ann}{Ann}
\DeclareMathOperator{\Zar}{Zar}


%%%%%%%%%%%%
%%%%% SGA4-2
\DeclareMathOperator{\Sex}{Sex}
\DeclareMathOperator{\vois}{vois}
\DeclareMathOperator{\Vois}{Vois}
\DeclareMathOperator{\Ext}{Ext}
\DeclareMathOperator{\ext}{ext}
\DeclareMathOperator{\Fer}{Fer}
\DeclareMathOperator{\Add}{Add}
\DeclareMathOperator{\Cocart}{Cocart}
\DeclareMathOperator{\cosk}{cosk}
\DeclareMathOperator{\plat}{plat}

\DeclareMathOperator{\faisc}{faisc}
\DeclareMathOperator{\sk}{sk}
\DeclareMathOperator{\ch}{ch}

\def\HH{\mathbf{H}}
\def\qc{\mathrm{qc}}
\def\et{\mathrm{et}}
\def\coh{\mathrm{coh}}

\def\Hvr#1{%
\begin{array}[b]{l}
{\kern2pt\scriptstyle r}\\[-8pt]
{\kern1pt\scriptstyle \vee}\\[-5pt]
H^{#1}
\end{array}}

\def\Hvv#1{%
\begin{array}[b]{l}
\kern2pt\rotatebox{90}{${\scriptscriptstyle\ll}$}\\[-6pt]
H^{#1}
\end{array}}

\def\Hvi#1{%
\begin{array}[b]{l}
{\kern0.2pt\scriptscriptstyle \infty}\\[-8pt]
{\kern1pt\scriptstyle \vee}\\[-5pt]
H^{#1}
\end{array}}

\def\Hv#1{%
\begin{array}[b]{l}
{\kern1pt\scriptstyle \vee}\\[-5pt]
H^{#1}
\end{array}}

\DeclareMathOperator{\cart}{cart}
\DeclareMathOperator{\Ffaisc}{\mathbf{Faisc}}


\def\Ev#1{%
\begin{array}[b]{l}
{\kern1pt\scriptstyle \vee}\\[-5pt]
E'_{#1}
\end{array}}

\DeclareMathOperator{\Cont}{Cont}
\DeclareMathOperator{\Ccont}{\mathscr{C}\!\mathit{ont}}
\DeclareMathOperator{\Site}{Site}
\newcommand\MM{\mathbf{M}}
\def\limtop#1{{\displaystyle{\varprojlim_{#1}\!\mbox{\rm top }}}}
\def\Limtop{\mathop{\varprojlim\!\mbox{\rm top}}}
\DeclareMathOperator{\spec}{spec}
\def\ss{\scriptstyle{\mathscr{S}}}
\DeclareMathOperator{\Gal}{Gal}
\def\smallm{\scriptstyle{\mathscr{M}}}
\DeclareMathOperator{\End}{End}
\DeclareMathOperator{\ass}{ass}
\DeclareMathOperator{\zar}{zar}
\DeclareMathOperator{\Ogg}{Ogg}
\DeclareMathOperator{\res}{res}
\DeclareMathOperator{\tr}{tr}
\DeclareMathOperator{\Supp}{Supp}
\DeclareMathOperator{\codim}{codim}
\DeclareMathOperator{\rk}{rk}
\DeclareMathOperator{\cur}{cur}
\DeclareMathOperator{\cdqc}{cdqc}
\newcommand\DD{\mathbf{D}}
\DeclareMathOperator{\cl}{cl}
\def\JJ{\mathbf{J}}
\def\FF{\mathbf{F}}
\def\BB{\mathbf{B}}
\DeclareMathOperator{\Alget}{Alget}
\DeclareMathOperator{\Ass}{Ass}
\DeclareMathOperator{\supp}{supp}
\DeclareMathOperator{\rad}{rad}
\DeclareMathOperator{\prof}{prof}
\DeclareMathOperator{\cint}{int}
\DeclareMathOperator{\sic}{sic}
\DeclareMathOperator{\Tor}{Tor}
\DeclareMathOperator{\tors}{tors}
\DeclareMathOperator{\Sym}{Sym}
\DeclareMathOperator{\parf}{parf}
\DeclareMathOperator{\tor}{tor}
\DeclareMathOperator{\defn}{defn}
\DeclareMathOperator{\Div}{Div}
\DeclareMathOperator{\Pol}{Pol}
\DeclareMathOperator{\Fex}{Fex}
\DeclareMathOperator{\Lie}{Lie}
\DeclareMathOperator{\tot}{tot}
\DeclareMathOperator{\Aut}{Aut}
\DeclareMathOperator{\HOM}{HOM}
\DeclareMathOperator{\pch}{pch}
\DeclareMathOperator{\Ch}{Ch}
\DeclareMathOperator{\EXT}{EXT}
\DeclareMathOperator{\PIC}{PIC}
\DeclareMathOperator{\Rf}{Rf}
\DeclareMathOperator{\inj}{inj}
\newcommand\VV{\mathbf{V}}

\def\Xv{%
\begin{array}[b]{l}
{\kern1pt\scriptstyle \vee}\\[-5pt]
X
\end{array}}

%%%%%%%
\newcommand\TT{\mathbf{T}}
\newcommand\XX{\mathbf{X}}
\newcommand\EE{\mathbf{E}}
%%%%%%

\def\qedsymbol{\textup{C.Q.F.D.}}
\let\CQFD\qed
\let\cqfd\qed

%%%%%\bigdot
\font\bigsymb=cmsy10 at 4pt
\def\bigdot{{\kern1.2pt\raise 1.5pt\hbox{\bigsymb\char15}}}
%%%%%

\def\Lright{\mathop{\vcenter{\hbox{\hbox{$\longrightarrow$}\lower8pt \hbox{\kern-16pt\hbox{$\longrightarrow$}}}}}}

\def\Aa{{\mathcal A}}
\def\Bb{{\mathcal B}}
\def\Cc{{\mathcal C}}
\def\Dd{{\mathcal D}}
\def\Ee{{\mathcal E}}
\def\Ff{{\mathcal F}}
\def\Gg{{\mathcal G}}
\def\Hh{{\mathcal H}}
\def\Ii{{\mathcal I}}
\def\Jj{{\mathcal J}}
\def\Kk{{\mathcal K}}
\def\Ll{{\mathcal L}}
\def\Mm{{\mathcal M}}
\def\Nn{{\mathcal N}}
\def\Oo{{\mathcal O}}
\def\Pp{{\mathcal P}}
\def\Qq{{\mathcal Q}}
\def\Rr{{\mathcal R}}
\def\Ss{{\mathcal S}}
\def\Tt{{\mathcal T}}
\def\Uu{{\mathcal U}}
\def\Vv{{\mathcal V}}
\def\Ww{{\mathcal W}}
\def\Xx{{\mathcal X}}
\def\Yy{{\mathcal Y}}
\def\Zz{{\mathcal Z}}

\newcommand{\CC}{\mathbf{C}}
\newcommand{\GG}{\mathbf{G}}
\newcommand{\KK}{\mathbf{K}}
\newcommand{\LL}{\mathbf{L}}
\newcommand{\NN}{\mathbf{N}}
\newcommand{\PP}{\mathbf{P}}
\newcommand\QQ{\mathbf{Q}}
\newcommand\RR{\mathbf{R}}
\newcommand\ZZ{\mathbf{Z}}

\def\to{\mathchoice{\longrightarrow}{\rightarrow}{\rightarrow}{\rightarrow}}
\def\from{\mathchoice{\longleftarrow}{\leftarrow}{\leftarrow}{\leftarrow}}
\def\mto{\mathchoice{\longmapsto}{\mapsto}{\mapsto}{\mapsto}}
\def\mfrom{\mathrel{\from\joinrel\mapstochar\,}}
\def\hto{\mathrel{\lhook\joinrel\to}}
\def\hfrom{\mathrel{\from\joinrel\rhook}}

\def\implies{\mathchoice{\Longrightarrow}{\Rightarrow}{\Rightarrow}{\Rightarrow}}
\def\implied{\mathchoice{\Longleftarrow}{\Leftarrow}{\Leftarrow}{\Leftarrow}}
\def\iff{\mathchoice{\Longleftrightarrow}{\Leftrightarrow}{\Leftrightarrow}{\Leftrightarrow}}

\def\To#1{\mathchoice{\xrightarrow{\labelstyle~~#1~}}{\stackrel{#1}{\longrightarrow}}{}{}}
\def\Mto#1{\mathrel{\mapstochar\To{#1}}}
\def\From#1{\mathchoice{\xleftarrow{\labelstyle~~#1~}}{\stackrel{#1}{\longleftarrow}}{}{}}

\newcommand{\isomto}{\overset{\sim}{\longrightarrow}}
\newcommand{\isomfrom}{\overset{\sim}{\longleftarrow}}

\newcommand{\IFF}{\unskip~$\iff$~\ignorespaces}
\newcommand{\THEN}{\unskip~$\implies$~\ignorespaces}

\let\leq\leqslant
\let\geq\geqslant
\let\le\leqslant
\let\ge\geqslant

\def\cf{{cf.}\kern.3em}
\def\Cf{{Cf.}\kern.3em}
\def\ie{{i.e.}\kern.3em}
\def\iev{{i.e.},\xspace}
%\def\loccito{loc.\kern2pt cit.\xspace}
%\def\loccitn{\textit{loc.\kern2pt cit.}\xspace}
%\def\loccit{\sisi{\loccito}{\loccitn}}
\def\EGA{\textup{EGA}\kern.2em}
\def\SGA{\textup{SGA}\kern.3em}

\let\oldS\S
\def\S{\oldS\kern .15em}
\def\ptbl{.\kern .2em}
\def\resp{resp.\kern.3em}
\def\Exp{Exp.\kern.3em}

\let\bar\overline
\let\hat\widehat
\let\tilde\widetilde
\let\dpl\displaystyle
\let\ldots\dots
\let\emptyset\varnothing
\let\skpt\skippointrait
\let\epsilon\varepsilon

\def\sfrac#1#2{{#1}/{#2}}

\DeclareMathOperator{\cd}{cd}
%\DeclareMathOperator{\Hhom}{\mathcal{H}\!\mathit{om}}
\DeclareMathOperator{\Hhom}{\mathbf{Hom}}
%\DeclareMathOperator{\Eext}{\mathcal{E}\!\mathit{xt}}
\DeclareMathOperator{\Eext}{\mathbf{Ext}}
\DeclareMathOperator{\Spec}{Spec}
\DeclareMathOperator{\Tr}{Tr}

\let\emptyset=\varnothing

%%%%%%%%%%%%%%%
%%% expose XVII
%%%%%%%%%%%%%%%

\makeatletter
\let\old@optexpose\optexpose
\newcommand{\empty@optexpose}[1]{}
\newcommand{\paren@optexpose}[1]{\old@optexpose{#1}(}
\newcommand{\refnoexpose}[1]{\let\optexpose\empty@optexpose\ref{#1}\let\optexpose\old@optexpose}%\refnoexpose{X.cor5.2} engendre "5.2" au lieu de "X 5.2"
\newcommand{\neweqref}[1]{\let\optexpose\paren@optexpose\textup{\ref{#1})}\let\optexpose\old@optexpose}%\neweqref{XII.eq4.1} engendre "XII (4.1)" au lieu de "(XII 4.1)"
\makeatother
\newcommand{\refvb}[2][VI${}_{\text{B}}$ ]{\sisi{#1\refnoexpose{#2}}{\ref{#2}}}

\newcommand{\leftexp}[2]{{\vphantom{#2}}^{#1}{#2}}
\newcommand{\AAA}{\mathbf{A}}
\newcommand{\Sf}{\mathfrak{S}}
\DeclareMathOperator*{\quotprojlim}{\text{``$\varprojlim$''}}
\DeclareMathOperator*{\quotinjlim}{\text{``$\varinjlim$''}}
\DeclareMathOperator{\Fl}{F\sisi{\ell}{l}}
\DeclareMathOperator{\Pl}{P\sisi{\ell}{l}}
%\DeclareMathOperator{\cl}{c\sisi{\ell}{l}}
\DeclareMathOperator{\tdf}{tdf}
\DeclareMathOperator{\rg}{rg}
\newcommand{\iInd}{\sisi{I}{i}nd}
\newcommand{\iind}{ind}
\newcommand{\ppro}{pro}
\newcommand{\tTor}{\sisi{T}{t}or}
\newcommand{\ttor}{tor}
\newcommand{\chap}{Chap.}
\newcommand{\Ffex}{\Ff\mathit{ex}}
\newcommand{\sheafk}{\underline{k}}
\newcommand{\sheafn}{\underline{n}}
\newcommand{\sheafGamma}{\underline{\Gamma}}

\sisi{
\DeclareMathOperator{\simarrow}{\rule{.15mm}{1.8mm}\kern -3.4pt\rightsquigarrow}
\newcommand{\sheafell}{\underline{\ell}}
\newcommand{\sheafZZ}{\underline{\ZZ}}
}{
\newcommand{\nderef}[1]{la N.D.E.~(\ref{#1}) page~\pageref{#1}}
}

%%%%%%%%%%%%%%%%%
%%%%% expose XVII
%%%%%%%%%%%%%%%%%




%\externaldocument{../01/01.tex}
\externaldocument{../02/01.tex}
\externaldocument{../03/03.tex}
\externaldocument{../04/04.tex}

\setboolean{original}{false}

\def\labelenumi{(\theenumi)}
\def\labelenumii{(\theenumii)}
\def\labelenumiii{(\theenumiii)}
\def\labelenumiv{(\theenumiv)}

% \expose{i}

\begin{document}

\begin{center}
I. Préfaisceaux\\
A. Grothendieck et J.-L. Verdier
\end{center}

\selectlanguage{english}version : \showgitstatus\selectlanguage{french}

\tableofcontents

\else
\chapter{Préfaisceaux}
%\addtocontents{toc}{par A. Grothendieck et J.-L. Verdier}
\begin{center}
A. Grothendieck et J.-L. Verdier
\end{center}
\fi




Dans\pageoriginale les numéros 0 à 5 de cet exposé, on
présente les
propriétés élémentaires, et le plus souvent bien connues, des
catégories de préfaisceaux\footnote[1]{Le lecteur pourra aussi
  consulter SGA 3 I \S\S 1 à 3.}. Ces propriétés sont
utilisées constamment dans la suite du séminaire et leur
connaissance est essentielle pour la compréhension des exposés
suivants. Les démonstrations sont immédiates; elles sont le plus
souvent omises. Dans les numéros 6 à 9 sont abordés quelques thèmes
utilisés à différentes reprises dans la suite. Le lecteur pressé %%
pourra les omettre en première lecture. Le numéro 10 fixe la %%
terminologie employée. L'appendice 11 est %%
dû à N. Bourbaki.


\setcounter{section}{-1}
\section{Univers}\etiquette[0]{I.sec0}

Un \emph{univers} est un ensemble non-vide\nde{Cette définition
contredit l'appendice de Bourbaki \textit{infra}. Toutefois
l'intérêt des univers vides semble discutable...} $𝒰$ qui
jouit des propriétés suivantes:
\begin{itemize}
\item[($𝒰$ 1)] Si $x\in 𝒰$ et si $y\in x$ alors $y\in 𝒰$.

\item[($𝒰$ 2)] Si $x$, $y\in 𝒰$, alors $\{x,y\}\in 𝒰$.

\item[($𝒰$ 3)] Si $x\in 𝒰$, alors $\Pp(x)\in 𝒰$.

\item[($𝒰$ 4)] Si $(x_i, i\in I\sisi{\in 𝒰}{})$ est une
  famille d'éléments de $𝒰$\sisi{}{ et $I\in𝒰$}, alors $\cup_{i\in I}x_i \in
  𝒰$.
\end{itemize}

Des\pageoriginale axiomes précédents on déduit facilement les
propriétés:
\begin{itemize}
\item Si $x\in 𝒰$, l'ensemble $\{x\}$ appartient à $𝒰$.

\item Si $x$ est un sous-ensemble de $y\in 𝒰$, alors $x\in 𝒰$.

\item Si $x$, $y\in 𝒰$, le couple $(x,y)=\{\{x,y\},x\}$ (définition
  de Kuratowski) est un élément de $𝒰$.

\item Si $x$, $y\in 𝒰$, la réunion $x\cup y$ et %%
le produit $x\times
   y$ sont des éléments de $𝒰$.

\item Si $(x_i, i\in I\in 𝒰)$ est une famille
  d'éléments de $𝒰$,
   le produit $\prod_{i\in I}x_i$ est un élément de $𝒰$.

\item Si $x\in 𝒰$, alors $\card(x)< \card(𝒰)$ (strictement). En
   particulier la relation $𝒰\in 𝒰$ n'est pas vérifiée.
\end{itemize}

On peut donc faire toutes les opérations usuelles de la théorie des
ensembles à partir des éléments d'un univers sans, pour cela, que le
résultat final cesse d'être un élément de l'univers.

La notion d'univers a pour premier intérêt de fournir une définition
des catégories usuelles: la catégorie des ensembles appartenant à
l'univers $𝒰$ ($𝒰$-$\Ens$), la catégorie des espaces
topologiques appartenant à l'univers $𝒰$, la catégorie des
groupes commutatifs appartenant à l'univers $𝒰$
($𝒰$-$\Ab$), la catégorie des catégories appartenant à
l'univers $𝒰$\ldots.

Cependant le seul univers connu est l'ensemble des symboles du type
$\{\{∅\}, \{\{∅\}, ∅\}\}$ \etc. (tous
les éléments de cet univers sont des ensembles finis et cet univers
est dénombrable). En particulier, on ne connaît pas d'univers qui
contienne un élément de cardinal infini. On est donc amené
à ajouter
aux axiomes de la théorie des ensembles l'axiome: $(𝒰 A)$
\emph{Pour tout ensemble $x$ il existe un univers $𝒰$ tel
que $x\in
  𝒰$.}

L'intersection\pageoriginale d'une famille d'univers étant un
univers, on en déduit immédiatement que tout ensemble est
élément
d'un plus petit univers. On peut montrer que l'axiome $(𝒰
A)$ est indépendant des axiomes de la théorie des ensembles.

On ajoutera aussi l'axiome:


\noindent \emph{$(𝒰 B)$ Soit $R\{x\}$ une relation et
$𝒰$ un univers. S'il
  existe un élément $y\in 𝒰$ tel que $R\{y\}$, alors
  $\tau_xR\{x\}\in 𝒰$.}

La non-contradiction des axiomes $(𝒰 A)$ et $(𝒰 B)$
par rapport aux autres axiomes de la théorie des ensembles n'est pas
démontrée ni démontrable, semble-t-il.

Soit $𝒰$ un univers et $c(𝒰)$ la borne supérieure
des cardinaux des éléments de $𝒰(c(𝒰)\leq
\card(𝒰))$. Le cardinal $c(𝒰)$ jouit des propriétés
suivantes:

$(FI)$ Si $a<c(𝒰)$, alors $2^a<c(𝒰)$.

$(FII)$ Si $(a_i, i\in I)$ est une famille de cardinaux strictement
inférieurs à $c(𝒰)$ et si $\card(I)$ est %%
strictement inférieur à $c(𝒰)$, $\Sigma_{i\in I} a_i <
c(𝒰)$.

Les cardinaux qui possèdent les propriétés $(FI)$ et $(FII)$ sont
appelés \emph{cardinaux fortement inaccessibles.}

Le cardinal 0 et le cardinal infini dénombrable sont fortement
inaccessibles.

L'axiome $(𝒰 A)$ implique :

\emph{$(𝒰 A')$ Tout cardinal est majoré strictement par un
  cardinal fortement inaccessible.}

On peut montrer (11) que réciproquement la non contradiction de
$(UA')$ implique la non contradiction de $(UA)$, et que la non
contradiction de ces axiomes entraîne celle de l'axiome
$(𝒰 B)$.

Appelons\pageoriginale \emph{ensemble artinien} tout ensemble $E$
tel qu'il n'existe pas de familles infinies $(x_n, n\in N)$ telle
que $x_\circ \in E$, $x_{n+1}\in x_n$. On peut alors montrer
[\loccit] qu'il y a une correspondance biunivoque entre les
cardinaux fortement inaccessibles et les univers artiniens, définie
ainsi: à tout cardinal fortement inaccessible $c$ on fait
correspondre l'unique univers artinien $U_c$ tel que
$$
\card(𝒰_c)=c{\quoi}.
$$
Soient $c<c'$ deux cardinaux fortement inaccessibles. On a:
$$
𝒰_c\in 𝒰'_{c'}{\quoi}.
$$
En particulier les univers artiniens de cardinaux inférieurs à un
cardinal donné forment un ensemble bien ordonné (pour la relation
d'appartenance). L'axiome $(𝒰 A')$ est équivalent à
l'axiome:

$(𝒰 A'')$. \emph{Tout ensemble artinien est élément d'un
univers
  artinien.}

Remarquons que tous les ensembles usuels $(\ZZ, \QQ,\ldots)$ sont
des ensembles artiniens.

\section{$𝒰$-catégories. Préfaisceaux d'ensembles}\etiquette[1]{I.sec1}

\setcounter{subsection}{-1}
\subsection{}\etiquette[1.0]{I.sec1.0}
\emph{Dans la suite du séminaire et sauf mention expresse du
  contraire, les univers considérés posséderont un élément
  de cardinal infini.} Soit $𝒰$ un univers. On dit qu'un ensemble
  est $𝒰$-\emph{petit} (ou, quand aucune confusion n'en résulte,
  \emph{petit}) s'il est \emph{isomorphe} à un élément de
  $𝒰$. On utilise aussi la terminologie: petit groupe, petit anneau,
  petite catégorie\nde{Une catégorie $C$ est vue comme un ensemble de flèches
  (muni du sous-ensemble des objets, vu comme ensemble des flèches
  identiques et des applications « source, but »). Ainsi, les expressions « $C$ est élément de $𝒰$ » ou
  « $C$ est $𝒰$-petite » font sens.}
  \ldots On supposera souvent, sans mention
  explicite, que les schémas, espaces topologiques, ensembles
  d'indices\ldots avec lesquels on travaille sont $\in 𝒰$,
où\pageoriginale tout au moins ont un cardinal $\in 𝒰$;
  cependant, de nombreuses catégories avec lesquelles on travaillera
  ne seront pas $\in 𝒰$.

\begin{enonce*}{Définition 1.1}\etiquette[1.1]{I.def1.1}
Soient $𝒰$ un univers et $C$ une catégorie. On dit que $C$
est une $𝒰$-catégorie si pour tout couple $(x,y)$ d'objets
de $C$, l'ensemble $\Hom_C(x,y)$ est $𝒰$-petit.
\end{enonce*}

\setcounter{subsection}{1}
\subsubsection{}\etiquette[1.1.1]{I.sec1.1.1}
Soient $C$ et $D$ deux catégories et $\text{Fonct}(C,D)$ la
catégorie des foncteurs de $C$ dans $D$. On vérifie immédiatement
les assertions suivantes:
\begin{enumerate}
\item[a)] Si $C$ et $D$ sont éléments d'un univers $𝒰$
  (\resp  $𝒰$-petites) la catégorie $\text{Fonct}(C,D)$ est un élément
  de $𝒰$ (\resp  est $𝒰$-petite).

\item[b)] Si $C$ est $𝒰$-petite et si $D$ est une $𝒰$-catégorie,
  $\text{Fonct}(C,D)$ est une $𝒰$-catégorie.
\end{enumerate}

\begin{enonce*}[remark]{Remarque 1.1.2}\etiquette[1.1.2]{I.rem1.1.2}
Soit $D$ une catégorie possédant les propriétés suivantes:
\begin{itemize}
\item[(C1)] L'ensemble $\ob(D)$ est contenu dans l'univers $𝒰$.

\item[(C2)] Pour tout couple $(x,y)$ d'objets de $D$, l'ensemble
  $\Hom_D(x,y)$ est un élément de $𝒰$.
\end{itemize}
(Les catégories usuelles construites à partir d'un univers
$𝒰$ possèdent ces deux propriétés: $𝒰$-$\Ens$,
$𝒰$-$\Ab$,\ldots). Soit $C$ une catégorie appartenant à
$𝒰$. Alors la catégorie $\text{Fonct}(C,D)$ ne possède pas
en général les propriétés (C1) et (C2). Par exemple la catégorie
$\text{Fonct}(C,𝒰$-$\Ens$) ne possède aucune des propriétés
(C1) et (C2). C'est ce qui justifie la définition adoptée de
$𝒰$-catégorie, de préférence à la notion plus restrictive
par les conditions (C1) et (C2) ci-dessus.
\end{enonce*}

\begin{enonce*}{Définition 1.2}\etiquette[1.2]{I.def1.2}
Soit\pageoriginale $C$ une catégorie. On appelle catégorie des %%
préfaisceaux d'ensembles sur $C$ relative à l'univers $𝒰$
(ou, lorsqu'aucune confusion n'en résulte, catégorie des
préfaisceaux sur $C$) la catégorie des foncteurs contravariants sur
$C$ à valeur dans la catégorie des $𝒰$-ensembles.
\end{enonce*}

On désigne par $\widehat{C}_{𝒰}$ (ou plus
simplement, lorsqu'aucune confusion n'en résulte, par
$\widehat{C}$\,) la catégorie des préfaisceaux
d'ensembles sur $C$ relative à l'univers $𝒰$. Les objets de
$\widehat{C}_{𝒰}$ sont appelés
$𝒰$-préfaisceaux (ou plus simplement préfaisceaux) sur $C$.
Lorsque $C$ est $𝒰$-petite, la catégorie
$\widehat{C}_{𝒰}$ est une
$𝒰$-catégorie. Lorsque $C$ est une $𝒰$-catégorie,
$\widehat{C}_{𝒰}$ n'est pas nécessairement une
$𝒰$-catégorie.


\begin{enonce*}[remark]{Construction-définition
    1.3}\etiquette[1.3]{I.condef1.3}
Soit $x$ un objet d'une $𝒰$-catégorie $C$. On appelle
$𝒰$-\emph{foncteur représenté par} $x$ le foncteur
$h_{𝒰}(x):C^\circ\rightarrow 𝒰$-$\Ens$ dont la
construction suit \footnote[1]{$C^\circ$ désignera toujours la
catégorie opposée
  à $C$.}. Soit $y$ un objet de $C$.
\begin{enumerate}
\item[a)] Si $\Hom_C(y,x)$ est un élément de $𝒰$, alors:
      $$
      h_{𝒰}(x)(y)=\Hom_C(y,x)
      $$

\item[b)] Supposons que $\Hom_C(y,x)$ ne soit pas un élément de $𝒰$
  et soit $R(Z,x,y)$ la relation: « L'ensemble $Z$ est but d'un
  isomorphisme $\Hom_C(y,x)\xrightarrow{\sim}Z$ ». On pose alors:
  $$
  h_{𝒰}(x)(y)=\tau_ZR(Z){\quoi}.
  $$
\end{enumerate}
(En vertu de l'axiome $(𝒰 B)$, $h_{𝒰}(x)(y)$ est un
élément de $𝒰$). Soit $R'(u,x,y)$ la relation: « $u$ est
une bijection
$$
u:\Hom_C(y,x)\xrightarrow{\sim} h_{𝒰}(x)(y)\mbox{
»}{\quoi}.
$$
On\pageoriginale pose alors:
$$
\varphi(y,x)=\tau_uR'(u){\quoi}.
$$

On remarquera qu'on $a$, dans les cas (a) et (b), un isomorphisme
canonique:
$$
\varphi(y,x):\Hom_C(y,x)\xrightarrow{\sim}
h_{𝒰}(x)(y){\quoi}.
$$
(Dans le cas a) $\varphi(y,x)$ est l'identité). Soit alors
$u:y\rightarrow y'$ une flèche de $C$. Le morphisme $u$ définit, par
la composition des morphismes, une application:
$$
\Hom_C(u,x):\Hom_C(y',x)\longrightarrow \Hom_C(y,x){\quoi}.
$$
On pose alors
$$
h_{𝒰}(x)(u)=\varphi(y,x)\Hom_C(x,u)\varphi(y,x)^{{}-1}{\quoi}.
$$

On vérifie immédiatement que $h_{𝒰}(x)$, ainsi défini, est
un foncteur $C^\circ\rightarrow 𝒰$-$\Ens$.
\end{enonce*}

\setcounter{subsection}{3} \setcounter{subsubsection}{0}
\subsubsection{}\etiquette[1.3.1]{I.sec1.3.1}
De même, on définit, à l'aide des isomorphismes $\varphi$, pour tout
morphisme $y:x\rightarrow x'$ un morphisme de foncteurs:
$$
h_{𝒰}(y):h_{𝒰}(x)\longrightarrow h_{𝒰}(x')
$$
et on vérifie immédiatement qu'on a défini ainsi un foncteur
$$
h_{𝒰}:C\longrightarrow
\widehat{C}_{𝒰}{\quoi}.
$$

\subsubsection{}\etiquette[1.3.2]{I.sec1.3.2}
Soit maintenant $𝒱$ un univers tel que $𝒰\subset
𝒱$. On a alors un foncteur canonique pleinement fidèle:
$$
𝒰\text{-}\Ens\hookrightarrow 𝒱\text{-}\Ens
$$
d'où un foncteur pleinement fidèle:
$$
\widehat{C}_{𝒰}\hookrightarrow
\widehat{C}_{𝒱}
$$
et\pageoriginale le diagramme:

%\def\objectstyle{\mathrm}--here
\[
\xymatrix@R=1.75cm@C=3.5cm@M=.2cm{
C\ar[r]^{h_{\underline{U}}} \ar[dr]_{h_{\underline{V}}}& C{\widehat{\phantom{i}}}_{\underline{U}}\ar@{_{(}->}[d]\\
& C\widehat{\phantom{i}}_{\underline{V}} }
\]
est commutatif à isomorphisme canonique près. Lorsque $C$ est un
élément de $𝒰$, cet isomorphisme canonique est l'identité.

\subsubsection{}\etiquette[1.3.3]{I.sec1.3.3}
\emph{Dans la pratique, l'univers $𝒰$ est fixé
  une fois pour toutes et n'est pas mentionné.} On utilise alors les
  notations $\widehat{C}$ (pour la catégorie des
  $𝒰$-préfaisceaux d'ensembles) et
$$
h:C\longrightarrow \widehat{C}{\quoi}.
$$
Pour tout objet $x$ de $C$, le préfaisceau $h(x)$ est appelé le
\emph{préfaisceau représenté par} $x$ et nous identifierons toujours
la valeur en $y\in \ob(C)$ du préfaisceau $h(x)$ avec $\Hom_C(y,x)$.

\begin{enonce*}{Proposition 1.4}\etiquette[1.4]{I.prop1.4}
Soient $C$ une $𝒰$-catégorie, $F$ un préfaisceau sur $C$ et
$X$ un objet de $C$. Il existe un isomorphisme fonctoriel en $X$ et
en $F$:
$$
i:\Hom_{\widehat{C}}(h(X),F)\xleftarrow{\sim}
F(X){\quoi},
$$
l'application $i$ n'étant autre, lorsque $F$ est de la forme $h(Y)$,
que l'application:
$$
h:\Hom(h(X), h(Y))\longleftarrow \Hom(X,Y){\quoi}.
$$
En particulier le foncteur $h$ est pleinement fidèle.
\end{enonce*}

\setcounter{subsection}{4} \setcounter{subsubsection}{0}
\subsubsection{}\etiquette[1.4.1]{I.sec1.4.1}
Cette proposition justifie les abus de langage habituels identifiant
un objet de $C$ et le foncteur contravariant correspondant. Un
préfaisceau\pageoriginale isomorphe à un objet image par $h$ (ou, en
utilisant l'abus de langage signalé ci-dessus, isomorphe à un objet
de $C$) est appelé \emph{préfaisceau représentable.}

\subsubsection{}\etiquette[1.4.2]{I.sec1.4.2}
Soit $𝒱$ un univers contenant un univers $𝒰$. La
catégorie des $𝒰$-préfaisceaux est une sous-catégorie
\emph{pleine} de la catégorie des $𝒱$-préfaisceaux et par
suite un $𝒰$-préfaisceau est représentable si et seulement
si son image dans la catégorie des $𝒱$-préfaisceaux est un
$𝒱$-préfaisceau représentable.

\section{Limites projectives et inductives}\etiquette[2]{I.sec2}

Soient $C$ une $𝒰$-catégorie et $I$ une petite catégorie.
Notons $\Fonct(I,C)$ la catégorie des foncteurs de $I$ dans $C$. La
catégorie $\Fonct(I,C)$ est une $𝒰$-catégorie. À tout objet
$X$ de $C$ associons la sous-catégorie $𝒳$ de $C$ ayant pour
seul objet l'objet $X$ et pour seule flèche l'identité de $X$.
Désignons par $i_X:𝒳\rightarrow C$ le foncteur d'inclusion.
Il existe un et un seul foncteur $e_X:I\rightarrow 𝒳$, et
nous désignerons par $k_X:I\rightarrow C$ le foncteur $i_X\circ
e_X$. (On dira que $k_X$ est le foncteur constant de valeur $X$). La
correspondance $X\mapsto k_X$ est visiblement fonctorielle en $X$,
ce qui nous permet de définir, pour tout foncteur $G:I\rightarrow
C$, le préfaisceau $X\mapsto \Hom_{\Fonct(I,C)}(k_X,G)$.

\begin{enonce*}{Définition 2.1}\etiquette[2.1]{I.def2.1}
On appelle limite projective de $G$ et on note $\varprojlim_I G$ le
préfaisceau:
$$
X\longmapsto \Hom_{\Fonct(I,C)}(k_X,G){\quoi}.
$$
Lorsque le préfaisceau $\varprojlim_I G$ est représentable, on
désigne encore par $\varprojlim_I G$ un\pageoriginale objet de $C$
qui le représente. L'objet $\varprojlim_I G$ n'est donc défini qu'à
isomorphisme près. Lorsqu'aucune confusion n'en résulte on emploie
la notation abrégée $\varprojlim G$.
\end{enonce*}

Pour $G$ variable, $\varprojlim G$ est un foncteur de la catégorie
$\Fonct(I,C)$ à valeurs dans $\widehat{C}$.

\setcounter{subsection}{1} \setcounter{subsubsection}{0}
\subsubsection{}\etiquette[2.1.1]{I.sec2.1.1}
On définit de même par symétrie (renversement du sens des flèches
dans $C$) la limite inductive d'un foncteur: c'est un foncteur
covariant sur $C$ à valeur dans la catégorie des
$𝒰$-ensembles. Nous emploierons les notations $\varinjlim_I
G$ ou bien $\varinjlim G$.

On notera que les produits, produits fibrés, noyaux sont des limites
projectives. De même les sommes, sommes amalgamées, conoyaux sont
des %%
limites inductives.

\begin{enonce*}{Définition 2.2}\etiquette[2.2]{I.def2.2}
Soient $I$ et $C$ deux catégories et $G:I\rightarrow C$ un foncteur.
On dit que la limite projective de $G$ est représentable s'il existe
un univers $𝒰$ tel que:
\begin{enumerate}
\renewcommand{\labelenumi}{\theenumi)}
\item La catégorie $I$ soit $𝒰$-petite.

\item La catégorie $C$ soit une $𝒰$-catégorie.

\item Le préfaisceau $\varprojlim G$ à valeurs dans la
  catégorie des $𝒰$-ensembles soit représentable.
\end{enumerate}
\end{enonce*}

Il résulte au ${\rm n^\circ}$ 1 que l'objet $\varprojlim G$
représentant le préfaisceau $\varprojlim G$ ne dépend pas, à
isomorphisme près, de l'univers $𝒰$. Notons aussi qu'il
existe un plus petit univers $𝒰'$ possédant les propriétés
(1) et (2), et que le préfaisceau $\varprojlim G$ est nécessairement
à valeurs dans la catégorie des $𝒰'$-ensembles.

\setcounter{subsection}{2} \setcounter{subsubsection}{0}
\subsubsection{}\etiquette[2.2.1]{I.sec2.2.1}
Soient\pageoriginale $I$ et $C$ deux catégories. On dit que
\emph{les $I$-limites projectives dans $C$ sont
  représentables} si pour tout foncteur $G:I\rightarrow C$, la
  limite projective de $G$ est représentable. Enfin soient $C$ une
  catégorie et $𝒰$ un univers. On dit que les
  $𝒰$-\emph{limites
  projectives dans $C$ sont représentables} si pour toute
  catégorie $I$ $𝒰$-petite et pour tout foncteur $G:I\rightarrow
  C$, la limite projective de $G$ est représentable.

\begin{enonce*}{Proposition 2.3}\etiquette[2.3]{I.prop2.3}
Soient $C$ une catégorie et $𝒰$ un univers\sisi{ non
vide}{}. Les assertions suivantes sont équivalentes:
\begin{enumerate}
\renewcommand{\theenumi}{\roman{enumi}}
\renewcommand{\labelenumi}{\theenumi)}
\item Les $𝒰$-limites projectives dans $C$ sont représentables.

\item Les produits, indexés par un petit ensemble sont
  représentables, et les noyaux de couples de flèches sont
  représentables.
\item Les produits, indexés par un petit ensemble sont
  représentables, et les produits fibrés sont représentables.
\end{enumerate}
\end{enonce*}

\begin{enonce*}[remark]{Preuve}
Il suffit de remarquer qu'il existe un isomorphisme fonctoriel en
$G$
$$
\varprojlim G=\Ker (\prod_{i\in \ob(I)} %%
G(i)\rightrightarrows \prod_{u\in \rm{F\ell}(I)} G(\but(u))){\quoi},
$$
le couple de flèches étant défini par les morphismes
\begin{align*}
\prod_{i\in \Ob(I)}  G(i) %%
 & \xrightarrow{\pr_{\but(u)}}{} G(\but(u))\\
\prod_{i\in \Ob(I)} G(i) %%
 &\xrightarrow{\pr_{\text{source}(u)}}{} G(\text{source
}(u))\xrightarrow{G(u)} G(\but(u)){\quoi}.
\end{align*}
De plus, il est clair que $\Ker (X  \displaystyle \Lright^{u}_{v}Y)=
X\prod_{X\times Y} X$ les deux morphismes de $X$ dans $X\times Y$
étant $\id_X \times u$ et $\id_X\times v$.
\end{enonce*}

\setcounter{subsection}{3} \setcounter{subsubsection}{0}
\subsubsection{}\etiquette[2.3.1]{I.sec2.3.1}
Il existe évidemment des définitions et assertions analogues pour
les limites inductives que nous n'expliciterons pas. De même pour
les limites projectives\pageoriginale et inductives finies (\ie
relatives à une %%
catégorie $I$ finie).


\begin{enonce*}{Corollaire 2.3.2}\etiquette[2.3.2]{I.cor2.3.2}
Désignons par $𝒰$-$\Ens$ la catégorie des
$𝒰$-ensembles et par $𝒰$-$\Ab$ la catégorie des
objets groupes abéliens de $𝒰$-$\Ens$. Les
$𝒰$-limites projectives et inductives dans
$𝒰$-$\Ens$ et dans $𝒰$-$\Ab$ sont représentables.
\end{enonce*}

\begin{enonce*}{Proposition 2.3.3}\etiquette[2.3.3]{I.prop2.3.3}
Soient $I$ une petite catégorie, $F:I\rightarrow 𝒰$-$\Ens$
un foncteur et \linebreak %%
$G\in\ob$ $I$ un petit ensemble d'objets de $I$ tel que pour tout
objet $X$ de $I$ il existe un objet $Y\in G$ et un morphisme
$Y\rightarrow X$ (\resp  $X\rightarrow Y$). Alors $\varprojlim_I F$
(\resp  $\varinjlim_I F$) est représentable et on a
$\card(\varprojlim_I F)\leq \Pi_{Y\in G}\card(F(Y))$ (\resp
$\card(\varinjlim_I F)\leq \Sigma_{Y\in G}\card(F(Y))$).
\end{enonce*}

\begin{enonce*}[remark]{Preuve}
On vérifie immédiatement que $\varprojlim_I F'$ (\resp $\varinjlim_I
F'$) est isomorphe à un sous-objet (\resp  un quotient de $\Pi_{Y\in
G}F(Y)$ (\resp  $\amalg_{Y\in G}F(Y)$); d'où la proposition.
\end{enonce*}

\begin{enonce*}{Définition 2.4.1}\etiquette[2.4.1]{I.def2.4.1}
Soient $C$ une catégorie où les limites projectives (\resp
inductives) finies soient représentables et $u:C\rightarrow C'$ un
foncteur. On dit que $u$ est exact à gauche (\resp  à droite) s'il
« commute » aux limites projectives (\resp  inductives) finies.
Un foncteur exact à gauche et à droite est appelé un foncteur exact.
\end{enonce*}

\setcounter{subsection}{4} \setcounter{subsubsection}{1}
\subsubsection{}\etiquette[2.4.2]{I.sec2.4.2}
Il résulte de 2.3 que, pour qu'un foncteur soit exact\sisi{}{ à gauche}, il faut et il %% étiquette !
suffit qu'il transforme l'objet final (= produit vide) en l'objet
final, le produit de deux objets en le produit des deux objets
images, et le noyaux des couples de deux flèches en le noyau des
couples images ou encore, qu'il transforme l'objet final en l'objet
final et les produits fibrés en produits fibrés\sisi{. (O}{ (o}n
suppose que dans $C$ les $\varprojlim$ finies sont
représentables\sisi{.)}{).}

\setcounter{subsection}{5} \setcounter{subsubsection}{-1}
\subsubsection{}\etiquette[2.5.0]{I.sec2.5.0}
Soient\pageoriginale $I$, $J$ et $C$ trois catégories, $G:I\times
J\rightarrow C$ un foncteur (\ie un foncteur de $I$ à valeur dans la
catégorie des foncteurs de $J$ dans $C$). Supposons que les limites
projectives des foncteurs
\begin{align*}
G_i:&J\longrightarrow C\qquad i\in \ob(I)\\
    & j\longmapsto G(i\times j)
\end{align*}
soient représentables, et que le foncteur:
$$
i\longmapsto \varprojlim G_i
$$
admette une limite projective représentable. Il est clair qu'alors
le foncteur $G$ admet une limite projective représentable et qu'on a
un isomorphisme canonique:
$$
\varprojlim_{I\times J} G\\
\xrightarrow{\sim}\varprojlim_I \varprojlim_J G_i{\quoi},
$$
et que par suite on a un isomorphisme canonique
$$
\varprojlim_I \varprojlim_J G_i\xrightarrow{\sim}\varprojlim_J
\varprojlim_I G_j{\quoi}.
$$
Nous dirons par la suite plus brièvement que les limites projectives
commutent aux limites projectives. On voit de même que les limites
inductives commutent aux limites inductives. Mais il n'est pas vrai
en général que les limites inductives commutent aux limites
projectives.

\setcounter{subsection}{5}
\begin{enonce*}{Définition 2.5}\etiquette[2.5]{I.def2.5}
Soient $C$ une catégorie à produits fibrés représentables, $I$ une
catégorie, $G:I\rightarrow C$ un foncteur, $g:G\rightarrow X$ un
morphisme de $G$ dans un objet de $C$ (\ie un morphisme de $G$ dans
le foncteur constant associé à $X$), $m:Y\rightarrow X$ un morphisme
de $C$. Soit\pageoriginale $\fprod{G}{Y}{X}:I\rightarrow C$ le
foncteur
$$
i\longmapsto \fprod{G(i)}{Y}{X} {\quoi}.
$$
On dit que la limite inductive de $G$ est universelle si pour tout
objet $X$, tout morphisme $g:G\rightarrow X$, tout morphisme
$m:Y\rightarrow X$,
\begin{enumerate}
\item[a)] la limite inductive du foncteur $\fprod{G}{Y}{X}$ est
  représentable,

\item[b)] le morphisme canonique $\varinjlim (\fprod{G}{Y}{X}
  )\rightarrow \fprod{(\varinjlim G)}{Y}{X}$ est un isomorphisme.
\end{enumerate}
\end{enonce*}

\begin{enonce*}{Proposition 2.6}\etiquette[2.6]{I.prop2.6}
Soit $𝒰$ un univers. Les limites inductives dans
$𝒰$-$\Ens$ qui sont représentables (en particulier, les
$𝒰$-limites inductives $(2.2.1)$ dans $𝒰$-$\Ens$) %% étiquette !
sont universelles.
\end{enonce*}
Nous utiliserons aussi un autre résultat de commutation entre
limites projectives et inductives que nous allons présenter
maintenant.

\begin{enonce*}{Définition 2.7}\etiquette[2.7]{I.def2.7}
Une catégorie $I$ est pseudo-filtrante lorsqu'elle possède les
propriétés suivantes:
\begin{itemize}
\item[PS 1)] Tout diagramme de la forme:
\[
\xymatrix@R=.5cm@C=1cm{
       & j\\
i \ar[ru]\ar[dr]     &  \\
       & j'
}
\]
peut être inséré dans un diagramme commutatif:
\[
\xymatrix@R=.5cm@C=1cm{
  & j \ar[rd]&\\
i\ar[ur]\ar[rd] &   & k\\
  & j'\ar[ru]&
}
\]

\item[PS 2)] Tout diagramme de la forme:
\[
\xymatrix@R=.75cm@C=1.5cm@M=.25cm{
i\ar@<2.5pt>[r]^{u}\ar@<-2.5pt>[r]_{v} & j }
\]
%%
  peut \pageoriginale être inséré dans un diagramme:
\[
\xymatrix@C=1.75cm@M=.25cm{ i\ar@<2.5pt>[r]^{u}\ar@<-2.5pt>[r]_{v} &
j\ar[r]%%
^{w} & k\\
}
\]
  tel que
  $$
  w\circ u=w\circ v{\quoi}.
  $$
\end{itemize}
Une catégorie $I$ est dite filtrante si elle est pseudo-filtrante,
non vide et connexe, \ie si deux objets quelconques de $I$ peuvent
être reliés par une suite de flèches (on n'impose aucune condition
sur le sens des flèches); cela signifie aussi, en présence de {\rm
PS 2}, que $I\neq ∅$ et que pour deux objets $a$, $b$ de
$I$, il existe toujours un objet $c$ de $I$ et des flèches $a\sisi{$
$}{\rightarrow}c$ et $b\sisi{$ $}{\rightarrow}c$. On dit aussi
qu'une catégorie $I$ est cofiltrante si $I^\circ$ est filtrante.
\end{enonce*}

\begin{enonce*}{Exemple 2.7.1}\etiquette[2.7.1]{I.exem2.7.1}
Si dans $I$ les sommes amalgamées (\resp  les sommes de deux objets)
et les conoyaux de doubles flèches sont représentables, alors $I$
est pseudo-filtrante (\resp  filtrante).
\end{enonce*}

\begin{enonce*}{Proposition 2.8}\etiquette[2.8]{I.prop2.8}
Soit $𝒰$ un univers. Les $𝒰$-limites inductives
filtrantes dans $𝒰$-$\Ens$ commutent aux limites projectives
finies.
\end{enonce*}

On se ramène immédiatement à démontrer que les $𝒰$-limites
filtrantes commutent aux produits fibrés. La démonstration est
laissée au lecteur. On pourra utiliser la description de la limite
donnée par le

\begin{enonce*}{Lemme 2.8.1}\etiquette[2.8.1]{I.lem2.8.1}
Soient $I$ une petite catégorie filtrante, $i\mapsto X_i$ un
foncteur de $I$ dans $𝒰$-$\Ens$. Sur l'ensemble somme
$\amalg_{i\in
 \ob I} X_i$, soit $R$ la relation:
\begin{itemize}
\item[(R)] Deux éléments $x_i\in X_i$ et $x_j\in X_j$ sont
  reliés s'il existe un objet $k\in \ob I$ et deux morphismes
  $u:i\rightarrow k$ et $v:j\rightarrow k$ tels que les images dans
  $X_k$ de $x_i$ et de $x_j$ par les applications de transition $u$ et
  $v$ respectivement soient égales.
\end{itemize}
Alors:\pageoriginale
\begin{enumerate}
\renewcommand{\labelenumi}{\theenumi)}
\item $R$ est une relation d'équivalence.

\item Le quotient $\amalg_{i\in \ob I}X_i/R$ est canoniquement isomorphe
  à $\varinjlim_I X_i$.

\item Deux éléments $x_i\in X_i$ et $x_j\in X_j$ sont
  respectivement équivalents suivant $R$ à deux éléments d'un
  même $X_k$.

\item Pour tout $i\in \ob I$, deux éléments $\alpha$ et $\beta$ de $X_i$ sont
  équivalents suivant $R$ si et seulement s'il existe un morphisme
  $u:i\rightarrow j$ tel que $u(\alpha)=u(\beta)$.
\end{enumerate}
\end{enonce*}

L'assertion 1) résulte de (PS 1) (2.7). Pour démontrer 2), on %% étiquette !
vérifie que $\amalg_{i\in \ob I}X_i/R$ possède la propriété
universelle de la limite inductive (on n'utilise que (PS 1)).
L'assertion 3) résulte du fait que $I$ est connexe. L'assertion 4)
résulte de (PS 2).

\begin{enonce*}{Corollaire 2.9}\etiquette[2.9]{I.cor2.9}
Soit $%%
\gamma$ une espèce de structure algébrique « définie par limites
projectives finies ». (Le lecteur est prié de donner un sens
mathématique à la phrase précédente. Notons seulement que les
structures de groupes, groupes abéliens, anneaux, modules, etc... %%
sont de telles structures). Désignons par $𝒰$-$\gamma$ la
catégorie des $\gamma$-objets de $𝒰$-$\Ens$. Le foncteur qui
à chaque objet de $𝒰$-$\gamma$ associe l'ensemble
sous-jacent, commute aux $𝒰$-limites filtrantes. Par suite
les $𝒰$-limites filtrantes dans $𝒰$-$\gamma$
commutent aux limites projectives finies.
\end{enonce*}

\begin{enonce*}{Corollaire 2.10}\etiquette[2.10]{I.cor2.10}
Les $𝒰$-limites pseudo-filtrantes dans $𝒰$-$\Ab$
commutent aux limites projectives finies.
\end{enonce*}

(On se ramène aux limites filtrantes en décomposant la catégorie
d'indices en composantes connexes).

Notons\pageoriginale maintenant un résultat que nous utiliserons
constamment: Soient $C$ et $C'$ deux $𝒰$-catégories, $u$ et
$v$ deux foncteurs $\xymatrix@C=1.25cm{ C  \ar@<-.6ex>[r]_u & C'
\ar@<-.6ex>[l]_v }$ $u$ \emph{adjoint à gauche} de $v$. Rappelons
que ceci veut dire qu'il existe un isomorphisme entre les
bifoncteurs à valeur dans $𝒰$-$\Ens$:
$$
\Hom_{C'}(u(X),X')\xrightarrow{\sim}\Hom_{C}(X,v(X')){\quoi}.
$$

\begin{enonce*}{Proposition 2.11}\etiquette[2.11]{I.prop2.11}
Le foncteur $u$ commute aux limites inductives représentables; le
foncteur $v$ commute aux limites projectives représentables.
\end{enonce*}

Cette assertion signifie que pour toute catégorie $I$ et tout
foncteur $G:I\rightarrow C'$ tel que la limite projective de $G$
soit représentable, le foncteur $v\circ G$ admet une limite
projective représentable et que l'on a un isomorphisme canonique:
$$
v(\varprojlim G)\longrightarrow \varprojlim(v\circ G){\quoi}.
$$
Le lecteur explicitera de lui-même l'assertion concernant le
foncteur $u$.

Signalons enfin un calcul de limites projectives dans le cas où la
catégorie d'indices admet des produits.

\begin{enonce*}{Proposition 2.12}\etiquette[2.12]{I.prop2.12}
Soient $I$ et $C$ deux catégories, $𝒱$ un univers. On
suppose que les $𝒱$-limites projectives dans $C$ sont
représentables et qu'il existe dans $I$ une famille d'objets
$(i_\alpha)_{\alpha\in A}$, où $A$ est un élément de $𝒱$,
telle que:
\begin{enumerate}
\renewcommand{\labelenumi}{\theenumi)}
\item les produits $i_\alpha\times i_\beta$ soient représentables
  pour tout couple $(\alpha,\beta)\in A\times A$,

\item tout objet de $I$ s'envoie dans un au moins des $i_\alpha$.
\end{enumerate}


Alors pour tout contra foncteur $F$ de $I$ dans $C$
$$
F:I^\circ\longrightarrow C{\quoi},
$$
la limite projective de $F$ existe et il existe un isomorphisme
fonctoriel en $F$:
$$
\varprojlim F\xrightarrow{\sim} \Ker \left(\prod_{\alpha\in
  A} F(i_\alpha)\rightrightarrows \prod_{(\alpha,\beta)\in A\times A}
F(i_\alpha \times i_\beta)\right){\quoi},
$$
les\pageoriginale deux flèches étant définies par les projections
des produits $i_\alpha\times i_\beta$ sur les facteurs.
\end{enonce*}

\section{Propriétés d'exactitude de la catégorie des préfaisceaux}\etiquette[3]{I.sec3}

Soient $C$ une $𝒰$-catégorie, $\widehat{C}$ la
catégorie des préfaisceaux sur $C$. Les propriétés d'exactitude de
$\widehat{C}$ se déduisent toutes de la proposition
suivante:

\begin{enonce*}{Proposition 3.1}\etiquette[3.1]{I.prop3.1}
Les $𝒰$-limites projectives et inductives dans
$\widehat{C}$ sont représentables. Pour tout objet $X$
de $C$, le foncteur sur $\widehat{C}$:
$$
F\longmapsto F(X)\qquad F\in \widehat{C}
$$
commute aux limites inductives et projectives.

En d'autres termes, « les limites inductives et projectives dans
$\widehat{C}$ se calculent argument par argument ».
Citons quelques corollaires.
\end{enonce*}

\begin{enonce*}{Corollaire 3.2}\etiquette[3.2]{I.cor3.2}
Soit $\gamma$ une structure algébrique définie par limites
projectives finies. La catégorie des foncteurs contravariants à
valeurs dans $𝒰$-$\gamma$ est équivalente à la catégorie des
$\gamma$-objets de $\widehat{C}$.
\end{enonce*}

\begin{enonce*}{Corollaire 3.3}\etiquette[3.3]{I.cor3.3}
Un morphisme de $\widehat{C}$ qui est à la fois un
monomorphisme et un épimorphisme est un isomorphisme. Un morphisme
de $\widehat{C}$ se factorise de manière unique en un
épimorphisme suivi d'un monomorphisme. Les limites inductives dans
$\widehat{C}$ qui sont représentables, sont
universelles. Les $𝒰$-limites inductives filtrantes dans
$\widehat{C}$ commutent aux limites projectives
finies. Le foncteur canonique $C\rightarrow
\widehat{C}$ commute aux limites projectives
représentables.
\end{enonce*}

Plus\pageoriginale généralement on peut dire que la catégorie
$\widehat{C}$ hérite de toutes les propriétés de la
catégorie des $𝒰$-ensembles faisant intervenir des limites
inductives et projectives.

\setcounter{subsection}{4} \setcounter{subsubsection}{-1}
\subsubsection{}\etiquette[3.4.0]{I.sec3.4.0}
Soit $F$ un objet de $\widehat{C}$. On désigne par
$C/F$ la catégorie suivante: Les objets de $C/F$ sont les couples
formés d'un objet $X$ de $C$ et d'un morphisme $u$ de $X$ dans $F$.
Soient $(X,u)$ et $(Y,v)$ deux objets. Un morphisme de $(X,u)$ dans
$(Y,v)$ est un morphisme $g$ de $X$ dans $Y$ tel que le diagramme
ci-après soit commutatif:
\[
\xymatrix@=1.25cm{
X \ar[rr]^{g}\ar[rd]_{u}&   & Y\ar[ld]^{v}\\
            & F &.
}
\]

\begin{enonce*}{Proposition 3.4}\etiquette[3.4]{I.prop3.4}
Avec les notations de $(3.4.0)$, le foncteur source $C/F\rightarrow
\widehat{C}$ admet une limite inductive représentable
dans $\widehat{C}$. Le morphisme canonique:
$$
\mathop{\underrightarrow{\text{\rm lim source}}}_{C/F}\;
(.)\longrightarrow F
$$
est un isomorphisme.
\end{enonce*}


\begin{enonce*}{Corollaire 3.5}\etiquette[3.5]{I.cor3.5}
Soient $F$ et $H$ deux objets de $\widehat{C}$. Il
existe un isomorphisme canonique:
$$
\Hom(F,H)\xrightarrow{\sim}\varprojlim_{(X,u)\in \ob(C/F)}
H(X){\quoi}.
$$
\end{enonce*}

\begin{enonce*}[remark]{Preuve}
Le corollaire se déduit immédiatement de \hyperref[I.prop3.4]{3.4}
et de \hyperref[I.def1.2]{1.2}.
\end{enonce*}

\begin{enonce*}{Proposition 3.6}\etiquette[3.6]{I.prop3.6}
Soient\pageoriginale $C$ une $𝒰$-catégorie, $𝒱$ un
univers contenant $𝒰$ et
$\nobreak{i:\widehat{C}_{𝒰}\rightarrow
\widehat{C}_{𝒱}}$ le foncteur d'injection
naturel des catégories de préfaisceaux correspondantes $(1.3)$. Le
foncteur $i$ commute aux limites inductives et projectives. De plus,
pour tout objet $F$ de $\widehat{C}_{𝒰}$, tout
sous-objet de $i(F)$ est isomorphe à l'image par $i$ d'un unique
sous-objet de $F$ (de sorte qu'on obtient une bijection de
l'ensemble des sous-objets de $F$ avec l'ensemble des sous-objets de
$i(F)$).
\end{enonce*}

\section{Cribles}\etiquette[4]{I.sec4}

\begin{enonce*}{Définition 4.1}\etiquette[4.1]{I.def4.1}
Soit $C$ une catégorie. On appelle crible de la catégorie $C$ une
sous-catégorie pleine $D$ de $C$ possédant la propriété suivante:
tout objet de $C$ tel qu'il existe un morphisme de cet objet dans un
objet de $D$ est dans $D$. Soit $X$ un objet de $C$; on appelle (par
abus de langage) cribles de $X$ les cribles de la catégorie $C/X$.
\end{enonce*}

Soit $𝒰$ un univers tel que $C$ soit une
$𝒰$-catégorie. Soit $\widehat{C}$ la catégorie
de préfaisceaux correspondante. À tout crible de $X$ on associe un
sous-objet de $X$ dans $\widehat{C}$ de la manière
suivante: à tout objet $Y$ de $C$, on fait correspondre l'ensemble
des morphismes $f:Y\rightarrow X$ tels que l'objet $(Y,f)$
appartienne au crible.

\begin{enonce*}{Proposition 4.2}\etiquette[4.2]{I.prop4.2}
L'application définie ci-dessus, établit une %%
bijection entre l'ensemble des cribles de $X$ et l'ensemble des
sous-objets de $X$ dans $\widehat{C}$.
\end{enonce*}

\begin{enonce*}[remark]{Preuve}
Montrons seulement \sisi{qu'elle}{quelle} est l'application inverse. À tout
sous-foncteur $R$ de $X$ on associe la catégorie $C/R$ des objets de
$C$ au-dessus de $R$ \sisi{(3.4)}{(\hyperref[I.prop3.4]{3.4})}. On vérifie
immédiatement que $C/R$ est un crible de $X$.
\end{enonce*}

\begin{enonce*}[remark]{Remarque 4.2.1}\etiquette[4.2.1]{I.rem4.2.1}
On\pageoriginale voit de même que les cribles de $C$ sont en
correspondance biunivoque canonique avec l'ensemble des
sous-foncteurs du « foncteur final » sur $C$ (objet « final » de $\widehat{C}$).
\end{enonce*}

\setcounter{subsection}{2}
\subsection{}\etiquette[4.3]{I.sec4.3}
Soit $C$ une $𝒰$-catégorie. Par abus de langage nous
appellerons aussi cribles de $X$, les sous-objets de $X$ dans la
catégorie $\widehat{C}$. Cet abus de langage nous
permet pour tout préfaisceau $F$ et tout crible $R$ de $X$ de
définir $\Hom_{\widehat{C}}(R,F)$ comme étant
l'ensemble des %%
morphismes du foncteur $R$ dans $F$. On a d'ailleurs un isomorphisme
canonique fonctoriel en $F$ (3.5):
$$
\Hom_{\widehat{C}}(R,F)\xrightarrow{\sim}\varprojlim_{C/R}
F(.){\quoi},
$$
ce qui permet d'en donner une définition directe\nde{Plus
simplement, $C/R$ s'identifie à $R$ de sorte qu'on a la formule $
\Hom_{\widehat{C}}(R,F)\xrightarrow{\sim}\varprojlim_{R}
F{\quoi}, $ où $F$ désigne abusivement le composé de $F$ et du
foncteur « source » tautologique $R\rightarrow C/X\rightarrow
C$.}. De même, la proposition 4.2 nous permet de transposer aux
cribles les opérations usuelles sur les foncteurs. Citons:

\subsubsection{}\etiquette[4.3.1]{I.4.3.1}
\emph{Changement de base.} Soit $R$ un crible de $X$ et
$f:Y\rightarrow X$ un morphisme d'objets de $C$. Le produit fibré
$\fprod{R}{Y}{X}$ est un crible de $Y$ qu'on appelle \emph{crible
déduit
  de $R$ par changement de base.} La sous-catégorie correspondante
de $C/Y$ est l'image inverse de la sous-catégorie de $C/X$ définie
par $R$ par le foncteur canonique $C/Y\rightarrow C/X$ défini par
$f$.

\subsubsection{}\etiquette[4.3.2]{I.sec4.3.2}
\emph{Relation d'ordre, intersection, réunion.} La relation
d'inclusion sur les sous-foncteurs de $X$ est une relation d'ordre.
On peut définir la réunion et l'intersection d'une famille de
cribles indexés par un ensemble quelconque comme étant la borne
supérieure et la borne inférieure de la famille de sous-préfaisceaux
correspondante.

\subsubsection{}\etiquette[4.3.3]{I.sec4.3.3}
\emph{Image, crible engendré.}\pageoriginale Soient
$(F_\alpha)_{\alpha\in A}$ une famille de préfaisceaux et pour
chaque $\alpha\in A$, un morphisme $f_\alpha:F_\alpha\rightarrow X$
où $X$ est un objet de $C$. On appelle \emph{image} de cette famille
de morphismes la réunion des images des $f_\alpha$. L'image de cette
famille est donc un crible de $X$. En particulier si tous les
$F_\alpha$ sont des objets de $C$, le crible image sera appelé le
crible engendré par les morphismes $f_\alpha$. La catégorie $C/R$
est la sous-catégorie pleine de $C/X$ formée des objets
$X'\rightarrow X$ au-dessus de $X$ tels qu'il existe un
$X$-morphisme de $X'$ dans un des $F_{\alpha}$.

Le lecteur pourra, à titre d'exercice, traduire en termes des
catégories $C/R$ les relations et opérations définies ici sur les
sous-foncteurs. Il constatera alors que ces relations et opérations
ne dépendent pas de l'univers $𝒰$ tel que $C$ soit une
$𝒰$-catégorie, et qu'elles sont par suite définies pour
toute catégorie, sans que l'on soit obligé de préciser l'univers
auquel les ensembles de morphismes appartiennent; ce qu'on pouvait
d'ailleurs prévoir \textit{a priori} grâce à 3.6.

\section{Fonctorialité des catégories de préfaisceaux}\etiquette[5]{I.sec5}

\setcounter{subsection}{-1}
\subsection{}\etiquette[5.0]{I.sec5.0}
Soient $C$, $C'$, $D$ trois catégories et $u:C\rightarrow C'$ un
foncteur. On désignera par $u^\ast$ le foncteur: \sisi{$$
u^\ast:\Hhom({C'}^\circ,D)\longrightarrow \Hhom(C^\circ, D),
G\longmapsto G\circ u{\quoi},
$$}{$$u^\ast:\left\{\begin{array}{ccc}
                      \Hhom({C'}^\circ,D) & \longrightarrow & \Hhom(C^\circ, D) \\
                      G & \longmapsto & G\circ u{\quoi}
                    \end{array}
\right.$$} obtenu en composant avec le foncteur $u$. Le foncteur
$u^\ast$ commute aux limites inductives et projectives.

\begin{enonce*}{Proposition 5.1}\etiquette[5.1]{I.prop5.1}
Supposons\pageoriginale que $C$ soit petite, et que, dans $D$, les
$𝒰$-limites inductives (\resp  projectives) soient
représentables. Le foncteur $u^\ast$ admet un adjoint à gauche
$u_{!}$ (\resp  à droite $u_\ast$). On a donc un isomorphisme:
\begin{gather*}
\Hom_{\Hhom(C^\circ,D)}(F,u^\ast G)\xrightarrow{\sim}
\Hom_{\Hhom({C'}^\circ,D)}(u_{!}F,G)\\
(\mbox{\resp  } \Hom_{\Hhom(C^\circ,D)}(u^\ast G,F)\simeq
\Hom_{\Hhom({C'}^\circ,D)}(G,u_\ast F){\quoi}).
\end{gather*}
\end{enonce*}

\begin{enonce*}[remark]{Preuve}
Nous n'indiquerons que la démonstration de l'existence du foncteur
adjoint à gauche. La partie \resp  de la proposition s'en déduira
alors formellement grâce aux isomorphismes:
$$
\Hhom(C^\circ,
D)\xrightarrow{\sim}\Hhom(C,D^\circ)^\circ\xrightarrow{\sim}\Hhom((C^\circ)^\circ,D^\circ)^\circ{\quoi}.
$$
\end{enonce*}

Soit $Y$ un objet de $C'$. Désignons par $I^Y_u$ la catégorie
suivante: Les objets de $I^Y_u$ sont les couples $(X,m)$ où $X$ est
un objet de $C$ et $m$ un morphisme $Y\rightarrow u(X)$. Soient
$(X,m)$ et $(X',m')$ deux objets de $I^Y_u$. Un morphisme de $(X,m)$
dans $(X',m')$ est un morphisme $\xi:X\rightarrow X'$ tel que
$m'=u(\xi)m$. La composition des morphismes se définit de la manière
évidente.

Soit $f:Y\rightarrow Y'$ un morphisme de $C'$. Le morphisme $f$
définit par composition un foncteur $I^f_u:I^Y_u\rightarrow
I^{Y'}_u$.

On a de plus un foncteur $\pr_Y:I^Y_u\rightarrow C$ qui, à l'objet
$(X,m)$ associe l'objet $X$. Notons qu'alors le diagramme:
\[\vcenter{
\xymatrix@=1.5cm{
       &I^Y_u \ar[r] \ar[rd]_{\pr_Y}& I^f_u\ar[r] & I^{Y'}_u \ar[ld]^{\pr_{Y'}}\\
       & &   C   & \\
}}\leqno{(*)}
\]
est commutatif.

Soit\pageoriginale maintenant $F$ un préfaisceau sur $C$ et posons:
$$
\displaylines{\rlap{(5.1.1)}\hfill u_{!}
  F(Y)=\varinjlim_{I^Y_u} F\sisi{}{\circ}\pr_Y(.){\quoi}.\hfill}
$$
La commutativité du diagramme ($^\ast$) et la fonctorialité de la
limite inductive font de $u_{!}F$ un préfaisceau sur $C'$. Montrons
que %%
le foncteur $u_{!}$ est un adjoint à gauche du foncteur $u^\ast$.
Pour cela montrons que pour tout préfaisceau $G$ sur $C'$, il existe
un isomorphisme fonctoriel
$$
\Hom(u_{!}F,G)\xrightarrow{\sim}\Hom(F,u^\ast G){\quoi}.
$$
Soit $\xi\in \Hom(u_{!}F,G)$. Pour tout objet $X$ de $C$, on a donc
un morphisme:
$$
\xi_X:u_{!}(u(X))\longrightarrow G(u(X)){\quoi}.
$$
Mais $(u(X),\id_{u(X)})$ est un objet de $I^{u(X)}_u$, et par
définition de la limite inductive, on a un morphisme canonique:
$$
F(X)\longrightarrow u_{!}F(u(X)){\quoi}.
$$
On en déduit pour tout objet $X$ de $C$ un morphisme:
$$
\eta_X:F(X)\longrightarrow G(u(X))
$$
qui est visiblement fonctoriel en $X$. D'où un morphisme
$$
\eta:F\longrightarrow u^\ast G{\quoi}.
$$

Réciproquement, soit $\eta\in \Hom(F,u^\ast G)$. On en déduit, pour
tout objet $Y$ de $C$, un morphisme de foncteur:
$$
\eta_Y:F\sisi{}{\circ}\pr_Y\longrightarrow (u^\ast G)\pr_Y{\quoi},
$$
d'où, en composant avec le morphisme évident du foncteur $(u^\ast
G)\pr_Y$ dans le foncteur constant $G(Y)$, un morphisme:
$$
\xi_Y:u_{!}F(Y)\longrightarrow G(Y)
$$
qui est fonctoriel en $Y$. D'où un morphisme:
$$
\xi:u_{!}F\longrightarrow G{\quoi}.
$$

Le\pageoriginale lecteur vérifiera que les deux applications ainsi
définies sont inverses l'une de l'autre, et achèvera ainsi la
démonstration.

\begin{enonce*}{Proposition 5.2}\etiquette[5.2]{I.prop5.2}
Supposons que dans $D$, les $𝒰$-limites inductives soient
représentables, les limites projectives finies soient représentables
et que les $𝒰$-limites inductives filtrantes commutent aux
limites projectives finies. Supposons de plus que dans $C$ les
limites projectives finies soient représentables et que le foncteur
$u$ soit exact à gauche $(\hyperref[I.cor2.3.2]{2.3.2})$. Alors les limites
projectives finies sont représentables dans $\Hhom(C^\circ, D)$ et
dans $\Hhom({C'}^\circ,D)$, et le foncteur $u_{!}$ est exact à
gauche.
\end{enonce*}

\begin{enonce*}[remark]{Preuve}
La première assertion est triviale. Démontrons la seconde. D'après
la démonstration de 5.1, pour tout préfaisceau $F$ sur $C$ et tout
objet $Y$ de $C'$ on a:
$$
u_{!}F(Y)\xrightarrow{\sim} \varinjlim_{I^Y_u}
F\sisi{}{\circ}\pr_Y{\quoi}.
$$
 Il suffit donc de montrer que la catégorie $(I^Y_u)^\circ$
 vérifie les axiomes (PS 1) et (PS 2) (2.7) et que cette catégorie
 est connexe. La vérification est laissée au lecteur.
\end{enonce*}

\setcounter{subsection}{2}
\subsection{}\etiquette[5.3]{I.sec5.3}
En particularisant ces résultats au cas où $D$ est la catégorie des
$𝒰$-ensembles, on obtient une suite de trois foncteurs:
$$
u_{!}, u^\ast, u_\ast {\quoi},
$$
qui est une « suite de foncteurs adjoints » dans le sens que
pour deux foncteurs consécutifs de la suite celui de droite est
adjoint à droite de l'autre. Leurs\pageoriginale propriétés
essentielles sont résumées dans la:

\begin{enonce*}{Proposition 5.4}\etiquette[5.4]{I.prop5.4}
Soient $C$ une petite catégorie, $C'$ une $𝒰$-catégorie, et
$u:C\rightarrow C'$ un foncteur.
\begin{enumerate}
\renewcommand{\labelenumi}{\theenumi)}
\item Le foncteur $u^\ast:\widehat{C'}\rightarrow
  \widehat{C}$ commute aux limites inductives et
  projectives.

\item Le foncteur $u_\ast:\widehat{C}\rightarrow
  \widehat{C'}$ commute aux limites projectives. Pour
  tout préfaisceau $F$ sur $C$ et tout objet $Y$ de $C'$, on a:
  $$
  u_\ast F(Y)\xrightarrow{\sim}
  \Hom_{\widehat{C}}(u^\ast(Y),F){\quoi}.
  $$

\item Le foncteur $u_{!}:\widehat{C}\rightarrow
  \widehat{C'}$ commute aux limites inductives. Le
  foncteur $u_{!}$ n'est défini qu'à isomorphisme près, mais on
  peut toujours le choisir tel que le diagramme
\[
\xymatrix@C=2.75cm@R=1.5cm{
  C \ar[r]^{u}\ar[d]_{h}& C' \ar[d]^{h^\prime}\\
C{\widehat{\phantom{\;i}}}\ar[r]^{u_{!}}&C'{\widehat{\phantom{\;i}}}}
\]
($h$ et $h'$ sont les foncteurs d'inclusion canonique) soit
commutatif.

Pour tout préfaisceau $F$ sur $C$, on a:
$$
u_{!}F\xrightarrow{\sim} \varinjlim_{C/F} h'\sisi{}{\circ}u{\quoi}.
$$

\item Si les limites projectives  finies sont représentables dans
  $C$ et si $u$ est exact à gauche $(2.3.2)$, le foncteur $u_{!}$
  est exact à gauche.
\end{enumerate}
\end{enonce*}

\begin{enonce*}[remark]{Preuve}
L'assertion (1) est triviale. L'assertion (2) se déduit du fait que
$u_\ast$ est un foncteur adjoint à droite par (\hyperref[I.prop2.11]{2.11}) et
(\hyperref[I.prop1.4]{1.4}). Il en est de même pour l'assertion (3) mais on
applique en plus (\hyperref[I.prop3.4]{3.4}). Enfin l'assertion (4) n'est
autre que \hyperref[I.prop5.2]{5.2} qu'on peut appliquer grâce à
(\hyperref[I.def2.7]{2.7}).
\end{enonce*}

\begin{enonce*}{Proposition 5.5}\etiquette[5.5]{I.prop5.5}
Soient\pageoriginale $C$ et $C'$ deux petites catégories et
$\xymatrix@C=1cm{
 C \ar@<-2.5pt>[r]_{u} & C'\ar@<-2.5pt>[l]_{v}
}$ un couple de foncteurs, où $v$ est adjoint à gauche de $u$. Il
existe alors des isomorphismes, compatibles avec les isomorphismes
d'adjonction:
\[
\xymatrix@R=.75cm{
v^{\ast} \ar[rr]^{\sim}&    & u_!      &\\
v_{\ast} \ar[rr]^{\sim}&    & u^{\ast} & . }
\]
\end{enonce*}

\begin{enonce*}[remark]{Preuve}
Il suffit d'exhiber un isomorphisme $v^\ast\xrightarrow{\sim}
u_{!}$; l'autre isomorphisme s'en déduira par adjonction. Soient $F$
un préfaisceau sur $C$ et $Y$ un objet de $C'$. On a alors:
$$
v^\ast F(Y)\xrightarrow{\sim} \Hom(v(Y),F){\quoi}.
$$
Puis en utilisant (\hyperref[I.prop3.4]{3.4}):
$$
\Hom(v(Y),F)\xrightarrow{\sim} \varinjlim_{C/F} \Hom(v(Y),.){\quoi}.
$$
Mais $v$ est adjoint à gauche de $u$ et par suite:
$$
\varinjlim_{C/F}\Hom(v(Y),.)\xrightarrow{\sim}\varinjlim_{C/F}\Hom(Y,u(.)){\quoi}.
$$
Utilisant alors (\hyperref[I.prop5.4]{5.4}.3)), il vient:
$$
\varinjlim_{C/F}\Hom(Y,u(.))\xrightarrow{\sim}
\Hom(Y,u_{!}F)\xrightarrow{\sim} u_{!}F(Y){\quoi}.
$$
On a donc déterminé, pour tout objet $Y$ de $C'$, un isomorphisme
$v^\ast F(Y)\xrightarrow{\sim} u_{!}F(Y)$ qui est visiblement
fonctoriel en $Y$ et en $F$, \cqfd
\end{enonce*}

\begin{enonce*}{Corollaire 5.5.1}\etiquette[5.5.1]{I.cor5.5.1}
Soit $u:C\rightarrow C'$ un foncteur qui admet un adjoint à gauche.
Le foncteur $u_{!}:\widehat{C}\rightarrow
\widehat{C'}$ commute aux limites projectives
(rappelons qu'il commute aux limites inductives par $(1.4.3)$).
\end{enonce*}

\begin{enonce*}[remark]{Remarque 5.5.2}\etiquette[5.5.2]{I.rem5.5.2}
On\pageoriginale trouve ainsi une « suite de \emph{quatre}
foncteurs adjoints » (\cf 5.3):
$$
v_{!}, v^\ast=u_{!}, v_\ast = u^\ast, u_\ast{\quoi},
$$
dont les trois premiers (\resp  derniers) commutent donc aux
$\varinjlim$ (\resp  $\varprojlim$).
\end{enonce*}

\begin{enonce*}{Proposition 5.6}\etiquette[5.6]{I.prop5.6}
Les hypothèses sont celles de $\hyperref[I.prop5.4]{5.4}$. Les conditions
suivantes sont équivalentes:
\begin{enumerate}
\renewcommand{\theenumi}{\roman{enumi}}
\renewcommand{\labelenumi}{\theenumi)}
\item Le foncteur $u$ est pleinement fidèle.

\item Le foncteur $u_{!}$ est pleinement fidèle.

\item Le morphisme d'adjonction
  $\id_{\widehat{C}}\rightarrow u^\ast u_{!}$ est un
  isomorphisme.

\item Le foncteur $u_\ast$ est pleinement fidèle.

\item Le morphisme d'adjonction $u^\ast u_\ast\rightarrow
  \id_{\widehat{C}}$ est un isomorphisme.
\end{enumerate}
\end{enonce*}

\begin{enonce*}[remark]{Preuve}
Il est clair que ii) $\Leftrightarrow$ iii) et iv) $\Leftrightarrow$
v) (propriétés générales des foncteurs adjoints) et que ii)
$\Rightarrow$ i) (\hyperref[I.prop5.4]{5.4}.3)). Montrons que i) $\Rightarrow$
iii). Les foncteurs $\id_{\widehat{C}}$, $u^\ast$ et
$u_{!}$ commutent aux limites inductives. D'après (3.4), il suffit
donc de démontrer que $H\rightarrow u^\ast u_{!} H$ est un
isomorphisme lorsque $H$ est représentable ce qui est évident.
Montrons que %%
iii) est équivalent à %%
v). Pour tout objet $H$ (\resp  $K$) de $\widehat{C}$
désignons par $\Phi(H):H\rightarrow u^\ast u_{!} H$ (\resp  par
$\Psi(K):u^\ast u_\ast\rightarrow K$) le morphisme d'adjonction. On
a alors un diagramme commutatif:
\[
\xymatrix@C=3.5cm@L=6pt{
\Hom_{\widehat{C}}(H,u^{\ast}u_{\ast}
K)\ar@{-}[d]_{\wr}\ar[r]^{\quad \Hom(H,\Psi(K))}   & \Hom_{\widehat{C}}(H,K)\\
\Hom_{\widehat{C}}(u_{!}H,u_{\ast}K)\ar@{-}[d]_{\wr} &  \\
\Hom_{\widehat{C}}(u^{\ast}u_{!}H,K)\ar[uru]_{\quad
   \Hom(\Phi(H),K)} & .
}
\]

Par\pageoriginale suite $\Phi(H)$ est un isomorphisme pour tout $H$
si et seulement si $\Psi(𝒦)$ est un isomorphisme pour tout
$K$, \cqfd
\end{enonce*}

\begin{enonce*}[remark]{Remarque 5.7}\etiquette[5.7]{I.rem5.7}
\begin{enumerate}
\renewcommand{\theenumi}{\alph{enumi}}
\renewcommand{\labelenumi}{\theenumi)}
\item Les équivalences ii) $\Leftrightarrow$ iii)
  $\Leftrightarrow$ iv) $\Leftrightarrow$ v) sont des résultats
  généraux sur les foncteurs adjoints.

\item La forme explicite de $u_{!}$ \resp  $u_\ast$ donnée dans la
  démonstration de 1.4 montre aussitôt que, sous les hypothèses
  générales de 1.1, si $u$ est pleinement fidèle, alors le
  morphisme d'adjonction $\id\rightarrow u^\ast u_{!}$ (\resp  $u^\ast
  u_\ast\rightarrow \id$) est un isomorphisme, \ie que $u_{!}$
  (\resp  $u_\ast$) est pleinement fidèle.
\end{enumerate}
\end{enonce*}

\setcounter{subsection}{8} \setcounter{subsubsection}{-1}
\subsubsection{}\etiquette[5.8.0]{I.sec5.8.0}
Soient $\gamma$ une espèce de structure algébrique définie par
limites projectives finies $𝒰$-$\gamma$-$\Ens$ la catégorie
des $\gamma$-objets de $𝒰$-$\Ens$, esj:
$𝒰$-$\gamma$-$\Ens\rightarrow 𝒰$-$\Ens$ le foncteur
ensemble sous-jacent (pour simplifier nous supposons que l'espèce de
structure envisagée a un seul ensemble de base). Soit $C$ une
catégorie. La composition avec esj fournit un foncteur noté
$$
\mbox{esj}^{\widehat{\phantom{,}}}:\Hhom(C^\circ,
𝒰\text{-}\gamma\text{-}\Ens)\longrightarrow
\widehat{C}{\quoi}.
$$
Comme dans $\widehat{C}$, les limites projectives se
calculent argument par argument, le foncteur
$\mbox{esj}^{\widehat{\phantom{,}}}$ se factorise en une
équivalence.
$$
\displaylines{\rlap{(5.8.1)}\hfill \Hhom(C^\circ,
  𝒰\text{-}\gamma\text{-}\Ens)\xrightarrow{\approx}
  \widehat{C}_\gamma{\quoi},\hfill}
$$
où $\widehat{C}_\gamma$ est la catégorie des
$\gamma$-objets de $\widehat{C}$, et un foncteur
encore noté
$$
\mbox{esj}^{\widehat{\phantom{,}}}:\widehat{C}_\gamma\longrightarrow
\widehat{C}{\quoi},
$$
et appelé le foncteur « \emph{préfaisceau d'ensembles
  sous-jacent} ».

\setcounter{subsubsection}{1}
\subsubsection{}\etiquette[5.8.2]{I.sec5.8.2}
Supposons\pageoriginale que le foncteur esj:
$𝒰$-$\gamma$-$\Ens\rightarrow 𝒰$-$\Ens$ admette un
adjoint à gauche $\Lib:𝒰$-$\Ens\rightarrow
𝒰$-$\gamma$-$\Ens$ \footnote[1]{Le foncteur $\Lib$ est le
foncteur « $\gamma$-objets libre
  engendré ». Exemple: groupe libre, groupe commutatif libre,
  $A$-module libre, \etc} (on peut montrer en fait que cette
condition est toujours satisfaite). La composition avec $\Lib$
fournit un foncteur
$$
\Lib^{\widehat{\phantom{,}}}:\widehat{C}\longrightarrow
\Hhom(C^\circ, 𝒰\text{-}\gamma\text{-}\Ens)
$$
et en composant avec l'équivalence (5.8.1), un foncteur encore noté
$$
\Lib^{\widehat{\phantom{,}}}:\widehat{C}\longrightarrow
\widehat{C}_\gamma
$$
et appelé le foncteur « préfaisceau de $\gamma$-objets libres
engendré ». Le foncteur $\Lib^{\widehat{\phantom{,}}}$ est adjoint
à gauche au foncteur $\esj^{\widehat{\phantom{,}}}$.

\begin{enonce*}{Proposition 5.8.3}\etiquette[5.8.3]{I.prop5.8.3}
Soient $\gamma$ une espèce de structure algébrique définie par
limites projectives finies telle que dans la catégorie des
$\gamma$-objets de $𝒰$-$\Ens$, les $𝒰$-limites
inductives soient représentables \footnote[2]{On peut montrer que
cette condition est
  toujours satisfaite.}, $C$ une catégorie appartenant à $𝒰$,
$C'$ une $𝒰$-catégorie, $u:C\rightarrow C'$ un foncteur.
Désignons par $\widehat{C}_\gamma$ (\resp
$\widehat{C'}_\gamma$) la catégorie des
$\gamma$-objets de $\widehat{C}$ (\resp
$\widehat{C'}$) et par ${u^\ast}^\gamma$ le foncteur
sur les $\gamma$-objets déduit du foncteur $u^\ast$. Il résulte de
\hyperref[I.prop5.1]{5.1}  et de l'équivalence 5.8.1 qu'il existe un
foncteur adjoint à gauche (\resp  à droite) au foncteur
${u^\ast}^\gamma$. Ce foncteur est noté $u_{!_\gamma}$ (\resp
$u_{\ast_\gamma}$).
\begin{enumerate}
\item Le foncteur ${u^\ast}^\gamma$ commute aux limites inductives et
  projectives. Le diagramme
\[\vcenter{
\xymatrix@C=3cm@R=1.5cm{
 \widehat{C'}_\gamma \ar[r]^{u^{\ast^{\gamma}}} \ar[d]_{\esj'\widehat{\phantom{a}}}&
\widehat{C}_\gamma  \ar[d]_{\esj^{\widehat{\phantom{a}}}}\\
 \widehat{C'} \ar[r]^{u^{\ast}}& \widehat{C}
}}\leqno{(*)}
\]
est\pageoriginale commutatif (${\esj'}^{\widehat{\phantom{,}}}$ et
$\esj^{\widehat{\phantom{,}}}$ désignent les foncteurs « ensemble
sous-jacent »),

\item Le foncteur $u_{\ast_\gamma}$ commute aux limites projectives. Le
  diagramme
\[\vcenter{
\xymatrix@C=3cm@R=1.5cm{
 \widehat{C}_\gamma \ar[r]^{u_{\ast_\gamma}} \ar[d]_{\esj\widehat{\phantom{a}}}&
  \widehat{C'}_\gamma\ar[d]_{\esj'\widehat{\phantom{a}}}\\
\widehat{C} \ar[r]^{u_{\ast}}&
\widehat{C'} }}\leqno{({*}{*})}
\]
est commutatif à isomorphisme près.

\item Le foncteur $u_{!_\gamma}$ commute aux limites inductives. Supposons
  que $\esj^{\widehat{\phantom{,}}}$
  (\resp  ${\esj'}^{\widehat{\phantom{,}}}$) admette un adjoint à
  gauche $\Lib^{\widehat{\phantom{,}}}$
  (\resp  ${\Lib'}^{\widehat{\phantom{,}}}$) $(5.8.2)$. Le diagramme
\[\vcenter{
\xymatrix@C=3.5cm@R=1.75cm{
 \widehat{C} \ar[r]^{u_{!}} \ar[d]_{\Lib\widehat{\phantom{a}}}&
  \widehat{C'} \ar[d]_{\Lib'\widehat{\phantom{a}}}\\
\widehat{C}_{\gamma} \ar[r]^{u_{!\gamma}} &
\widehat{C'}_{\gamma} }}\leqno{({*}{*}{*})}
\]
est commutatif à isomorphisme près.
\end{enumerate}

Supposons que $u_{!}$ commute aux limites projectives finies
(\hyperref[I.prop5.4]{5.4}~4) et \hyperref[I.prop5.6]{5.6}). Alors le diagramme
\[\vcenter{
\xymatrix@C=3.5cm@R=1.5cm{
 \widehat{C}_\gamma \ar[r]^{u_{!\gamma}} \ar[d]_{\esj\widehat{\phantom{a}}}&
  \widehat{C'}_\gamma\ar[d]_{\esj'\widehat{\phantom{a}}}\\
\widehat{C} \ar[r]^{u_{!}}& \widehat{C'}
}}\leqno{({*}{*}{*}{*})}
\]
est commutatif à isomorphisme près, et $u_{!\gamma}$ commute aux
limites projectives finies.
\end{enonce*}

\begin{enonce*}[remark]{Preuve}
L'assertion (1) est évidente. L'assertion (2) aussi car $u_\ast$ est
un adjoint à droite et par suite (\hyperref[I.prop2.11]{2.11}) commute aux limites
projectives et, en particulier, aux limites projectives finies; d'où
la commutativité du diagramme $({\ast}{\ast})$. La commutativité du
diagramme $({\ast}{\ast}{\ast})$ se déduit de l'unicité,
à\pageoriginale isomorphisme près du foncteur adjoint à gauche, et
enfin la commutativité du diagramme $({\ast}{\ast}{\ast}{\ast})$ se
déduit immédiatement du fait que $u_{!}$ commute aux limites
projectives finies.
\end{enonce*}

\begin{enonce*}[remark]{Notation 5.9}\etiquette[5.9]{I.not5.9}
Par abus de notation, les foncteurs ${u^\ast}^\gamma$ et
$u_{\ast\gamma}$ seront souvent notés $u^\ast$ et $u_\ast$, ce qui
ne risque pas d'apporter des confusions en vertu de la commutativité
des diagrammes $(\ast)$ et $({\ast}{\ast})$. En revanche, lorsque
$u_{!}$ ne commute pas aux limites projectives finies, le diagramme
$({\ast}{\ast}{\ast}{\ast})$ n'est pas commutatif à isomorphisme
près, et les notations $u_{!}$ et $u_{!\gamma}$ devront être
employées pour éviter des confusions.
\end{enonce*}

\setcounter{subsection}{9}
\subsection{}\etiquette[5.10]{I.sec5.10}
Soit $C$ une petite catégorie. Pour tout objet $X$ de
$\widehat{C}$, on désigne par $C/X$ la catégorie des
flèches de but $X$ et de source un objet de $C$. Le foncteur source
définit un foncteur $j_X:C/X\rightarrow C$. Soit $m:Y\rightarrow X$
un morphisme de $\widehat{C}$. La composition des
morphismes définit un foncteur $j_m:C/Y\rightarrow C/X$. Le
diagramme:
\[
\xymatrix@=1.5cm{
C/Y \ar[rr]^{j_m} \ar[rd]_{j_{Y}}&   & C/X\ar[ld]^{j_{X}}\\
    & C &
}
\]
est commutatif. Il résulte de \hyperref[I.prop5.1]{5.1} que pour tout objet $F$ de
$(C/X)^{\widehat{\phantom{,}}}$ et tout objet $Y$ de $C$, on a:
$$
\displaylines{\rlap{(5.10.1)}\hfill j_{X!}F(Y)=\coprod_{u\in
    \Hom_{\widehat{C}}(Y,X)} F(u){\quoi}.\hfill}
$$
La formule (5.10.1) permet de définir $j_{X!}$ lorsque $C$ est une
$𝒰$-catégorie, et on vérifie que le foncteur $j_{X!}$ ainsi
défini est toujours adjoint à droite au\pageoriginale foncteur
$j^\ast_X:\widehat{C}\rightarrow
(C/X)^{\widehat{\phantom{,}}}$ (\hyperref[I.sec5.0]{5.0}).

\begin{enonce*}{Proposition 5.11}\etiquette[5.11]{I.prop5.11}
Soit $C$ une $𝒰$-catégorie, $X$ un préfaisceau sur $C$.
\begin{enumerate}
\renewcommand{\labelenumi}{\theenumi)}
\item Le foncteur
$$
j_{X!}:(C/X)^{\widehat{\phantom{,}}}\longrightarrow
\widehat{C}
$$
se factorise par la catégorie $\widehat{C}/X$:
$$
(C/X)^{\widehat{\phantom{,}}}\xrightarrow{e_X}
\widehat{C}/X\longrightarrow
\widehat{C}{\quoi}.
$$
Le foncteur $e_X$ est une équivalence de catégories.

\item Le foncteur $e_{X}\circ j^\ast_X:\widehat{C}
\to  \widehat{C}/X$ est canoniquement isomorphe au
foncteur
  $H\mapsto (H\times X \xrightarrow{\pr_2} X)$.
\end{enumerate}
\end{enonce*}

\begin{enonce*}[remark]{Preuve}
\begin{enumerate}
\renewcommand{\labelenumi}{\theenumi)}
\item Soit $f$ l'objet final de $(C/X)^{\widehat{\phantom{,}}}$. On a
  un isomorphisme canonique $f\xrightarrow{\sim} \varinjlim_{Y\in
  \ob C/X}$ $Y$ et par suite $j_{X!}(f)=\varinjlim_{Y\in
  \ob C/X} j_X(Y)$ (5.4). Or $j_{X!}(f)\simeq X$; d'où la
  factorisation. Pour montrer que $e_X$ est une équivalence nous
  nous contenterons d'exhiber un foncteur quasi-inverse: A tout objet
  $H\rightarrow X$ de $\widehat{C}/X$ on associe le
  préfaisceau sur $C/X$:
$$
(Y\rightarrow X)\longmapsto
  \Hom_{\widehat{C}/X}((Y\rightarrow X), (H\rightarrow
  X)){\quoi}.
$$

\item Le foncteur $e_X\circ j^\ast_X$ est adjoint à droite au
  foncteur d'oubli et par suite le foncteur $e_X\circ j^\ast$ est
  canoniquement isomorphe au foncteur $H\mapsto (H\times
  X\xrightarrow{\pr_2}X)$.
\end{enumerate}
\end{enonce*}

\setcounter{subsection}{11}
\subsection{}\etiquette[5.12]{I.sec5.12}
Soit $m:Y\rightarrow X$ un morphisme de $\widehat{C}$.
D'après (5.11), le morphisme $m$ est canoniquement isomorphe à
l'image par $e_X$ d'un objet de $(C/X)^{\widehat{\phantom{,}}}$ que
nous noterons $[m]$. Le foncteur $e_X$ définit, par restriction aux
sous-catégories, une équivalence
$$
(C/X)/[m]\xrightarrow{e_m}C/Y{\quoi}.
$$
Le\pageoriginale diagramme
\[
\xymatrix{
  (C/X)/[m]\ar[rr]^{e_m} \ar[rd]_{j_{[m]}} &       & C/Y \ar[ld]^{j_m}\\
          & C/X   &
}
\]
est commutatif à isomorphisme canonique près.

\subsection{}\etiquette[5.13]{I.sec5.13}

Signalons un résultat qui nous sera utile dans Exp. VI. Soit
$u:C\rightarrow C'$ un foncteur entre petites catégories. Pour tout
objet $H$ de $\widehat{C}$ désignons par
$u/H:C/H\rightarrow C'/u_{!}H$ le foncteur qui associe à tout
morphisme $m:X\rightarrow H$ le morphisme
$u_{!}X\xrightarrow{u_{!}m} u_{!}H$ (on sait (\hyperref[I.prop5.4]{5.4}) qu'on
peut toujours poser $u_{!}X=uX$). Le diagramme ci-après est
commutatif:
\[\vcenter{
\xymatrix@C=3cm@R=1.75cm{
C/H \ar[r]^{u/H}\ar[d]_{J_H} & C'/u_{!}H \ar[d]_{j_{u_{!}H}} \\
C\ar[r]^{u}   & C'\quad. }}\leqno{(5.13.1)}
\]
On a donc un diagramme commutatif à isomorphisme près:
%\[\vcenter{
%\xymatrix@C=2.5cm@R=1.5cm{ (C/H){\widehat{\phantom{a}}}
%\ar[r]^{(u/H)_{!}} \ar[d]_{(j_H)_!}&
%  (C'/u^·H){\widehat{\phantom{a}}} \ar[d]_{(j_{u!H})_{!}}\\
%  C{\widehat{\phantom{\;a}}} \ar[r]^{u_!}& C'{\widehat{\phantom{\;a}}}\quad,
%}}\leqno{(5.13.2)}
%\]
et comme $(u/H)_{!}$ transforme l'objet final de
$(C/H)^{\widehat{\phantom{,}}}$ en l'objet final de
$(C'/u_{!}H)^{\widehat{\phantom{,}}}$, le diagramme:
\[\vcenter{
\xymatrix@C=3.5cm@R=1.5cm{ (C/H)\widehat{\phantom{a}}
\ar[r]^{(u/H)_!} \ar[d]_{e_H}&
  (C'/u_!H)\widehat{\phantom{a}} \ar[d]_{e_{u_{!}H}}\\
C\widehat{\phantom{a}}/H  \ar[r]^{u_{!}/H} &
C'\widehat{\phantom{a}}/u_{!}H }}\leqno{(5.13.3)}
\]
est\pageoriginale commutatif à isomorphisme près (6.1).

\begin{enonce*}{Proposition 5.14}\etiquette[5.14]{I.prop5.14}
Soit $u:C\rightarrow C'$ un foncteur entre petites catégories. On
suppose que $u$ possède la propriété suivante:

{\rm(PPF)} Pour tout objet $X$ de $C$ le foncteur
$u/X:C/X\rightarrow C'/uX$ est pleinement fidèle.

Alors:
\begin{enumerate}
\renewcommand{\labelenumi}{\theenumi)}
\item Soit $f$ l'objet final de $\widehat{C}$. Le
  foncteur $u$ se factorise en
$$
C=C/f\xrightarrow{u/f}C'/u_{!}f\xrightarrow{j_{u_{!}f}} C'{\quoi}.
$$
Le foncteur $u/f$ est pleinement fidèle.

\item Le foncteur $u_{!}:\widehat{C}\rightarrow
  \widehat{C'}$ se factorise en
$$
\widehat{C}=(C/f)^{\widehat{\phantom{,}}}\xrightarrow{(u/f)_{!}}
(C'/u_{!}f)^{\widehat{\phantom{,}}}\xrightarrow{e_{u_{!}f}}\widehat{C'}/u_{!}f\longrightarrow
\widehat{C'}{\quoi},
$$
où le foncteur $(u/f)_{!}$ est pleinement fidèle, le foncteur
$e_{u_{!}f}$ une équivalence, et le foncteur
$\widehat{C'}/u_{!}f\rightarrow
\widehat{C'}$ le foncteur d'oubli.

\item En particulier le foncteur $u_{!}$ est fidèle et par suite le
  morphisme d'adjonction $\id\xrightarrow{\Phi}u^\ast u_{!}$ est un
  monomorphisme. De plus, pour tout morphisme $\alpha:H\rightarrow K$
  de $\widehat{C}$, le diagramme:
\[
\xymatrix@C=3.5cm@R=1.5cm@M=.2cm{
H \ar@{^{(}->}[r]^{\Phi(H)}\ar[d]_{\alpha} & u^{*}u_{!}H \ar[d]^{u.u^{\bigdot}(\alpha)}\\
K \ar@{^{(}->}[r]^{\Phi(K)} &  u^{*}u_!K }
\]
est cartésien.
\end{enumerate}
\end{enonce*}

\begin{enonce*}[remark]{Preuve}
\begin{enumerate}
\renewcommand{\labelenumi}{\theenumi)}
\item La\pageoriginale factorisation provient du diagramme
  (5.13.1). Le foncteur $u$ est fidèle. Donc $u/f$ est
  fidèle. Montrons qu'il est pleinement fidèle. Soient $X$ et $Y$
  deux objets de $C$, $\can_X:uX\rightarrow u_{!}f$
  (\resp  $\can_Y:uY\rightarrow u_{!}f$) les morphismes canoniques et
\[
\xymatrix@=1.25cm{
uX \ar[rr]^{m} \ar[rd]_{\can_X}&               & uY \ar[ld]^{\can_Y}\\
   & u^{\bigdot}f &
}
\]
un morphisme de $C'/u_{!}f$. On a $u_{!}f=\varinjlim_{Z\in
  \ob C} uZ$ et par suite (3.1):
$$
\Hom_{\widehat{C'}}(uX,u_{!}f)=\varinjlim_{Z\in
  \ob C}\Hom_{C'}(uX, uZ){\quoi}.
$$
Par définition de la limite inductive, dire que $\can_Y m=\can_X$
équivaut à dire qu'il existe
\begin{enumerate}
\renewcommand{\theenumii}{\alph{enumii}}
\renewcommand{\labelenumii}{\theenumii)}
\item une suite finie d'objets de $C$, $X_i$, $i\in [0,n]$,
  $X_\circ=X$, $X_n=Y$,

\item pour tout $i$, un morphisme $m_i:uX\rightarrow
  uX_i(m_\circ=\id_{X}, m_n=m)$,

\item pour tout couple $(i,i+1)$ un morphisme $f_i:X_i\rightarrow
  X_{i+1}$ ou bien $f_i:X_{i+1}\rightarrow X_i$,
\end{enumerate}
tels que les diagrammes
\begin{displaymath}
\xymatrix@R=1.25cm@C=1.25cm{
   &  uX \ar[rd]^{\;\;m_{i+1} \qquad\quad\text{ ou bien}} \ar[ld]_{m_i}&\\
  uX_i\ar[rr]^{u(f_i)} &          &  uX_{i+1}
}\qquad \xymatrix@R=1.25cm@C=1.25cm{
      &  uX \ar[rd]^{\;\;m_{i+1}} \ar[ld]_{m_i}&\\
 uX_i &                                        &
   uX_{i+1}\ar[ll]_{u(f_i)}\quad ,
}
\end{displaymath}
soient commutatifs.

On démontre alors immédiatement, par récurrence sur $i$ et en
utilisant la propriété (PPF), que $m_i$ est de la forme $u(p_i)$. En
particulier $m=u(p)$ et par suite $u/f$ est pleinement fidèle.

\item La\pageoriginale factorisation est immédiate. Le foncteur
  $(u/f)_{!}$ est pleinement fidèle en vertu de \hyperref[I.prop5.6]{5.6}. Les autres
  assertions résultent de \hyperref[I.prop5.11]{5.11}.

\item Le foncteur $u_{!}$ est composé du foncteur d'oubli qui est
  fidèle, et de foncteurs pleinement fidèles. Il est par suite
  fidèle. Il en résulte, d'après les propriétés
  générales des foncteurs adjoints, que le morphisme d'adjonction
  $\id\rightarrow u^\ast u_{!}$ est un monomorphisme. D'après 2) le
  foncteur $u_{!}$ apparaît comme le composé d'un foncteur
  pleinement fidèle $v:\widehat{C}\to
  \widehat{C'}/u_{!}f$ et du foncteur d'oubli. Le
  foncteur $u^\ast$, adjoint à droite de $u_{!}$, est donc le
  composé du foncteur « produit par $u_{!}f$ », adjoint à
  droite du foncteur d'oubli, et d'un foncteur $w$ adjoint à droite
  de $v$. De plus, $v$ étant pleinement fidèle, le morphisme
  d'adjonction $\id\rightarrow wv$ est un isomorphisme. La dernière
  assertion en résulte aisément.
\end{enumerate}
\end{enonce*}

\section{Foncteurs fidèles et foncteurs
  conservatifs}\etiquette[6]{I.sec6}\pageoriginale

\begin{enonce*}{Définition 6.1}\etiquette[6.1]{I.def6.1}
Soient $E$ une catégorie, $(\varphi_i:E\rightarrow F_i)_{i\in I}$
une famille de foncteurs
$$
\varphi_i:E\longrightarrow F_i{\quoi}.
$$
On dit que la famille de foncteurs $(\varphi_i)$ est fidèle si pour
tout couple d'objets $X$, $Y$, de $E$, et tout couple de flèches
$u$, $v:X\rightrightarrows Y$, la relation
$\varphi_i(u)=\varphi_i(v)$ pour tout $i\in I$ implique $u=v$ (en
d'autres termes, si l'application $\Hom(X,Y)\rightarrow
\Pi_i\Hom(\varphi_i(X),\varphi_i(Y))$ définie par $(\varphi_i)$ est
injective). On dit que la famille de foncteurs $(\varphi_i)$ est
conservative si toute flèche $u$ de $E$, telle que $\varphi_i(u)$
soit un isomorphisme pour tout $i\in I$, est un isomorphisme. On dit
que $(\varphi_i)$ est conservative pour les monomorphismes (\resp
pour les épimorphismes, \resp \ldots) si la condition précédente est
vérifiée chaque fois que $u$ est un monomorphisme (\resp  un
épimorphisme, \resp \ldots).
\end{enonce*}

\setcounter{subsection}{1} \setcounter{subsubsection}{0}
\subsubsection{}\etiquette[6.1.1]{I.sec6.1.1}
Si on introduit le foncteur unique
$$
\varphi:E\longrightarrow F=\prod_{i\in I} F_i
$$
défini par la famille de foncteurs $(\varphi_i)$, il est clair que
celle-ci est fidèle (\resp  conservative, \resp  conservative pour
les monomorphismes, \resp\ldots) si et seulement si le foncteur
$\varphi$ est fidèle (\resp  conservatif, \resp\ldots) (par quoi on
entend que la famille réduite au seul objet $\varphi$ est fidèle,
\resp  conservative, \resp\ldots). On pourrait donc sans
inconvénient majeur nous borner par la suite au cas d'une famille
réduite à un seul foncteur. Pour la commodité des futures
références, nous donnerons néanmoins les énoncés suivants pour les
familles.

Les\pageoriginale notions de 6.1 sont surtout utiles lorsque les
$\varphi_i$ satisfont à des propriétés d'exactitude convenables, et
dans ce cas ont une tendance à coïncider:

\begin{enonce*}{Proposition 6.2}\etiquette[6.2]{I.prop6.2}
Les notations sont celles de \hyperref[I.prop5.1]{5.1}.
\begin{enumerate}
\renewcommand{\theenumi}{\roman{enumi}}
\renewcommand{\labelenumi}{(\theenumi)}
\item Si les noyaux de doubles flèches, ou les conoyaux de doubles
  flèches, sont représentables dans $E$, et si les $\varphi_i$ $y$
  commutent, alors on a l'implication
$$
(\varphi_i)\mbox{ conservative } \Longrightarrow (\varphi_i)\mbox{
    fidèle}{\quoi}.
$$

\item Supposons que les produits fibrés (\resp  les sommes
  amalgamées) soient représentables dans $E$, et que les
  $\varphi_i$ $y$ commutent. Supposons $(\varphi_i)$ fidèle ou
  conservative; alors pour toute flèche $u$ de $E$, $u$ est un
  monomorphisme (\resp  un épimorphisme) si et seulement si pour tout
  $i\in I$, il en est ainsi pour $\varphi_i(u)$.

\item Supposons que dans $E$ les produits fibrés et les sommes
  amalgamées sont représentables et que les $\varphi_i$ $y$
  commutent, et que toute flèche dans $E$ qui est un bimorphisme
  (\ie un monomorphisme et un épimorphisme) soit un
  isomorphisme. Alors on a l'implication
$$
(\varphi_i)\mbox{ fidèle } \Longrightarrow (\varphi_i)\mbox{
    conservative}{\quoi}.
$$

\item Supposons que dans $E$ les produits fibrés (\resp  les sommes
  amalgamées) soient représentables, et que les $\varphi_i$ $y$
  commutent. Alors, si $(\varphi_i)$ est conservative pour les
  monomorphismes (\resp  pour les épimorphismes) alors $(\varphi_i)$
  est même conservative.

\item Soit $D$ un type de diagramme, $F:d\mapsto F(d)$ un diagramme de
  type $D$ dans $E$, $X$ un objet de $E$ et $u=(u_d)_{d\in D}$ une
  famille de flèches
$$
X\longrightarrow F(d)\quad (\mbox{\resp  } F(d)\longrightarrow
X){\quoi}.
$$
Supposons\pageoriginale que $(\varphi_i)$ soit conservative, que les
limites projectives (\resp  inductives) de type $D$ soient
représentables dans $E$, et que les $\varphi_i$ $y$ commutent.
Alors, pour que $u$ fasse de $X$ une limite projective (\resp
inductive) de $F$ dans $E$, il faut et il suffit que pour tout $i\in
I$, $\varphi_i(u)$ fasse de $\varphi_i(X)$ une limite projective
(\resp  inductive) de $\varphi_i(D)$ dans $E_i$.
\end{enumerate}
\end{enonce*}

\noindent \emph{Démonstration.}
\begin{enumerate}
\renewcommand{\theenumi}{\roman{enumi}}
\renewcommand{\labelenumi}{(\theenumi)}
\item Pour l'énoncé non respé, il suffit, pour une double
  flèche donnée $u$, $v:X\rightrightarrows Y$, d'exprimer
  l'égalité $u=v$ par la condition que l'inclusion
  $\Ker(u,v)\rightarrow X$ est un isomorphisme. Ici et par la suite,
  on se dispense de répéter l'argument dual pour l'énoncé
  dual.

\item Si $(\varphi_i)$ est fidèle, on exprime la condition que
  $u:X\rightarrow Y$ soit un monomorphisme par l'égalité
  $\pr_1=\pr_2$ pour le produit fibré $\fprod{X}{X}{Y}$. Si $(\varphi_i)$
  est conservatif, on l'exprime par la condition que le morphisme
  diagonal $\delta:X\rightarrow \fprod{X}{X}{Y}$ soit un isomorphisme.

\setcounter{enumi}{3}
\item Comme dans ce dernier argument, le morphisme $\delta$ est un
  monomorphisme, on voit qu'il suffisait en fait de supposer
  $(\varphi_i)$ conservative pour les monomorphismes. Mais ceci
  implique alors que $(\varphi_i)$ est conservative tout court. En
  effet, si $u\in %%
F\ell\,E$ est telle que les $\varphi_i(u)$ soient des
  isomorphismes, on en conclut que ce sont des monomorphismes
  d'après ce qui précède, donc des isomorphismes d'après
  l'hypothèse sur $(\varphi_i)$.

\setcounter{enumi}{2}
\item Est une conséquence triviale de (ii).

\setcounter{enumi}{4}
\item Est une conséquence triviale des définitions.
\end{enumerate}

Notons la conséquence suivante de (i) (ii) (iv):

\begin{enonce*}{Corollaire 6.3}\etiquette[6.3]{I.cor6.3}
Supposons\pageoriginale que dans $E$ les produits fibrés et les
sommes amalgamées soient représentables et que les $\varphi_i$ $y$
commutent, et que les noyaux de double flèches ou les conoyaux de
double flèches soient représentables et que les $\varphi_i$ $y$
commutent. (Il suffit par exemple que les limites projectives finies
et les limites inductives finies soient représentables dans $E$, et
que les $\varphi_i$ soient des foncteurs exacts.) Alors les
conditions suivantes sont équivalentes:
\begin{itemize}
\item[a)] $(\varphi_i)$ est fidèle.

\item[b)] $(\varphi_i)$ est conservative.

\item[c)] $(\varphi_i)$ est conservative pour les monomorphismes.

\item[${\rm c')}$] $(\varphi_i)$ est conservative pour les épimorphismes.
\end{itemize}
\end{enonce*}
Signalons aussi pour mémoire:

\begin{enonce*}{Proposition 6.4}\etiquette[6.4]{I.prop6.4}
Soient $\varphi:E\rightarrow F$ un foncteur admettant un adjoint à
droite $\psi$ (donc $\Hom(\varphi(X),Y)\simeq \Hom(X,\psi(Y))$).
Pour que $\varphi$ (\resp  $\psi$) soit fidèle, il faut et il suffit
que pour tout élément $X$ de $E$ (\resp  tout élément $Y$ de $F$),
le morphisme d'adjonction $X\rightarrow \psi\varphi(X)$ soit un
monomorphisme. Pour que $\varphi$ (\resp  $\psi$) soit pleinement
fidèle, il faut et il suffit que le morphisme d'adjonction précédent
soit un isomorphisme.
\end{enonce*}
En effet, si $X'$, $X$ sont deux objets de $E$, l'application
\begin{equation*}
\Hom(X',X)\longrightarrow
  \Hom(\varphi(X'),\varphi(X))\tag{$\ast$}
\end{equation*}
s'identifie à l'application déduite de
\begin{equation*}
X\longrightarrow \psi
  \varphi (X)\tag{${\ast}{\ast}$}
\end{equation*}
par application du foncteur $\Hom(X',-)$. Donc pour que $(\ast)$
soit un monomorphisme\pageoriginale (\resp  un isomorphisme) pour
tout $X'$, $X$ étant fixé, il faut et il suffit que le morphisme
d'adjonction $({\ast}{\ast})$ soit un monomorphisme (\resp  un
isomorphisme).

\begin{enonce*}{Proposition 6.5}\etiquette[6.5]{I.prop6.5}
Soit $p:E'\rightarrow E$ un foncteur. Les conditions suivantes sont
%%
équivalentes:
\begin{enumerate}
\renewcommand{\theenumi}{\roman{enumi}}
\renewcommand{\labelenumi}{(\theenumi)}
\item $p$ est fidèle, conservatif et fibrant {\rm(SGA 1 VI 6.1).}

\item $p$ est un foncteur fibrant à fibres des catégories
  discrètes.

\item Pour tout $X'\in \ob E'$, le foncteur $E'_{/X'}\rightarrow
  E_{/p(X)}$ induit par $p$ est une équivalence de catégories,
  surjective sur les objets.

\item (Lorsque $E$ est une $𝒰$-catégorie). Il existe un
  élément $F\in \ob\widehat{C}$ et une équivalence de
  catégories sur $E$ {\rm(SGA 1 VI 4.3)} $E'\xrightarrow{\approx}
  E_{/F}$ (où $E_{/F}$ est la sous-catégorie pleine de
  $\widehat{E}_{/F}$ formée des flèches $X\rightarrow F$ dont la
  source est dans $E$).
\end{enumerate}
\end{enonce*}

\setcounter{subsection}{5} \setcounter{subsubsection}{0}
\subsubsection{}\etiquette[6.5.1]{I.sec6.5.1}
Rappelons qu'une catégorie $C$ dite est \emph{discrète} si c'est un
\emph{groupoïde} (\ie toute flèche y est inversible) et si elle
est \emph{rigide} (\ie le groupe des automorphismes de tout objet
est réduit au groupe unité); il revient au même de dire que la
catégorie est équivalente à la catégorie $C'$ définie par un
ensemble $I$ (avec $\text{ob}C'=I$, et comme seules flèches les
flèches identiques). Quand on suppose déjà que $C$ est un
groupoïde, alors dire que $C$ est discrète revient à dire que pour
deux objets $X$, $Y$ de $C$, il existe au plus une flèche de $X$
dans $Y$, \ie que $C$ est isomorphe à la catégorie définie par un
ensemble préordonné.

L'équivalence\pageoriginale des conditions (i) et (ii) de
\hyperref[I.prop6.5]{6.5} est une conséquence immédiate des rappels précédents
et du

\begin{enonce*}{Lemme 6.5.2}\etiquette[6.5.2]{I.lem6.5.2}
Soit $p:E'\rightarrow E$ un foncteur fibrant. Alors:
\begin{enumerate}
\renewcommand{\theenumi}{\roman{enumi}}
\renewcommand{\labelenumi}{(\theenumi)}
\item Pour que $p$ soit conservatif, il faut et il suffit que ses
  catégories fibres soient des groupoïdes.

\item Pour que $p$ soit fidèle, il faut et il suffit que ses
  catégories fibres soient des catégories ordonnées.
\end{enumerate}
\end{enonce*}

\noindent \emph{Démonstration de} 6.5.2.
\begin{enumerate}
\renewcommand{\theenumi}{\roman{enumi}}
\renewcommand{\labelenumi}{(\theenumi)}
\item Supposons $p$ conservatif. Pour toute flèche $u$ d'une fibre
  $E'_X$, $p(u)=\id_X$ est un isomorphisme, donc $u$ est un
  isomorphisme dans $E'$, donc aussi dans $E'_X$ (car un inverse de
  $u$ dans $E'$ sera évidemment un inverse dans $E'_X$). Donc $E'_X$
  est un groupoïde. Inversement, supposons les $E'_X$ des
  groupoïdes, et soit $u'$ une flèche de $E'$ telle que $p(u')$
  soit un isomorphisme, prouvons que $u$ est un isomorphisme. Pour
  ceci on note que, $p$ étant fibrant, on peut factoriser
  $u':X'\rightarrow Y'$ en un composé $X'\rightarrow
  u^\ast(Y')\rightarrow Y'$, où la première flèche est un
  $X$-morphisme (N.B. $X=p(X')$, $u=p(u')$) et la deuxième est un
  morphisme cartésien au-dessus de $u$. La première flèche est
  un isomorphisme puisque $E'_X$ est un groupoïde, et la deuxième
  l'est, car un morphisme \emph{cartésien} %%
d'une catégorie fibrée
  est évidemment un isomorphisme dès que sa projection l'est.

\item Supposons $p$ fidèle, et soient $X'$, $Y'$ deux objets d'une
  catégorie fibre $E_X$. Alors deux flèches de $X'$ dans $Y'$ sont
  au-dessus de la même flèche $\id_X$ de $E$, donc sont identiques,
  donc $E_X$ est ordonnée. Inversement,\pageoriginale supposons les
  catégories fibres ordonnées, et prouvons que $p$ est
  fidèle. Soient donc $u'$, $v':X'\rightrightarrows Y'$ des
  flèches de $E'$ au-dessus d'une même flèche $u:X\rightarrow Y$
  de $E$. Elles se factorisent alors en $X'\rightrightarrows
  u^\ast(Y)\rightarrow Y'$, où les deux flèches
  $X'\rightrightarrows u^\ast(Y')$ sont des flèches de $E'_X$ de
  même source et même but; celles-ci sont donc égales, donc
  $u'=v'$, \cqfd
\end{enumerate}

Revenons à la démonstration de \hyperref[I.prop6.5]{6.5}. On a prouvé (i)
$\Leftrightarrow$ (ii). D'autre part (ii) %%
$\Leftrightarrow$ (iv) est assez claire: en effet, d'une part la
catégorie $E_{/F}$ est fibrée sur $E$ à catégories fibres les
catégories discrètes définies par les ensembles $F(X)$, comme il
résulte aussitôt des définitions; d'autre part, si $p$ est comme
dans (ii), alors en vertu du sorite SGA 1 VI 8 la catégorie fibrée
$E'$ sur $E$ est $E$-équivalente à la catégorie scindée sur $E$
définie par le foncteur $E\rightarrow (\Cat)$ définie par le
foncteur $F:E\rightarrow (\Ens)$, associant à tout $X\in \ob E$
l'ensemble des classes d'isomorphie d'objets de $E'_X$. Or cette
catégorie scindée est $E$-isomorphe à la catégorie $E_{/F}$. Comme
(iv) $\Rightarrow$ (iii) est claire, il reste à prouver (iii)
$\Rightarrow$ (i). Or il est clair que pour que $p$ soit fidèle
(\resp  conservatif) il faut et il suffit que les foncteurs induits
$E'_{/X'}\rightarrow E_{/p(X)}$ le soient, \textit{a fortiori} il
suffit que ceux-ci soient pleinement fidèles; donc il reste à
prouver seulement que (iii) implique que $p$ est fibrant. Mais on
voit encore qu'un foncteur $p$ est fibrant si et seulement si les
foncteurs induits $E'_{/X'}\rightarrow E_{/p(X')}$ le sont. Il en
est en particulier ainsi si ce sont des équivalences de catégories
surjectives sur les objets.

\section{Sous-catégories génératrices et
  cogénératrices}\etiquette[7]{I.sec7}\pageoriginale

\begin{enonce*}{Définition 7.1}\etiquette[7.1]{I.def7.1}
Soient $E$ une catégorie, $C$ une sous-catégorie pleine de $E$. On
dit que $C$ est une sous-catégorie de $E$ \emph{génératrice par
épimorphismes stricts (\resp  par épimorphismes)} si pour tout objet
$X$ de $E$, la famille des flèches de $E$ de but $X$, de source
$X'\in \ob C$, est épimorphique (\resp  épimorphique stricte)
(\hyperref[I.condef1.3]{1.3}). On dit que $C$ est une sous-catégorie de $E$
\emph{génératrice (\resp  génératrice pour les monomorphismes, \resp
génératrice pour les monomorphismes stricts)} si pour toute flèche
$u:Y\rightarrow X$ (\resp  tout monomorphisme de $E$, \resp  tout
monomorphisme strict de $E$ $(10.5)$), telle que pour tout $X'\in
\ob C$, l'application correspondante $\Hom(X',Y)\rightarrow
\Hom(X',X)$ soit bijective, $u$ est un isomorphisme. Enfin, on dit
qu'une famille $(X_i)$ d'objets de $E$ est génératrice par
épimorphismes stricts (\resp \ldots) si la sous-catégorie pleine $C$
de $E$ engendrée par cette famille est génératrice par épimorphismes
stricts (\resp \ldots).
\end{enonce*}

\setcounter{subsection}{1} \setcounter{subsubsection}{0}
\subsubsection{}\etiquette[7.1.1]{I.sec7.1.1}
Notons qu'en termes de la famille $(h_{X'})_{X'\in \ob C}$ des
foncteurs
$$
h_{X'}:E\longrightarrow (\Ens)\; (X'\in \ob C)
$$
représentés par les $X'\in \ob C$, on peut exprimer la condition que
$C$ soit génératrice (\resp  génératrice pour les
monomorphismes,
\resp  génératrice pour les monomorphismes stricts) par celle que la
famille $(h_{X'})$ soit \emph{conservative} (\resp  conservative
pour les monomorphismes, \resp  conservative pour les monomorphismes
stricts) (\hyperref[I.def6.1]{6.1}). Il résulte également immédiatement des
définitions que $C$ est génératrice par épimorphismes
si\pageoriginale et seulement si la famille $(h_{X'})_{X'\in \ob C}$
est \emph{fidèle} (\hyperref[I.def6.1]{6.1}). On donnera aussi ci-dessous
(\hyperref[I.prop7.2]{7.2} (i)) une interprétation analogue pour la condition
sur $C$ d'être génératrice par épimorphismes stricts.

\subsubsection{}\etiquette[7.1.2]{I.sec7.1.2}
Comme pour les notions introduites dans \hyperref[I.def6.1]{6.1}, les notions
de \hyperref[I.def7.1]{7.1} sont surtout utiles lorsque $E$ possède des
propriétés d'exactitude convenables, auquel cas les diverses notions
introduites ont une nette tendance à être toutes équivalentes
(\hyperref[I.cor7.3]{7.3}). C'est pourquoi la question de savoir laquelle de
ces notions \hyperref[I.def7.1]{7.1} doit être considérée comme la plus
importante ne se pose guère; dans les cas les plus importants, ces
notions coïncident et le terme « sous-catégorie
génératrice »
peut donc être interprété indifféremment comme se
rapportant à
n'importe laquelle des propriétés envisagées dans \hyperref[I.def7.1]{7.1}
(par exemple la première, qui est la plus forte de toute comme nous
allons voir (\hyperref[I.prop7.2]{7.2} (ii))).

\subsubsection{}\etiquette[7.1.3]{I.sec7.1.3}
Supposons que $E$ soit une $𝒰$-catégorie, et considérons le
foncteur canonique
\begin{equation*}
\varphi:E\longrightarrow
  \widehat{C}=\Hhom(C, 𝒰\text{-}\Ens)\tag{7.1.3.1}
\end{equation*}\etiquette[7.1.3.1]{I.sec7.1.3.1}
composé des foncteurs $E\rightarrow \widehat{E}\rightarrow
\widehat{C}$, où le premier foncteur est
le foncteur canonique (\hyperref[I.sec1.3.3]{1.3.3}), et
le deuxième
le foncteur restriction à $C$. Notons qu'il est évident qu'il
revient au même de dire que le foncteur précédent $\varphi$ est
conservatif (\resp  fidèle), ou de dire que la famille des foncteurs
$h_{X'}:X\rightarrow \Hom(X',X)=\varphi(X)(X')$, pour $X'\in\ob C$
variable, est une famille conservative (\resp  fidèle), c'est-à-dire
aussi (\hyperref[I.sec7.1.2]{7.1.2}) que $C$ est génératrice (\resp
génératrice par
épimorphisme). De même $\varphi$ est\pageoriginale conservative pour
les monomorphismes (\resp  pour les monomorphismes stricts) si et
seulement la famille des $h_{X'}(X'\in \ob C)$ est conservative pour
les monomorphismes (\resp  pour les monomorphismes stricts), \ie si
et seulement si la sous-catégorie $C$ de $E$ est génératrice pour
les monomorphismes (\resp  pour les monomorphismes stricts).

\begin{enonce*}{Proposition 7.2}\etiquette[7.2]{I.prop7.2}
Soient $E$ une $𝒰$-catégorie, $C$ une sous-catégorie pleine.
\begin{enumerate}
\renewcommand{\theenumi}{\roman{enumi}}
\renewcommand{\labelenumi}{(\theenumi)}
\item Les conditions suivantes sont équivalentes:
\begin{enumerate}
\renewcommand{\theenumii}{\alph{enumii}}
\renewcommand{\labelenumii}{\theenumii)}
\item $C$ est une sous-catégorie génératrice par épimorphismes
  stricts.

\item Pour tout $X\in \ob E$, désignant par $C_{/X}$ la
  sous-catégorie pleine de $E_{/X}$ formée des flèches
  $X'\rightarrow X$ de source $X'\in \ob C$, la flèche naturelle du
  foncteur d'inclusion $j:C_{/X}\rightarrow E$ dans le foncteur
  constant sur $C_{/X}$ de valeur $X$ fait de $X$ une limite inductive
  de $j$:
$$
X\xleftarrow{\sim} \lim_{C_{/X}}X'{\quoi}.
$$

\item Le foncteur canonique $\varphi$ de (\hyperref[I.sec7.1.3.1]{7.1.3.1}) est pleinement
  fidèle.
\end{enumerate}

\item On a entre les notions de \hyperref[I.def7.1]{7.1} les implications
suivantes~:
\[
\xymatrix@C=-.7cm@R=.45cm{
  &  \txt{{\rm 1)} C {%%
génératrice par épim.}\\{ stricts}\\ \emph{($\varphi$
pleinement fidèle)}} \ar@{=>}[rd]&\\
 \txt{{\rm 2)} C %%
génératrice par\\ épimorphismes\\ \emph{($\varphi$
 fidèle)}} &    &  \txt{{\rm 3)} C %%
génératrice \phantom{aaa}\\
  \emph{($\varphi$ conservatif)}} \ar@{=>}[d]\\
    &      &  \txt{{\rm 4)} C %%
génératrice pour mono-\\ morphismes \\
  \emph{($\varphi$ conservatif pour mon.)}}\ar@{=>}[d]\\
    &      &  \txt{{\rm 5)} C {%%
génératrice pour mon.}\\ stricts\\
  \emph{($\varphi$ conservatif pour mon. stricts).}}\\
}
\]

\item On\pageoriginale a les implications conditionnelles suivantes:
\begin{enumerate}
\renewcommand{\theenumii}{\alph{enumii}}
\renewcommand{\labelenumii}{\theenumii)}
\item Si dans $E$ les familles épimorphiques de flèches sont
  épimorphiques strictes, on a $2)\Rightarrow 1)$. Si dans $E$ les
  monomorphismes sont stricts, on a $5)\Rightarrow 4)$.

\item Si dans $E$ les noyaux de couples de flèches (\resp  les
  produits fibrés) sont représentables, alors on a $3)\Rightarrow
  2)$ (\resp  $4)\Rightarrow 3)$).

\item Si dans $E$ toute famille de morphismes $X_i\rightarrow X$ de
  même but $X$ se factorise en une famille épimorphique stricte
  (\resp  épimorphique) $X_i\rightarrow Y$ suivie d'une monomorphisme
  (\resp  d'une monomorphisme strict) $Y\rightarrow X$, alors on a
  $4)\Rightarrow 1)$ (\resp  $5)\Rightarrow 2)$).
\end{enumerate}
\end{enumerate}
\end{enonce*}
Signalons tout de suite le

\begin{enonce*}{Corollaire 7.3}\etiquette[7.3]{I.cor7.3}
Toutes les notions envisagées dans $6.1$ (et reprises dans le
diagramme d'implications de {\rm(ii)} ci-dessus) sont équivalentes
dans chacun des deux cas suivants:
\begin{enumerate}
\renewcommand{\theenumi}{\roman{enumi}}
\renewcommand{\labelenumi}{(\theenumi)}
\item Dans $E$, les noyaux de doubles flèches et les produits
  fibrés sont représentables, les monomorphismes sont stricts et
  les familles épimorphiques de flèches sont épimorphiques
  strictes.

\item Dans $E$, toute famille $(X_i\rightarrow X)_{i\in I}$ de
  flèches de même but $X$ se factorise en une famille
  épimorphique $(X_i\rightarrow Y)$ suivie d'un monomorphisme
  $Y\rightarrow X$, tout monomorphisme de $E$ est strict, et toute
  famille épimorphique de flèches de $E$ est stricte.
\end{enumerate}
\end{enonce*}

En\pageoriginale effet, dans le cas (i) on a 4) $\Rightarrow$ 3) et
3) $\Rightarrow$ 2) grâce à b), et 5) $\Rightarrow$ 4) et 2)
$\Rightarrow$ 1) grâce à a). Dans le cas (ii) on a grâce à c), les
implications 4) $\Rightarrow$ 1) et 5) $\Rightarrow$ 2). On conclut
donc grâce au diagramme d'implications \hyperref[I.prop6.2]{6.2} (ii).

\medskip
\noindent \emph{Démonstration de} \hyperref[I.prop7.2]{7.2}.
\begin{enumerate}
\renewcommand{\theenumi}{\roman{enumi}}
\renewcommand{\labelenumi}{(\theenumi)}
\item L'implication b) $\Rightarrow$ a) résulte aussitôt des
  définitions. Prouvons a) $\Rightarrow$ b). Donc sous
  l'hypothèse a), il faut prouver que pour tout $X$, $Y\in \ob E$,
  tout système de flèches
$$
u_{X'}:X'\longrightarrow Y
$$
indexé par les $X'\in \ob C_{/X}$, telle que l'on ait
$u_{X'}f=u_{X''}$ pour toute flèche $f:X''\rightarrow X'$ dans
$C_{/X}$, se factorise par une flèche (nécessairement unique par
l'hypothèse a)) $X\rightarrow Y$. D'après l'hypothèse a), il suffit
de vérifier que pour tout objet $Z$ de $E_{/X}$ et tout couple de
morphismes $v':Z\rightarrow X'$, $v'':Z\rightarrow X''$ dans
$E_{/X}$, avec $X'$ et $X''$ dans $C_{/X}$, on a
$u_{X'}v'=u_{X''}v''$. Or, grâce à l'hypothèse a), la famille des
flèches $w:X'''\rightarrow Z$, avec $X'''\in \ob C$, est
épimorphique, et il suffit donc de vérifier que pour toute telle
$w$, on a $(u_{X'}v')w=(u_{X''}v'')w$, ce qui s'écrit aussi
$u_{X'}(v'w)=u_{X''}(v''w)$ et résulte aussitôt de l'hypothèse faite
sur la famille des $u$.

Prouvons maintenant l'équivalence des conditions b) et c). Pour ceci
notons que pour tout $X\in \ob E$, l'objet $\varphi(X)$ de
$\widehat{C}$ est limite inductive dans $\widehat{C}$ du foncteur
canonique $\widehat{C}_{/\varphi(X)}\rightarrow \widehat{C}$ (\hyperref[I.prop3.4]{3.4});
or $\widehat{C}_{/\varphi(X)}$ est canoniquement isomorphe à
$C_{/X}$, de sorte qu'on a dans $\widehat{C}$
$$
\varphi(X)=\varinjlim_{C_{/X}}X'{\quoi},
$$
où\pageoriginale on identifie l'objet $X'$ de $C$ avec le foncteur
$\in \widehat{C}$ qu'il représente. On a par suite, pour un deuxième
objet $Y$ de $E$, un isomorphisme canonique
\begin{align*}
\Hom(\varphi(X),\varphi(Y))\simeq\varprojlim_{C_{/X}}
\Hom(X',\varphi(Y))&\simeq
\varprojlim_{C_{/X}}\varphi(Y)(X')(1.4)\\
&\simeq \varprojlim_{C_{/C}} \Hom(X'%%
,Y){\quoi}.\tag{$\ast$}
\end{align*}
Ceci posé, l'application
$$
\Hom(X,Y)\longrightarrow \Hom(\varphi(X),\varphi(Y))
$$
définie par $\varphi$ n'est autre, via l'isomorphisme entre les
membres extrêmes de $(\ast)$, que l'application
$$
\Hom(X,Y)\longrightarrow \varprojlim_{C_{/X}} \Hom(X',Y)
$$
déduite du système inductif de flèches $X'\rightarrow X$ indexé par
$C_{/X}$ envisagé dans \hyperref[I.prop7.2]{7.2} (i) b). Donc la première
application est bijective pour tout $Y$ ($X$ étant fixé) si et
seulement si $X$ est une limite inductive du foncteur d'inclusion
$j:C_{/X}\rightarrow E$, ce qui prouve l'équivalence de b) et c).

\item Les implications 1) $\Rightarrow$ 2) et 3) $\Rightarrow$ 4)
  $\Rightarrow$ 5) sont triviales en vertu des
  définitions. L'implication 1) $\Rightarrow$ 3) s'obtient en
  interprétant la propriété 1) par la pleine fidélité de
  $\varphi$ grâce à (i), et en observant qu'un foncteur pleinement
  fidèle est conservatif. Or on a déjà observé
  (\hyperref[I.sec7.1.1]{7.1.1}) que 3)
  signifie que $\varphi$ est conservatif.

\item L'assertion\pageoriginale a) est une tautologie. L'assertion
  b) résulte de \hyperref[I.prop6.2]{6.2} (i) (\resp  \hyperref[I.prop6.2]{6.2} (iv)) appliqué au foncteur
  $\varphi$, compte tenu que ce dernier est exact à gauche. Prouvons
  enfin c). Considérons, pour un $X\in \ob E$, la famille de tous
  les morphismes $X_i\rightarrow X$, avec $X_i\in \ob C$; par
  hypothèse sur $E$ elle se factorise en une famille $X_i\rightarrow
  Y$ épimorphique effective (\resp  épimorphique) suivie d'un
  monomorphisme (\resp  d'un monomorphisme strict) $Y\rightarrow
  X$. Alors l'hypothèse 4) (\resp  5)) implique que $Y\rightarrow
  X$ est un isomorphisme, donc la famille envisagée est
  épimorphique stricte (\resp  épimorphique), \cqfd
\end{enumerate}

\begin{enonce*}{Proposition 7.4}\etiquette[7.4]{I.prop7.4}
Soient $E$ une catégorie, $C$ une sous-catégorie pleine de $E$, $X$
un objet de $E$. On suppose $C$ génératrice dans $E$ (\resp $C$
génératrice pour les monomorphismes, et que le produit fibré de deux
sous-objets de $X$ sur $X$ est représentable dans $E$). Alors un
sous-objet strict $(10.11)$ (\resp  un sous-objet) $X'$ de $X$ est
connu %%
quand on connaît, pour tout $T\in \ob C$, la partie de $\Hom(T,X)$
image de $\Hom(T,X')$. Par suite, le cardinal de l'ensemble des
sous-objets stricts (\resp  de l'ensemble des sous-objets) de $X$
est majoré par $\prod_{T\in \ob C} 2^{\card
  \Hom(T,X)}$, et si $X\in \ob C$, il est majoré par $2^{\card%%
F\ell\,C}$.
\end{enonce*}

Prouvons d'abord l'assertion respée. Soient $X'$, $X''$ deux
sous-objets de $X$ tels que pour tout $T\in \ob C$, les images de
$\Hom(T,X')$ et $\Hom(T,X'')$ dans $\Hom(T,X)$ soient égales. Elles
sont donc aussi égales à l'image de $\Hom(T,X''')$, où $X'''$ est le
produit fibré de $X'$ et $X''$ sur $X$. Comme $C$ est génératrice
pour les monomorphismes, il s'ensuit que les
monomorphismes\pageoriginale $X'''\rightarrow X'$ et
$X'''\rightarrow X''$ sont des isomorphismes, donc $X'$ et $X''$
sont égaux, étant séparément égaux au sous-objet $X'''$ de $X$.

Prouvons l'assertion non respée. Par définition de la notion de
sous-objet strict, il suffit de vérifier que la connaissance de la
partie $\Hom(T,X')$ de $\Hom(T,X)$ pour tout $T\in\ob C$ implique la
connaissance de celles des doubles flèches $\xymatrix{
X\ar@<2pt>[r]^{u,v}\ar@<-2pt>[r]& T }$ telles que $ui=vi$, où
$i:X'\rightarrow X$ est l'injection canonique. Or comme $C$ est
génératrice, la relation $ui=vi$ équivaut à la relation
$(ui)f=(vi)f$ pour tout $f\in \Hom(T,X')$, \ie à $ug=vg$ pour tout
$g\in \Hom(T,X)$ provenant de $\Hom(T,X')$ (\ie de la forme if, avec
$f\in \Hom(T,X')$), ce qui prouve notre assertion.

\begin{enonce*}{Corollaire 7.5}\etiquette[7.5]{I.cor7.5}
Soient $E$ une catégorie, $C$ une sous-catégorie pleine
génératrice,
$X$ un objet de $C$. Alors un quotient strict (\hyperref[I.sec10.8]{10.8}) $X'$ de $X$
est connu quand on connaît, pour tout $T\in \ob C$, la partie de
$\Hom(T,X)^2$ formée des couples $(u,v)$ tels que $qu=qv$, où
$q:X\rightarrow X'$ est le morphisme canonique. Donc le cardinal de
l'ensemble des quotients stricts de $X$ est majoré par $\prod_{T\in
\ob C} 2^{\card\Hom(T,X)^2}$.
\end{enonce*}

En effet, par %%
définition, un quotient strict $X'$ de $X$ est connu quand on
connaît, pour tout objet $Y$ de $E$, la partie de $\Hom(Y,X)\times
\Hom(Y,X)$ formée des couples $(u,v)$ tels que $qu=qv$, où
$q:X\rightarrow X'$ est le morphisme canonique. Or la relation
$qu=qv$ équivaut à la relation $(qu)f=(qv)f$ pour tout morphisme
$f:T\rightarrow Y$ de source $T$ dans $C$, puisque $C$ est
génératrice. Cette relation s'écrit encore $q(uf)=q(vf)$, ce qui
prouve la première assertion de \hyperref[I.cor7.5]{7.5}. La seconde en
résulte aussitôt.

\setcounter{subsection}{5} \setcounter{subsubsection}{0}
\subsubsection{}\etiquette[7.5.1]{I.sec7.5.1}
On\pageoriginale peut généraliser \hyperref[I.cor7.5]{7.5}, en
introduisant, pour une
famille d'objets $(X_i)_{i\in I}$ de $E$, la notion de
\emph{quotient
  strict dans $E$ de la famille}, par quoi on entend une famille
épimorphique stricte $(p_i:X_i\rightarrow X')_{i\in I}$ de
morphismes de $E$, --- étant entendu, comme pour les quotients
ordinaires, qu'on identifie deux telles familles
$(p_i:X_i\rightarrow X')$ et $(q_i:X_i\rightarrow X'')$ si on peut
trouver un isomorphisme $v:X'\rightarrow X''$ (nécessairement
unique) tel que l'on ait $fp_i=q_i$ pour tout $i\in I$. Avec cette
terminologie, la démonstration de \hyperref[I.cor7.5]{7.5} s'applique {\it ne
varietur} pour donner la

\begin{enonce*}[remark]{Variante 7.5.2}\etiquette[7.5.2]{I.var7.5.2}
Soient $E$, $C$ comme dans \hyperref[I.cor7.5]{7.5}, et $(X_i)_{i\in I}$ une famille
d'objets de $E$. Alors un quotient strict $X'$ de $(X_i)_{i\in I}$
dans $E$ (\hyperref[I.sec7.5.1]{7.5.1}) est connu quand on connaît, pour tout
$T\in \ob C$,
et tout couple $(i,j)\in I\times I$, la partie de $\Hom(T,X_i)\times
\Hom(T,X_j)$ formée des couples $(u,v)$ tels que $p_iu=p_jv$, où
pour tout $i\in I$, $p_i:X_i\rightarrow X'$ désigne le morphisme
canonique. Par suite, le cardinal de l'ensemble des quotients
stricts de $(X_i)_{i\in I}$ dans $E$ est majoré par $ \prod_{T\in
\ob C} \prod_{i, j\in I}2^{\card \Hom(T,X_i)\times
  \card\Hom(T,X_j)}{\quoi}.
$
\end{enonce*}

\setcounter{subsubsection}{2}
\subsubsection{}\etiquette[7.5.3]{I.sec7.5.3}
On voit tout de suite que la conclusion analogue est vraie si on
suppose seulement que $C$ est génératrice pour les monomorphismes
stricts, pourvu que l'on suppose que les produits $X_i\times X_j$
sont représentables et que l'on se borne aux quotients
\emph{effectifs} de la famille $(X_{i})_{i\in I}$, \ie aux quotients
stricts $X'$ tels que les produits fibrés $\fprod{X_i}{X_{j}}{X'}$
soient représentables dans $E$ (ce qui n'est pas une restriction si
$E$ est stable par produits fibrés).

\begin{enonce*}{Proposition 7.6}\etiquette[7.6]{I.prop7.6}
Soient\pageoriginale $E$ une catégorie, $C$ une sous-catégorie
pleine de $E$ génératrice par épimorphismes
(\hyperref[I.def7.1]{7.1}), $\Pi_\circ=$ un
cardinal infini $\geq \card  F\ell  C$, $\Pi\geq \Pi_\circ$ un
cardinal, $(u_i:X_i\rightarrow X)_{i\in I}$ une famille épimorphique
universelle (\hyperref[I.sec10.3]{10.3}) dans $E$ formée de morphismes quarrables
(\hyperref[I.sec10.7]{10.7}) à sources $X_i\in\ob C$, et telle que $\card I\leq \Pi$.
Alors $\card F\ell  C_{/X}\leq \Pi^{\Pi_\circ}$.
\end{enonce*}

Soit $T\in \ob C$, il suffira de prouver qu'on a
\begin{equation*}
\card\Hom(T,X)\leq
  \Pi^{\Pi_\circ}{\quoi},\tag{7.6.1}
\end{equation*}\label{I.eq7.6.1}
car on en conclura successivement
$$
\card\ob C_{/X}\leq(\Pi^{\Pi_\circ})\times\card \ob C_{/X}\leq
(\Pi^{\Pi_\circ})\times \Pi_\circ=\Pi^{\Pi_\circ}{\quoi},
$$
enfin $\card F\ell C_{/X}\leq ((\Pi^{\Pi_\circ})^2)\times
\Pi_\circ=\Pi^{\Pi_\circ}$, puisque pour deux objets $S$, $T$ de
$C_{/X}$, on a $\card \Hom_{C_{/X}}(S,T)\leq
\car\Hom_C(S,T)\leq\Pi_\circ$. Pour prouver \ref{I.eq7.6.1}, notons d'abord le

\begin{enonce*}{Lemme 7.6.2}\etiquette[7.6.2]{I.lem7.6.2}
Soit $J$ un ensemble tel que $\card J=\Pi_\circ$. Pour tout objet
$T$ de $C$, et tout homomorphisme $f:T\rightarrow X$, il existe une
famille $(v_j:S_j\rightarrow T)_{j\in J}$ épimorphique, à sources
$S_j\in\ob C$, et pour tout $j\in J$ un $i(j)\in I$ et un
$g_j:S_j\rightarrow X_{i(j)}$ tels que l'on ait $u_{i(j)}g_j=fv_j$.
\end{enonce*}

En effet, la famille des $\fprod{X_i}{T}{X}\rightarrow$ est
épimorphique par hypothèse, d'autre part, comme $C$ est
génératrice
par épimorphismes stricts, pour tout $i$, la famille des flèches
$S\rightarrow \fprod{X_i}{T}{X}$ de source $S\in \ob C$ est
épimorphique, donc par transitivité la famille des flèches
$v:S\rightarrow T$ de source dans $\ob C$ qui se factorisent par un
des $\fprod{X_i}{T}{X}$ est épimorphique. Or ce sont les flèches
$v:S\rightarrow T$ de source dans $C$ pour lesquelles il existe un
$i\in I$ et un $g:S\rightarrow X_i$ tel que $u_ig=fv$. Comme
l'ensemble de\pageoriginale ces flèches $v:S\rightarrow T$ est
contenu dans $F\ell$ $C$, donc de cardinal $\leq \Pi_\circ$, il peut
s'indexer par l'ensemble d'indices $J$, d'où le lemme.

Notons maintenant qu'un morphisme $f:T\rightarrow X$ est connu quand
on connaît les $fv_j$, qui sont connus quand on connaît les $g_j$.
donc $\card\Hom(T,X)$ est majoré par le cardinal de l'ensemble des
familles $(i(j),S_j,v_j,g_j)_{j\in J}$; comme le cardinal de
l'ensemble des applications $j\rightarrow I$ est $\leq
\Pi^{\Pi_\circ}$, et comme pour une telle application $j\mapsto
i(j)$ fixée, le cardinal de l'ensemble des familles correspondantes
$S_j$, $f_j$, $v_j$ est majoré par celui de l'ensemble des
applications de $J$ dans $F\ell  C\times F\ell  C$, qui est $\leq
\Pi^{\Pi_\circ}_\circ$ puisque $\card(F\ell C\times F\ell
C)=\Pi_\circ^2=\Pi_\circ$, on trouve que le premier membre de
(7.6.1) est majoré par $\Pi^{\Pi_\circ}\times
\Pi_\circ^{\Pi_\circ}=\Pi^{\Pi_\circ}$, ce qui achève la
démonstration de \hyperref[I.prop7.6]{7.6}.

\begin{enonce*}{Proposition 7.7}\etiquette[7.7]{I.prop7.7}
Soient $E$ une catégorie où les produits fibrés sont représentables,
$(X_\alpha)_{\alpha\in A}$ une famille d'objets de $E$ génératrice,
$(\varphi_i)_{i\in I}$ une famille de foncteurs
$\varphi_i:E\rightarrow E_i$ commutant aux produits fibrés. Pour que
$(\varphi_i)$ soit conservative (\hyperref[I.prop5.1]{5.1}), il faut et il
suffit que pour tout $\alpha \in A$ et pour tout sous-objet $X'$ de
$X_\alpha$ distinct de $X_\alpha$, il existe un $i\in I$ tel que
$\varphi_i(X')\rightarrow \varphi_i(X_\alpha)$ ne soit pas un
isomorphisme. Dans ce cas, si $E$ est une $𝒰$-catégorie et
si $A$ est $𝒰$-petit, il existe une partie
$𝒰$-petite $J$ de $I$ telle que $(\varphi_j)_{j\in J}$ soit
déjà une famille conservative de foncteurs.
\end{enonce*}

La nécessité de la condition envisagée de
conservativité étant
évidente, prouvons sa suffisance. En vertu de \hyperref[I.def6.1]{6.1} (iv),
il suffit de
prouver que $(\varphi_i)$ est conservative pour les monomorphismes.
Soit $u:Y'\rightarrow Y$\pageoriginale un monomorphisme dans $E$ qui
n'est pas un isomorphisme, il faut prouver qu'il existe $i\in I$ tel
que $\varphi_i(u)$ n'est pas un isomorphisme. Par hypothèse sur
$(X_\alpha)$, il existe un $\alpha$ et un morphisme
$X_\alpha\rightarrow Y$ qui ne se factorise pas par $Y'$, en
d'autres termes, tel que l'image inverse $X'$ de $Y'$ soit un
sous-objet de $X_\alpha$ distinct de $X_\alpha$. Par hypothèse, il
existe un $i\in I$ tel que $\varphi_i(X')\rightarrow \varphi(X)$ ne
soit pas un isomorphisme. Comme $\varphi_i(X')\simeq
\fprod{\varphi(X_\alpha)}{\varphi_i(Y')}{\varphi_{i(Y)}}$, il
s'ensuit bien que $\varphi_i(Y')\rightarrow \varphi_i(Y)$ n'est pas
un isomorphisme.

La deuxième assertion de \hyperref[I.prop7.7]{7.7} résulte aussitôt de la
première, compte tenu de 7.4.

\begin{enonce*}{Corollaire 7.7.1}\etiquette[7.7.1]{I.cor7.7.1}
Soient $E$ une $𝒰$-catégorie où les produits fibrés sont
représentables, et admettant une famille génératrice d'objets qui
soit $𝒰$-petite. Alors pour toute famille génératrice
$(Y_i)_{i\in I}$ de $E$, il existe une sous-famille génératrice
$𝒰$-petite $(Y_j)_{j\in J}(\card J\in 𝒰)$.
\end{enonce*}

Il suffit en effet d'appliquer \hyperref[I.prop7.7]{7.7} à une
$𝒰$-petite famille génératrice $(X_\alpha)_{\alpha\in A}$ de
$E$ et à la famille des foncteurs $\varphi_i$ $(i\in I)$ représentés
par les $Y_i$.

\begin{enonce*}{Proposition 7.8}\etiquette[7.8]{I.prop7.8}
Soient $E$ une catégorie, $C$ une sous-catégorie pleine génératrice
par épimorphismes stricts, $D$ une catégorie, $\Hhom'(E,D)$ la
sous-catégorie pleine de $\Hhom(E,D)$ formée des foncteurs qui
commutent aux limites inductives du type $C_{/X}$, où $X$ est un
objet quelconque de $E$. Alors le foncteur $F\rightarrow F| C$
induit un foncteur pleinement fidèle
$$
\Hhom'(E,D)\longrightarrow \Hhom(C,D){\quoi}.
$$
Par\pageoriginale suite, si $\Hhom(C,D)$ est une
$𝒰$-catégorie $(1.1.)$, par exemple {\rm(1.1.1. b))} si $C$
est $𝒰$-petite et $D$ est une $𝒰$-catégorie, alors
$\Hhom'(E,D)$ est également une $𝒰$-catégorie.
\end{enonce*}

Soient $F$, $G:E\rightarrow D$ deux foncteurs dans le premier
membre, et $u:F\rightarrow G$ un homomorphisme. Si $X=\varinjlim
X_i$ dans $E$ et si $F$ et $G$ commutent à la limite inductive
envisagée, alors $u(X):F(X)\rightarrow G(X)$ s'identifie à la limite
des morphismes $u(X_i):F(X_i)\rightarrow G(X_i)$, et est donc connu
quand on connaît les $u(X_i)$. Ceci montre que le foncteur envisagé
dans \hyperref[I.prop7.8]{7.8} est fidèle, compte tenu de l'implication a)
$\Rightarrow$ b) dans \hyperref[I.prop7.2]{7.2} (i). Soit inversement $v:F|
\sisi{D}{C}\rightarrow G|\sisi{D}{C}$ un homomorphisme, prouvons qu'il provient d'un
homomorphisme $u:F\rightarrow G$. On définira, pour tout $X\in \ob
E$,
$$
u(X):F(X)=\varinjlim_{C_{/X}}F(X_i)\longrightarrow
G(X)=\varinjlim_{C_{/X}}G(X_i)
$$
comme la limite inductive des $v(X_i)$. Il est immédiat que l'on
obtient bien un homomorphisme fonctoriel en $X$, donc un
$u:F\rightarrow G$, \sisi{en}{et} enfin que le morphisme induit par $u$ de
$F|C\sisi{'}{}$ dans $G|C\sisi{'}{}$ est $v$, ce qui achève la démonstration.


\setcounter{subsection}{8}
\subsection{Familles et sous-catégories cogénératrices}\etiquette[7.9]{I.sec7.9}

Soient $E$ une catégorie, $C$ une sous-catégorie pleine de $E$. On
dit que $C$ est cogénératrice par monomorphismes stricts (\resp
cogénératrice par monomorphismes, \resp cogénératrice, \resp
cogénératrice pour les épimorphismes, \resp cogénératrice pour les
épimorphismes stricts) si la sous-catégorie pleine $C^\circ$ de
$E^\circ$ est génératrice par épimorphismes stricts (\resp\ldotsn).
Terminologie analogue pour les familles. Tous les résultats du
présent numéro concernant\pageoriginale la notion de sous-catégorie
génératrice et ses variantes (\hyperref[I.def7.1]{7.1}), redonnent donc
des résultats
correspondants pour les notions duales, que nous laissons au lecteur
le soin de formuler pour sa satisfaction personnelle.

\begin{enonce*}{Proposition 7.10}\etiquette[7.10]{I.prop7.10}
Soient $E$ une catégorie, $C$ une sous-catégorie pleine
génératrice
$(7.1)$, $D$ une sous-catégorie pleine de $E$. Pour que $D$ soit
cogénératrice $(7.9)$, il faut et il suffit que pour toute double
flèche $\xymatrix{ T\ar@<2pt>[r]^{u,v}\ar@<-2pt>[r]& X }$ dans $E$
de source $T\in \ob C$, avec $u\neq v$, il existe une flèche
$w:X\rightarrow I$ de but $I\in \ob D$, telle que $wu\neq wv$.
\end{enonce*}

Par définition, dire que $D$ est cogénératrice signifie que pour
toute double flèche $\xymatrix{ Y\ar@<2pt>[r]^{u,v}\ar@<-2pt>[r]& X
}$ dans $E$ telle que $u\neq v$, il existe une flèche
$w:X\rightarrow I$, avec $I\in\ob D$, telle que $wu\neq wv$. Donc
\hyperref[I.prop7.10]{7.10} signifie simplement qu'il suffit de tester cette
propriété lorsque $Y\in \ob C$. Or comme $C$ est génératrice,
l'hypothèse $u\neq v$ implique qu'il existe $f:T\rightarrow Y$ telle
que $uf\neq vf$, d'où par hypothèse l'existence d'une
$w:X\rightarrow I$ de but $I\in\ob D$ telle que $w(uf)\neq w(vf)$,
d'où $wu=wv$, \cqfd

\begin{enonce*}{Corollaire 7.11}\etiquette[7.11]{I.cor7.11}
Les notations étant celles de \hyperref[I.prop7.10]{7.10}, supposons que les
objets $I$ de $D$ sont des objets injectifs de $E$, \ie tels que
pour tout monomorphisme $X\rightarrow Y$ dans $E$, l'application
$\Hom(Y,I)\rightarrow \Hom(X,I)$ correspondante soit surjective.
Supposons de plus que toute double flèche $\xymatrix{
T\ar@<2pt>[r]^{u,v}\ar@<-2pt>[r]& X }$ dans $E$ se factorise en une
double flèche épimorphique (\resp épimorphique effective)
$\xymatrix{ T\ar@<2pt>[r]^{u',v'}\ar@<-2pt>[r]& X' }$ suivie d'un
monomorphisme $i:X'\rightarrow X$. Alors dans le critère
\hyperref[I.prop7.10]{7.10} pour que $D$ soit cogénératrice, on peut se borner
aux doubles flèches $(u,v)$ qui sont épimorphiques (\resp
épimorphiques effectives).
\end{enonce*}

On\pageoriginale en conclut:

\begin{enonce*}{Corollaire 7.12}\etiquette[7.12]{I.cor7.12}
Soit $E$ une $𝒰$-catégorie admettant une petite
sous-catégorie génératrice $C$, et telle que tout objet de $E$ soit
source d'un monomorphisme dans un objet injectif de $E$. Supposons
de plus que pour tout $T\in\ob C$, la somme $T\amalg T$ dans $E$
soit représentable, et que tout morphisme de source $T\amalg T$ se
factorise en un épimorphisme effectif suivi d'un monomorphisme.
Alors $E$ admet une petite sous-catégorie pleine $D$ cogénératrice.
De façon précise, on peut prendre $D$ telle que $\card \ob D\leq
\prod_{T,T'\in \ob C}2^{\card(\Hom(T',T\coprod T))^2}$.
\end{enonce*}

En effet, en vertu de \hyperref[I.cor7.11]{7.11}, il suffit pour tout $T\in \ob C$ et pour
tout quotient effectif $X$ de $T\amalg T$, de choisir un plongement
de $X$ dans un objet injectif $I$ de $E$, et de prendre pour $D$ la
sous-catégorie pleine de $E$ engendrée par ces $I$. La conclusion
résulte alors de 7.5.

\begin{enonce*}[remark]{Exemples 7.13}\etiquette[7.13]{I.exem7.13}
Pour construire de petites sous-catégories cogénératrice en termes
de petites sous-catégories génératrices, on est donc amené à
chercher des conditions pour qu'une $𝒰$-catégorie $E$
admette « suffisamment d'injectifs », \ie que tout objet se
plonge dans un objet injectif (par un monomorphisme). Il est bien
connu [Tohoku] que cette condition est satisfaite dans une
$𝒰$-catégorie abélienne à (petites) limites inductives
filtrantes exactes (axiome $AB$ 5 de \loccit) admettant une petite
famille génératrice. La construction de \loccit n'est d'ailleurs pas
liée de façon essentielle aux catégories abéliennes, et marche
également dans la catégorie des faisceaux de $𝒰$-ensembles
sur un espace topologique $X\in 𝒰$. Nous n'énoncerons pas
ici les propriétés d'exactitude qui font marcher la construction en
question, et nous bornerons à signaler que dans\pageoriginale le cas
particulier de la catégorie des faisceaux d'ensembles sur $X$, on se
ramène immédiatement au cas de \loccit de la façon suivante. On note
que si $\mathcal{O}_X$ est un Anneau sur $X$, alors tout objet
injectif $I$ de la catégorie des $\mathcal{O}_X$-Modules est aussi
injectif en tant qu'objet de la catégorie des faisceaux d'ensembles.
En effet, si $F$ est un %%
faisceau d'ensembles, et $\mathcal{O}_X[F]$ le « $\mathcal{O}_X$-Module libre engendré par $F$ », on a par
définition un homomorphisme de faisceaux d'ensembles
\begin{equation*}
F\longrightarrow \mathcal{O}_X[F]\tag{$\ast$}
\end{equation*}
donnant lieu à un isomorphisme, fonctoriel en $F$
$$
\Hom_{\mathcal{O}_X}(\mathcal{O}_X[F], I)\xrightarrow{\sim} \Hom
(F,I){\quoi}.
$$
Comme le foncteur $F\mapsto \mathcal{O}_X[F]$ transforme
manifestement monomorphisme en monomorphisme, il s'ensuit aussitôt
que si $I$ est un $\mathcal{O}_X$-Module injectif, c'est aussi un
faisceau d'ensembles injectif. Il s'ensuit que si le morphisme
$(\ast)$ est un monomorphisme (ce qui est le cas si on prend pour
$\mathcal{O}_X$ un Anneau constant de valeur un anneau $A\neq 0$),
alors un plongement de $\mathcal{O}_X[F^\ast]$ dans un
$\mathcal{O}_X$-Module injectif $I$ donne un plongement de $F$ dans
le faisceau d'ensembles injectif sous-jacent à $I$.
\end{enonce*}

Le même argument s'applique, sans changement, au cas de la catégorie
des faisceaux de $𝒰$-ensembles sur un $𝒰$-site, qui
sera introduite dans l'exposé suivant.

L'intérêt de l'existence d'une petite sous-catégorie
cogénératrice
réside surtout dans le critère de représentabilité
\hyperref[I.cor8.12.7]{8.12.7} plus
bas.

\setcounter{section}{7}
\section{Ind-objets et pro-objets}\etiquette[8]{I.sec8}\pageoriginale

\subsection{Foncteurs cofinaux et sous-catégories
  cofinales}\etiquette[8.1]{I.sec8.1}

\begin{enonce*}{Définition 8.1.1}\etiquette[8.1.1]{I.def8.1.1}
Soit
$$
\displaylines{\rlap{$(8.1.1.1)$}\hfill \varphi:I\longrightarrow
I'\hfill}
$$
un foncteur. On dit que $\varphi$ est un foncteur cofinal si pour
toute catégorie $C$ et pour tout foncteur $u:I'\rightarrow C$,
considérant $\varinjlim u$ et $\varinjlim u\circ \varphi$ comme des
objets de $\Hhom(G(𝒱\text{-}\Ens))^\circ$ (\hyperref[I.def2.1]{2.1},
\hyperref[I.sec2.3.1]{2.3.1}) (où
$𝒱$ est un univers tel que $I$, $I'\in 𝒱$, et que
$C$ soit une $𝒱$-catégorie), le morphisme canonique
$$
\displaylines{\rlap{$(8.1.1.2)$}\hfill \varinjlim
  u\circ\varphi\longrightarrow \varinjlim u\hfill}
$$
est un isomorphisme. Lorsque $\varphi$ est le foncteur d'inclusion
d'une sous-catégorie de $I'$, on dit que $I$ est une sous-catégorie
cofinale si $\varphi$ est cofinal.
\end{enonce*}

\setcounter{subsubsection}{1}
\subsubsection{}\etiquette[8.1.2]{I.sec8.1.2}
Remarquons que pour $\varphi$, $u$ données, la bijectivité de
(8.1.1.2) ne dépend pas du choix de l'univers $𝒱$.

Revenant à la définition des termes figurant dans
(8.1.1.2) (\hyperref[I.def2.1]{2.1}, \hyperref[I.sec2.3.1]{2.3.1}), on voit que dire que
$\varphi$ est
cofinal signifie aussi
que pour tout univers $𝒱$ tel que $I$ et $I'$ soient
$𝒱$-petits, et tout foncteur
$$
v:{I'}^\circ\longrightarrow (𝒱\text{-}\Ens){\quoi},
$$
l'homomorphisme canonique
$$
\displaylines{\rlap{$(8.1.2.1)$}\hfill \varprojlim v\longrightarrow
  \varprojlim v\circ\varphi\hfill}
$$
est un isomorphisme \ie une bijection. Pour voir que cette condition
est bien nécessaire, on observera, posant $u=v^\circ$, qu'elle
signifie aussi que la condition de la définition 8.1.1 est remplie
quand on y fait $C=(𝒱\text{-}\Ens)^\circ$ (ce\pageoriginale
qui implique que les limites inductives envisagées dans (8.1.1.2)
sont représentables dans $C$).

\setcounter{paragraph}{1}
\paragraph{}\etiquette[8.1.2.2]{I.sec8.1.2.2}
Il est immédiat sur la définition que \emph{le composé de deux
  foncteurs cofinaux est cofinal.}

\begin{enonce*}{Proposition 8.1.3}\etiquette[8.1.3]{I.prop8.1.3}
Soit $\varphi:I\rightarrow I'$ un foncteur.
\begin{enumerate}
\renewcommand{\theenumi}{\alph{enumi}}
\renewcommand{\labelenumi}{\theenumi)}
\item Pour que $\varphi$ soit cofinal, il faut que $\varphi$
  satisfasse la condition:

\begin{itemize}
\item[F 1)] Pour tout objet $i'$ de $I'$, il existe un objet $i$ de
  $I$ tel que $\varphi(i)$ majore $i'$ (\ie tel que
  $\Hom(i',\varphi(i))\neq ∅$).
\end{itemize}

\item Supposons $I$ filtrante. Pour que $\varphi$ soit cofinal, il
  faut et il suffit que $\varphi$ satisfasse la condition {\rm (F 1)}
  ainsi que la condition suivante:
\begin{itemize}
\item[F 2)] Pour tout objet $i$ de $I$ et toute double flèche
$\xymatrix{ i'\ar@<2pt>[r]^{f',g'}\ar@<-2pt>[r]& \varphi(i) }$ dans
$I'$ de but $\varphi(i)$, il existe une flèche $h:i\rightarrow j$
dans $I$ telle que $\varphi(h)f'=\varphi(h)g'$.
\end{itemize}
De plus, si $\varphi$ est cofinal, $I'$ est filtrante.

\item Supposons $I'$ filtrante, et $\varphi$ pleinement fidèle. Pour
  que $\varphi$ soit cofinal, il faut et il suffit qu'il satisfasse la
  condition {\rm F 1)} de {\rm a)}. Cela implique que $I$ est filtrante.
\end{enumerate}
\end{enonce*}

\noindent \emph{Démonstration.} Pour la nécessité dans a) et b), on
utilisera la définition \hyperref[I.def8.1.1]{8.1.1} dans le seul cas où $u$ est le
foncteur d'inclusion canonique (\hyperref[I.sec1.3.3]{1.3.3})
$$
u:I'\hookrightarrow \widehat{I}'
$$
(un univers $𝒰$ tel que $I$, $I'$ soient $𝒰$-petits
étant choisi). On peut alors interpréter (8.1.1.2)
comme une flèche
de $\widehat{I}'$ (\hyperref[I.prop3.1]{3.1}), dont le but est le foncteur\pageoriginale
\emph{final} sur $I'$ \hyperref[I.prop3.4]{3.4}). On voit immédiatement que
la condition
     F 1) exprime que cette flèche est un épimorphisme de
$\widehat{I}'$, \ie est surjective sur chaque argument (compte tenu
que les limites inductives dans $\widehat{I}'$ se calculent « argument par argument »). Cela prouve a). Supposons maintenant $I$
filtrante et $\varphi$ cofinal. Alors $I$ est filtrante: en effet,
la condition que deux objets de $I'$ soient majorés par un troisième
résulte aussitôt de la même condition sur $I$, et de F 1), et il
faut seulement encore prouver la condition PS 2) de \hyperref[I.def2.7]{2.7}, \ie que
pour toute double flèche
$$
f',g':i'\rightrightarrows j'
$$
dans $I'$, il existe une flèche $h':j'\rightarrow k'$ de $I'$ telle
que $h'f'=h'g'$. Or en vertu de F~1) on peut supposer $k'$ de la
forme $\varphi(i)$. Mais comme par l'hypothèse $\varphi$ cofinal,
pour tout objet $i'$ de $I'$\; $\varinjlim_i\Hom(i',\varphi(i))$ est
l'ensemble réduit à un élément, il résulte de la description
standard des limites inductives \emph{filtrantes} dans (Ens) (\hyperref[I.lem2.8.1]{2.8.1})
que cela implique l'existence d'une flèche $h:i\rightarrow j$ dans
$I$ telle que $\varphi(h)f'=\varphi(h)g'$. Cela achève donc de
prouver que $I'$ est filtrant, et prouve en même temps la condition
F~2).

Pour prouver que les conditions énoncées dans b) sont suffisantes
pour que $\varphi$ soit cofinal, on utilise la forme \hyperref[I.sec8.1.2]{8.1.2} de la
définition. On constate aussitôt que la condition F 1) implique que
(8.1.2.1) est un monomorphisme \ie est injectif sur chaque argument,
tandis que la condition F 2) (jointe à F 1) assure qu'il est
bijectif (compte tenu du calcul des limites projectives dans
$\widehat{{I'}^\circ}$ argument par argument). Cela prouve donc b).

Enfin,\pageoriginale si $I'$ est filtrante et $\varphi$ pleinement
fidèle, alors il est immédiat que la condition F 1) implique que $I$
est filtrante, et implique la condition F 2). Donc c) résulte de a)
et de b).

\setcounter{subsubsection}{3}
\subsubsection{}\etiquette[8.1.4]{I.sec8.1.4}
Lorsque $I'$ est une catégorie filtrante, $I$ une sous-catégorie
pleine de $I'$, on voit donc par 8.1.3 c) que la condition que $I$
soit cofinale dans $I'$ ne dépend que de la partie $\Ob$ $I$ de
l'ensemble préordonné $OB$ $I'$ (pour la relation de préordre
« $x\leq y$ » $\Leftrightarrow$ « $\Hom(x,y)\neq ∅$ »). Si $J'$ est un ensemble préordonné, $J$ une partie de $J'$, on
%%
dira parfois que $J$ est une \emph{partie cofinale de} $J'$ lorsque
tout élément de $J'$ est majoré par un élément de $J$. Lorsque $J'$
est filtrante, cela signifie donc que le foncteur d'inclusion
$𝒥\hookrightarrow 𝒥'$ pour les catégories associées
est cofinal.

\subsubsection{}\etiquette[8.1.5]{I.sec8.1.5}
Dans la suite, nous n'utiliserons la notion de foncteur cofinal que
dans les cas où les catégories $I$ et $I'$ sont filtrantes.
Classiquement, on se bornait même à des catégories associées à des
ensembles préordonnés (\ie dans lesquelles il existe au plus une
flèche de source et de but donnés). Il apparaît cependant que cette
restriction est gênante dans les applications, les catégories
filtrantes « naturelles » qui s'introduisent dans de nombreuses
applications n'étant \emph{pas} des catégories ordonnées. Le
résultat suivant, %%
dû à P. \textsc{Deligne}, montre cependant qu'il n'y a pas de
différence essentielle entre les deux points de vue:

\begin{enonce*}{Proposition 8.1.6}\etiquette[8.1.6]{I.prop8.1.6}
Soit\pageoriginale $I$ une petite catégorie filtrante. Alors il
existe un petit ensemble ordonné $E$, et un foncteur cofinal
$\varphi:ℰ\rightarrow I$, où $ℰ$ désigne la
catégorie associée à $E$.
\end{enonce*}

Supposons d'abord que l'ensemble préordonné $\Ob I$ n'ait pas de
plus grand élément. Appelons sous-diagramme de $I$ un couple
$D=(\mathcal{O},F)$ formé d'une partie $F$ de $F\ell(I)$ et d'une
partie $\mathcal{O}$ de $\Ob I$, tel que pour toute $f\in F$, la
source et le but de $f$ appartiennent à $\mathcal{O}$. Un élément
$e$ de $\mathcal{O}$ est appelé \emph{objet final} du sous-diagramme
$D$ si pour tout $x\in D'$, l'ensemble $\Hom(x,e)\cap F$ des flèches
de $D$ de source $x$ et de but $y$ a exactement un élément $f_x$, si
pour toute flèche $g:x\rightarrow y$ de $D$, on a $f_x=f_y\circ g$
et si $f_e=\id_e$. Soit $E$ l'ensemble des sous-diagrammes
\emph{finis} $D$ de $I$ ayant un élément final \emph{unique}
$\varphi(D)$, et ordonnons $E$ par inclusion. Si $D$, $D'$ sont deux
éléments de $E$ et $D\subset D'$, alors il existe une unique flèche
$\varphi(D',D)$ de $D'$ de source $\varphi(D)$, de but
$\varphi(D')$, et si on a des inclusions $D\subset D'\subset D''$
dans $E$, on a évidemment
$\varphi(D'',D)=\varphi(D'',D')\varphi(D',D)$; on a également
$\varphi(D,D)=\id_{\varphi(D)}$. On trouve donc un foncteur
$\varphi:ℰ\rightarrow I$, et tout revient à prouver que $E$
est filtrant et que $\varphi$ est cofinal. Compte tenu de
\hyperref[I.prop8.1.3]{8.1.3} b), on doit faire trois vérifications:
\begin{enumerate}
\renewcommand{\labelenumi}{\theenumi)}
\item Condition F 1) de 8.1.3: pour tout $i\in I$, on prend pour $D$
  le sous-diagramme de $I$ réduit à l'objet $i$ et sa flèche
  identique, alors $i\leq \varphi(D)=i$. Cette condition F 1) montre
  en même temps que $E\neq ∅$.

\item Condition F 2): Pour tout $i\in \Ob I$, tout $D\in E$ et toute
  double flèche $i\rightrightarrows \varphi(D)$, trouver un
  diagramme $D'\in E$, contenant $D$, tel que
  $\varphi(D',D):\varphi(D)\rightarrow \varphi(D')$ \emph{égalise}
  la double flèche. Or comme $I$ est filtrant, il\pageoriginale
  existe un objet $j$ de $I$ et une flèche
$f:\varphi(D)\rightarrow j$ qui égalise la double flèche. Comme $I$
n'a pas de plus grand élément, on peut supposer que $j$ est un
majorant \emph{strict} de $\varphi(D)$, ce qui implique qu'il est
distinct de tous les autres objets de $D$. Soit alors $D'$ le
sous-diagramme de $I$ dont l'ensemble d'objets est
$\mathcal{O}\cup\{j\}$, et l'ensemble des flèches est formé des
flèches de $D$, plus les composés
$x\xrightarrow{f_x}\varphi(D)\xrightarrow{f}j$, où $x$ est un objet
de $D$ et $f_x$ est l'unique flèche de $D$ de source $x$ et de but
$\varphi(D)$, et $\id_j$. Il est clair alors que $D'$ est un
sous-diagramme fini de $I$, qu'il admet $j$ comme unique objet
final, donc $D'\in E$, et $D'$ satisfait à la condition voulue.

\renewcommand{\labelenumi}{$3^\circ$}
\item Deux sous-diagrammes $D$, $D'\subset E$ sont contenus dans un
  même $D''\in E$. On peut en effet trouver un majorant strict $j$
  de $\varphi(D)$ et de $\varphi(D')$,
\[
\xymatrix@C=1.5cm@R=1cm{
\varphi(D)\ar[rd]^f && \\
           &j &\\
\varphi(D')\ar[ru]_{f'}&\quad , }
\]
et on prend pour $D''$ le sous-diagramme dont l'ensemble des objets
est la réunion de l'ensemble des objets de $D$, de $D'$ et de
$\{j\}$, et l'ensemble des flèches est la réunion de l'ensemble des
flèches de $D$, de $D'$, de l'ensemble des composés $f\circ f_x$
($x$ objet de $D$) et $f'\circ f'_{x'}$ ($x'$ objet de $D'$), et
$\{\id_j\}$. On obtient bien ainsi un sous-diagramme fini de $E$,
montrons que, quitte à remplacer $j$, $f$, $f'$ par $j'$, $gf$,
$gf'$ avec $g:j\rightarrow j'$ convenable, on peut obtenir que $j$
soit un objet final de $D'$. Il revient au même de dire que pour
tout $x$ qui est à la fois objet de $D$ et de $D'$, on a
$ff_x=f'f'_{x'}$. Or c'est là d'un ensemble \emph{fini} de
doubles-flèches qu'il s'agit\pageoriginale d'égaliser par un
$g:j\rightarrow j'$, ce qui est possible puisque $I$ est filtrante.
\end{enumerate}

Cela achève la démonstration dans le cas envisagé. On ramène le cas
général à celui-ci, en introduisant la catégorie filtrante $\NN$
associée à l'ensemble ordonné $\NN$ des entiers naturels, et en
notant que la catégorie $\NN\times I$ est filtrante, qu'elle n'a pas
de plus grand élément, et que la projection $\NN\times I\rightarrow
I$ est un foncteur cofinal.

\begin{enonce*}{Corollaire 8.1.7}\etiquette[8.1.7]{I.cor8.1.7}
Soit $I$ une $𝒰$-catégorie. Pour qu'il existe un petit
ensemble ordonné filtrant $E$, et un foncteur cofinal
$ℰ\rightarrow I$, il faut et il suffit que $I$ soit
filtrante et que $\Ob I$ admette une petite partie cofinale $I'$
(\hyperref[I.sec8.1.4]{8.1.4}).
\end{enonce*}

C'est nécessaire en vertu de 8.1.3 a), et suffisant en vertu de
8.1.6 appliqué à la sous-catégorie pleine de $I$
définie par $I'$.

\begin{enonce*}{Définition 8.1.8}\etiquette[8.1.8]{I.def8.1.8}
Soit $I$ une catégorie filtrante. On dit que $I$ est essentiellement
petite si $I$ est une $𝒰$-catégorie, et si elle satisfait
aux conditions équivalentes de \hyperref[I.cor8.1.7]{8.1.7}.
\end{enonce*}

\setcounter{subsection}{1}
\subsection{Ind-objets et foncteurs ind-représentables}\etiquette[8.2]{I.sec8.2}

\subsubsection{}\etiquette[8.2.1]{I.sec8.2.1}
Nous fixons dans la suite une $𝒰$-catégorie $C$, que nous
considérons toujours comme plongée dans la catégorie $\widehat{C}$
des $𝒰$-préfaisceaux à l'aide du foncteur canonique (1.3.1)
$$
\displaylines{\rlap{$(8.2.1.1)$}\hfill h:C\hookrightarrow
  \widehat{C}{\quoi}.\hfill}
$$
Nous appellerons \emph{ind-objet de} $C$ (ou \emph{système inductif}
de $C$), tout foncteur
$$
\displaylines{\rlap{$(8.2.1.2)$}\hfill \varphi:I\longrightarrow
  C{\quoi},\hfill}
$$
où\pageoriginale $I$ est une catégorie \emph{filtrante}
(\hyperref[I.def2.7]{2.7}),
appelée la \emph{catégorie d'indices} du ind-objet. ($I$ est par
ailleurs quelconque, et nous n'exigeons pas, notamment, que $I$ soit
une $𝒰$-catégorie.) Il sera souvent commode de désigner un
ind-objet $\varphi$ par la notation indicielle $(X_i)_{i\in \ob I}$,
ou même (par nouveau abus de notation) $(X_i)_{i\in I}$, ou
$(X_i)_i$ ou $(X_i)$, où $X_i=\varphi(i)$ pour $i\in \ob I$. Cette
notation n'est ni plus ni moins abusive que la notation utilisée
classiquement pour les systèmes inductifs indexés par les ensembles
ordonnés filtrants (où la mention des morphismes de transition est
également absente). Les $X_i$ s'appellent les \emph{objets
  composants} du ind-objet $𝒳$. On fera attention que dans le cas
général envisagé ici, les « \emph{morphismes de transition} »
$\varphi(f)$ du système inductif sont indexés par l'ensemble $F\ell$
$I$ des flèches de $I$, qui ne s'identifie pas en général à une
partie de $\Ob I\times \Ob I$; en d'autres termes, pour $i$ et $j$
donnés, $j$ majorant $i$, il peut exister plus d'un morphisme de
transition de $X_i$ dans $X_j$.

\subsubsection{}\etiquette[8.2.2]{I.sec8.2.2}
Les ind-objets $𝒳=(X_i)_{i\in I}$ de $C$ les plus utiles
sont ceux pour lesquels la catégorie d'indices $I$ est
essentiellement petite (8.1.8); un tel ind-objet est appelé
\emph{essentiellement
  petit}. Si $(X_i)_{i\in I}$ est ainsi, utilisant le fait que dans
$\widehat{C}$ les petites limites inductives sont représentables
(3.1), on peut considérer
$$
\displaylines{\rlap{$(8.2.2.1)$}\hfill
L(𝒳)\mathrel{\mathop=^{\rm def}}
  \varinjlim h\cdot 𝒳 = \varinjlim_i X_i \in \ob
  \widehat{C}{\quoi},\hfill}
$$
où la dernière limite est prise \emph{dans} $\widehat{C}$. Donc
$L(𝒳)$ est le préfaisceau
$$
\displaylines{\rlap{$(8.2.2.2)$}\hfill L(𝒳):Y\longmapsto
  \varinjlim_i \Hom(Y,X_i)\hfill}
$$
sur\pageoriginale $C$. On dit que ce préfaisceau est le
\emph{préfaisceau
  ind-représenté par le ind-objet $𝒳=(X_i)_{i\in I}$ de $C$.} Un
préfaisceau $F$ sur $C$ est appelé un \emph{préfaisceau
  ind-représentable} s'il est isomorphe à un préfaisceau
$L(𝒳)$ ind-représenté par un ind-objet essentiellement
petit. Il résulte d'ailleurs de \hyperref[I.prop8.1.6]{8.1.6} que dans cette
définition, on peut supposer que $I$ est une petite catégorie, et
même que $I$ est associée à un ensemble ordonné $\in 𝒰$.

\subsubsection{}\etiquette[8.2.3]{I.sec8.2.3}
Considérons un ind-objet $𝒳=(X_i)_{i\in I}$, donné par un
foncteur $\varphi:I\rightarrow C$, et soit
$$
u:I'\longrightarrow I
$$
un foncteur, où $I'$ est une deuxième catégorie filtrante. Alors le
foncteur $\varphi\circ u$ est un ind-objet $𝒳'$ de $C$ de
catégorie d'indices $I'$, qui en notation indicielle s'écrit
$(X_{u(i')})_{i'\in I'}$, qu'on appelle le ind-objet de $C$
\emph{déduit de $(X_i)_{i\in I}$ par changement de catégories
  d'indices} à l'aide du foncteur $u$. Si $I$ et $I'$ (\ie $𝒳$ et
$𝒳'$) sont essentiellement petits, on trouve un morphisme
canonique
$$
\displaylines{\rlap{$(8.2.3.1)$}\hfill
  L(𝒳')=\varinjlim_{i'}X_{u(i')}\longrightarrow L(𝒳)=\varinjlim_i
  X_i\hfill}
$$
entre les préfaisceaux ind-représentés par $𝒳'$ et
$𝒳$ respectivement. Lorsque $u$ est un foncteur cofinal
(\hyperref[I.def8.1.1]{8.1.1}), alors $𝒳$ est essentiellement petit si et seulement
si $𝒳'$ l'est (8.1.7), et l'homomorphisme précédent
(8.2.3.1) est un isomorphisme
$$
\displaylines{\rlap{$(8.2.3.2)$}\hfill L(𝒳')\simeq
  L(𝒳){\quoi}.\hfill}
$$

\subsubsection{}\etiquette[8.2.4]{I.sec8.2.4}
Soient\pageoriginale
$$
\displaylines{\rlap{$(8.2.4.1)$}\hfill 𝒳=(X_i)_{i\in
    I},𝒴=(Y_j)_{j\in I}\hfill}
$$
deux ind-objets de $C$ essentiellement petits, indexés par des
catégories (filtrantes et essentiellement petites) $I$ et $J$. On
appelle \emph{morphisme} du ind-objet $𝒳$ dans le ind-objet
$𝒴$ tout morphisme
$$
\displaylines{\rlap{$(8.2.4.2)$}\hfill L(𝒳)\longrightarrow
  L(𝒴)\hfill}
$$
entre les préfaisceaux qu'ils ind-représentent. On pose
$$
\displaylines{\rlap{$(8.2.4.3)$}\hfill
\Hom_{\text{ind}-\ob}(𝒳,𝒴)=\Hom_{\widehat{C}}(L(𝒳),L(𝒴)){\quoi}.
   \hfill}
$$
(On supprime l'indice ind-ob au $\Hom$ si on ne craint pas de
confusion). On définit la composition des ind-objets de $𝒞$
comme la composition des morphismes des préfaisceaux qu'ils
ind-représentent. Si $𝒱\supset 𝒰$ est un univers
contenant $𝒰$, et si on se borne aux ind-objets
essentiellement petits qui sont $\in 𝒱$, ceux-ci forment
donc l'ensemble d'objets d'une catégorie, dont l'ensemble des
flèches est formé des triples $(𝒳,𝒴,u)$, où
$𝒳$ et $𝒴$ sont des ind-objets essentiellement
petits $\in 𝒱$ et où $u:𝒳\rightarrow 𝒴$ est
un morphisme de ind-objets, \ie un morphisme
$L(𝒳)\rightarrow L(𝒴)$. La catégorie ainsi obtenu
est notée
$$
\displaylines{\rlap{$(8.2.4.4)$}\hfill
  \Ind_{𝒱}(C,𝒰){\quoi}.\hfill}
$$
Lorsque $𝒱=𝒰$, on la note simplement
$$
\displaylines{\rlap{$(8.2.4.5)$}\hfill \Ind_{𝒰}(C)\text{ ou
}
  \Ind(C){\quoi},\hfill}
$$
et on l'appelle \emph{catégorie des ind-objets de} $C$ relativement
à $𝒰$, en supprimant la mention de $𝒰$ quand aucune
confusion n'est à craindre.

Évidemment,\pageoriginale si $𝒱\subset 𝒱'$ sont deux
univers contenant $𝒰$, $\Ind_{𝒱}(C,𝒰)$ est
une sous-catégorie pleine de $\Ind_{𝒱'}(,𝒰)$; il
résulte alors de 8.2.3 que le foncteur d'inclusion est une
\emph{équivalence de catégories}:
$$
\displaylines{\rlap{$(8.2.4.6)$}\hfill
\Ind_{𝒱}(C,𝒰)\xrightarrow{\approx}\Ind_{𝒱'}(C,𝒰){\quoi}.\hfill}
$$
Cela montre en particulier que \emph{tout ind-objet essentiellement
  petit de $C$ définit un élément de $\Ind(C)$, à isomorphisme
  unique près.} Cette constatation justifie l'abus de langage, assez
courant en pratique, consistant à identifier tout ind-objet
essentiellement petit de $C$ à un objet de $\Ind(C)$.

Il résulte aussitôt des définitions que pour tout univers
$𝒱$, nous avons un foncteur canonique
$$
\displaylines{\rlap{$(8.2.4.7)$}\hfill
  L:\Ind_{𝒱}(C,𝒰)\longrightarrow \widehat{C}\hfill}
$$
qui est pleinement fidèle. Ces foncteurs sont évidemment connus à
isomorphisme unique près, grâce à (8.2.4.6), lorsqu'on
connaît l'un
d'eux, et en particulier lorsqu'on connaît le foncteur canonique
$$
\displaylines{\rlap{$(8.2.4.8)$}\hfill L:\Ind(C)\longrightarrow
  \widehat{C}{\quoi}.\hfill}
$$
Comme ce foncteur est pleinement fidèle, on l'utilise fréquemment
pour identifier un objet du premier membre avec le foncteur qu'il
pro-représente, voire pour identifier $\Ind(C)$ avec son image
essentielle dans $\widehat{C}$, formée des préfaisceaux
ind-représentables. On fera attention cependant que cette
identification présente nettement plus d'inconvénients que
l'identification analogue de $C$ à une sous-catégorie pleine de
$\widehat{C}$ par le foncteur $h$ (8.2.1.1), car contrairement à ce
dernier, le foncteur (8.2.4.8) n'est pas en général
injectif sur les objets.

\subsubsection{}\etiquette[8.2.5]{I.sec8.2.5}
Explicitons\pageoriginale la définition (8.2.4.3) en termes des
expressions indicielles (8.2.4.1) des ind-objets considérés. On
trouve, compte tenu de la définition (8.2.2.1):
$$
\Hom(𝒳,𝒴)\simeq \varprojlim_i
\Hom(X_i,𝒴)\mathrel{\mathop=^{\rm dfn}} \varprojlim_i
\Hom(X_i, L(𝒴))=\varprojlim_i\varinjlim_j \Hom_C(X_i,
Y_j){\quoi},
$$
la dernière égalité provenant de (8.2.2.2)
appliqué à $𝒴$.
On trouve donc une bijection canonique
$$
\displaylines{\rlap{$(8.2.5.1)$}\hfill \Hom((X_i)_{i\in
I},(Y_j)_{j\in
    J})\simeq \varprojlim_{i\in I} \varinjlim_{j\in
    J}\Hom_C(X_i,Y_j){\quoi}.\hfill}
$$
Nous laissons au lecteur le soin d'expliciter la composition des
morphismes de ind-objets sur cette formule. Notons qu'il résulte
aussitôt de cette formule que l'ensemble $\Hom(𝒳,𝒴)$
d'homomorphismes de ind-objets est $𝒰$-\emph{petit}, \ie que
les catégories (8.2.4.4) sont des
$𝒰$-catégories. Il revient
au même de dire (en vertu de (8.2.4.6)) que la \emph{catégorie
  $\Ind(C)$ est une $𝒰$-catégorie,} \ie que le deuxième membre
de (8.2.5.1) est petit lorsque $I$, $J\in 𝒰$, ce qui résulte
aussitôt du fait que $C$ est une $𝒰$-catégorie donc les
$\Hom_C(X_i,Y_j)$ sont petits.

\subsubsection{}\etiquette[8.2.6]{I.sec8.2.6}
Soit $I$ une catégorie filtrante essentiellement petite. Alors on
voit sur (8.2.5.1), ou sur la fonctorialité de la limite inductive
d'un foncteur $I\rightarrow \widehat{C}$ par rapport au dit
foncteur, qu'on a un foncteur canonique
$$
\displaylines{\rlap{$(8.2.6.1)$}\hfill \Hhom(I,C)\longrightarrow
  \Ind(C){\quoi}.\hfill}
$$
On fera attention que ce foncteur \emph{n'est pas} en général
pleinement fidèle, ni même fidèle.

\begin{enonce*}[remark]{Remarque 8.2.7}\etiquette[8.2.7]{I.sec8.2.7}
Les\pageoriginale définitions de \hyperref[I.sec8.2.4]{8.2.4} rendent claire la notion
d'\emph{isomorphie} de deux ind-objets (essentiellement petits), qui
signifie aussi l'isomorphie des préfaisceaux qu'ils
ind-représentent. On fera attention que c'est une notion
d'isomorphie très faible, si on la compare à la notion d'isomorphie
dans des catégories de la forme $\Hhom(I,C)$. Ainsi, un grand nombre
de relations assez naturelles qu'on peut imposer à un ind-objet (tel
que celle d'avoir ses composantes dans une sous-catégorie
strictement pleine donnée, ou celle d'être strict
(\hyperref[I.sec8.12]{8.12} ci-dessous))
\emph{ne sont pas} stables par isomorphisme. Il faudra donc, si on
tient à travailler avec des notions stables par isomorphisme de
ind-objets, « saturer » les notions envisagées en passant à la
sous-catégorie strictement pleine de $\Ind(C)$ engendrée par la
sous-catégorie de $\Ind(C)$ formée des objets satisfaisant à la
condition envisagée. Voir par exemple la notion de ind-objet
essentiellement constant, introduite plus bas (\hyperref[I.sec8.4]{8.4}).
\end{enonce*}

\begin{enonce*}[remark]{Exercice 8.2.8}\etiquette[8.2.8]{I.exer8.2.8}
Soient $𝒰\subset 𝒱$ deux univers et $C$ un
$𝒰$-catégorie qui appartienne à $𝒱$. On désigne par
$\SysInd(C)$ la catégorie suivante:
\begin{enumerate}
\renewcommand{\theenumi}{\alph{enumi}}
\renewcommand{\labelenumi}{\theenumi)}
\item Les objets de $\SysInd(C)$ sont les foncteurs
$\varphi:I\rightarrow C$ où $I$ est une catégorie filtrante
essentiellement $𝒰$-petite appartenant à $𝒱$.

\item Soient $(I,\varphi)$ et $(J,\psi)$ deux objets de
$\SysInd(C)$. Un morphisme de $\SysInd(C)$ de $(I,\varphi)$ dans
$(J,\psi)$ est un couple $(m,u)$, où $m:I\rightarrow J$ est un
foncteur et où $u:\varphi\rightarrow \psi\circ m$ est un morphisme
de foncteurs. La composition des morphismes dans $\SysInd(C)$ se
définit de manière évidente.
\end{enumerate}
\end{enonce*}

Notons\pageoriginale $S$ l'ensemble des morphismes
$(m,u):(I,\varphi)\rightarrow (J,\psi)$ de $\SysInd(C)$ tel que $m$
soit cofinal et $u$ un isomorphisme. Soit $(I,\varphi)$ un objet de
$\SysInd(C)$. La catégorie $F\ell(I)$ des flèches de $I$ s'envoie
par deux foncteurs naturels dans $I$: la source et le but. De plus
ces deux foncteurs sont liés par le morphisme canonique de foncteurs
$v$: source $\rightarrow$ but. D'où deux morphismes dans
$\SysInd(C)$ de $(F\ell(I),\varphi\circ\text{ source})$ dans
$(I,\varphi):p_1(I,\varphi)=(\text{source, } \id)$,
$p_2(I,\varphi)=(\mbox{but, } \varphi^\ast v:\varphi\circ
\text{source }\rightarrow \varphi\circ{\text{ but}})$. Notons
$p:\SysInd(C)\rightarrow\Ind(C)$ le foncteur évident. Soit $B$ une
catégorie. Montrer que le foncteur $F\mapsto F\circ p$:
$$
\Hhom(\Ind(C),B)\longrightarrow \Hhom(\SysInd(C),B)
$$
est pleinement fidèle, et qu'un foncteur $G:\SysInd(C)\rightarrow B$
appartient à l'image essentielle si et seulement s'il possède les
deux propriétés suivantes:
\begin{enumerate}
\renewcommand{\labelenumi}{\theenumi)}
\item Pour tout $s\in S$, $F(s)$ est un isomorphisme de $B$.

\item Pour tout objet $(I,\varphi)$ de $\SysInd(C)$, $F(p_1(I,\varphi))=F(p_2(I,\varphi))$.
\end{enumerate}

\subsection{Caractérisation des foncteurs ind-représentables}\etiquette[8.3]{I.sec8.3}

\begin{enonce*}{Proposition 8.3.1}\etiquette[8.3.1]{I.prop8.3.1}
Soit $F$ un préfaisceau ind-représentable sur $C$. Alors $F$ est
exact à gauche, \ie (\hyperref[I.cor2.3.2]{2.3.2}) pour toute catégorie finie
$J$ et tout
foncteur $\varphi:J\rightarrow C$ tel que $\varinjlim \varphi$ soit
représentable (\ie $\varprojlim \varphi^\circ$ soit représentable),
l'application canonique
$$
F(\varinjlim \varphi)\longrightarrow \varprojlim F\circ\varphi
$$
est bijective.
\end{enonce*}

En\pageoriginale effet, les préfaisceaux représentables étant
évidemment exacts à gauche, il résulte aussitôt
de \hyperref[I.prop2.8]{2.8} qu'il en est
de même de toute limite inductive filtrante de tels foncteurs, \cqfd

\setcounter{subsubsection}{1}
\subsubsection{}\etiquette[8.3.2]{I.sec8.3.2}
$F$ étant un préfaisceau sur $C$, nous aurons à travailler souvent
avec la catégorie
$$
\displaylines{\rlap{$(8.3.2.1)$}\hfill C_{/F}\hookrightarrow
\widehat{C}_{/F}{\quoi},\hfill}
$$
qui désigne la sous-catégorie pleine du deuxième membre formée des
flèches $X\rightarrow F$ \emph{de source dans} $C$. Il résulte de
1.4 que les objets de cette catégorie s'identifient aux couples
$(X,u)$, où $X$ est un objet de $C$ et $u\in F(X)$. Un morphisme de
$(X,u)$ dans $(X',u')$ s'interprète alors comme une flèche
$f:X\rightarrow X'$ dans $C$ telle que l'on ait $F(f)(u')=u$. Le
foncteur « oubli du but » de $\widehat{C}_{/F}$ dans
$\widehat{C}$ induit un foncteur (qualifié de \emph{canonique}).
$$
\displaylines{\rlap{$(8.3.2.2)$}\hfill C_{/F}\longrightarrow
C{\quoi},\hfill}
$$
déjà considéré dans \hyperref[I.prop3.4]{3.4}, où on prouve que
la limite inductive de ce
foncteur existe dans $\widehat{C}$ (sans condition de petitesse sur
$C_{/F}$ !) et est isomorphe canoniquement à $F$.


\setcounter{paragraph}{2}
\paragraph{}\etiquette[8.3.2.3]{I.sec8.3.2.3}
Notons que si dans $C$ les limites inductives finies sont
représentables, et si $F$ est exact à gauche (\ie les transforme en
limites projectives finies de $(\Ens)$), alors il en est de même
dans $C_{/F}$, et fortiori (\hyperref[I.exem2.7.1]{2.7.1}) $C_{/F}$ est filtrante. Si,
plus généralement, dans $C$ les sommes de deux objets et les
conoyaux de doubles flèches sont représentables, et si $F$ les
transforme en produits \resp en noyaux, alors\pageoriginale $C_{/F}$
est stable sous les mêmes types de limites inductives finies, donc
elle est filtrante si et seulement si elle est non vide, \ie si et
seulement si le foncteur $F$ n'est pas le foncteur constant de
valeurs $∅$.

\begin{enonce*}{Théorème 8.3.3}\etiquette[8.3.3]{I.them8.3.3}
Soient $C$ une $𝒰$-catégorie, et $F\in \ob \widehat{C}$,
$F:C^\circ\rightarrow (𝒰\text{-}\Ens)$ un
$𝒰$-préfaisceau sur $C$. Les conditions suivantes sont
équivalentes:
\begin{itemize}
\item[(i)] $F$ est ind-représentable (\hyperref[I.sec8.2.2]{8.2.2}).

\item[(ii)] La catégorie $C_{/F}$ (\hyperref[I.sec8.3.2]{8.3.2}) est filtrante et
essentiellement petite.

\item[(iii)] (si dans $C$ les limites inductives finies sont
représentables.) Le foncteur $F$ est exact à gauche, et $\ob C_{/F}$
admet une petite partie cofinale (\hyperref[I.sec8.1.4]{8.1.4}).

\item[(iii bis)] (Si dans $C$ les sommes de deux objets et les
conoyaux de doubles flèches sont représentables.) Le foncteur $F$
transforme somme de deux objets de $C$ en produits, conoyaux de
doubles flèches de $C$ en noyaux, la catégorie $C_{/F}$ est non vide
\ie le foncteur $F$ n'est pas le foncteur constant de valeur
$∅$, enfin il existe une petite famille d'objets de
$C_{/F}$ telle que tout objet $X$ de $C_{/F}$ soit majoré par un
objet $X_i$, \ie admette un $F$-morphisme $X_i\rightarrow X$.

\item[(iv)] (si la catégorie $C$ est équivalente à une petite
catégorie.) La catégorie $C_{/F}$ est filtrante.

\item[(v)] (Si la catégorie $C$ est équivalente à une petite
catégorie, et si les limites inductives filtrantes $y$ sont
représentables.) Le foncteur $F$ est exact à gauche.
\end{itemize}
\end{enonce*}

\noindent \emph{Démonstration.} (i)\pageoriginale $\Rightarrow$
(ii). Supposons $F$ ind-représenté par $(X_i)_{i\in I}$, avec $I$
petit. Prouvons que $C_{/F}$ est filtrante. Soient $X$, $X'$ deux
objets de $C_{/F}$, \ie des objets de $C$ munis de morphismes
$X\rightarrow F$, $X'\rightarrow F$, prouvons qu'ils sont majorés
par un troisième objet de $C_{/F}$. Or les morphismes $X\rightarrow
F$, $X'\rightarrow F$ proviennent de morphismes $X\rightarrow X_i$,
$X'\rightarrow X'_i$, et quitte à remplacer $i$, $i'$ par un
majorant commun dans $\Ob I$, on peut supposer $i=i'$, et on prend
comme majorant commun de $X$, $X'$ l'objet $X_i$ muni du morphisme
canonique $X_i\rightarrow F$. Soit maintenant $\xymatrix{
X\ar@<2pt>[r]^{f,g}\ar@<-2pt>[r]& X' \ar@<2pt>[r]^{h} & F }$
 une double flèche dans $C_{/F}$, prouvons qu'elle est
égalisée par une flèche $X'\rightarrow X''$ de $C_{/F}$. Or
$h:X'\rightarrow F$ est donné par un morphisme $h_i:X'\rightarrow
X_i$, et la condition $hf=hg$ signifie qu'il existe
$\alpha:i\rightarrow j$ dans $I$ tel que
$\varphi(\alpha)(h_if)=\varphi(\alpha)(h_ig)$, \ie quitte à
remplacer $h_i$ par $\varphi(\alpha)h_i$, on égalise $f$ et $g$.
Cela prouve que $C_{/F}$ est filtrante. Il est alors immédiat
qu'elle est essentiellement petite puisque les objets
$(X_i\rightarrow F)_{i\in \Ob I}$ forment une petite famille dans
$\Ob C_{/F}$ qui est cofinale.

(ii) $\Rightarrow$ (i). Posons $I=C_{/F}$, et considérons le
foncteur canonique (8.3.2.1)
$$
\varphi:I=C_{/F}\longrightarrow F{\quoi}.
$$
On sait que le préfaisceau représenté par cet ind-objet de $C$ est
$F$ (\hyperref[I.prop3.4]{3.4}), donc $F$ est ind-représentable par
définition (\hyperref[I.sec8.2.2]{8.2.2}).

(ii) $\Leftrightarrow$ (iii). En effet, (ii) $\Rightarrow$ (iii) car
un foncteur ind-représentable est exact à gauche (8.3.1), et (iii)
$\Rightarrow$ (ii) car on a signalé (\hyperref[I.prop8.3.1]{8.3.1}), que $F$
exact à gauche implique que $C_{/F}$ est filtrante, et\pageoriginale
on applique la définition \hyperref[I.def8.1.8]{8.1.8} pour conclure que cette
catégorie est essentiellement petite.

(ii) $\Leftrightarrow$ (iii bis) se prouve de même que (ii)
$\Leftrightarrow$ (iii).

Les équivalences (ii) $\Leftrightarrow$ (iv) et (iii)
$\Leftrightarrow$ (v) sont triviales, puisque pour $C$ équivalente à
une petite catégorie, $C_{/F}$ est évidemment équivalente également
à une petite catégorie et \textit{a fortiori} elle est
automatiquement essentiellement petite dès qu'elle est filtrante.
Cela achève la démonstration de 8.3.3.

\begin{enonce*}[remark]{Remarque 8.3.4}\etiquette[8.3.4]{I.rem8.3.4}
Soit $F:C'\rightarrow C''$ un foncteur entre $𝒰$-catégories.
Il apparaît que la notion d'exactitude à gauche
(\hyperref[I.cor2.3.2]{2.3.2}) ne présente
guère d'intérêt en pratique que si dans $C'$ les limites
projectives
finies sont représentables. Dans le cas où
$C''=(𝒰\text{-}\Ens)$, et où $C'$ est $𝒰$-petit,
écrivant $C'$ sous la forme $C^\circ$ (donc $C={C'}^\circ)$, il
convient de considérer que la « bonne notion » qui remplace,
dans le cas général, celle d'exactitude à gauche est celle de
ind-représentabilité de $F$ (considéré comme
préfaisceau sur
$C^\circ$; \ie celle de pro-représentabilité de $F$,
considéré comme
foncteur sur $C'$, \cf \hyperref[I.sec8.10]{8.10}). Cette notion coïncide
bien avec celle
d'exactitude à gauche lorsque dans $C'$ les limites projectives
finies sont représentables %%
 \textit{i.e.} dans $C$ les limites inductives finies sont
                   %%
représentables (\hyperref[I.them8.3.3]{8.3.3}). Lorsque l'on ne suppose plus $C'$
$𝒰$-petit ou tout au moins équivalente à une catégorie
$𝒰$-petite, alors les deux notions pour un foncteur $F$,
d'exactitude à gauche et de pro-représentabilité, ne coïncident plus
nécessairement, même si dans $C'$ les limites projectives finies
sont représentables (\cf \hyperref[I.exer8.12.9]{8.12.9}). En fait, il semble que
dans\pageoriginale ce cas, les foncteurs exacts à gauche non
ind-représentables doivent être considérés comme étant de nature
pathologique, les « bons » objets restant les foncteurs
ind-représentables; comparer \hyperref[I.cor8.13.3]{8.13.3}. Pour le cas des foncteurs
$F:C'\rightarrow C''$, avec $C'$ et $C''$ à nouveau des
$𝒰$-catégories quelconques, nous développerons plus bas
(8.11.5), plus généralement, une notion qui « améliore » celle
d'exactitude à gauche, (savoir celle d'un foncteur admettant un
foncteur pro-adjoint).
\end{enonce*}

\subsection{Ind-objets constants, essentiellement constants}\etiquette[8.4]{I.sec8.4}

Choisissons une catégorie finale $\sheaf{e}$, \ie telle que $\Ob
% 2013-5-24 problème de compilation avec \sheaf \XXX
\sheaf{e}$ et $F\ell$ $\sheaf{e}$ soient réduits à un élément (par
exemple, on peut prendre pour $\sheaf{e}$ la sous-catégorie pleine
de $(\Ens)$ formée par l'ensemble vide, si on veut un choix
canonique). C'est évidemment une catégorie filtrante et petite. Si
$X$ est un objet de $C$, on lui associe le ind-objet constant indexé
par $\sheaf{e}$, de valeur $X$, soit $c(X)$. Il est clair que pour
$X$ variable, on trouve ainsi un foncteur pleinement fidèle
$$
\displaylines{\rlap{$(8.4.1)$}\hfill c:C\longrightarrow
\Ind(C){\quoi},\hfill}
$$
d'ailleurs injectif sur les objets, et par lequel nous identifierons
$C$ à une sous-catégorie pleine de $\Ind(C)$. On notera que le
foncteur composé
$$
\displaylines{\rlap{$(8.4.2)$}\hfill
C\xrightarrow{c}\Ind(C)\xrightarrow{L}\widehat{C}\hfill}
$$
est le foncteur canonique $h$ (8.2.1.1).

\setcounter{subsubsection}{2}
\subsubsection{}\etiquette[8.4.3]{I.sec8.4.3}
Plus\pageoriginale généralement, soit $I$ une catégorie filtrante.
Le foncteur constant sur $I$ de valeur sur un objet $X$ de $C$ est
aussi appelé \emph{le ind-objet constant de valeur $X$ indexé par
$I$.} Il est clair, si $I$ est essentiellement petit, que cet
ind-objet ind-représente le foncteur $h(X)$ représenté par $X$, donc
cet ind-objet est isomorphe à $c(X)$.

\subsubsection{}\etiquette[8.4.4]{I.sec8.4.4}
Un ind-objet $𝒳$ de $C$ est dit \emph{essentiellement
constant} s'il est isomorphe (dans une catégorie
$\Ind_{𝒱}(C,𝒱')$, pour $𝒱$, $𝒱'$
convenables) à un ind-objet de la forme $c(X)$, $X\in \ob C$.
L'objet $X$, qui est alors déterminé à isomorphisme canonique près,
est appelé la \emph{valeur} du ind-objet essentiellement constant
envisagé. Donc par définition, le foncteur $c$ de (8.4.1) induit une
équivalence de $C$ avec la sous-catégorie pleine de $\Ind(C)$ formée
des ind-objets qui sont essentiellement constants, cette
sous-catégorie étant l'image essentielle du foncteur $c$.

Évidemment un ind-objet constant est essentiellement constant,
l'inverse n'étant vrai que dans le seul cas, trivial, où $C$ est
vide ou une catégorie ponctuelle.

\subsection{Limites inductives filtrantes dans $\Ind(C)$}\etiquette[8.5]{I.sec8.5}

\begin{enonce*}{Proposition 8.5.1}\etiquette[8.5.1]{I.prop8.5.1}
Soit $C$ une $𝒰$-catégorie. Dans $\Ind(C)$, les petites
limites inductives filtrantes sont représentables, et le foncteur
canonique $(8.2.4.8)$
$$
L:\Ind(C)\longrightarrow \widehat{C}
$$
y commute.
\end{enonce*}

Compte\pageoriginale tenu du fait que $L$ est pleinement fidèle, les
assertions de la proposition équivalent à la suivante:

\begin{enonce*}{Corollaire 8.5.2}\etiquette[8.5.2]{I.cor8.5.2}
Toute petite limite inductive filtrante (dans $\widehat{C}$) de
préfaisceaux ind-représentables est ind-représentable.
\end{enonce*}

Cela résulte facilement du critère \hyperref[I.them8.3.3]{8.3.3} (ii). Le
détail de la vérification est laissée au lecteur.

\setcounter{subsubsection}{2}
\subsubsection{}\etiquette[8.5.3]{I.sec8.5.3}
Le cas « tautologique » de limites inductives filtrantes dans
$\Ind(C)$ est celui où on part d'un ind-objet essentiellement petit
$𝒳=(X_i)_{i\in I}$ de $C$. On a alors, dans $\widehat{C}$
donc aussi dans $\Ind(C)$ (ou dans $\widehat{C}$):
$$
\displaylines{\rlap{$(8.5.3.1)$}\hfill 𝒳=\varinjlim_I
X_i{\quoi}.\hfill}
$$
On fera attention, en écrivant cette formule, qu'il ne s'agit pas
d'une limite inductive dans $C$, et que même lorsque cette dernière
existe, elle n'est pas isomorphe dans $\Ind(C)$ à $𝒳$: en
effet, le foncteur canonique (8.4.1) $c:C\rightarrow \Ind(C)$
\emph{ne commute pas en général aux limites inductives filtrantes.}
(\Cf \hyperref[I.exer8.5.5]{8.5.5} ci-dessous.)

Pour obvier à cette possibilité de confusion dans l'écriture de
(8.5.3.1), « certains auteurs » (= {P. Deligne}) préfèrent
l'écrire
$$
\displaylines{\rlap{$(8.5.3.2)$}\hfill
%𝒳=\lq\lq{\varinjlim_I}\textrm{ »\,}
𝒳=\textrm{« }{\varinjlim_I}\textrm{ »\,}X_i{\quoi},\hfill}
$$
le rôle des guillemets étant d'indiquer que la limite inductive est
prise dans une catégorie de Ind-objets $\Ind(C)$. Par extension, il
y aurait lieu alors de dénoter par $\textrm{« }{\lim}\textrm{ »}$ toute
opération de limite inductive dans $\Ind(C)$ (sans que les
composants du système inductif envisagé dans $\Ind(C)$ soient
nécessairement dans l'image, ou l'image essentielle, de
$c:C\rightarrow \Ind(C)$).

\subsubsection{}\etiquette[8.5.4]{I.sec8.5.4}
Pour\pageoriginale calculer la limite inductive ou projective d'un
foncteur
$$
\displaylines{\rlap{$(8.5.4.1)$}\hfill \varphi:J\longrightarrow
\Ind(C){\quoi},\hfill}
$$
où $J$ est maintenant une catégorie quelconque, il est souvent
commode d'expliciter $\varphi$ à l'aide d'un foncteur
$$
\displaylines{\rlap{$(8.5.4.2)$}\hfill \Phi:J\times I\longrightarrow
C{\quoi},\hfill}
$$
où $I$ est une catégorie filtrante essentiellement petite, ou de
préférence petite. Un tel $\Phi$ définit en effet un foncteur
$$
j\longmapsto (i\longmapsto \Phi (j,i)):J\longrightarrow
\Hhom(I,C){\quoi},
$$
d'où, en composant avec la flèche canonique (8.2.6.1)
$\Hhom(I,C)\rightarrow \Ind(C)$, un foncteur
$$
\displaylines{\rlap{$(8.5.4.3)$}\hfill j\longmapsto (i\longmapsto
\Phi(j,i)):J\longrightarrow \Ind(C){\quoi}.\hfill}
$$
On pourra dire que $\Phi$ est une \emph{expression indicielle de
$\varphi$, indexée par la catégorie (filtrante) $I$}, si le foncteur
correspondant (8.5.4.3) est isomorphe à $\varphi$. Nous étudierons
plus bas (\hyperref[I.sec8.8]{8.8}) des conditions générales d'existence
d'une expression
indicielle pour un foncteur $\varphi$ donné, et nous bornerons ici à
partir d'un foncteur $\varphi$ donné sous forme indicielle
(8.5.4.3), dans le cas où $J$ est une petite catégorie
filtrante,
pour indiquer dans ce cas le calcul de
$\textrm{«}{\varinjlim}\textrm{»\,}\varphi=\textrm{«}{\varinjlim_j}\textrm{»\,}\varphi(j)$.
La formule (8.5.3.2), et la formule d'associativité des limites
inductives (\hyperref[I.sec2.5.0]{2.5.0}) donne alors immédiatement le « \emph{calcul
tautologique} » de
$\textrm{«}{\varinjlim}\textrm{»\,}\varphi$:
$$
\displaylines{\rlap{$(8.5.4.4)$}\hfill
\textrm{«}{\varinjlim_J}\textrm{»\,}\varphi=\textrm{«}{\varinjlim_{J\times
I}}\textrm{»\,}\Phi{\quoi},\hfill}
$$
\ie\pageoriginale
$$
\displaylines{\rlap{$(8.5.4.5)$}\hfill
\textrm{«}{\varinjlim_j}\textrm{»\,}(i\longmapsto \Phi
(j,i))=\textrm{«}{\varinjlim_{j,i}}\textrm{»\,}\Phi(j,i){\quoi},\hfill}
$$
Ainsi, le système inductif cherché n'est autre que $\Phi$ lui-même
(la catégorie d'indices étant $J\times I$, qui est bien filtrante
puisque $J$ et $I$ le sont).

\begin{enonce*}[remark]{Exercice 8.5.5}\etiquette[8.5.5]{I.exer8.5.5}
Soit $C$ une $𝒰$-catégorie
\begin{enumerate}
\renewcommand{\theenumi}{\alph{enumi}}
\renewcommand{\labelenumi}{\theenumi)}
\item Prouver que les conditions suivantes sont équivalentes:

\begin{itemize}
\item[(i)] Dans $C$ les petites limites inductives sont représentables,
et le foncteur $c:C\rightarrow \Ind(C)$ y commute.

\item[(iii)] Le foncteur $c$ est une équivalence de catégories.

\item[(ii)] Dans $C$ les petites limites inductives sont
représentables, et pour tout $X\in \Ob C$, le foncteur $Y\mapsto
\Hom(X,Y)$ commute aux dites limites.

\item[(iv)] Les petites limites inductives filtrantes sont
représentables dans $C$, et le foncteur $\lim c:\Ind C\rightarrow C$
(\cf 8.7.1.5) est pleinement fidèle (ou encore, une
équivalence).
\end{itemize}

\item Si $C$ est une catégorie finie, prouver que les conditions
précédentes équivalent à la suivante: pour tout
projecteur dans $C$
(\hyperref[I.sec10.6]{10.6}), l'image est représentable dans $C$.
\end{enumerate}
\end{enonce*}

\subsection{Extension d'un foncteur aux
  ind-objets}\etiquette[8.6]{I.sec8.6}

\subsubsection{}\etiquette[8.6.1]{I.sec8.6.1}
Soit
$$
\displaylines{\rlap{$(8.6.1.1)$}\hfill f:C\longrightarrow C'\hfill}
$$
un foncteur entre $𝒰$-catégories. Pour tout ind-objet
$$
\varphi:I\longrightarrow C
$$
de\pageoriginale $C$ ($I$ une catégorie filtrante), le foncteur
composé $f\circ \varphi$ est un ind-objet de $C'$, noté aussi
$\Ind(f)(\varphi)$. En notations indicielles, si $\varphi$ est écrit
sous la forme $𝒳=(X_i)_{i\in I}$, on obtient
$$
\displaylines{\rlap{$(8.6.1.2)$}\hfill \Ind(f)(X_i)_{i\in
I}=(f(X_i))_{i\in I}{\quoi}.\hfill}
$$
Soit $𝒴=(Y_j)_{j\in J}$ un deuxième système inductif de $C$.
On trouve alors une application évidente, également notée
$u\rightarrow \Ind(f)(u)$:
$$
\displaylines{\rlap{$(8.6.1.3)$}\hfill
\Hom(𝒳,𝒴)=\varprojlim_i \varinjlim_j \Hom
(X_i,Y_j)\rightarrow\hfill\cr \hfill
\Hom(\Ind(f)(𝒳),\Ind(f)(𝒴))\simeq \varprojlim_i
\varinjlim_j \Hom(f(X_i),f(Y_j)){\quoi}.\hfill}
$$
On constate immédiatement que ces applications sont compatibles avec
la composition des morphismes de ind-objets, en se référant à la
formule non écrite de composition des morphismes de ind-objets
(\hyperref[I.sec8.2.5]{8.2.5}). On a ainsi obtenu, pour tout univers $𝒱$, un
foncteur
$$
\displaylines{\rlap{$(8.6.1.4)$}\hfill
\Ind(f):\Ind_{𝒱}(C,𝒰)\longrightarrow
\Ind_{𝒱}(C',𝒰)\hfill}
$$
(notations de (8.2.4.5)), et en particulier un foncteur
$$
\displaylines{\rlap{$(8.6.1.5)$}\hfill
\Ind(f):\Ind(C)\longrightarrow \Ind(C'){\quoi}.\hfill}
$$
Si on a un deuxième foncteur
$$
g:C'\longrightarrow C''{\quoi},
$$
on a évidemment une identité de foncteurs entre catégories $\Ind$
(ou $\Ind_{𝒱}$, au choix):
$$
\displaylines{\rlap{$(8.6.1.6)$}\hfill \Ind(gf)=\Ind(g)\circ
\Ind(f){\quoi},\hfill}
$$
et\pageoriginale pour $C'=C$, on a la sempiternelle formule
$$
\displaylines{\rlap{$(8.6.1.7)$}\hfill
\Ind(\id_C)=\id_{\Ind(C)}{\quoi}.\hfill}
$$
De façon imagée, on peut donc dire que la catégorie $\Ind(C)$
\emph{dépend fonctoriellement de $C$} (de façon covariante)
\footnote[1]{Pour une dépendance contravariante, sous certaines
conditions, \cf 8.11.}. Nous laissons au lecteur le soin
d'expliciter même une dépendance 2-\emph{fonctorielle}, en
définissant, pour tout couple de $𝒰$-catégories $C$, $C'$,
un foncteur canonique
$$
\displaylines{\rlap{$(8.6.1.8)$}\hfill
\Iind(f):\Hhom(C,C')\longrightarrow \Hhom(\Ind(C),
\Ind(C')){\quoi},\hfill}
$$
et en précisant la nature fonctorielle des isomorphismes identiques
(8.6.1.6) et (8.6.1.7).

\subsubsection{}\etiquette[8.6.2]{I.sec8.6.2}
Reprenons le foncteur (8.6.1.1), et considérons le foncteur
correspondant (5.1)
$$
\displaylines{\rlap{$(8.6.2.1)$}\hfill
f_{!}:\widehat{C}\longrightarrow \widehat{C'}{\quoi},\hfill}
$$
et le diagramme de foncteurs
\[\vcenter{
\xymatrix@C=3.25cm@R=1.25cm{
\Ind(C) \ar[r]^{\ind(f)} \ar[d]_{L_C}& \Ind(C')\ar[d]^{L_{C'}}\\
\hat{C} \ar[r]^{f_!}& \hat{C}' }}\leqno{(8.6.2.2)}
\]
où les flèches verticales désignent les foncteurs canoniques $L$ de
(8.2.4.8). Comme le foncteur $f_{!}$ « prolonge $f$ »
et commute aux limites inductives (5.4.3), et que les
foncteurs $L_C$ et $L_{C'}$ commutent aux petites limites inductives
filtrantes (\hyperref[I.prop8.5.1]{8.5.1}), il résulte de (8.5.3.2) que
pour tout ind-objet $(X_i)_{i\in I}$ de $C$, on a dans
$\widehat{C'}$ un isomorphisme canonique
$$
f_{!}(L_C((X_i)_{i\in I}))\simeq \varinjlim_i f_{!}L_C(X_i)\simeq
\varinjlim_i L_{C'}f(X_i)=L_{C'}(f(X_i)_{i\in I})
$$
\ie\pageoriginale un isomorphisme canonique
$$
f_{!}L_C(𝒳)\simeq L_{C'}(\Ind(f)(𝒳)){\quoi}.
$$
On laisse au lecteur le soin de vérifier que ce dernier est
fonctoriel, \ie qu'il correspond à un isomorphisme canonique
$$
\displaylines{\rlap{$(8.6.2.3)$}\hfill f_{!}L_C\simeq L_{C'}\circ
\Ind(f){\quoi}.\hfill}
$$
En d'autres termes, le \emph{carré (8.6.2.2) est commutatif à
isomorphisme canonique près.} Compte tenu du fait que $f_{!}$
commute aux petites limites inductives, et de 8.5.1, on conclut de
ceci:

\begin{enonce*}{Proposition 8.6.3}\etiquette[8.6.3]{I.prop8.6.3}
Soit $f:C\rightarrow C'$ un foncteur entre $𝒰$-catégories.
Alors le foncteur
$$
\Ind(f):\Ind(C)\longrightarrow \Ind(C')
$$
commute aux limites inductives filtrantes. De plus, il rend
commutatif le diagramme suivant
\[
\xymatrix@C=2.5cm@R=1.5cm{
C       \ar[r]^{f}       \ar[d]_{c_{C}}& C' \ar[d]^{c_{C'}}\\
\Ind(C) \ar[r]^{\Ind(f)} & \Ind(C')\quad , }
\]
où les flèches verticales $c_C$, $c_{C'}$ sont les foncteurs
canoniques $(8.4.1)$.
\end{enonce*}
La dernière assertion est triviale sur les définitions, et a
été mis
pour la commodité des références. Noter d'ailleurs que les
propriétés énoncées dans 8.6.3 \emph{caractérisent le foncteur}
$\Ind(f)$ à isomorphisme unique près,\pageoriginale comme étant
induit par le foncteur $f_{!}$ (8.6.2.1), comme il résulte de la
démonstration qu'on vient de donner de (8.6.2.3).

\begin{enonce*}{Proposition 8.6.4}\etiquette[8.6.4]{I.prop8.6.4}
Les notations sont celles de $8.6.3$.
\begin{enumerate}
\renewcommand{\theenumi}{\alph{enumi}}
\renewcommand{\labelenumi}{\theenumi)}
\item Pour que $\Ind(f)$ soit fidèle (\resp pleinement fidèle), il
faut et il suffit que $f$ le soit.

\item Pour que $\Ind(f)$ soit une équivalence de catégories, il
faut et il suffit que $f$ soit pleinement fidèle, et que tout objet
de $C'$ soit isomorphe à un facteur direct $(10.6)$ d'un objet dans
l'image de $f$.
\end{enumerate}
\end{enonce*}

\noindent \emph{Démonstration.}
\begin{enumerate}
\renewcommand{\theenumi}{\alph{enumi}}
\renewcommand{\labelenumi}{\theenumi)}
\item La nécessité résulte évidemment du fait que $f$ est
induit par $\Ind(f)$. Pour la suffisance, il suffit d'utiliser la
forme (8.6.1.3) de $\Ind(f)$ sur des ensembles
$\Hom(𝒳,𝒴)$, en se rappelant que les limites
inductives filtrantes d'ensembles, et les limites projectives
quelconques, transforment monomorphismes en monomorphismes,
isomorphismes en isomorphismes.

\item On peut supposer déjà $g$ donc $\Ind(f)$ pleinement
fidèle. Comme tout objet de $\Ind(C')$ est une petite limite
inductive filtrante d'objets de $C'$, la pleine fidélité de
$\Ind(f)$ implique que pour ce foncteur soit essentiellement
surjectif, il revient au même que tout objet $X'$ de $C'$ soit dans
l'image essentielle. Or si on a un isomorphisme
$$
X'\xrightarrow{\sim} \textrm{« }{\varinjlim}\textrm{ »\,}f(X_i){\quoi},
$$
cet isomorphisme se factorise par un des $f(X_i)$, ce qui montre que
$X'$ est un facteur direct de $f(X_i)$, ce qui prouve la nécessité
dans b). Pour la suffisance, utilisant la pleine fidélité de
$\Ind(f)$, elle résulte aussitôt du
\end{enumerate}

\begin{enonce*}{Corollaire 8.6.5}\etiquette[8.6.5]{I.cor8.6.5}
Dans\pageoriginale $\Ind(E)$ les images de projecteurs (\hyperref[I.sec10.6]{10.6}) sont
représentables, et le foncteur $\Ind(f):\Ind(C)\rightarrow \Ind(C')$
commute à la formation desdites images.
\end{enonce*}

Cela résulte en effet du fait que dans $\Ind(C)$ les petites limites
inductives filtrantes sont représentables (\hyperref[I.prop8.5.1]{8.5.1}) et que
$\Ind(f)$ y commute (\hyperref[I.prop8.6.3]{8.6.3}), compte tenu que l'image
d'un projecteur $p:X\rightarrow X$ s'interprète comme la limite
inductive d'un système inductif filtrant indexé par $\NN$
$$
X\xrightarrow{P}X\xrightarrow{P}X\longrightarrow\ldots
$$
ou au choix, comme limite inductive du foncteur qu'on devine sur la
catégorie filtrante $P$ ayant un seul objet, et une flèche non
identique $\Pi$ telle que $\Pi^2=\Pi$.

\subsection{}\etiquette[8.7]{I.sec8.7}
\emph{Le foncteur $\varinjlim_C:\Ind(C)\rightarrow C$.
Caractérisations universelles de la catégorie $\Ind(C)$.}

\subsubsection{}\etiquette[8.7.1]{I.sec8.7.1}
Reprenons une $𝒰$-catégorie $C$, et le foncteur canonique
$$
\displaylines{\rlap{$(8.7.1.1)$}\hfill c:C\longrightarrow
\Ind(C){\quoi}.\hfill}
$$
Soit $𝒳=(X_i)_{i\in I}$ un objet de $\Ind(C)$. Il résulte
immédiatement de la définition des limites inductives
représentables
(\hyperref[I.def2.1]{2.1}, \hyperref[I.sec2.1.1]{2.1.1}) que $\varinjlim_i X_i$ \emph{est
  représentable dans $C$
si et seulement si l'adjoint à gauche de $c$ est défini en
$𝒳$}, et que dans le cas $\varinjlim_i X_i$ (calculé dans
$C$) n'est autre que la valeur en $𝒳$ dudit adjoint à
gauche:
$$
\displaylines{\rlap{$(8.7.1.2)$}\hfill \Hom_C(\varinjlim_i X_i,
Y)\simeq \Hom_{\Ind(C)}((X_i)_{i\in I}, c(Y)){\quoi}.\hfill}
$$
Si\pageoriginale on désigne par
$$
\displaylines{\rlap{$(8.7.1.3)$}\hfill \Ind(C)'\subset
\Ind(C)\hfill}
$$
la sous-catégorie pleine de $\Ind(C)$ formée des ind-objets de $C$
qui admettent une limite inductive \emph{dans} $C$, il résulte de
l'observation précédente que cette sous-catégorie est
\emph{strictement pleine}, et que l'on a un foncteur canonique
$$
\displaylines{\rlap{$(8.7.1.4)$}\hfill
\textstyle{\varinjlim}_C:\Ind(C)'\longrightarrow C{\quoi},\hfill}
$$
dont la valeur en chaque objet $(X_i)_{i\in I}$ de $\Ind(C)'$ est sa
limite inductive dans $C$. Dans le cas particulier où dans $C$ les
petites limites inductives filtrantes sont représentables, on
obtient donc un foncteur naturel
$$
\displaylines{\rlap{$(8.7.1.5)$}\hfill
{\varinjlim}_C:\Ind(C)\longrightarrow C{\quoi}.\hfill}
$$

\setcounter{paragraph}{5}
\paragraph{}\etiquette[8.7.1.6]{I.sec8.7.1.6}
Bien entendu, on peut définir les foncteurs précédents
également sur
des catégories du type $\Ind_{𝒱}(C,𝒰)'$,
$\Ind_{𝒱}(C,𝒰)$, mais, compte tenu de $(8.2.4.6)$,
ils sont déjà déterminés (à isomorphisme unique
près) par la
connaissance des foncteurs précédents, correspondants au cas
$𝒱=𝒰$.

\paragraph{}\etiquette[8.7.1.7]{I.sec8.7.1.7}
De la construction précédente de $\varinjlim_C$ comme un foncteur
adjoint à gauche, il résulte immédiatement que ce foncteur
\emph{commute aux limites inductives quelconques}, et en particulier
aux petites limites inductives filtrantes (ces dernières étant
représentables dans $\Ind(C)$ (\hyperref[I.prop8.5.1]{8.5.1})). Notons aussi que
$\Ind(C)'$ contient toujours l'image essentielle de $c$ (formée des
ind-objets essentiellement constants), et que l'on a un isomorphisme
canonique fonctoriel en $X\in \ob \Ind(C)'$:
$$
\displaylines{\rlap{$(8.7.1.8)$}\hfill {\varinjlim}_C c(X)\simeq
X{\quoi},\hfill}
$$
\ie,\pageoriginale dans le cas favorable où $C$ est stable par
petites limites inductives, on a un isomorphisme canonique
$$
\displaylines{\rlap{$(8.7.1.9)$}\hfill {\varinjlim}_C \circ c\simeq
\id_C{\quoi}.\hfill}
$$

\subsubsection{}\etiquette[8.7.2]{I.sec8.7.2}

Considérons maintenant un foncteur
$$
\displaylines{\rlap{$(8.7.2.1)$}\hfill f:C\longrightarrow
E{\quoi},\hfill}
$$
où $E$ est une $𝒰$-catégorie où les petites limites
inductives sont représentables, de sorte qu'on a un foncteur
$$
{\varinjlim}_E:\Ind(E)\longrightarrow E{\quoi}.
$$
Comme on a défini également (8.6)
$$
\Ind(f):\Ind(C)\longrightarrow \Ind(E){\quoi},
$$
on peut considérer le composé
$$
\displaylines{\rlap{$(8.7.2.2)$}\hfill \bar{f}={\varinjlim}_E\circ
\Ind(f):\Ind(C)\longrightarrow E{\quoi},\hfill}
$$
qu'on appelle parfois le \emph{prolongement canonique de $f$ aux
ind-objets} (mais qu'on se gardera de confondre avec $\Ind(f) !$).
Comme composé de deux foncteurs commutant aux petites limites
inductives filtrantes (\hyperref[I.prop8.6.3]{8.6.3}, \hyperref[I.sec8.7.1.7]{8.7.1.7}), ce
foncteur lui-même commute aux petites limites inductives filtrantes.
De plus, il résulte de l'isomorphisme (8.7.1.9) appliqué à
$E$, et de (8.6.2.3), que $\bar{f}$ \emph{prolonge} $f$ à
isomorphisme près, \ie qu'on a un isomorphisme canonique
$$
\displaylines{\rlap{$(8.7.2.3)$}\hfill \bar{f}\circ c_C\simeq
f{\quoi}.\hfill}
$$
Il\pageoriginale est assez clair d'ailleurs, compte tenu du fait que
tout objet de $\Ind(C)$ est petite limite inductive filtrante
d'objets de $C$, que les deux propriétés précédentes
\emph{caractérisent} encore $\bar{f}$, à isomorphisme unique près.
On peut préciser ce point pour obtenir en même temps une
caractérisation universelle, à équivalence près, de $\Ind(C)$ parmi
les $𝒰$-catégories $E$ où les petites limites inductives
filtrantes sont représentables. Si $E$ et $F$ sont deux telles
$𝒰$-catégories, désignons par
$$
\displaylines{\rlap{$(8.7.2.4)$}\hfill \Hhom(E,F)'\subset
\Hhom(E,F)\hfill}
$$
la sous-catégorie pleine de $\Hhom(E,F)$ formée des foncteurs qui
commutent aux petites limites inductives filtrantes. On a alors:

\begin{enonce*}{Proposition 8.7.3}\etiquette[8.7.3]{I.prop8.7.3}
Soit $C$ une $𝒰$-catégorie, et utilisons la notation
ci-dessus (8.7.2.4). Alors le foncteur canonique
$$
c_C:C\longrightarrow \Ind(C)
$$
est $2$-universel parmi les foncteurs de source $C$ et de but une
$𝒰$-catégorie à petites limites inductives filtrantes
représentables. En d'autres termes, $\Ind(C)$ est une telle
catégorie (\hyperref[I.sec8.2.5]{8.2.5}, \hyperref[I.prop8.5.1]{8.5.1}), et si $E$ est une
telle catégorie, le foncteur
$$
\displaylines{\rlap{$(8.7.3.1)$}\hfill g\longmapsto g\circ
c_C:\Hhom(\Ind(C),E)'\longrightarrow \Hhom(C,E)\hfill}
$$
est une équivalence de catégories.
\end{enonce*}

Pour prouver ce dernier point, il suffit de vérifier que
l'application $f\mapsto \bar{f}$ définie par (8.7.2.2) peut se
préciser par un \emph{foncteur}
$$
\displaylines{\rlap{$(8.7.3.2)$}\hfill f\longmapsto
\bar{f}={\varinjlim}_E\cdot \Ind(f):\Hhom(C,E)\longrightarrow
\Hhom(\Ind(C),E)'{\quoi},\hfill}
$$
et\pageoriginale que ce dernier est aussi-inverse de (8.7.3.1). Le
détail de la vérification, essentiellement triviale, est laissé au
lecteur. On peut aussi invoquer 7.8 pour conclure d'abord que
(8.7.3.1) est pleinement fidèle, et (8.7.2.3) pour
conclure qu'il
est essentiellement surjectif.

\setcounter{subsubsection}{3}
\subsubsection{}\etiquette[8.7.4]{I.sec8.7.4}
Reprenons un foncteur entre $𝒰$-catégories
$$
\displaylines{\rlap{$(8.7.4.1)$}\hfill f:C\longrightarrow
E{\quoi},\hfill}
$$
où $E$ est une $𝒰$-catégorie où les petites limites
inductives filtrantes sont représentables, d'où un foncteur
$$
\displaylines{\rlap{$(8.7.4.2)$}\hfill \bar{f}:(X_i)_{i\in I}
\longmapsto \mathop{\lim}_i f(X_i):\Ind(C)\longrightarrow
E{\quoi}.\hfill}
$$
Nous nous proposons d'étudier des conditions sur $f$ qui assurent
que $\bar{f}$ est pleinement fidèle, \resp une équivalence de
catégories. Comme $f$ est isomorphe au composé $\bar{f}\circ c_C$,
et que $c_C:C\rightarrow \Ind(C)$ est pleinement fidèle, on voit que
pour que $\bar{f}$ soit pleinement fidèle, il \emph{faut} que $f$ le
soit. Quitte à remplacer $C$ par son image essentielle dans $E$, on
voit donc qu'on ne perd pas en généralité, essentiellement, en
supposant que $f$ est le foncteur d'inclusion d'une sous-catégorie
$C$ de $E$, ce que nous supposerons par la suite, pour simplifier
les notations.

\begin{enonce*}{Proposition 8.7.5}\etiquette[8.7.5]{I.prop8.7.5}
Les notations sont celles de \hyperref[I.sec8.7.4]{8.7.4}.
\begin{enumerate}
\renewcommand{\theenumi}{\alph{enumi}}
\renewcommand{\labelenumi}{\theenumi)}
\item Pour que le foncteur $\bar{f}$ soit pleinement fidèle, il faut
et il suffit que l'on ait:
\begin{itemize}
\item[(i)] Tout objet $X$ de $C$ satisfait à la condition suivante:

{\rm(PF)} Pour tout petit système inductif filtrant $(Y_i)_{i\in I}$
dans $C$, de\pageoriginale limite inductive $\varinjlim_i Y_i$ dans
$E$, l'application canonique
$$
\displaylines{\rlap{$(8.7.5.1)$}\hfill \varinjlim_i \Hom
(X,Y_i)\longrightarrow \Hom(X,\varinjlim_i Y_i)\hfill}
$$
est bijective.
\end{itemize}

\item Pour que le foncteur $\bar{f}$ soit une équivalente de
catégories, il faut et il suffit que $C$ satisfasse la condition
{\rm(i)}, et les deux conditions suivantes:
\begin{itemize}
\item[(ii)] $C$ est une sous-catégorie de $E$ génératrice par
épimorphismes stricts (\hyperref[I.def7.1]{7.1}).

\item[\rm(iii)] Pour tout objet $X$ de $E$, la sous-catégorie pleine $C_{/X}$
de $E_{/X}$, formée des objets $X'$ au-dessus de $X$ de source dans
$\Ob C$, est filtrante et essentiellement petite.
\end{itemize}
\end{enumerate}

La condition {\rm(iii)} est vérifiée en particulier si $C$ est
équivalente à une petite catégorie, et si les limites inductives
finies dans $C$ sont représentables et le foncteur d'inclusion
$f:C\rightarrow E$ $y$ commute (\ie pour toute catégorie finie $J$
et tout foncteur $\varphi:J\rightarrow C$ la limite inductive de
$f\varphi$ est représentable et est isomorphe dans $E$ à un objet
de $C$).
\end{enonce*}


\noindent \emph{Démonstration}.
\begin{enumerate}
\renewcommand{\theenumi}{\alph{enumi}}
\renewcommand{\labelenumi}{\theenumi)}
\item Avec les notations de la condition (i), si $𝒴$ désigne le
ind-objet $(Y_i)$, $𝒳$ le ind-objet constant défini par $X$,
alors (8.7.5.1) n'est autre que l'application canonique
$$
\Hom(𝒳, 𝒴)\longrightarrow \Hom
(\bar{f}(𝒳),\bar{f}(𝒴)){\quoi},
$$
donc si $\bar{f}$ est pleinement fidèle cette application est bien
bijective. Réciproquement, si $𝒳=(X_j)_{j\in J}$ et
$𝒴=(Y_i)_{i\in I}$ sont deux ind-objets quelconques de $E$,
alors l'application $u:\Hom (𝒳, 𝒴)\rightarrow \Hom
(\bar{f}(𝒳), \bar{f}(𝒴))$ est\pageoriginale la
limite projective sur $j$ des applications $u_j:\Hom(𝒳_j,
𝒴)\rightarrow \Hom(\bar{f}(𝒳_j),𝒴)$, où
$𝒳_j$ est le ind-objet constant de valeur $X_j$; donc $u$
est bijective si les $u_j$ le sont, et (i) implique donc que
$\bar{f}$ est pleinement fidèle.

\item Supposons (i), (ii), (iii) vérifiées, et prouvons que
$\bar{f}$ est une équivalence. Il reste à prouver qu'il est
essentiellement surjectif, donc que tout $X\in \ob E$ est dans
l'image essentielle. Or l'hypothèse (ii) signifie que
$\varinjlim_{C_{/X}} Z\rightarrow X$ est un isomorphisme, et (iii)
que la catégorie $C_{/X}$ est filtrante et essentiellement petite,
donc $X$ est l'image du ind-objet de $E$ défini par le foncteur
naturel $C_{/X}\rightarrow E$. Inversement, supposons que $\bar{f}$
est une équivalence, donc en vertu de a), il reste à vérifier (ii)
et (iii). On peut supposer que $E=\Ind(C)$, $C$ étant identifiée à
la sous-catégorie $c(C)$ de $E$. Mais on sait (8.3.3 (ii)) que pour
tout $X\in \Ind(C)$, identifié si on le désire au foncteur $F$ sur
$C^\circ$ qu'il ind-représente, $C_{/X}=C_{/F}$ est une catégorie
filtrante essentiellement petite (ce qui prouve (iii)) et que la
limite inductive dans $\widehat{C}$ du foncteur $C_{/F}\rightarrow
\widehat{C}$ est $F$. A fortiori, il en est ainsi dans la
sous-catégorie pleine $E$ de $\widehat{C}$, puisque $F=X$ est dans
$E$, ce qui prouve (ii). Reste à prouver la dernière assertion de
b). Or l'hypothèse faite sur $C$ implique évidemment que dans la
catégorie $C_{/X}$, les limites inductives finies sont
représentables, \textit{a fortiori} la catégorie $C_{/X}$ est
filtrante.
\end{enumerate}

\begin{enonce*}[remark]{Remarque 8.7.5.2}\etiquette[8.7.5.2]{I.rem8.7.5.2}
Le\pageoriginale raisonnement qu'on vient de faire montre plus
généralement que lorsque la condition (i) est vérifiée, alors
l'image essentielle du foncteur pleinement fidèle $\bar{f}$
(8.7.4.2) est formée des $X\in \ob E$ tels que la
catégorie $C_{/X}$
soit filtrante et essentiellement petite (condition automatiquement
satisfaite si $C$ satisfait les conditions énoncées à la fin de b)),
\emph{et} que la limite inductive du foncteur canonique
$C_{/X}\rightarrow E$ soit $X$.
\end{enonce*}

\begin{enonce*}{Corollaire 8.7.6}\etiquette[8.7.6]{I.cor8.7.6}
Soit $E_{PF}$ la sous-catégorie pleine de $E$ formée des objets
satisfaisant la condition $(PF)$ de \hyperref[I.prop8.7.5]{8.7.5} {\rm a)}.
Alors pour tout foncteur $J\rightarrow E_{PF}$, $J$ une catégorie
finie, qui admet une limite inductive $X$ dans $E$, on a $X\in \ob
E_{PF}$. Le foncteur canonique $(X_i)_{i\in I}\rightarrow
\displaystyle{\mathop{\varinjlim}_i} {}_EX_i$ déduit de l'inclusion
$g:E_{PF}\hookrightarrow E$
$$
\displaylines{\rlap{$(8.7.6.1)$}\hfill
\bar{g}:\Ind(E_{PF})\longrightarrow E\hfill}
$$
est pleinement fidèle, et si dans $E$ les limites inductives finies
sont représentables et si $E_{PF}$ est équivalente à une petite
catégorie, alors l'image essentielle du foncteur $\bar{f}$ est
formée des objets $X$ de $E$ tels que le morphisme canonique
$\varinjlim_{(E_{PF})_{/X}}Y\rightarrow X$ soit un isomorphisme. Si
on suppose de plus que la sous-catégorie $E_{PF}$ de $E$ est
génératrice par épimorphismes stricts (\hyperref[I.def7.1]{7.1}),
alors le foncteur
(8.7.6.1) est une équivalence de catégories.
\end{enonce*}

Tous les faits sont évidents, dans l'ordre où ils sont donnés,
compte tenu de \hyperref[I.prop8.7.5]{8.7.5}, en utilisant \hyperref[I.prop2.8]{2.8} pour la
première assertion.

\begin{enonce*}{Corollaire 8.7.7}\etiquette[8.7.7]{I.cor8.7.7}
Supposons\pageoriginale que dans $E$ les limites inductives finies
soient représentables. Soit $C$ une sous-catégorie pleine de $E$,
équivalente à une petite catégorie, et considérons le foncteur
$(X_i)_{i\in I}\rightarrow
{\displaystyle{\mathop{\varinjlim}_i}{}_E} X_i$:
$$
\displaylines{\rlap{$(8.7.7.1)$}\hfill \bar{f}:\Ind(C)\rightarrow
E{\quoi}.\hfill}
$$
Soit $E_{PF}$ comme dans $8.7.6$. Les conditions suivantes sont
équivalentes:
\begin{itemize}
\item[(i)] Le foncteur $\bar{f}:\Ind(C)\rightarrow E$ est une
équivalence.

\item[(ii)] La catégorie $C$ dans $E$ est génératrice par
épimorphismes stricts, contenue dans $E_{PF}$, et tout objet de
$E_{PF}$ est isomorphe dans $E$ (ou dans $E_{PF}$, cela revient au
même (\hyperref[I.cor8.7.6]{8.7.6})) à un facteur direct d'un objet de $C$.

Lorsque la sous-catégorie $C$ de $E$ est stable par facteurs directs
(\hyperref[I.sec10.6]{10.6}), les conditions précédentes équivalent encore aux
suivantes:

\item[(ii bis)] $C=E_{PF}$, et $C$ est génératrice dans $E$ par
épimorphisme stricts.

\item[(iii)] La sous-catégorie $C$ de $E$ est génératrice par
épimorphismes stricts, contenue dans $E_{PF}$, et les limites
inductives finies $y$ sont représentables.

Lorsque de plus dans $E$ les limites inductives finies sont
représentables, ces conditions équivalent aussi à

\item[(iii bis)] La sous-catégorie $C$ de $E$ est génératrice
par épimorphismes stricts, contenue dans $E_{PF}$, et stable dans
$E$ par limites inductives finies (ou ce qui revient au même, par
sommes finies, et par conoyaux de doubles flèches).
\end{itemize}
\end{enonce*}

On\pageoriginale sait déjà (\hyperref[I.prop8.7.5]{8.7.5}) que la condition (i)
implique que $C\subset E_{PF}$, et que $C$ est génératrice par
épimorphismes stricts; prouvons aussi qu'alors tout objet $X$ de
$E_{PF}$ est isomorphe à un facteur direct d'un objet de $C$. En
effet on sait que $X$ est une limite inductive filtrante
$\varinjlim_i X_i$ dans $E$ d'objets de $C$, et par définition de
$E_{PF}$, on voit que pour $i$ convenable l'homomorphisme canonique
$X_i\rightarrow X$ admet un inverse à gauche, de sorte que $X$ est
bien facteur direct de l'objet $X_i$ de $C$. Cela prouve que (i)
$\Rightarrow$ (ii) (sans hypothèse sur $C$ ni $E$, d'ailleurs).
Prouvons (ii) $\Rightarrow$ (i). Comme $C$ est équivalente à une
petite catégorie et tout objet de $E_{PF}$ est facteur direct d'un
objet de $C$, on voit que $E_{PF}$ est également équivalente à une
petite catégorie, donc en vertu de \hyperref[I.cor8.7.6]{8.7.6} le foncteur
(8.7.6.1) est une équivalence de catégories, et on conclut
grâce à \hyperref[I.prop8.6.4]{8.6.4} b) appliqué à l'inclusion
$E\hookrightarrow E_{PF}$.

Lorsque tout facteur direct dans $E$ d'un objet de $C$ est dans $C$,
il est clair que (ii) $\Leftrightarrow$ (ii bis). D'autre part (ii
bis) implique (iii) \resp (iii bis) en vertu de 8.7.6. Enfin, (iii)
(et \textit{a fortiori} (iii bis)) implique (i) en vertu de
\hyperref[I.prop8.7.5]{8.7.5}. Cela achève la démonstration.

\begin{enonce*}[remark]{Exercice 8.7.8}\etiquette[8.7.8]{I.exer8.7.8}
(\emph{Enveloppes de Karoubi}). Soient $C$ une
$𝒰$-catégorie, $E=\Ind(C)$, et $E_{PF}\supset C$ %%
la sous-catégorie pleine de $E$ définie dans 8.7.6.
\begin{enumerate}
\renewcommand{\theenumi}{\alph{enumi}}
\renewcommand{\labelenumi}{\theenumi)}
\item Prouver que dans $E_{PF}$ les images de projecteurs sont
représentables, et que tout objet de $E_{PF}$ est facteur direct
d'un objet de $C$.

\item Appelons \emph{karoubienne} une catégorie $F$ dans laquelle
les images de projecteurs sont représentables. Soit alors
$$
\varphi:C\longrightarrow K
$$
un\pageoriginale foncteur pleinement fidèle, tel que $K$ soit
karoubienne et que tout objet de $K$ soit facteur direct d'un objet
de l'image de $C$. Prouver que $\varphi$ est 2-\emph{universel}
parmi les foncteurs de $C$ à valeurs dans des catégories
karoubiennes, de façon précise: pour toute catégorie karoubienne
$F$, le foncteur
$$
g\longmapsto g\circ \varphi :\Hhom(K,F)\longrightarrow \Hhom(K,F)
$$
est une équivalence de catégories. (On utilisera le fait que tout
foncteur commute aux images de projecteurs.) La catégorie $K$ munie
de $\varphi$, déterminée à équivalence près (elle-même déterminée à
isomorphisme unique près) par les propriétés précédentes, s'appelle
l'\emph{enveloppe de Karoubi} $\widetilde{C}$ de $C$. (Comparer
aussi IV \hyperref[I.cor7.5]{7.5}.)

\item Déduire de a) et b) que $\Ind(C)_{PF}$, munie du foncteur
$C\rightarrow \Ind(C)_{PF}$ induit par $c:C\rightarrow \Ind(C)$,
fait de $\Ind(C)_{PF}$ une enveloppe de Karoubi de $C$.

\item Montrer que tout foncteur $f:C\rightarrow C'$ se prolonge de
façon essentiellement unique en un foncteur
$\widetilde{f}:\widetilde{C}\rightarrow\widetilde{C'}$ des
enveloppes de Karoubi, et que si on prend ces enveloppes de Karoubi
comme dans c), $\widetilde{f}$ est le foncteur induit par
$\Ind(f):\Ind(C)\rightarrow \Ind(C')$. Montrer que les conditions de
\hyperref[I.sec8.5.4]{8.5.4} %%
b) %%
équivalent encore à la suivante: $\widetilde{f}$ est une équivalence
de catégories.
\end{enumerate}
\end{enonce*}

\begin{enonce*}[remark]{Exercice 8.7.9}\etiquette[8.7.9]{I.exer8.7.9}
Soit $E$ une $𝒰$-catégorie.
\begin{enumerate}
\renewcommand{\theenumi}{\alph{enumi}}
\renewcommand{\labelenumi}{\theenumi)}
\item Montrer que les conditions sont équivalentes:

\begin{itemize}
\item[(i)] $E$ est équivalente à une catégorie de la forme
$\Ind(C)$, avec $C$ une $𝒰$-catégorie.

\item[(i bis)] Comme\pageoriginale (i), avec de plus $C$ karoubienne
(\hyperref[I.exer8.7.8]{8.7.8} b)).

\item[(ii)] La sous-catégorie $E_{PF}$ de $E$ (\hyperref[I.cor8.7.6]{8.7.6}) est
génératrice par épimorphismes stricts, et pour tout objet $X$ de
$E$, $(E_{PF})_{/X}$ est une catégorie filtrante essentiellement
petite.
\end{itemize}

Montrer que si $C$ est une $𝒰$-catégorie \emph{karoubienne}
telle que $E$ soit équivalente à $\Ind(C)$, alors $C$ est
équivalente à $E_{PF}$ (par une équivalence déterminée à
isomorphisme unique près). Établir une 2-équivalence entre la
2-catégorie formée des catégories $E$ satisfaisant aux conditions
précédentes et les foncteurs entre icelles \emph{commutant aux
petites limites inductives filtrantes,} et la 2-catégorie formée des
$𝒰$-catégories \emph{karoubiennes} et les foncteurs
quelconques entre icelles.

\item Montrer que les conditions suivantes sont équivalentes:

\begin{itemize}
\item[(i)] $E$ est équivalente à une catégorie de la forme
$\Ind(C)$, où $C$ est une $𝒰$-catégorie où les limites
inductives finies sont représentables.

\item[(i bis)] Comme (i), mais avec de plus $C$ karoubienne.

\item[(ii)] Les limites inductives finies dans $E$ sont
représentables, la sous-catégorie $E_{PF}$ de $E$ est génératrice
par épimorphismes stricts, et pour tout objet $X$ de $E$, il existe
une petite partie $I$ de $\Ob E_{PF/X}$ telle que tout objet de
$E_{PF/X}$ soit majoré par un objet de $I$. (Utiliser 8.9.5 b).)
\end{itemize}
\end{enumerate}
\end{enonce*}

\subsection{Représentation indicielle d'un foncteur $J\rightarrow
\Ind(C)$}\etiquette[8.8]{I.sec8.8}\pageoriginale

\subsubsection{}\etiquette[8.8.1]{I.sec8.8.1}
Soit
$$
\varphi:J\longrightarrow \Ind(C)\text{ \ie } \varphi\in \Ob
\Hhom(J,\Ind(C))
$$
un foncteur, $C$ étant une $𝒰$-catégorie. Rappelons (8.5.4)
qu'une \emph{représentation indicielle de} $\varphi$ est par
définition un ind-objet de $\Hhom(J,C)$
$$
\psi:I\longrightarrow \Hhom(J,C){\quoi},
$$
I catégorie filtrante essentiellement petite, dont la limite
inductive dans $\Hhom(J,\Ind(C))$ soit isomorphe à $\varphi$ par un
isomorphisme donné. Donc $\varphi$ admet une représentation
indicielle si et seulement si il est isomorphe à une limite
inductive filtrante essentiellement petite dans $\Hhom(J,\Ind(C))$
d'objets de la sous-catégorie pleine $\Hhom(J,C)$. Nous dirons que
la catégorie $J$ est \emph{admissible pour} $C$ si tout foncteur
$\varphi:J\rightarrow \Ind(C)$ admet une représentation indicielle.

Comme dans $\Ind(C)$ les petites limites inductives filtrantes sont
représentables (\hyperref[I.prop8.5.1]{8.5.1}), il en est de même dans
$\Hhom(J,\Ind(C))$, et elles se calculent « argument par argument ».
Par suite, le foncteur d'inclusion
$$
\displaylines{\rlap{$(8.8.1.1)$}\hfill \Hhom(J,C)\hookrightarrow
\Hhom(J,\Ind(C))\hfill}
$$
se prolonge canoniquement en un foncteur
$$
\displaylines{\rlap{$(8.8.1.2)$}\hfill
\Ind(\Hhom(J,C))\longrightarrow \Hhom(J,\Ind(C))\hfill}
$$
(8.7.2). Les éléments de l'image essentielle de ce foncteur sont
       alors précisément les $\varphi$ admettant une
       représentation indicielle; donc dire que $J$ est admissible
       pour $C$ signifie aussi que le foncteur (8.8.1.2) est
       essentiellement surjectif. Signalons à ce propos:

\begin{enonce*}{Proposition 8.8.2}\etiquette[8.8.2]{I.prop8.8.2}
Si\pageoriginale la catégorie $J$ est équivalente à une
catégorie
finie, le foncteur canonique $(8.8.1.2)$ est pleinement fidèle; donc
$J$ est admissible pour $C$ si et seulement si c'est une équivalence
de catégories.
\end{enonce*}

En vertu de \hyperref[I.prop8.7.5]{8.7.5} a) tout revient à prouver que pour
tout petit système inductif filtrant $(\varphi_i)_{i\in I}$ dans
$\Hhom(J,C)$ et tout objet $\varphi$ de $\Hhom(C,J)$, l'application
canonique
$$
\displaylines{\rlap{$(8.8.2.1)$}\hfill \varinjlim_i \Hom(\varphi,
\varphi_i)\longrightarrow \Hom(\varphi, \varinjlim_i
\varphi_i)\hfill}
$$
est bijective, où $\varinjlim_i \varphi_i$ désigne la limite
inductive prise dans $\Hom(J,\Ind(C))$. Or, si $\varphi$, $\psi$
sont deux foncteurs $J\rightarrow C$, on a un diagramme exact
d'ensembles, fonctoriel en $\varphi$ et $\psi$:
$$
\Hom(\varphi,\psi)\longrightarrow \prod_{X\in \Ob J}
\Hom(\varphi(X),\psi(X))\rightrightarrows \prod_{\substack{f\in
F\ell J\\ f:X\rightarrow Y}} \Hom(\varphi(X), \psi (Y)){\quoi}.
$$
Comme les limites inductives filtrantes commutent aux noyaux de
doubles flèches et aux produits finis, la bijectivité de
(8.8.2.1)
s'ensuit quand $J$ est finie. Le cas où $J$ est équivalente à une
catégorie finie se ramène aussitôt au cas précédent.

\begin{enonce*}{Proposition 8.8.3}\etiquette[8.8.3]{I.prop8.8.3}
Soit $\varphi:J\rightarrow \Ind(C)$ un foncteur, avec $J$
équivalente à une petite catégorie. Pour que $\varphi$ admette une
représentation indicielle, il suffit qu'il satisfasse aux deux
conditions suivantes, et cela est également nécessaire lorsque $J$
est équivalente à une catégorie finie:
\begin{enumerate}
\renewcommand{\theenumi}{\alph{enumi}}
\renewcommand{\labelenumi}{\theenumi)}
\item La catégorie $\Hhom(J,C)_{/\varphi}$ (formée des
flèches de $\Hhom(J,\Ind(C))$ de but $\varphi$ et de source dans
$\Hhom(J,C)$) est filtrante.

\item Pour\pageoriginale tout objet $j$ de $J$, le foncteur
$$
\displaylines{\rlap{$(8.8.3.1)$}\hfill \psi\longmapsto
\psi(j):\Hhom(J,C)_{/\varphi}\longrightarrow C_{/\varphi(j)}\hfill}
$$
\end{enumerate}
est cofinal. \ie satisfaisant aux conditions {\rm %%
F 1), F 2)} de $8.1.3$.
\end{enonce*}

\noindent \emph{Démonstration}. Supposons a) et b) vérifiées,
prouvons que $\varphi$ admet une représentation indicielle. On peut
supposer évidemment $J$ petite. Soit
$$
\displaylines{\rlap{$(8.8.3.2)$}\hfill
I=\Hhom(J,C)_{/\varphi}{\quoi},\hfill}
$$
qui est une catégorie filtrante par hypothèse, et supposons d'abord
$I$ essentiellement petite, de sorte que « l'inclusion » de $I$
dans $\Hhom(J,C)$ définit un ind-objet de $\Hhom(J,C)$, $i\mapsto
\psi_i$. De plus, on a un homomorphisme canonique
$$
\displaylines{\rlap{$(8.8.3.3)$}\hfill \varinjlim_i
\psi_i\longrightarrow \varphi{\quoi},\hfill}
$$
donné argument par argument par l'homomorphisme
$$
\varinjlim_I \psi_i (j)\longrightarrow
\varphi(j)=\varinjlim_{C_{/\varphi(j)}}X
$$
déduit du foncteur (8.8.3.1). Comme %%
ce dernier est cofinal, on en conclut que (8.8.3.3) est un
isomorphisme, d'où la conclusion. Dans le cas où on ne suppose pas
$I$ essentiellement petite, il suffit de construire une
sous-catégorie pleine essentiellement petite $I'$, telle que
(8.8.3.3) reste un isomorphisme en prenant $\varinjlim_{I'}$ au lieu
de $\varinjlim_I$. Pour ceci, il suffit que les foncteurs composés
$$
I'\longrightarrow I\longrightarrow C_{/\varphi(j)}
$$
induits\pageoriginale par les foncteurs (8.8.3.1) soient tous
cofinaux. On conclut alors par le lemme suivant, dont la
démonstration est immédiate grâce au critère (8.1.3 b)), et est
laissée au lecteur:

\begin{enonce*}{Lemme 8.8.3.4}\etiquette[8.8.3.4]{I.lem8.8.3.4}
Soient $I$ une $𝒰$-catégorie filtrante, et
$$
f_j:I\longrightarrow I_j\quad (j\in J)
$$
une petite famille de foncteur cofinaux de $I$ dans des catégories
filtrantes essentiellement petites. Alors il existe une
sous-catégorie pleine $I'$ de $I$ qui est filtrante et
essentiellement petite, et telle que les foncteurs induits par les
$f_j$ soient cofinaux.
\end{enonce*}

Prouvons enfin la nécessité dans \hyperref[I.prop8.8.3]{8.8.3} lorsqu'on suppose
$J$ équivalente à une catégorie finie. Une représentation indicielle
de $\varphi$ à l'aide d'une catégorie d'indices filtrante
essentiellement petite $I$ définit un foncteur $\psi$
\[
\xymatrix@C=1.25cm@R=2cm{
 I \ar[rr]^(.4){\psi} \ar@{-->}[rrd]^{\psi_j}&&  \Hhom(J,C)_{/\varphi} \ar@{-->}[d] &\\
                        &     &   C_{/\varphi(j)}\quad ,\\
}\tag{8.8.3.5}
\]
et il suffit de prouver que $\psi$ et chacun des foncteurs $\psi_j$
qui s'en déduit sont cofinaux: en vertu de \hyperref[I.prop8.1.3]{8.1.3} b) il
s'ensuivra bien que $\Hhom(J,C)_{/\varphi}$ est filtrante, et par
définition \hyperref[I.def8.1.1]{8.1.1} que la flèche verticale de (8.8.3.5) est
également cofinale. En vertu de 8.8.2 on peut identifier $\varphi$ à
un objet de $\Ind(E)$, où $E=\Hhom(J,C)$, et le fait que le foncteur
$$
\psi:I\longrightarrow E_{/\varphi}{\quoi},
$$
déduit\pageoriginale du ind-objet $\varphi$ de $E$ indexé par $I$,
est cofinal est un fait général, qui se vérifie immédiatement à
l'aide des critères (F 1) et (F 2) de \hyperref[I.prop8.1.3]{8.1.3}. La même
raison (où $E$, $\varphi$ sont remplacés par $C$, $\varphi(j)$)
montre que $\psi_j$ est cofinal, ce qui achève la démonstration.

\begin{enonce*}[remark]{Remarques 8.8.4}\etiquette[8.8.4]{I.rem8.8.4}
\begin{enumerate}
\renewcommand{\theenumi}{\alph{enumi}}
\renewcommand{\labelenumi}{\theenumi)}
\item Dans le cas où $J$ est équivalente à une catégorie
finie, si $\varphi$ admet une représentation indicielle, on a vu que
(8.8.3.5) est un foncteur cofinal, ce qui implique que la catégorie
$\Hhom(J,C)_{/\varphi}$ est elle-même essentiellement petite.

\item Supposons que dans $C$ les limites inductives finies soient
représentables. Alors il en est de même dans $E=\Hhom(J,C)$, et
celles-ci se calculent argument par argument, et ce sont également
des limites inductives dans $\Hhom(J,\widehat{C})$ et \textit{a
fortiori} dans $\Hhom(J,\Ind(C))$. Il s'ensuit aussitôt que la
condition a) de \hyperref[I.prop8.8.3]{8.8.3} est alors automatiquement
satisfaite, et tout revient à regarder la condition b). J'ignore si
elle est automatiquement satisfaite lorsque de plus $J$ est supposée
petite.
\end{enumerate}
\end{enonce*}
Nous en arrivons au résultat principal du présent numéro:

\begin{enonce*}{Proposition 8.8.5}\etiquette[8.8.5]{I.prop8.8.5}
Soit $J$ une catégorie équivalente à une catégorie finie, et
supposons de plus $J$ rigide \ie que pour tout $j\in \Ob J$, tout
endomorphisme de $j$ soit l'identité. Alors $J$ est admissible pour
$C$ (quelle que soit la $𝒰$-catégorie $C$), \ie tout
foncteur $J\rightarrow \Ind(C)$ admet une représentation indicielle.
\end{enonce*}

Quitte \pageoriginale %%
à remplacer $J$ par une catégorie équivalente, nous pouvons supposer
que $J$ est \emph{réduite} \ie que deux objets isomorphes de $J$
sont identiques. Alors $J$ est même finie. Nous procédons par
récurrence sur $\text{card}\Ob J$, le cas où ce nombre est $\leq 0$
étant trivial; nous le supposerons donc $\geq 1$. Soit $j_\circ$ un
objet \emph{maximal} de $J$, \ie tel que pour toute flèche
$j_\circ\rightarrow j$, il existe une flèche $j\rightarrow j_\circ$;
compte tenu du fait que $J$ est rigide, cela implique que
$j_\circ\rightarrow j$ est un isomorphisme, et comme $J$ est
réduite, cela implique $j_\circ=j$. Soit donc $J'$ la sous-catégorie
pleine déduite de $J$ en lui enlevant l'objet $j_\circ$, et $J''$ la
sous-catégorie pleine réduite à l'objet $j_\circ$. La proposition
résulte alors du

\begin{enonce*}{Lemme 8.8.5.1}\etiquette[8.8.5.1]{I.lem8.8.5.1}
Soient $J$ une catégorie finie, $J'$ et $J''$ deux sous-catégories
pleine telles que $\Ob J=\Ob J'\cup \Ob J''$, et que pour tout
$j'\in \Ob J'$ et $J''\in \Ob J''$, on ait
$\Hom(j'',j')=∅$. Si $J'$ et $J''$ sont admissibles pour
$C$, il en est de même de $J$.
\end{enonce*}

Tout revient à prouver les critères a) et b) de \hyperref[I.prop8.8.3]{8.8.3},
ce qui revient à faire six vérifications élémentaires, savoir les
deux dernières conditions des catégories filtrantes pour
$\Hhom(J,C)_{/\varphi}$, et les conditions
F 1) et
F 2) pour les foncteurs (8.8.3.1), dans le cas $j\in \Ob J'$ et
$j\in \Ob J''$ successivement; ces dernières conditions impliquant
d'ailleurs que $\Hhom(J,C)_{/\varphi}$ est non vide. La vérification
assez fastidieuse n'offre pas de difficulté, et est laissée au
lecteur. (Le rédacteur a vainement cherché une démonstration
élégante qui court-circuiterait ces déplaisantes
vérifications.)

\begin{enonce*}[remark]{Exercice 8.8.6}\etiquette[8.8.6]{I.exer8.8.6}
\begin{enumerate}
\renewcommand{\theenumi}{\alph{enumi}}
\renewcommand{\labelenumi}{\theenumi)}
\item Soient\pageoriginale $J$, $J'$ deux catégories admissibles
pour $C$, l'une d'elles étant finie. Prouver que $J\times J'$ est
admissible pour $C$. (Utiliser 8.8.2.)

\item Prouver qu'une petite catégorie discrète est admissible pour
toute catégorie $C$.

\item Soit $J$ une catégorie satisfaisant les conditions suivantes:
(i) $J$ est rigide, (ii) $\Ob J$ est dénombrable, (iii) Pour deux
objets quelconques $j$, $j'$ de $J$, $\Hom(j,j')$ est fini, (iv)
Pour tout objet $j$ de $J$, l'ensemble des objets de $J$ qui sont
majorés par $j$ \ie qui sont source d'une flèche de but $j$, est
fini. Montrer que $J$ est admissible pour toute
$𝒰$-catégorie $C$. (Montrer d'abord qu'on peut écrire $J$
comme réunion filtrante de sous-catégories pleines finies $J_n$,
telles que tout objet de $J$ majoré par un objet de $J_n$ soit dans
$J_n$. Vérifier ensuite les critères a), b) de 8.8.3.)
\end{enumerate}
\end{enonce*}

\begin{enonce*}[remark]{Exercice 8.8.7}\etiquette[8.8.7]{I.exer8.8.7}
Nous identifions dans les notations un ensemble ordonné et la
catégorie qu'il définit.
\begin{enumerate}
\renewcommand{\theenumi}{\alph{enumi}}
\renewcommand{\labelenumi}{\theenumi)}
\item Soit $C$ un ensemble ordonné. Soit $C^{\widetilde{\phantom{,}}}$
l'ensemble des parties $A$ de $C$ qui sont filtrantes et telles que
$x\in A$ et $y\leq x$ implique $y\in A$. Soit $L_\circ:C\rightarrow
C^{\widetilde{\phantom{,}}}$ l'application qui associe à tout $x\in
C$ l'ensemble $L_\circ(x)$ des $y\in A$ tels que $y\leq x$. Montrer
que $L$ est une application injective et que l'ordre de $C$ est
induit par celui de $C^{\widetilde{\phantom{,}}}$. Pour tout $A\in
C^{\widetilde{\phantom{,}}}$, considérons l'inclusion
$f(A):A\rightarrow C$, %%
c'est un ind-objet de $C$; pour tout ind-objet $\varphi:I\rightarrow
C$ de $C$, soit $g(\varphi)\in C^{\widetilde{\phantom{,}}}$ la
partie de $C$ formée des éléments de $C$ majorés par un élément de
la forme $\varphi(i)$. Montrer\pageoriginale qu'on obtient ainsi
deux équivalences quasi-inverses l'une de l'autre
\[
\xymatrix@C=1.5cm{ \Ind(C)
\ar@{<->}[r]^(.5){\displaystyle\mathop{f,g}_{\approx}}&
C^{\widetilde{\phantom{a}}}, }
\]
transformant le foncteur $L_\circ$ en le foncteur canonique
$L:C\rightarrow \Ind(C)$.

\item Supposons que pour tout $x\in C$, l'ensemble des éléments
majorant (\resp minorant) $x$ est fini. Montrer que tout ind-objet
(\resp pro-objet) de $C$ est essentiellement constant, et même que
pour tout tel $\varphi:I\rightarrow C$ (\resp
$\varphi:I^\circ\rightarrow C$) il existe un $i_\circ\in \Ob I$ tel
que le foncteur induit sur $I_{/i_\circ}$ soit constant (\ie
transforme toute flèche en un isomorphisme).

\item Prenons $C\subset\NN^\circ\times
\NN$ formé des couples d'entiers naturels $(j,i)$ tels que $i\geq
j$. Montrer que $\NN^{\widetilde{\phantom{,}}}$ est isomorphe à
l'ensemble ordonné déduit de $\NN$ (identifié à un sous-ensemble
ordonné de $\NN^{\widetilde{\phantom{,}}}$ comme dans a)) en lui
ajoutant un plus grand élément $\infty$. Montrer que
$C^{\widetilde{\phantom{,}}}$ est isomorphe à $\NN^\circ\times
\NN^{\widetilde{\phantom{,}}}$. Considérons le foncteur
$$
\varphi:\NN^\circ\longrightarrow C^{\widetilde{\phantom{,}}}\quav
\varphi(j)=(j,\infty){\quoi}.
$$
\end{enumerate}
\end{enonce*}

\noindent Montrer que:
\begin{enumerate}
\renewcommand{\labelenumi}{\theenumi)}
\item la limite projective de $\varphi$ n'est pas représentable dans
$C^{\widetilde{\phantom{,}}}$, bien que les limites projectives
filtrantes dans $C$ soient représentables (et soient essentiellement
constantes, en vertu de b)).

\item Pour tout foncteur $\psi:\NN^\circ\rightarrow
C$, on a $\Hom(\psi,\varphi)=∅$, et \textit{a fortiori} le
foncteur $\varphi$ n'admet pas de représentation sous forme
indicielle.
\end{enumerate}

\begin{enonce*}[remark]{Exercice 8.8.8}\etiquette[8.8.8]{I.exer8.8.8}
Soient\pageoriginale $X=\NN\amalg \NN$ l'ensemble somme de deux
copies de $\NN$, $s$ la symétrie de $X$, et pour tout $i\in \NN$,
soit $X_i$ la partie $\Delta_{i+1}\amalg\Delta_i$ de $\NN$ (où
$\Delta_i$ désigne la partie de $\NN$ formée des $j$ tels que $j\leq
i$). Soit $C$ la sous-catégorie de la catégorie des sous-ensembles
de $X$, dont les objets sont les $X_i$, et les flèches les
applications entre des $X_i$ qui sont induites par l'identité de $X$
ou par sa symétrie $s$, et $C^{\widetilde{\phantom{,}}}$ la
catégorie définie de façon analogue, mais où on admet de plus
l'objet $X$.

\begin{enumerate}
\renewcommand{\theenumi}{\alph{enumi}}
\renewcommand{\labelenumi}{\theenumi)}
\item Montrer que $\Ind(C)$ est équivalente à
$C^{\widetilde{\phantom{,}}}$, le foncteur canonique $C\rightarrow
\Ind(C)$ correspondant à l'inclusion $C\rightarrow
C^{\widetilde{\phantom{,}}}$. Montrer qu'il n'existe pas de couple
$(Y,f)$ d'un objet $Y$ de $C$ et d'un endomorphisme $f$ de $Y$, et
de morphisme d'objets à endomorphisme de $\Ind(C)$ de $(Y,f)$ dans
$(X,s)$.

\item En conclure que si $J$ est la catégorie ayant un seul objet,
et en plus de la flèche identique une seule flèche d'ordre 2, alors
le foncteur $J\rightarrow \Ind(C)$ défini par $(X,s)$ n'admet pas de
représentation indicielle.
\end{enumerate}
\end{enonce*}

\subsection{Propriétés d'exactitude de $\Ind(C)$}\etiquette[8.9]{I.sec8.9}

\begin{enonce*}{Proposition 8.9.1}\etiquette[8.9.1]{I.prop8.9.1}
Soit $C$ une $𝒰$-catégorie.
\begin{enumerate}
\renewcommand{\theenumi}{\alph{enumi}}
\renewcommand{\labelenumi}{\theenumi)}
\item Les foncteurs canoniques $L$ (8.2.4.8) et $c$ (8.4.1)
$$
C\xrightarrow{c}\Ind(C)\xrightarrow{L}\widehat{C}
$$
commutent aux limites projectives.

\item Supposons que dans $C$ les limites inductives finies soient
représentables, et que $C$ soit équivalente à une petite
catégorie.
Alors dans $\Ind(C)$ les petites limites projectives sont
représentables.

\item Si\pageoriginale dans $C$ les petites limites projectives
(\resp les limites projectives finies) sont représentables, alors il
en est de même dans $\Ind(C)$. Soit $J$ une catégorie qui est finie
et rigide, ou discrète; si les limites projectives de type $J$ sont
représentables dans $C$, ce même type de limites projectives est
représentables dans $\Ind(C)$.

\item Les petites limites inductives filtrantes dans $\Ind(C)$ (elles
sont représentables en vertu de \hyperref[I.prop8.5.1]{8.5.1}) sont exacts à
gauche, \ie commutent aux limites projectives finies.
\end{enumerate}
\end{enonce*}

\noindent \emph{Démonstration}.
\begin{enumerate}
\renewcommand{\theenumi}{\alph{enumi}}
\renewcommand{\labelenumi}{\theenumi)}
\item Le fait que $L$ commute aux limites projectives résulte du
fait qu'il est pleinement fidèle et du calcul des limites
projectives dans $\widehat{C}$ « argument par argument », qui
implique que pour vérifier qu'un système projectif de morphismes
$F\rightarrow F_i$ $(i\in I)$ dans $\widehat{C}$ fait de $F$ une
limite projective des $F_i$, il suffit de vérifier que l'assertion
analogue est vraie pour les %%
systèmes projectifs d'applications ensemblistes
$\Hom(X,F)\rightarrow \Hom(X,F_i)$, pour tout $X\in \Ob C$. Or
$C\subset \Ind(C)$.

Le fait que $c$ commute aux limites projectives résulte formellement
du fait que $L$ y commute, ainsi que le composé $L\circ
c:C\rightarrow \widehat{C}$.

\item Compte tenu de a), l'assertion revient à dire que toute
limite projective de préfaisceaux ind-représentables est
ind-représentable, ce qui résulte aussitôt du critère
\hyperref[I.them8.3.3]{8.3.3} (v),
compte tenu qu'une limite projective de foncteurs exacts à gauche
est exact à gauche (ce qui résulte du fait que « les limites
projectives commutent aux limites projectives » (\hyperref[I.sec2.5.0]{2.5.0})).

\item Résulte\pageoriginale formellement de la propriété
  analogue de
$\widehat{C}$ (\hyperref[I.cor3.3]{3.3}), et du fait que $L$ est conservatif (étant
pleinement fidèle) et commute aux limites du type envisagé (a) et
8.5.1).

\item Il est bien connu (\cf \hyperref[I.prop2.3]{2.3}) que la
  représentabilité des
limites projectives finies équivaut à celle des limites projectives
des types

\begin{center}
vide $, \ldots ,\quad \vcenter{ \xymatrix@C=.25cm@R=.5cm{
                     \bigdot \ar[rd]&       &\bigdot\ar[ld]\\
                             &\bigdot&
}}$
\end{center}

\noindent (correspondants à des ensembles ordonnés finis
particuliers), et celle des petites limites projectives revient à
celle des produits et des limites projectives finies. Donc la
deuxième assertion faite dans (c) implique la première.
\end{enumerate}

\emph{Supposons d'abord $I$ fini.} Utilisant le résultat
\hyperref[I.prop8.8.5]{8.8.5} sur la représentabilité des %%
foncteurs $\varphi:J\rightarrow \Ind(C)$ sous forme indicielle
(\hyperref[I.sec8.5.4]{8.5.4})
$$
\displaylines{\rlap{$(8.9.1.1)$}\hfill \phi:J\times I\longrightarrow
C{\quoi},\hfill}
$$
l'existence des $\varprojlim \varphi$ est donc un cas particulier du
résultat plus précis et plus général:

\begin{enonce*}{Corollaire \sisi{8.8.2}{8.9.2}}\etiquette[8.9.2]{I.cor8.9.2}
Considérons un foncteur $\varphi:J\rightarrow \Ind(C)$ donné sous
forme indicielle $\Phi$ (8.9.1.1), $J$ étant une
catégorie finie.
Si pour tout $i\in \Ob I$, le foncteur partiel $j\rightarrow
\Phi(j,i)$ a une limite projective (\resp inductive) représentable
dans $C$, alors $\varphi$ a une limite projective (\resp inductive)
représentable dans $\Ind(C)$, et on a un isomorphisme canonique
$$
\displaylines{\rlap{$(8.9.2.1)$}\hfill
\varprojlim\varphi\simeq\textrm{« }{\varinjlim_i}\textrm{ »} \varprojlim_j
\Phi(j,i)\hfill}
$$
(\resp\pageoriginale
$$
\displaylines{\rlap{$(8.9.2.2)$} \hfill \varinjlim \varphi \simeq
\textrm{« }{\varinjlim_i}\textrm{ »} \varprojlim_j\Phi(j,i)){\quoi},\hfill}
$$
où $\varprojlim_j$ (\resp. $\varinjlim_j$) est calculé dans $C$.
\end{enonce*}

Pour la première formule, on utilise simplement que dans $\Ind(C)$
les limites inductives filtrantes commutent à $\varprojlim_J$, et
que $c$ commute à $\varprojlim_J$ (8.9.1 d) et a)):
$$
\varprojlim \varphi=\varprojlim_j \textrm{« }{\varinjlim_i}\textrm{ »\,} \Phi
(j,i) \xleftarrow{\sim} \textrm{« }{\varinjlim_i}\textrm{ »\,}
{\varprojlim_j}^{\Ind(C)} \Phi(j,i)\simeq \textrm{« }{\varinjlim_i}\textrm{ »}%%
\varprojlim_j{^C \ \Phi(j,i){\quoi}}.
$$
La seconde se prouve de même, en utilisant la commutation des
limites inductives de type $I$ aux limites inductives de type $J$
(\hyperref[I.sec2.5.0]{2.5.0}) et le fait que $c$ commute aux limites inductives de type
$J$, prouvé dans \hyperref[I.rem8.9.4]{8.9.4} a) ci-dessous.

\emph{Cas $J$ discret}. Ce cas est contenu dans l'assertion plus
précise suivante:

\begin{enonce*}{Corollaire 8.9.3}\etiquette[8.9.3]{I.cor8.9.3}
Soient $I_\alpha$ des catégories filtrantes essentiellement petites
indexées par un petit ensemble $A$, et pour tout $\alpha \in A$,
soit $X(\alpha)=(X(\alpha)_{i_\alpha})_{i_\alpha \in I_\alpha}$ un
ind-objet de $C$ indexé par $I_\alpha$. Soit $I$ la catégorie
produit des $I_\alpha$, et supposons que pour tout
$\sheaf{i}=(i_\alpha)_{\alpha\in A}\in I$, le produit
$\prod_{\alpha\in A}%%
X (\alpha)_{i_\alpha}$ soit représentable dans $C$. Alors le produit
$\prod_{\alpha \in A} %%
X (\alpha)$ est représentable dans $\Ind(C)$, et il est
canoniquement isomorphe à l'ind-objet indexé par $I$ donné par la
formule
$$
\displaylines{\rlap{$(8.9.3.1)$}\hfill \prod_{\alpha \in A}\;
\mathop{\textrm{« }\varinjlim\textrm{ »}}_{i_\alpha\in I_\alpha}\;
X(\alpha)_{i_\alpha}\simeq
\mathop{\textrm{« }\varinjlim\textrm{ »}}_{\sheaf{i}=(i_\alpha)_{\alpha\in
A}\in I}\; \left(\prod_{\alpha\in
A}A(\alpha)_{i_\alpha}\right){\quoi},\hfill}
$$
où\pageoriginale le produit du deuxième membre est le produit
calculé dans $C$. (On notera que $I$ est filtrant et essentiellement
petit, les $I_\alpha$ l'étant.)
\end{enonce*}

En effet, il est bien connu (et immédiat par réduction au cas où on
travaille dans la catégorie des ensembles) que la formule envisagée
est valable quand on calcule les limites dans $\widehat{C}$. La
conclusion résulte alors du fait que $L$ commute aux limites
envisagées (\hyperref[I.prop8.9.1]{8.9.1} a) et \hyperref[I.prop8.5.1]{8.5.1}).

\begin{enonce*}[remark]{Remarques 8.9.4}\etiquette[8.9.4]{I.rem8.9.4}
La démonstration donnée de 8.9.2, 8.9.3 montre, plus
généralement,
que si pour tout $i\in \Ob I$, $\varprojlim_j \Phi(j,i)$ (\resp
$\varinjlim_j \Phi(j,i)$) \emph{calculé dans} $\Ind(C)$ est
représentable, alors il en est de même de $\varprojlim_i\varphi$
(\resp de $\varinjlim_j\varphi$). Ceci et l'argument de c) montre
que pour une catégorie %%
donnée $J$ provenant d'un ensemble ordonné fini ou discret (ou plus
généralement, qui est $C$-admissible (\hyperref[I.prop8.3.1]{8.3.1})), les
limites projectives (\resp inductives) de type $J$ sont
représentables dans $\Ind(C)$ si (et seulement si) pour tout
foncteur $\varphi:J\rightarrow C$, la limite projective (\resp
inductive) de $\varphi$ \emph{calculée dans} $\Ind(C)$ est
représentable. Dans le cas non respé, cela signifie aussi, en vertu
de (a), que toute limite projective de type $J$ de préfaisceaux
représentables sur $C$ est ind-représentable.
\end{enonce*}

\begin{enonce*}{Proposition 8.9.5}\etiquette[8.9.5]{I.prop8.9.5}
Soit $C$ une $𝒰$-catégorie.
\begin{enumerate}
\renewcommand{\theenumi}{\alph{enumi}}
\renewcommand{\labelenumi}{\theenumi)}
\item Le foncteur canonique
$$
c:C\longrightarrow \Ind(C)
$$
est exact à droite (donc exact, compte tenu de $8.9.1$ {\rm a)}).

\item Si les limites inductives finies (\resp les sommes finies) sont
représentables dans $C$, alors les petites limites inductives (\resp
les petites sommes) sont représentables dans $\Ind(C)$. Soit $J$ un
ensemble préordonné\pageoriginale fini ou discret; si les limites
inductives de type $J$ sont représentables dans $C$, il en est de
même dans $\Ind(C)$.
\end{enumerate}
\end{enonce*}

\noindent \emph{Démonstration.}
\begin{enumerate}
\renewcommand{\theenumi}{\alph{enumi}}
\renewcommand{\labelenumi}{\theenumi)}
\item Supposons que dans $C$ on ait $X\simeq \varinjlim_J X_j$, où
$J$ est une catégorie finie, $X$ et les $X_j$ dans $C$. On a alors,
pour tout ind-objet $𝒴=(Y_i)_{i\in I}$ de $C$:
\begin{gather*}
\Hom(X,𝒴)\simeq \varinjlim_i \Hom(X,Y_i)\simeq \varinjlim_i
\varprojlim_j \Hom(X_j,Y_i)\\
\mathrel{\mathop\simeq^{2.8}}\varprojlim_j \varinjlim_i \Hom(X_j,
Y_i) \simeq \varprojlim_j \Hom (X_j,𝒴){\quoi},
\end{gather*}
et la conclusion voulue résulte de la comparaison des termes
extrêmes.

\item On a déjà noté dans \hyperref[I.rem8.8.4]{8.8.4} que les
  assertions faites
résultent de a) et de \hyperref[I.cor8.9.2]{8.9.2}, du moins pour les limites
inductives finies. Pour prouver les conclusions faites dans les cas
infinis, on est ramené au cas des sommes, qui peut s'interpréter
comme une limite inductive filtrante de sommes finies, et on conclut
donc grâce à \hyperref[I.prop8.5.1]{8.5.1}.
\end{enumerate}

\begin{enonce*}[remark]{Remarque 8.9.6}\etiquette[8.9.6]{I.rem8.9.6}
On a déjà observé que le foncteur $c$ ne commute pas en général aux
limites inductives filtrantes, donc pas non plus aux sommes
infinies, contrairement à ce qui a lieu pour $L:\Ind(C)\rightarrow
\widehat{C}$. Par contre, à l'exception de la commutation aux
limites inductives \emph{filtrantes} (\hyperref[I.prop8.5.1]{8.5.1}), $L$ n'a
pratiquement jamais de propriétés de commutation à quelque autre
type de limites inductives (objet initial, somme de deux objets
conoyaux de doubles flèches). Ainsi, si $C$ admet un objet initial
$∅_C$, qui est donc objet initial de $\Ind(C)$ en vertu de
a), $∅_C$ n'est jamais objet initial de $\widehat{C}$ (\ie
identique au préfaisceau constant de valeur $∅$), puisque
$\Hom(∅_C,∅_C)\neq ∅$!
\end{enonce*}

\begin{enonce*}{Proposition 8.9.7}\etiquette[8.9.7]{I.prop8.9.7}
Soit\pageoriginale
$$
f:C\longrightarrow C'
$$
un foncteur entre $𝒰$-catégories, et soit $J$ une catégorie
finie et rigide $(8.8.5)$. Si les limites inductives (\resp
projectives) de type $J$ sont représentables dans $C$ et si $f$ y
commute, alors $\Ind(f):\Ind(C)\rightarrow \Ind(C')$ commute
également à ce type de limites.
\end{enonce*}

Cela résulte aussitôt de \hyperref[I.prop8.8.5]{8.8.5} et du calcul
\hyperref[I.cor8.9.2]{8.9.2} des limites finies dans une catégorie de
ind-objets, pour un foncteur représenté sous forme indicielle.

\begin{enonce*}{Corollaire 8.9.8}\etiquette[8.9.8]{I.cor8.9.8}
Si dans $C$ les limites inductives (\resp projectives) finies sont
représentables, et si $f$ est exact à droite (\resp à gauche) alors
il en est de même de $\Ind(f)$; dans le cas non respé, $\Ind(f)$
commute même aux petites limites inductives quelconques.
\end{enonce*}

La première assertion résulte de \hyperref[I.prop8.9.7]{8.9.7}. La deuxième
résulte de la première, compte tenu que $\Ind(f)$ commute aux
limites inductives filtrantes (\hyperref[I.prop8.6.3]{8.6.3}).

\begin{enonce*}[remark]{Exercice 8.9.9}\etiquette[8.9.9]{I.exer8.9.9}
Soit $C$ une $𝒰$-catégorie.
\begin{enumerate}
\renewcommand{\theenumi}{\alph{enumi}}
\renewcommand{\labelenumi}{\theenumi)}
\item Supposons que dans $C$ les sommes finies sont
représentables. Montrer que si dans $C$ les sommes \emph{finies}
sont disjointes (\resp universelles) (\cf II 4.5), alors dans
$\Ind(C)$ les \emph{petites} sommes sont disjointes (\resp
universelles).

\item Supposons que dans $C$ tout morphisme se factorise en un
épimorphisme suivi d'un monomorphisme (\resp en un épimorphisme
effectif suivi\pageoriginale d'un monomorphisme, \resp en un
épimorphisme suivi d'un monomorphisme effectif, \resp un
épimorphisme effectif suivi d'un monomorphisme effectif). Montrer
que $\Ind(C)$ satisfait à la même propriété. Dans le dernier cas
respé, on conclut donc que tout épimorphisme de $\Ind(C)$ est
effectif, tout monomorphisme de $\Ind(C)$ est effectif, tout
bimorphisme de $\Ind(C)$ est un isomorphisme; montrer que dans ce
cas tout épimorphisme (\resp monomorphisme) de $\Ind(C)$ peut se
représenter par un système inductif d'épimorphismes (\resp de
monomorphismes) de $C$. Si on suppose de plus que dans $C$ tout
épimorphisme (\resp tout monomorphisme) est \emph{universel}, la
même propriété est vraie dans $\Ind(C)$.

\item Supposons $C$ additive (\resp abélienne), alors $\Ind(C)$
l'est également.

\item Soient $X=\textrm{« }\varinjlim_I\textrm{ »\,}X_i$ un ind-objet de $C$,
$R\rightrightarrows X$ une relation d'équivalence dans $X$. Pour
tout $i\in \Ob I$, soit $R_i\rightrightarrows X_i$ la relation
d'équivalence dans $X_i$ (considéré comme objet de $\Ind(C)$)
induite par $R$ via $X_i\rightarrow X$. (On suppose que le produit
fibré $R_i=R\times_{X\times X}X_i \times X_i$ dans $(\Ind(C))^{\wedge}$ est
représentable par un objet de $\Ind(C)$, ce qui est le cas si dans
$C$ les limites projectives finies sont représentables.) Montrer que
pour que $R$ soit effective, il suffit que les $R_i$ le soient.

\item Soient $X$ un objet de $C$, $R$ une relation d'équivalence
dans $c(X)$ \ie dans $X$ considéré comme objet de $\Ind(C)$.
Supposons que dans $C$ les produits fibrés soient représentables.
Pour que $R$ soit effective, il faut que $R$ soit de la forme
$\textrm{« }\varinjlim_i\textrm{ »\,} R_i$, où les $R_i$ sont des relations
d'équivalence dans $X$ (regardé comme un objet de $C$), et cette
condition est\pageoriginale suffisante si on suppose que dans $C$
les relations d'équivalence sont effectives.

\item Supposons que dans $C$ tout morphisme se factorise en un
épimorphisme effectif suivi d'un monomorphisme effectif, et que les
limites projectives finies soient représentables. Soit $X$ un objet
de $C$, et $R$ un sous-objet de $X\times X$, regardé comme élément
de $\Ind(C)$. Écrivons $R$ sous la forme $\textrm{« }\varinjlim\mbox{''
}Z_i$, où $(Z_i)_{i\in I}$ est une famille filtrante croissante de
sous-objets de $X\times X$ \emph{dans} $C$ (c'est possible grâce a
b)). Montrer que pour que $R$ soit une relation d'équivalence dans
$X$, il faut et il suffit que pour tout $i\in \Ob I$, existe un
$j\in \Ob I$ qui contienne $s(Z_i)$ et $Z_i\circ Z_i$, où $s$ est la
symétrie de $X\times X$.

\item Supposons que dans $C$ les limites projectives finies ainsi que
les sommes finies soient représentables, que toute relation
d'équivalence y soit effective, et tout morphisme s'y factorise en
épimorphisme effectif suivi par un monomorphisme effectif. Montrer
que dans $\Ind(C)$ les relations d'équivalence sont universelles si
et seulement si $C$ satisfait à la condition suivante:
\begin{itemize}
\item[%%
ST)] Pour tout objet $X$ de $C$ et tout sous-objet $Z$ de $X\times
X$ dans $C$, si on définit par récurrence la suite de sous-objets
$Z_n$ $(n\geq 0)$ de $X\times X$ dans $C$ par $Z_\circ=Z$,
$Z_n=\Sup(Z_n,s(Z_n),Z_n\circ Z_n)$ (le $\Sup$ pris dans l'ensemble
des sous-objets de $X\times X$, qui existe grâce aux hypothèses
faites sur C), alors la suite $(Z_n)_{n\geq 0}$ est
\emph{stationnaire.}
\end{itemize}

\setcounter{enumi}{10}
\item Montrer que dans $\Ind((\Ens))$, les relations d'équivalence
ne sont pas nécessairement effectives.

$^\ast$ NB. Pour des critères pour que $\Ind(C)$ soit un topos, \cf
VI 8.9.9.${}_\ast$
\end{enumerate}
\end{enonce*}

\begin{enonce*}[remark]{Exercice 8.9.10}\etiquette[8.9.10]{I.exer8.9.10}
Soit\pageoriginale $C$ une $𝒰$-catégorie. Posons
$\Pro(C)=(\Ind(C^\circ))^\circ$ (\cf 8.11).
\begin{enumerate}
\renewcommand{\theenumi}{\alph{enumi}}
\renewcommand{\labelenumi}{\theenumi)}
\item Supposons que dans $C$ les sommes finies (\resp les petites
sommes) sont représentables. Montrer que si elles sont disjointes
(\cf II 4.5), alors $\Pro(C)$ satisfait la même %%
condition.

\item Soit $J$ un petit ensemble, tel que les sommes indexées par
$J$ soient représentables dans $C$, donc aussi dans $\Pro(C)$. Soit
$(X(\alpha))_{\alpha\in J}$ une famille d'éléments de $\Pro(C)$,
avec $X(\alpha)=\textrm{« }\varprojlim_{I_\alpha}\textrm{ »\,}X_{i_\alpha}$.
Montrer que pour que la somme des $X(\alpha)$ dans $\Ind(C)$ soit
universelle, il suffit qu'il en soit de même pour chacune des
familles $\varprojlim_J X_{i_\alpha}$, où $(i_\alpha)_{\alpha\in I}\in
\prod_{\alpha\in J}I_\alpha$. En conclure que si dans $C$ les sommes
de type $J$ sont universelles, pour qu'il en soit de même dans
$\Pro(C)$, il faut et il suffit que pour toute famille
$(X(\alpha))_{\alpha\in J}$ comme dessus, avec $I_\alpha=I$ pour
tout $\alpha\in J$, l'homomorphisme canonique dans $\Pro(C)$
$$
\textrm{« }\varprojlim_I\textrm{ »\,}\amalg_\alpha X(\alpha)_i\longleftarrow
\textrm{« }\varprojlim_{I^J}\textrm{ »\,} \amalg_{\alpha\in
J}X(\alpha)_{i_\alpha}
$$
soit un isomorphisme. En conclure que dans $\Pro((\Ens))$ les sommes
de type $J$ sont universelles si et seulement si $J$ est fini.
\end{enumerate}
\end{enonce*}

\begin{enonce*}[remark]{Exercice 8.9.11}\etiquette[8.9.11]{I.exer8.9.11}
Soit $C$ une $𝒰$-catégorie.
\begin{enumerate}
\renewcommand{\theenumi}{\alph{enumi}}
\renewcommand{\labelenumi}{\theenumi)}
\item Soit $(X_i\rightarrow X)_{i\in I}$ une famille de morphismes
dans $C$. Montrer que pour qu'elle soit épimorphique dans $\Ind(C)$,
il suffit qu'il existe une sous-famille \emph{finie} qui soit
épimorphique dans $C$. Prouver que cette condition est également
nécessaire lorsqu'on suppose que dans $C$ toute famille finie de
morphismes de but $X_i$ se factorise en une famille
épimorphique,\pageoriginale suivie d'un monomorphisme effectif
(10.5). (Pour cette dernière assertion, soit $P$ l'ensemble des
      parties finies de $I$, ordonné par inclusion, et soit
      $X'=(X'_J)_{J\in P}$ le ind-objet formé des images des
      sous-familles finies de la famille donnée, enfin soit
      $X''=X\amalg_{X'}X=\textrm{« }\lim\textrm{ »\,}X\amalg_{X'_J}X$. Montrer
      que les deux morphismes canoniques $X\rightrightarrows X''$ sont
      distincts, mais coïncident sur les $X'_J$.)

\item Supposons que dans $C$ toute famille finie de morphismes de
      même but se factorise en une famille épimorphique, suivie
      d'un monomorphisme effectif. Prouver que $\Ind(C)$ admet une
      petite sous-catégorie génératrice par épimorphismes
      (\hyperref[I.def7.1]{7.1}) si et seulement si $C$ admet une petite
      sous-catégorie
      $C'$ telle que tout objet de $C$ soit but d'une famille
      épimorphique \emph{finie} de source dans $C'$. Lorsqu'on
      suppose que $C$ admet une petite sous-catégorie
      génératrice (\hyperref[I.def7.1]{7.1}), alors $\Ind(C)$ admet une petite
      sous-catégorie génératrice par épimorphismes si et
      seulement si $C$ est équivalente à une petite
      catégorie. (Pour ce dernier énoncé, utiliser
      \hyperref[I.var7.5.2]{7.5.2}).

\item $\Ind((\Ens))$ n'admet pas de petite sous-catégorie
  génératrice.
\end{enumerate}
\end{enonce*}

\begin{enonce*}[remark]{Exercice 8.9.12}\etiquette[8.9.12]{I.8.9.12}
Soit $C$ une $𝒰$-catégorie, et considérons
$\Pro(C)=\Ind(C^\circ)^\circ$.
\begin{enumerate}
\renewcommand{\theenumi}{\alph{enumi}}
\renewcommand{\labelenumi}{\theenumi)}
\item Montrer que si $\Pro(C)$ admet une petite sous-catégorie $P'$
génératrice par épimorphismes (7.1), alors il existe une petite
sous-catégorie $C'$ de $C$ qui est génératrice par épimorphismes
dans $\Pro(C)$. (Prendre la catégorie des composants des objets de
$P'$.)

\item Montrer\pageoriginale que la catégorie $\Pro((\Ens))$ n'a pas de petite
sous-catégorie génératrice par épimorphismes. (Si $C'$ est comme
dans a), choisir un ensemble $X$ dont le cardinal majore strictement
les cardinaux des éléments de $C'$, et considérer le pro-ensemble
$𝒳$ formé par les complémentaires dans $X$ des parties de
cardinal $<\card(X)$. Montrer que pour tout ensemble $Y$ non vide de
cardinal $<\card(X)$, on a $\Hom(Y,𝒳)=∅$, mais que
$𝒳$ n'est pas isomorphe au pro-ensemble constant
$∅$.)
\end{enumerate}
\end{enonce*}

\subsection{Notions duales: pro-objets, foncteurs pro-représentables}\etiquette[8.10]{I.sec8.10}

Soit $C$ une $𝒰$-catégorie. On appelle \emph{pro-objet} de
$C$ tout foncteur
$$
\displaylines{\rlap{$(8.10.1)$}\hfill \varphi:I^\circ\longrightarrow
C{\quoi},\hfill}
$$
où $I$ est une catégorie filtrante essentiellement petite (appelée
la catégorie d'indices), et comme d'habitude $I^\circ$ désigne la
catégorie opposée. Soit $𝒱$ un univers tel que
$𝒱\supset 𝒰$. On fera attention que les pro-objets
de $C$ indexés par des $I\in 𝒱$ sont en correspondance
biunivoque avec les ind-objets de $C^\circ$ indexés par des $I\in
𝒱$, en associant à tout tel ind-objet $\psi:I\rightarrow
C^\circ$ le pro-objet $\varphi=\psi^\circ:I^\circ\rightarrow C$.
S'inspirant de cette correspondance, on définit « par transport de
structure et renversement des flèches » la notion de morphisme
entre pro-objets de $C$ à partir de la notion analogue (8.2.4.2),
(8.2.5.1) pour les in-objets, et on définit en conséquence la
catégorie des pro-objets de $C$ indexés par des $I\in 𝒱$, et
un isomorphisme canonique:
$$
\displaylines{\rlap{$(8.10.2)$}\hfill
\Pro_{𝒱}(C,𝒰)\simeq \Ind_{𝒱}(C^\circ,
𝒰)^\circ{\quoi};\hfill}
$$
pour\pageoriginale $𝒱=𝒰$ on parle simplement de la
\emph{catégorie des pro-objets de} $C$, notée $\Pro_{𝒰}(C)$
ou simplement $\Pro(C)$, qui est donc définie par
$$
\displaylines{\rlap{$(8.10.3)$}\hfill
\Pro(C)=\Pro_{𝒰}(C,𝒰)\simeq
\Ind(C^\circ)^\circ{\quoi}.\hfill}
$$
Si deux pro-objets sont donnés sous \emph{forme indicielle}
\begin{equation}
𝒳=(X_i)_{i\in I}\quav 𝒴=(Y_j)_{j\in
J}{\quoi},\tag{8.10.4}
\end{equation}
la formule (8.2.5.1) prend ici la forme
\begin{equation}
\Hom(𝒳,𝒴)\simeq \varprojlim_j \varinjlim_i
\Hom(X_i,Y_j){\quoi}.\tag{8.10.5}
\end{equation}
Le foncteur (8.4.1) appliqué à $C^\circ$ donne, par passage aux
catégories opposées, un foncteur canonique (dénoté par
la même
lettre $c$ s'il n'y a pas risque de confusion), permettant
d'identifier $C$ à une sous-catégorie pleine de $\Pro(C)$:
\begin{equation}
c:C\longrightarrow\Pro(C){\quoi};\tag{8.10.6}
\end{equation}
c'est pour avoir un tel foncteur, et non $C\rightarrow
\Pro(C)^\circ$, qu'on a « renversé les flèches » dans la
définition des morphismes de pro-objets à partir de la définition
analogue pour les ind-objets. Les pro-objets dans l'image
essentielle de (8.10.6) s'appellent encore \emph{pro-objets
essentiellement constants.}

Posons
\begin{equation}
\check{C}=(C^\circ)^{\widehat{\phantom{,}}}=\Hhom(C,(\Ens)){\quoi},\tag{8.10.7}
\end{equation}
alors le foncteur (8.2.4.7) peut être considéré
comme un foncteur
canonique pleinement fidèle
\begin{equation}
L:\Pro_{𝒱}(C,𝒰)\longrightarrow
\check{C}^\circ{\quoi},\tag{8.10.8}
\end{equation}\pageoriginale
et on trouve en particulier
\begin{equation}
L:\Pro(C)\hookrightarrow \check{C}^\circ{\quoi}.\tag{8.10.9}
\end{equation}

Nous laisserons au lecteur le soin de traduire, au fur et mesure des
besoins, les résultats des sections précédentes et de celles qui
suivent du langage des ind-objets dans celui des pro-objets. Suivant
le contexte mathématique, c'est l'un ou l'autre langage qui est le
plus utile, sans compter les cas où les deux notions s'introduisent
simultanément, par exemple lorsqu'il y a lieu de considérer des
catégories complexes comme $\Pro(\Ind(C))$ ou $\Ind(\Pro(C))$. Nous
nous contentons de donner quelques indications supplémentaires, pour
fixer la terminologie et les notations, et préciser le « yoga ».


\setcounter{subsubsection}{9}
\subsubsection{}\etiquette[8.10.10]{I.sec8.10.10}
Un foncteur \emph{covariant $F:C\rightarrow (\Ens)$} est dit
\emph{foncteur pro-représentable} s'il est dans l'image essentielle
de (8.10.9) (ou 8.10.8, cela revient au même); la
sous-catégorie
pleine de $\check{C}$ formée de ces foncteurs est donc équivalente à
$\Pro(C)^\circ$. On préférera généralement travailler dans la
catégorie opposée, image essentielle \emph{en tant que
sous-catégorie de} (8.10.9), qui est donc équivalente à $\Pro(C)$
lui-même. En accord avec cet usage, il est souvent préférable de
regarder les foncteurs covariants
$$
F:C\longrightarrow (\Ens)
$$
comme des objets de $\Hhom(C,(\Ens))^\circ=\check{C}^\circ$, et
d'écrire en conséquence la catégorie des homomorphismes
$$
X\longrightarrow F\mbox{ dans } \check{C}^\circ{\quoi},
$$
\ie\pageoriginale des couples $(X,u)$ d'un objet $X$ de $C$ et d'un
$u\in F(X)$, comme
\begin{equation}
{}_{\displaystyle F}\backslash C{\quoi}  %%
\ .\tag{8.10.11}
\end{equation}
Avec cette convention, on trouvera donc un foncteur \emph{covariant}
(foncteur source)
\begin{equation}
{}_{\displaystyle F}\backslash C\rightarrow C{\quoi},\tag{8.10.12}
\end{equation}
et 3.4 se récrit sous la forme
\begin{equation}
F\xrightarrow{\sim}%%
\textrm{« }\varprojlim_{({}_{\displaystyle F}\backslash C)^\circ}\textrm{ »\,}
X\tag{8.10.13}
\end{equation}

\setcounter{subsubsection}{13}
\subsubsection{}\etiquette[8.10.14]{I.sec8.10.14}
Le critère \hyperref[I.them8.3.3]{8.3.3} appliqué à $C^\circ$ devient
un critère de
proreprésentabilité pour $F:$ il faut et il suffit que
$({}_F\backslash C)^\circ$ soit filtrante et essentiellement petite,
cette dernière condition étant superflue si $C$ est équivalente à
une petite catégorie; si de plus dans $C$ les limites projectives
finies sont représentables, il revient au même de dire que $F$ y
commute \ie est \emph{exact à gauche.}

\subsubsection{}\etiquette[8.10.15]{I.sec8.10.15}
Dans $\Pro(C)$ les petites limites projectives filtrantes sont
représentables et le foncteur (8.10.9) y commute; en d'autres
termes, ce foncteur transforme limites projectives en $\Pro(C)$ en
limites \emph{inductives} de $\check{C}=\Hhom(C,(\Ens))$, tout comme
son composé $C\rightarrow \check{C}^\circ$ avec (8.10.6), qui
partage avec lui la fâcheuse propriété d'être
contravariant quand on
le considère à valeurs dans $\check{C}$. On fera attention que les
foncteurs pro-représentables\pageoriginale de $C$ dans $(\Ens)$ sont
les foncteurs qui sont des petites limites \emph{inductives}
filtrantes de foncteurs représentables (et non des limites
projectives, comme la terminologie pourrait éventuellement le
suggérer).

\subsection{Ind-adjoints et pro-adjoints}\etiquette[8.11]{I.sec8.11}

\subsubsection{}\etiquette[8.11.1]{I.sec8.11.1}
Soit
\begin{equation}
f:C\longrightarrow C'\tag{8.11.1.1}
\end{equation}
un foncteur entre $𝒰$-catégories, d'où un foncteur
$F \sisi{→}{↦} F\circ f$
\begin{equation}
f^\ast:\widehat{C}'\longrightarrow\widehat{C}{\quoi}.\tag{8.11.1.2}
\end{equation}
On dit que $f$ \emph{admet un ind-adjoint} si le foncteur précédent
transforme foncteurs ind-représentables en foncteurs
ind-représentables. Comme $f^\ast$ commute aux limites inductives,
et que la sous-catégorie pleine de $C'$ formée des foncteurs
ind-représentables est stable par petites limites inductives
filtrantes, il revient au même de dire que $f$ admet un ind-adjoint,
ou que $f^\ast$ applique objets de $C$ (identifié à une
sous-catégorie pleine de $\widehat{C}$) dans des foncteurs
ind-représentables sur $C$, \ie
\begin{equation}
X\longmapsto \Hom(f(X),Y')\mbox{ est ind-représentable pour tout }
Y'\in \Ob C'{\quoi}.\tag{8.11.1.3}
\end{equation}
Une autre façon d'exprimer la condition que $f$ admette un
ind-adjoint est de dire qu'il existe un foncteur
\begin{equation}
g:\Ind(C')\longrightarrow \Ind(C)\tag{8.11.1.4}
\end{equation}
qui soit essentiellement induit par $f^\ast$, \ie tel qu'on ait un
isomorphisme de bifoncteurs
\begin{equation}
\Hom_{\Ind(C)}(X,g(𝒴'))\simeq
\Hom_{\Ind(C')}(f(X),𝒴'){\quoi},\tag{8.11.1.5}
\end{equation}
où\pageoriginale $X\in \Ob C$ et $Y'\in \Ob \Ind(C')$. En fait, il
suffit même d'avoir un foncteur
\begin{equation}
g_\circ:C'\longrightarrow \Ind(C)\tag{8.11.1.6}
\end{equation}
et un isomorphisme de bifoncteurs
\begin{equation}
\Hom_{\Ind(C)}(X,g_\circ(Y'))\simeq \Hom_{C'}(f(X),Y')\tag{8.11.1.7}
\end{equation}
en $X\in \Ob C$, $Y'\in \Ob C'$. Le foncteur $g$ (\resp $g_\circ$)
est évidemment déterminé à isomorphisme unique près par $f$, et
inversement on reconstitue $f$ à isomorphisme canonique près par la
connaissance de ce $g$ (\resp $g_\circ$). Il est clair sur
(8.11.1.5) que le foncteur $g$ commute aux limites inductives, et
qu'il « prolonge » le foncteur $g_\circ$; cela le détermine
donc à isomorphisme unique près en termes de $g_\circ$
(\hyperref[I.sec8.7.2]{8.7.2}). Le
foncteur $g$, et parfois aussi le foncteur $g_\circ$ qu'il prolonge,
est appelé le \emph{foncteur ind-adjoint de} $f$ (inutile ici de
préciser: à droite, car l'autre, s'il existe, s'appellera le
pro-adjoint, \cf \hyperref[I.sec8.11.5]{8.11.5} plus bas).

Bien entendu, lorsque $f$ admet un adjoint à droite
$$
f^{\ad}:C'\longrightarrow C{\quoi},
$$
il admet un ind-adjoint $g$, et celui-ci est isomorphe canoniquement
au prolongement canonique de $f$ aux ind-objets:
\begin{equation}
g\simeq\Ind(f^{\ad}){\quoi}.\tag{8.11.1.8}
\end{equation}
La notion de ind-adjoint est donc une généralisation naturelle de la
notion d'adjoint à droite.

Considérons\pageoriginale maintenant le prolongement canonique
\begin{equation}
\Ind(f):\Ind(C)\longrightarrow\Ind(C'){\quoi}.\tag{8.11.1.9}
\end{equation}
On a alors:

\begin{enonce*}{Proposition 8.11.2}\etiquette[8.11.2]{I.prop8.11.2}
Pour que le foncteur $f:C\rightarrow C'$ admette un ind-adjoint, il
faut et il suffit que le foncteur $\Ind(f)$ $(8.11.1.9)$ admette un
adjoint à droite. Ce dernier est canoniquement isomorphe au
ind-adjoint $(8.11.1.4)$.
\end{enonce*}

La suffisance et la dernière assertion sont triviales sur la formule
de ind-adjonction (8.11.1.5). Pour la nécessité, on
note que par
passage à la limite projective sur cette formule sur des $X_i$, on
déduit un isomorphisme d'adjonction en les pro-objets
$𝒳=(X_i)_{i\in I}$ et $𝒴'$:
\begin{equation}
\Hom_{\Ind(C)}(𝒳, g(𝒴'))\simeq
\Hom_{\Ind(C')}(\Ind(f)(𝒳),𝒴'){\quoi},\tag{8.11.2.1}
\end{equation}
\cqfd

\begin{enonce*}{Corollaire 8.11.3}\etiquette[8.11.3]{I.cor8.11.3}
Si $f$ admet un ind-adjoint, alors $\Ind(f)$ commute aux limites
inductives, et le ind-adjoint $g$ commute aux limites projectives.
\end{enonce*}

\begin{enonce*}{Proposition 8.11.4}\etiquette[8.11.4]{I.prop8.11.4}
Pour que le foncteur $f:C\rightarrow C'$ admette un ind-adjoint, il
faut que $f$ soit exact à \sisi{gauche}{droite}, et cette condition est suffisante
lorsque $C$ est équivalente à une petite catégorie.
\end{enonce*}

Cela résulte du critère (8.11.1.3), et de \hyperref[I.prop8.3.1]{8.3.1}
et \hyperref[I.them8.3.3]{8.3.3} (iv).

Pour un autre critère en termes de la notion de foncteur accessible,
\cf 8.13.3.

\setcounter{subsubsection}{4}\etiquette[8.11.5]{I.sec8.11.5}
\subsubsection{}
Considérons\pageoriginale maintenant le foncteur $F\mapsto F\circ f$
\begin{equation}
f^{\circ\ast}:\check{C}'\longrightarrow \check{C}\tag{8.11.5.1}
\end{equation}
induit par \sisi{$F$}{$f$}. On dira que $f$ \emph{admet un pro-adjoint} si le
foncteur précédent applique foncteur pro-représentable en foncteur
pro-représentable, \ie s'il existe un foncteur (appelé
\emph{foncteur pro-adjoint de $f$})
\begin{equation}
g:\Pro(\sisi{C}{C'})\longrightarrow \Pro(\sisi{C'}{C}){\quoi},\tag{8.11.5.2}
\end{equation}
et un isomorphisme de bifoncteurs
\begin{equation}
\Hom_{\Pro(C)}(g(𝒴'),𝒳)\simeq
\Hom_{\Pro(C')}(𝒴'\sisi{)}{},f(𝒳)){\quoi}.\tag{8.11.5.3}
\end{equation}
Bien entendu, dire que $f$ admet un pro-adjoint signifie que
$f^\circ$ admet un ind-adjoint, de sorte que les notions et
résultats pour les ind-adjoints se traduisent trivialement en termes
de pro-adjoints. Signalons seulement que $f$ admet un pro-adjoint si
et seulement si $\Pro(f):\Pro(C)\rightarrow \Pro(C')$ admet un
adjoint \emph{à \sisi{droite}{gauche}}, et que dans ce cas le foncteur $g$
précédent est un tel adjoint à \sisi{droite}{gauche} de
$\Pro(\sisi{C}{f})$; et qu'il faut
pour ceci que $f$ soit exact à \emph{gauche}, cette condition étant
également suffisante lorsque $C$ est équivalente à une petite
catégorie. Dans ce cas, $f$ est donc exact si et seulement si il
admet à la fois un ind-adjoint et un pro-adjoint.

\begin{enonce*}[remark]{Exemple 8.11.6}\etiquette[8.11.6]{I.exem8.11.6}
Considérons le cas d'un foncteur
$$
f:C\longrightarrow (\Ens){\quoi}.
$$
Si ce foncteur admet un pro-adjoint, il est pro-représentable, et la
réciproque est vraie si et seulement si la sous-catégorie pleine de
$\check{C}$ formée des foncteurs pro-représentables est stable par
petits produits; c'est\pageoriginale le cas en particulier si $C$
est équivalente à une petite catégorie (\hyperref[I.sec8.10.14]{8.10.14}) ou si
dans $C$ les petits produits sont représentables (\hyperref[I.prop8.9.5]{8.9.5}
b) appliquée à $C^\circ$). À peu de choses près, on peut donc dire
que pour un foncteur $f:C\rightarrow C'$, la notion d'existence d'un
pro-adjoint est la \emph{généralisation naturelle de la notion de
pro-représentabilité de $f$}, qui est définie lorsque $C'=(\Ens)$.
\end{enonce*}

\subsection{Ind-objets et pro-objets stricts. Application à un
critère de représentabilité}\etiquette[8.12]{I.sec8.12}

\subsubsection{}\etiquette[8.12.1]{I.sec8.12.1}
Soit
$$
𝒳=(X_i)_{i\in I}
$$
un ind-objet de la $𝒰$-catégorie $C$, et soit $F$ le
préfaisceau qu'il ind-représente. On voit alors aussitôt qu'il
revient au même de dire que les morphismes canoniques
$$
X_i\longrightarrow F=\textrm{« }\varinjlim_i\textrm{ »\,}X_i
$$
sont des monomorphismes de $\Ind(C)$ (ou, ce qui revient au même, de
$\widehat{C}$, \ie des monomorphismes de foncteurs « argument par
argument »), ou de dire que pour toute flèche $i\rightarrow j$ de
$I$, la flèche de transition correspondante
$$
X_i\longrightarrow X_j
$$
est un monomorphisme. Lorsque ces conditions sont remplies, et si de
plus $I$ est une catégorie ordonnée, on dit que $𝒳$ est un
\emph{ind-objet strict}. On\pageoriginale notera que cette condition
n'est pas invariante par isomorphisme de ind-objets; un ind-objet
sera appelé \emph{essentiellement strict} s'il est isomorphe à un
ind-objet strict. Un préfaisceau sera appelé \emph{strictement
ind-représentable} s'il est ind-représentable par un ind-objet
strict; donc $F=\textrm{« }\varinjlim_i\textrm{ »\,}X_i$ est strictement
ind-représentable si et seulement si $𝒳=(X_i)_{i\in I}$ est
essentiellement strict.

\paragraph{}\etiquette[8.12.1.1]{I.8.12.1.1}
Soit $F$ un préfaisceau sur $C$, et considérons la sous-catégorie
pleine de $C_{/F}$ formée des flèches $X\rightarrow F$ de source
dans $C$ qui sont des monomorphismes. On l'appellera la
\emph{catégorie des sous-foncteurs représentables de $F$}; c'est la
catégorie associée à l'ensemble ordonné des sous-foncteurs
représentables de $F$, ordonné par l'ordre induit de celui de
l'ensemble des sous-objets de $F$. Il résulte alors aussitôt des
définitions:

\begin{enonce*}{Proposition 8.12.2}\etiquette[8.12.2]{I.prop8.12.2}
Pour que le préfaisceau $F$ sur $C$ soit strictement
ind-représentable, il faut et il suffit que la catégorie (ordonnée)
$I$ des sous-foncteurs représentables de $F$ soit filtrante et
essentiellement petite et que l'on ait
\begin{equation*}
\varinjlim_I X_i \xrightarrow{\sim}F{\quoi}.\tag{8.12.2.1}
\end{equation*}
Lorsque pour tout objet $X$ de $C$, l'ensemble des sous-objets de
$X$ dans $C$ est petit (par exemple si $C$ admet une petite
sous-catégorie génératrice $(7.4)$), cela implique que la catégorie
des sous-foncteurs représentables est même petite.
\end{enonce*}

\setcounter{subsubsection}{2} \setcounter{paragraph}{1}
\paragraph{}\etiquette[8.12.2.2]{I.sec8.12.2.2}
Si\pageoriginale $F$ est strictement ind-représentable, il y a donc
une façon privilégiée de le ind-représenter par un ind-objet, et ce
dernier est même un ind-objet strict: on prend la représentation
(8.12.2.1). Un ind-objet strict $𝒴$ est dit
\emph{saturé} s'il est isomorphe à un ind-objet de la forme $(X_i)$
figurant dans (8.12.2.1), pour $F$ convenable; donc il
existe à isomorphisme unique près, un seul ind-objet strict saturé
isomorphe au ind-objet strict donné, savoir celui envisagé dans
\hyperref[I.prop8.12.2]{8.12.2}, en prenant $F=\varinjlim 𝒴$ (limite dans
$\widehat{C}$).

\paragraph{}\etiquette[8.12.2.3]{I.sec8.12.2.3}
Supposons qu'on sache déjà que l'on puisse trouver une petite partie
cofinale dans l'ensemble $\Ob C_{/F}$, ce qui est le cas en
particulier si $F$ est ind-représentable; alors il s'ensuit que la
même condition est vérifiée dans la sous-catégorie pleine $I$
envisagée dans \hyperref[I.prop8.12.2]{8.12.2}, donc on peut dans le critère
\hyperref[I.prop8.12.2]{8.12.2} omettre la condition que $I$ soit essentiellement
petite.

\subsubsection{}\etiquette[8.12.3]{I.sec8.12.3}
Soit $F$ un préfaisceau sur $C$, et considérons un objet
$$
u:X\longrightarrow F
$$
de $C_{/F}$. On dit que $u$ (ou le couple $(X,u)$) est
\emph{minimal} si pour toute factorisation de $u$ en
\begin{equation}
X\xrightarrow{p}X'\xrightarrow{u'}F{\quoi},\tag{8.12.3.1}
\end{equation}
avec $p$ un épimorphisme strict (\hyperref[I.sec10.2]{10.2}), $p$ est un isomorphisme.
Considérant $u$ comme un objet de $F(X)$ (1.4), dire que $u$ est
minimal signifie donc que\pageoriginale tout épimorphisme strict
$p:X\rightarrow X'$ tel que $u\in \Iim(F(p):F(X')\rightarrow F(X))$
est un isomorphisme. Cette notion s'éclaire par la partie a) du
lemme suivant:

\begin{enonce*}{Lemme 8.12.4}\etiquette[8.12.4]{I.lem8.12.4}
Soit $F$ un préfaisceau sur $C$.
\begin{enumerate}
\renewcommand{\theenumi}{\alph{enumi}}
\renewcommand{\labelenumi}{\theenumi)}
\item Supposons que $F$ transforme conoyaux en noyaux et considérons
un morphisme
$$
u:X\longrightarrow F{\quoi},
$$
avec $X\in \Ob C$. Pour que $u$ soit un monomorphisme, il suffit,
lorsque dans $C$ les conoyaux de doubles flèches sont
représentables, que $u$ soit minimal; cette condition est également
nécessaire si dans $C$ les produits fibrés sont représentables.

\item Supposons que dans $C$ les $\varinjlim$ finies soient
représentables. Pour que %%
la sous-catégorie des sous-foncteurs représentables de $F$
(\hyperref[I.8.12.1.1]{8.12.1.1}) soit filtrante (pas nécessairement petite) et ait
pour limite inductive $F$, il suffit que $F$ soit exact à gauche et
que tout morphisme $u:X\rightarrow F$, avec $X\in \Ob C$, se
factorise en
$$
X\rightarrow X'\xrightarrow{u'}F{\quoi},
$$
avec $u'$ minimal; cette condition est également nécessaire si dans
$C$ les produits fibrés sont représentables.

\item Supposons que $C$ admette une petite sous-catégorie
génératrice (\hyperref[I.def7.1]{7.1}), et que $I$ la catégorie des sous-foncteurs
représentables de $F$ soit filtrante, alors $I$ est petite.
\end{enumerate}
\end{enonce*}

\noindent \emph{Démonstration}.\pageoriginale
\begin{enumerate}
\renewcommand{\theenumi}{\alph{enumi}}
\renewcommand{\labelenumi}{\theenumi)}
\item Supposons $u$ minimal, prouvons que $u$ est un monomorphisme,
\ie que pour toute double-flèche $v$, $v':Y\rightrightarrows X$
telle que $uv=uv'$, on a $v=v'$. En effet, si $X'=\Coker(v,v')$,
alors $u$ se factorise en $X\xrightarrow{p}X'\xrightarrow{u'}F$ ($F$
transformant conoyaux en noyaux), et comme $p$ est un épimorphisme
strict par construction, il s'ensuit que $p$ est un isomorphisme \ie
$v=v'$. Supposons que $u$ est un monomorphisme, prouvons qu'il est
minimal. Considérons une factorisation (8.12.3.1); comme $p$ est un
épimorphisme strict, c'est le conoyaux de la double flèche canonique
$v$, $v':\fprod{X}{X}{X'}\rightrightarrows X$, et comme
$(u'p)v=(u'p)v'$ et que $u'p=u$ est un monomorphisme, on a $v=v'$
donc $p$ est un isomorphisme.

\item Suffisance: Comme $F$ est exact à gauche, $C_{/F}$ est
filtrante. Comme on sait que $F=\varinjlim_{C_{/F}}X$, on est ramené
par \hyperref[I.prop8.1.3]{8.1.3} c) à prouver que la sous-catégorie pleine $I$
de $C_{/F}$ des sous-objets de $F$ est cofinale dans $C_{/F}$; or en
vertu du « il suffit » dans a), c'est ce qu'assure l'hypothèse
que tout objet de $C_{/F}$ est majoré par un objet « minimal ».
Nécessité: Comme toute limite inductive filtrante de foncteurs
exacts à droite est itou, la première condition est trivialement
nécessaire. La deuxième résulte alors du « il faut » dans a).

\item Un sous-foncteur représentable $X\hookrightarrow F$ de $F$ est
connu quand on connaît la sous-catégorie pleine $C'_{/X}$ de
$C'_{/F}$, où $C$ est une petite sous-catégorie génératrice fixée de
$C$. (Utiliser l'hypothèse $I$ filtrante.) Comme $C'_{/F}$ est
petite, l'ensemble de ses sous-catégories pleine est petit, d'où la
conclusion.
\end{enumerate}

\begin{enonce*}{Proposition 8.12.5}\etiquette[8.12.5]{I.prop8.12.5}
Soit\pageoriginale $C$ une $𝒰$-catégorie où les limites
inductives finies sont représentables, et admettant une petite
sous-catégorie génératrice (\hyperref[I.def7.1]{7.1}). Soit $F$ un
préfaisceau sur $C$.
Pour que $F$ soit strictement ind-représentable, il suffit qu'il
satisfasse les deux conditions suivantes, et celles-ci sont
également nécessaires si dans $C$ les produits fibrés sont
représentables:
\begin{enumerate}
\renewcommand{\theenumi}{\alph{enumi}}
\renewcommand{\labelenumi}{\theenumi)}
\item $F$ est exact à gauche.

\item Tout couple $(X,u)$, avec $X\in \Ob C$ et $u\in F(X)$, est
majoré dans $C_{/F}$ par un couple minimal $(8.12.3)$, \ie il existe
un couple minimal $(X',u')$ $(X'\in \Ob C, u'\in F(X))$ et un
morphisme $f:X\rightarrow X'$ tel que $u=F(f)(u')$.
\end{enumerate}
\end{enonce*}

La suffisance résulte de \hyperref[I.prop8.12.2]{8.12.2} et de \hyperref[I.lem8.12.4]{8.12.4}
b), c), la nécessité de \hyperref[I.prop8.3.1]{8.3.1} et de \hyperref[I.lem8.12.4]{8.12.4}
a).

\begin{enonce*}[remark]{Remarque 8.12.6}\etiquette[8.12.6]{I.rem8.12.6}
Lorsque $F$ transforme sommes amalgamées de $C$ en produits fibrés,
on voit aussitôt que pour un couple $(X,u)$ donné comme dans b),
l'ensemble des quotients stricts $X'$ de $X$ tels que $u\in
\Iim(F(X')\rightarrow F(X))$ est filtrant décroissant, ce qui
implique que si l'ensemble des quotients stricts de $X$ est
artinien, alors la condition b) de \hyperref[I.prop8.12.5]{8.12.5} est
automatiquement satisfaite. Si donc la condition précédente sur $X$
est satisfaite pour tout objet $X$ de $C$ (on dit aussi alors,
parfois, que les objets de $C^\circ$ sont \emph{artiniens}), alors
il résulte de 8.12.5 que $F$ est strictement ind-représentable si et
seulement si $F$ est exact à gauche; dans ce cas, $F$ est donc
strictement ind-représentable dès qu'il est ind-représentable
(\hyperref[I.prop8.3.1]{8.3.1}).
\end{enonce*}

\begin{enonce*}{Corollaire 8.12.7}\etiquette[8.12.7]{I.cor8.12.7}
Soit\pageoriginale $C$ une $𝒰$-catégorie où les petites
limites projectives sont représentables, et admettant une petite
sous-catégorie cogénératrice (\hyperref[I.sec7.9]{7.9},
\hyperref[I.exem7.13]{7.13}). Alors un foncteur
$F:C\rightarrow (\Ens)$ est représentable si et seulement si $F$
commute aux petites limites projectives.
\end{enonce*}
La nécessité est claire. Pour la suffisance, on applique
\hyperref[I.prop8.12.5]{8.12.5}  à la catégorie opposée $C^\circ$. Pour prouver
d'abord que $F$ est (strictement) pro-représentable, on est ramené à
prouver que tout couple $(X,u)$, $X\in \Ob C$ et $u\in F(X)$, est
majoré par un couple « minimal ». Or la famille $(X_i)_{i\in I}$
des sous-objets \emph{stricts} de $X$ tels que $u\in
\Iim(F(X_i)\rightarrow F(X))$ est petite (\hyperref[I.cor7.5]{7.5} sous forme
duale). Grâce au fait que $F$ est exact à gauche, elle est
co-filtrante (\hyperref[I.rem8.12.6]{8.12.6}), et grâce au fait que $F$ commute
aux petites limites projectives on voit que $X'=\varprojlim_i X_i$
est un plus petit objet de cette famille. (Utiliser le fait que, $F$
étant exact à gauche, transforme monomorphismes en monomorphismes.)
Si $u'\in F(X')$ est l'unique élément dont l'image dans $F(X)$ est
$u$, on voit alors que $(X',u')$ est un couple minimal majorant
$(X,u)$. Cela prouve que $F$ est pro-représentable, et la conclusion
résulte alors du

\begin{enonce*}{Lemme 8.12.8.1}\etiquette[8.12.8.1]{I.lem8.12.8.1}
Soit $F:C\rightarrow(\Ens)$ un foncteur, où $C$ est une
$𝒰$-catégorie où les petites limites projectives sont
représentables. Pour que $F$ soit représentable, il faut et il
suffit qu'il soit pro-représentable et qu'il commute aux petites
limites projectives.
\end{enonce*}

La\pageoriginale nécessité est claire, prouvons la suffisance. Dire
que $F$ est représentable signifie évidemment que ${}_F\backslash C$
admet un élément initial (en fait, les objets initiaux de
${}_F\backslash C$ sont précisément les isomorphismes $F\rightarrow
X$, \ie les données de représentation pour $F$). Or $F$ étant
proreprésentable, $({}_F\backslash C^\circ)$ est filtrante et
équivalente à une petite catégorie, donc
$X=\varprojlim_{({}_F\backslash C)^\circ} X$ est représentable dans
$C$, et comme $F$ commute à la limite envisagée, il s'ensuit que $X$
est un objet initial de ${}_F\backslash C$, \cqfd

\begin{enonce*}{Corollaire 8.12.8}\etiquette[8.12.8]{I.cor8.12.8}
Les hypothèses sur $C$ étant celles de $8.12.7$, soit
$f:C\rightarrow C'$ un foncteur de $C$ dans une
$𝒰$-catégorie $C'$. Pour que $f$ admette un adjoint à
gauche, il faut et il suffit que $f$ commute aux petites limites
projectives.
\end{enonce*}

Cela se ramène en effet trivialement à \hyperref[I.cor8.12.7]{8.12.7}, en appliquant cet
énoncé aux foncteurs composés de la forme $X\mapsto \Hom(%%
Y',f(X))$.

\begin{enonce*}[remark]{Exemples 8.12.9}\etiquette[8.12.9]{I.exer8.12.9}
Comme on a signalé dans \hyperref[I.exem7.13]{7.13}, les hypothèses sur
$E$ de \hyperref[I.cor8.12.7]{8.12.7} et
\hyperref[I.cor8.12.8]{8.12.8} sont vérifiées si $C$ est la catégorie des
$𝒰$-faisceaux d'ensembles sur un espace topologique $X\in
𝒰$ (ou plus généralement, sur un $𝒰$-site (II
(3.0.2)). Donnons un exemple instructif \footnote[1]{(dû à H.
BASS).} qui montre que l'hypothèse d'existence d'une petite
sous-catégorie cogénératrice $D$ de $C$ n'est pas surabondante dans
\hyperref[I.cor8.12.7]{8.12.7}. Prenons pour $C$ la catégorie des groupes
éléments de
$𝒰$. Soit $J$ l'ensemble des classes d'isomorphie de groupes
simples $\in C$, choisissons pour tout $j\in J$ un groupe simple
$G_j$ dans la classe de $j$, et soit $I$ l'ensemble ordonné filtrant
des parties $𝒰$-petites de $J$, et pour $i\in I$, soit
$X_i=\amalg_{j\in i}G_j$. Les $X_i$ forment alors un\pageoriginale
système projectif $(X_i)_{i\in I}$ dans $E$, à morphismes de
transition des épimorphismes, et le foncteur correspondant
\begin{equation}
F(X)=\varinjlim_i \Hom(X_i,X)\tag{$\ast$}
\end{equation}
prend ses valeurs dans $(𝒰\text{-}\Ens)$ (bien que
l'ensemble d'indices $I$ n'ait évidemment pas un cardinal $\in
𝒰$); plus précisément, montrons que pour toute petite
sous-catégorie pleine $C_\circ$ de $C$, il existe un $i_\circ\in I$
tel que la restriction de $F$ à $C_\circ$ soit représentable par
$X_{i_\circ}$, ce qui prouvera à la fois que $F$ est à valeurs dans
$𝒰\text{-}\Ens$, et qu'il commute aux petites limites
projectives. Pour prouver notre assertion, il suffit de noter que
pour tout $X\in \Ob C_\circ$, le cardinal de l'ensemble $J(X)$ des
$j\in J$ tels qu'il existe un homomorphisme non trivial de $G_j$
dans $X$ est nécessairement petit, puisque un tel morphisme est
nécessairement un monomorphisme ($G_j$ étant simple); par suite, si
$i_\circ$ est la partie de $J$ réunion des $J(X)$ pour $X\in \ob
C_\circ$, $i_\circ$ est petit \ie $i_\circ\in I$, et il fait
l'affaire. D'autre part il est clair que $F$ n'est pas
représentable, puisque on a $\card I\not\in 𝒰$. De ceci et
de 8.12.7 on conclut donc que \emph{la catégorie $C$ des groupes
$\in 𝒰$ n'admet pas une petite sous-catégorie pleine
cogénératrice.} Comme l'objet $\ZZ$ de $C$ est d'autre part un
générateur, il résulte alors de la démonstration de
\hyperref[I.cor7.12]{7.12} qu'il
existe un groupe $G$ à deux générateurs qui ne se plonge
pas dans un
objet injectif de la catégorie $C$ des groupes $\in 𝒰$. Il
semble d'ailleurs plausible que $C$ n'admette pas d'autre objet
injectif que les groupes unités.
\end{enonce*}

\subsection{Foncteurs proreprésentables et foncteurs
  accessibles}\etiquette[8.13]{I.sec8.13}\pageoriginale

\subsubsection{}\etiquette[8.13.1]{I.sec8.13.1}
Dans le présent numéro, nous utilisons quelques notions et résultats
du paragraphe suivant, et notamment \hyperref[I.prop9.11]{9.11} et
\hyperref[I.prop9.13]{9.13}, pour obtenir un critère de proreprésentabilité que
nous utiliserons (incidemment) dans IV 9.16. $C$ désigne par la
suite une $𝒰$-catégorie satisfaisant la condition $L$ de
\hyperref[I.def9.1]{9.1} b), cette condition étant remplie par exemple si dans
$C$ les petites limites inductives filtrantes sont représentables.

\begin{enonce*}{Proposition 8.13.2}\etiquette[8.13.2]{I.prop8.13.2}
Soient $C$ comme ci-dessus, et
$$
f:C\longrightarrow(\Ens)
$$
un foncteur.
\begin{enumerate}
\renewcommand{\theenumi}{\alph{enumi}}
\renewcommand{\labelenumi}{\theenumi)}
\item Supposons que chaque objet de $C$ est accessible (\hyperref[I.def9.3]{9.3}). Si $f$
est pro-représentable, $f$ est accessible (\hyperref[I.def9.2]{9.2}) et
exact à gauche.

\item Supposons que dans $C$ les limites projectives finies soient
représentables, et que $C$ admette une filtration cardinale
(\hyperref[I.def9.12]{9.12}). Si $f$ est accessible et exact à gauche, alors $f$ est
proreprésentable.
\end{enumerate}
\end{enonce*}

\noindent \emph{Démonstration.}
\begin{enumerate}
\renewcommand{\theenumi}{\alph{enumi}}
\renewcommand{\labelenumi}{\theenumi)}
\item L'hypothèse sur $C$ signifie que les foncteurs covariants
représentables de $C$ dans $(\Ens)$ sont accessibles. Il en est donc
de même de toute petite limite inductive de tels foncteurs
(\hyperref[I.prop9.6]{9.6} (i)), donc aussi de tout foncteur
pro-représentable.

\item En vertu de \hyperref[I.them8.3.3]{8.3.3} (iii), il reste à prouver que dans
$\Ob({}_F\backslash C)^\circ$ il y a une petite sous-catégorie
cofinale. Or par hypothèse il existe un cardinal $\Pi$ tel que $f$
soit $\Pi$-accessible. Soit alors $(X,u)$,\pageoriginale $u\in
F(X)$, un objet de ${}_F\backslash C$. Avec les notations de 9.12 c)
on a alors $X=\varinjlim_i X_i$, avec $I$ filtrant grand devant
$\Pi$ et les $X_i$ dans $C'=\Filt^{\Pi}(C)$, d'où
$F(X)\xleftarrow{\sim}\varinjlim_i F(X_i)$. Cela montre que la
petite sous-catégorie $({}_F\backslash C')^\circ$ est cofinale dans
$({}_F\backslash C)^\circ$, et achève la démonstration.
\end{enumerate}

\begin{enonce*}{Corollaire 8.13.3}\etiquette[8.13.3]{I.cor8.13.3}
Soit $C$ une $𝒰$-catégorie satisfaisant aux conditions
suivantes:
\begin{enumerate}
\renewcommand{\theenumi}{\alph{enumi}}
\renewcommand{\labelenumi}{\theenumi)}
\item Dans $C$ les limites projectives finies sont représentables.

\item Dans $C$ les petites limites inductives filtrantes sont
représentables.

\item Le foncteur $\Ker$ sur la catégorie des doubles flèches de
$C$ est accessible (par exemple, il commute aux petites $\varinjlim$
filtrantes).

\item Tout épimorphisme strict de $C$ est strict universel $(10.2)$.

\item Il existe une petite sous-catégorie de $C$ génératrice par
épimorphismes stricts $(7.1)$.
\end{enumerate}
Sous ces conditions, un foncteur $f:C\rightarrow (\Ens)$ est
proreprésentable si et seulement si il est exact à gauche et
accessible (\hyperref[I.def9.2]{9.2}).
\end{enonce*}

En effet, les conditions de \hyperref[I.prop8.13.2]{8.13.2} a) et b) sur $C$ sont
vérifiées, en vertu de \hyperref[I.prop9.11]{9.11} et \hyperref[I.prop9.13]{9.13}
respectivement.


\begin{enonce*}{Corollaire 8.13.4}\etiquette[8.13.4]{I.cor8.13.4}
Soit $f:C\rightarrow C'$ un foncteur entre $𝒰$-catégories
satisfaisant aux conditions {\rm a)} à {\rm e)} de \hyperref[I.cor8.13.3]{8.13.3}. Pour
que $f$ admette un pro-adjoint, il faut et il suffit que $f$ soit
exact à gauche et accessible.
\end{enonce*}

Comme\pageoriginale les foncteurs représentables
$h_{Y'}:X'\rightarrow \Hom(Y',X'):C'\rightarrow(\Ens)$ sont exacts à
gauche et accessibles (\hyperref[I.prop9.11]{9.11}), si $f$ a ces mêmes
propriétés, il en est de même de ses composés avec les foncteurs
précédents, qui sont donc proreprésentables par \hyperref[I.cor8.13.3]{8.13.3},
\ie $f$ admet un proadjoint. Inversement, supposons que $f$ admet un
proadjoint, et soient $C'_1$ une petite sous-catégorie génératrice
de $C$, $\Pi$ un cardinal tel que les $Y'\in \Ob C'$ soient
$\Pi$-accessibles; pour prouver que $f$ est $\Pi$-accessible, il
suffit donc de prouver qu'il en est ainsi de ses composés avec les
$h_{Y'}$, $Y'\in \Ob C'$; or par hypothèse ces composés sont
proreprésentables, donc ils sont accessibles (\hyperref[I.cor8.13.3]{8.13.3}),
donc $\Pi$-accessibles pourvu qu'on prenne $\Pi$ assez grand, \cqfd

\section{Foncteurs accessibles, filtrations cardinales et construction
de petites sous-catégories génératrices}\etiquette[9]{I.sec9}

Le présent paragraphe, de nature plus technique que les autres
paragraphes de cet exposé, ne servira dans ce séminaire que dans IV
9 et dans VI 4, qui ne sont pas utilisés ailleurs dans le Séminaire.
Il s'impose donc d'omettre la lecture du présent paragraphe, du
moins en première lecture!

\setcounter{subsection}{-1}
\subsection{}\etiquette[9.0]{I.sec9.0}
Toutes les catégories envisagées dans le présent numéro sont
supposées être des $𝒰$-catégories. Sauf pour les petites
catégories d'indices $I$, $J\ldots$ que nous aurons à utiliser, les
développements qui suivent s'appliqueront surtout à des « grosses »
catégories $E$, $F\ldots$ qui sont stables par petites limites
inductives filtrantes. Il suffira cependant le\pageoriginale plus
souvent qu'une condition un peu plus faible soit vérifiée (condition
$L$ dans \hyperref[I.def9.1]{9.1} ci-dessous). Tous les cardinaux envisagés dans le
présent numéro sont supposés $\in 𝒰$.

Suivant une suggestion de P. \textsc{Deligne}, nous allons étudier,
pour un foncteur $f:E\rightarrow F$ entre grosses catégories, une
condition de commutation de $f$ à certains types de limites
inductives filtrantes, condition remarquablement stable, et qui sera
vérifiée pour les foncteurs les plus importants qu'on rencontre dans
la nature. Les applications que nous avons en vue, pour notre
séminaire, sont 9.13.3, 9.13.4 (utilisés dans VI
4) et surtout \hyperref[I.cor9.25]{9.25},
qui donne, dans un cas non trivial, l'existence d'une petite famille
génératrice dans une catégorie de sections d'une
catégorie fibrée;
ce résultat sera utilisé dans IV 9.16.

\begin{enonce*}{Définition 9.1}\etiquette[9.1]{I.def9.1}
\begin{enumerate}
\renewcommand{\theenumi}{\alph{enumi}}
\renewcommand{\labelenumi}{\theenumi)}
\item Soient $I$ un ensemble préordonné, $\pi$ un cardinal. On
dit que $I$ est grand devant $\pi$ si $I$ est filtrant, et si toute
partie de $I$ de cardinal $\leq \pi$ admet un majorant dans $I$.

\item Soit $E$ une catégorie. Si $\pi$ est un cardinal, on dit que
$E$ satisfait la condition $L_{\pi}$ si pour tout petit ensemble
ordonné $I$ grand devant $\pi$, $E$ est stable par les limites
inductives de type $I$. On dit que $E$ satisfait à la condition $L$
s'il existe un cardinal $\pi\in 𝒰$ tel que $E$ satisfasse à
la condition $L_{\pi}$.
\end{enumerate}
\end{enonce*}

\setcounter{subsection}{1}
\subsubsection{}\etiquette[9.1.1]{I.sec9.1.1}
Lorsque dans \hyperref[I.def9.1]{9.1} a) on a $\pi \geq 2$, la deuxième condition
énoncée
implique déjà que $I$ est filtrant, et si $\pi$ est fini, $I$ grand
devant $\pi$ signifie simplement que $I$ est filtrant. Nous ne nous
intéresserons guère par la suite qu'au cas où $\pi$ est infini.
Notons que si $\pi$, $\pi'$ sont deux cardinaux tels que $\pi'\geq
\pi$, alors $I$ grand devant $\pi'$ implique évidemment $I$ grand
devant $\pi$.

\subsubsection{}\etiquette[9.1.2]{I.sec9.1.2}
Comme\pageoriginale annoncé dans \hyperref[I.sec9.0]{9.0}, les conditions
$L_{\pi}$, $L$
doivent être considérées comme des variantes techniques de la
condition plus forte de stabilité par petites limites inductives
filtrantes. Il est clair que si $\pi$, $\pi'$ sont des cardinaux
tels que $\pi'\geq \pi$, alors la condition $L_{\pi}$ implique la
condition $L_{\pi'}$.

\begin{enonce*}{Définition 9.2}\etiquette[9.2]{I.def9.2}
Soit $f:E\rightarrow F$ un foncteur. Si $\Pi$ est un cardinal, on
dit que $f$ est $\Pi$-préaccessible (\resp $\Pi$-accessible) si $E$
satisfait $L_{\Pi}$ (\hyperref[I.def9.1]{9.1}) et si pour tout ensemble ordonné $I\in
𝒰$ grand devant $\Pi$, et tout système inductif $(X_i)_{i\in
I}$ dans $E$ de type $I$, le morphisme canonique
$$
\varinjlim f(X_i)\longrightarrow f(\varinjlim X_i)
$$
est un monomorphisme (\resp un isomorphisme). On dit que $f$ est
préaccessible \\ (\resp accessible) (relativement à l'univers
$𝒰$) s'il existe un cardinal $\Pi\in 𝒰$ tel que $f$
soit $\Pi$-préaccessible (\resp $\Pi$-accessible).
\end{enonce*}

La catégorie des foncteurs $\Pi$-accessibles (\resp accessibles) de
$E$ dans $F$ sera noté $\Hhom(E,F)_{\Pi}$ \resp $\Hhom(E,F)_{\acc}$.

\setcounter{subsection}{2} \setcounter{subsubsection}{0}
\subsubsection{}\etiquette[9.2.1]{I.sec9.2.1}
Évidemment, un foncteur commutant aux petites $\varinjlim$
filtrantes (p. ex. un foncteur admettant un adjoint à droite) est
$\Pi$-accessible pour tout cardinal $\Pi\geq 2$.

\begin{enonce*}{Définition 9.3}\etiquette[9.3]{I.def9.3}
Soient $E$ une catégorie, $X$ un objet de $E$,
$$
h^\circ_X:E\longrightarrow (𝒰\text{-}\Ens)
$$
le foncteur covariant qu'il représente, $\Pi\in 𝒰$ un
cardinal. On dit que $X$ est un \emph{objet $\Pi$-préaccessible
(\resp $\Pi$-accessible) de} $E$ si le foncteur
$h^\circ_X$\pageoriginale est $\Pi$-préaccessible (\resp
$\Pi$-accessible); on dit que $X$ est \emph{préaccessible (\resp
accessible)} s'il existe un cardinal $\Pi\in 𝒰$ tel que $X$
soit un objet $\Pi$-préaccessible (\resp $\Pi$-accessible) de $E$.
\end{enonce*}

\setcounter{subsection}{3} \setcounter{subsubsection}{0}
\subsubsection{}\etiquette[9.3.1]{I.sec9.3.1}
On désigne par $E_{\Pi}$ la sous-catégorie pleine de $E$ formée des
objets $\Pi$-accessibles de $E$.

Lorsqu'on applique la définition 9.3 à une catégorie de la forme
$\Hhom(C,F)$, la terminologie introduite présente \textit{a priori}
une ambiguïté avec la terminologie analogue introduite dans 9.2,
lorsqu'on interprète les objets de $\Hhom(C,F)$ comme des foncteurs;
il ne semble pas cependant qu'il y ait un risque de confusion
sérieux.

\begin{enonce*}{Définition 9.4}\etiquette[9.4]{I.def9.4}
Soient $E$ une catégorie, $\pi\in 𝒰$ un cardinal. On dit que
$E$ est une catégorie $\pi$-préaccessible (\resp $\pi$-accessible)
s'il existe dans $E$ une petite sous-catégorie pleine $C$ qui est
génératrice (\hyperref[I.def7.1]{7.1}) et dont les objets sont
$\pi$-préaccessibles
(\resp $\pi$-accessibles) (\hyperref[I.def9.3]{9.3}). On dit que $E$ est une
catégorie
pré-accessible (\resp. accessible) s'il existe un cardinal $\pi\in
𝒰$ tel que $E$ soit $\pi$-préaccessible (\resp
$\pi$-accessible).
\end{enonce*}

Pour des exemples importants, \cf \hyperref[I.rem9.11.3]{9.11.3} plus bas.

\begin{enonce*}{Proposition 9.5}\etiquette[9.5]{I.prop9.5}
Soit $f:E\rightarrow F$ un foncteur entre catégories telles que $F$
soit accessible et que tout objet de $E$ soit accessible. Alors, si
$f$ admet un adjoint à gauche, $f$ est accessible.
\end{enonce*}

En\pageoriginale effet, par hypothèse, $F$ admet une petite famille
conservative de foncteurs représentables $F\rightarrow
(𝒰\text{-}\Ens)$ qui sont accessibles, donc on est ramené
aussitôt à montrer que le composé de $f$ avec chacun des foncteurs
précédents est accessible. Or comme $f$ admet un adjoint à gauche,
ces composés sont des foncteurs $E\rightarrow
(𝒰\text{-}\Ens)$ \emph{représentables,} donc accessibles
d'après l'hypothèse sur $E$.

\begin{enonce*}{Proposition 9.6}\etiquette[9.6]{I.prop9.6}
Soient $E$ et $F$ deux catégories, $\pi\in 𝒰$ un cardinal et
considérons la sous-catégorie pleine $\Hhom(E,F)_{\Pi}$ de
$\Hhom(E,F)$ formée des foncteurs $\Pi$-accessibles $(9.2)$ de $E$
dans $F$.
\begin{enumerate}
\renewcommand{\theenumi}{\roman{enumi}}
\renewcommand{\labelenumi}{(\theenumi)}
\item Cette sous-catégorie est stable par tout type de $\varinjlim$
qui est représentable dans~$F$.

\item Supposons $F$ $\pi$-accessible (\hyperref[I.def9.4]{9.4}), et soit $J$ une
catégorie telle que $\card  F\ell \  J\leq \pi$ et que les
$\varprojlim$ de type $J$ soient représentables dans $F$, donc les
$\varprojlim$ de type $J$ sont représentables dans $\Hhom(E,F)$.
Alors la sous-catégorie $\Hhom(E,F)_{\Pi}$ est stable par les
$\varprojlim$ de type $J$.
\end{enumerate}
\end{enonce*}

\begin{enonce*}{Corollaire 9.7}\etiquette[9.7]{I.cor9.7}
Soient $E$ et $F$ deux catégories, avec $F$ accessible (\hyperref[I.def9.4]{9.4}).
Alors la sous-catégorie pleine $\Hhom(E,F)_{\acc}$ de $\Hhom(E,F)$
formée des foncteurs accessibles est stable par tout type de limite
inductive ou projective, relative à une petite catégorie d'indices
$J$, qui est représentable dans $F$ (donc dans $\Hhom(E,F)$).
\end{enonce*}

Preuve de \hyperref[I.prop9.6]{9.6}. L'assertion (i) résulte trivialement de
la commutation du foncteur $\varinjlim$ aux limites inductives
quelconques. Pour (ii), soit $(f_j)_{j\in J}$ un système projectif
de foncteurs $\Pi$-accessibles $E\rightarrow F$, $f=\varprojlim
f_j$\pageoriginale sa limite projective, qui se calcule « argument
par argument », prouvons que $f$ est $\Pi$-accessible, \ie que
pour tout ensemble ordonné $I\in 𝒰$ grand devant $\Pi$, et
tout système inductif $(X_i)_{i\in I}$ dans $E$, le morphisme
canonique
$$
\varinjlim_i ((\varprojlim_j f_j)(X_i))\longrightarrow
(\varprojlim_j f_j) (\varinjlim_i X_i)
$$
est un isomorphisme. Or le calcul « argument par argument » du
foncteur $\varprojlim_j f_j$ nous permet d'identifier le morphisme
canonique précédent au morphisme canonique
$$
\varinjlim_i \varprojlim_j f_j (X_i)\longrightarrow \varprojlim_j
\varinjlim_i f_j (X_i)
$$
associé au bifoncteur
$$
(i,j)\longmapsto f_j(X_i):I\times J\longrightarrow F{\quoi}.
$$
Donc \hyperref[I.prop9.6]{9.6} est une conséquence de l'assertion plus
générale:

\begin{enonce*}{Corollaire 9.8}\etiquette[9.8]{I.cor9.8}
Soient $F$ une catégorie $\Pi$-accessible, $I$ un ensemble ordonné
grand devant $\Pi$, $J$ une petite catégorie telle que $\card F\ell
\leq \Pi$, et que les $\varprojlim$ de type $J$ soient
représentables dans $F$, $h:I\times J\rightarrow F$ un foncteur;
alors le morphisme canonique
\begin{equation}
\varinjlim_i \varprojlim_j h(i,j)\longrightarrow \varprojlim_j
\varinjlim_i h(i,j)\tag{9.8.1}
\end{equation}
est un isomorphisme. En d'autres termes, le foncteur
\begin{equation}
\varinjlim_i:\Hhom(I,F)\longrightarrow F\tag{9.8.2}
\end{equation}
commute aux limites projectives de type $J$ (pour toute petite
catégorie $J$\pageoriginale telle que $\card  F\ell\, (J)\leq \Pi$
et telle que les limites projectives de type $J$ soient
représentables dans $F$). Ou encore, pour toute $J$ comme ci-dessus,
le foncteur
\begin{equation}
\varprojlim_j:\Hhom(J,F)\longrightarrow F\tag{9.8.3}
\end{equation}
est $\Pi$-accessible.
\end{enonce*}

Comme par hypothèse $F$ admet une famille conservative de foncteurs
covariants représentables $\Pi$-accessibles $g:F\rightarrow
(𝒰\text{-}\Ens)$, on voit aussitôt qu'on est ramené à
prouver 9.8 dans le cas où $F=𝒰\text{-}\Ens$. Nous
prouverons alors l'assertion sous la forme de la commutation de
(9.8.2) aux $\varprojlim$ de type $J$, avec $\card F\ell\,  J\leq
\Pi$. Comme $I$ est filtrant, nous savons que (9.8.2) commute aux
limites projectives finies (2.8). On est donc ramené par un argument
standard (\cf 2.3) à prouver qu'il commute aux produits indexés par
un ensemble $J$ tel que $\card J\leq \Pi$. Cela nous ramène à
prouver la bijectivité de (9.8.1) lorsque $J$ est discrète. Prouvons
l'injectivité: considérons deux éléments $a$, $b$ du
premier membre,
ils proviennent donc de $\prod_j h(i_\circ, j)$ %%
pour $i_\circ \in I$ convenable, soient $\prod_j a(i_\circ, j)$ et
$\prod_j b(i_\circ, j)$; supposons que les éléments $\prod_j
h(\infty, j)$ du deuxième membre qu'ils définissent soient égaux,
\ie que pour tout $j$, il existe $i(j)\geq i_\circ$ tel que
$a(i_\circ, j)$ et $b(i_\circ, j)$ aient même image dans $h(i,j)$.
Comme $I$ est grand devant $\card J$, il s'ensuit qu'il existe un
majorant commun $i_1\in I$ de tous les $i(j)$, ce qui implique que
les deux éléments envisagés de $\prod_j h(i_\circ, j)$ ont même
image dans $\prod_j h(i_1,j)$, donc définissent %%
le même élément du premier membre de (9.8.1), ce qui établit
l'injectivité.\pageoriginale Pour la surjectivité, soit $\prod_j
a(\infty, j)$ un élément du second membre; donc pour tout $j\in J$,
$a(\infty, j)$ provient d'un élément de $h(i(j),j)$, pour un
$i(j)\in I$ convenable. Comme précédemment, on peut trouver un
majorant commun $i\in I$ des $i(j)$, donc $\prod_j a(\infty, j)$
provient d'un élément $\prod_j a(i,j)$ de $\prod_j h(i,j)$, donc est
dans l'image de (9.8.1). Cela achève la démonstration.

\begin{enonce*}{Corollaire 9.9}\etiquette[9.9]{I.cor9.9}
Soit $F$ une catégorie $\Pi$-accessible (\resp accessible), alors
pour toute catégorie $J$ telle que $\card F\ell\,  J\leq \Pi$ (\resp
toute petite catégorie $J$) et pour tout foncteur $(X_j)_{j\in
J}:J\rightarrow F$ dont la limite inductive dans $F$ est
représentable, si les $X_j$ sont $\Pi$-accessibles (\resp
accessibles) il en est de même de $\varprojlim X_j$.
\end{enonce*}

En effet, le foncteur $F\rightarrow(𝒰\text{-}\Ens)$
représente par $\varinjlim X_j$ est la limite projective des
foncteurs représentés par les $X_j$, et on applique \hyperref[I.prop9.6]{9.6}
(ii) au système projectif formé par ces foncteurs.

\begin{enonce*}[remark]{Remarque 9.10}\etiquette[9.10]{I.rem9.10}
Dans les énoncés \hyperref[I.prop9.6]{9.6}, \hyperref[I.cor9.7]{9.7}, \hyperref[I.cor9.8]{9.8} et
\hyperref[I.cor9.9]{9.9} on peut remplacer partout les mots « $\Pi$-accessibles », « accessibles » par
« $\Pi$-préaccessible », « préaccessibles ». La démonstration
donnée prouve en effet également cette variante des énoncés
précédents.
\end{enonce*}

\begin{enonce*}{Proposition 9.11}\etiquette[9.11]{I.prop9.11}
Soit $E$ une catégorie, satisfaisant la condition $L$ (\hyperref[I.def9.1]{9.1}), $C$
une petite sous-catégorie pleine génératrice (\hyperref[I.def7.1]{7.1}). On suppose
satisfaite la condition:
\begin{enumerate}
\renewcommand{\theenumi}{\alph{enumi}}
\renewcommand{\labelenumi}{\theenumi)}
\item $E$ est stable par noyaux de doubles flèches, et le foncteur
$\Ker$ sur la catégorie des doubles flèches de $E$ est accessible
(par exemple, il\pageoriginale commute aux petites limites
inductives filtrantes).

Alors tout objet de $E$ est préaccessible (\hyperref[I.def9.3]{9.3}), \textit{a
fortiori} $E$ est préaccessible. Supposons que $C$ soit même
génératrice par épimorphisme stricts (\hyperref[I.def7.1]{7.1}), et
supposons que $E$
satisfasse aux conditions:

\item $E$ est stable par produits fibrés.

\item Toute petite famille épimorphique stricte dans $E$ est
épimorphique stricte universelle (\hyperref[I.sec10.3]{10.3}), ou dans $E$ les petites
sommes directes sont représentables, et tout épimorphisme strict de
$E$ est un épimorphisme strict universel.
\end{enumerate}

Alors tout objet de $E$ est accessible, \textit{a fortiori} $E$ est
accessible.
\end{enonce*}

Le fait que, moyennant (a), tout objet de $E$ soit accessible,
résulte du fait que pour tout objet $X$ de $E$, l'ensemble des
sous-objets stricts de $E$ est petit (\hyperref[I.prop7.4]{7.4}), et du

\begin{enonce*}{Lemme 9.11.1}\etiquette[9.11.1]{I.lem9.11.1}
Sous les conditions de (\hyperref[I.prop9.11]{9.11}a), soient $X\in \ob E$ et
$\Pi$ un cardinal tels que $E$ satisfasse $L_{\Pi}$
(\hyperref[I.def9.1]{9.1}), que le foncteur $\Ker$ dans (\hyperref[I.prop9.11]{9.11}a)
soit $\Pi$-accessible, et que l'ensemble des sous-objets stricts de
$X$ soit de cardinal majoré par $\Pi$. Alors $X$ est
$\Pi$-préaccessible.
\end{enonce*}

En effet, si $I$ est un ensemble grand devant $I$, $(Y_i)_{i\in I}$
un système inductif dans $E$ de limite inductive $Y$, et $(u_i,
v_i:X\rightrightarrows Y_i)$ une double flèche telle que la double
flèche composée $ \xymatrix{ X \ar@<2pt>[r]^{u,v} \ar@<-2pt>[r]& Y
}$ satisfasse $u=v$, prouvons qu'il existe $j\geq i$ dans $I$ tel
que $u_j=v_j$. Pour ceci,\pageoriginale considérons le système
inductif des doubles flèches $(u_j,v_j:X\rightrightarrows
Y_j)_{j\geq i}$, dont la limite inductive est $(u,v)$. Par hypothèse
on a $X=\Ker(u,v)=\varinjlim_j \Ker (u_j,v_j)=\varinjlim_j X_j$, où
$X_j=\Ker(u_j,v_j)$. Or les $X_j$ sont des sous-objets stricts de
$X$, donc il existe une partie $I'$ de $I$ formée d'indices $j\geq
i$, telle que $\card J\leq\Pi$ et que tout $X_j$ soit égal à un des
$X_{i'}$ $(i'\in I')$. Comme $I$ est grand devant $\Pi$, il existe
un majorant $j$ de $I'$ dans $I$. Alors $X_j$ contient tous les
$X_{j'}$ pour $j'\geq i$, donc $X_j\rightarrow X$ est un
épimorphisme (puisque la famille des $X_{j'}\rightarrow X$ est
épimorphique), donc un isomorphisme puisque c'est un monomorphisme
strict. Donc on a $u_j=v_j$, ce qui prouve \hyperref[I.lem9.11.1]{9.11.1}.

La deuxième assertion de \hyperref[I.prop9.11]{9.11} résulte de la première, et
du

\begin{enonce*}{Lemme 9.11.2}\etiquette[9.11.2]{I.lem9.11.2}
Sous les conditions de (\hyperref[I.prop9.11]{9.11}a), b), c), soient $X$ un
objet de $E$, et $\Pi$ un cardinal infini tels que $E$ satisfasse
$L_{\Pi}$ $(9.1)$, que l'on ait $\card \ob C_{/X}\leq \Pi$, que pour
deux objets $X'\rightarrow X$ et $X''\rightarrow X$ de $C_{/X}$,
$\fprod{X'}{X''}{X}$ soit $\Pi$-préaccessible, enfin que $X$ soit
$\Pi$-préadmissible. Alors $X$ est un $\Pi$-accessible.
\end{enonce*}

Avec les notations de la démonstration de \hyperref[I.lem9.11.1]{9.11.1}, il suffit de
prouver que tout morphisme $u:X\rightarrow Y$ provient d'un
morphisme $u_i:X\rightarrow Y_i$. Or considérons la famille des
morphismes $Y_i\rightarrow Y$, qui est épimorphique stricte; grâce à
b) et c), la famille des $\fprod{Y_i}{X}{Y} \rightarrow X$ est
également épimorphique stricte; d'autre part, pour tout $i$, comme
$C$ est génératrice par épimorphismes stricts, la famille des
flèches
$$
T\longrightarrow \fprod{Y_i}{X}{Y}
$$
de\pageoriginale source dans $C$ est épimorphique stricte. Dans la
deuxième alternative envisagée dans c), on peut trouver une telle
flèche épimorphique stricte, de source, une (petite) somme d'objets
de $C$. En vertu de la transitivité de la notion de famille
épimorphique stricte universelle (II 2.5) il s'ensuit alors que la
famille des flèches $T_\alpha\xrightarrow{f_\alpha}X$ de source dans
$C$ qui se factorisent par un des $\fprod{Y_i}{X}{Y}$ est
épimorphique stricte. Soit $J$ l'ensemble d'indices de cette
famille, qui est de cardinal majoré par $\card \ob C_{/X}\leq \Pi$.
Choisissons pour tout $\alpha\in J$ un $i=i(\alpha)\in I$ et un
$X$-morphisme $T_\alpha\rightarrow \fprod{Y_i}{X}{Y}$, ou ce qui
revient au même, un $v_\alpha:T_\alpha\rightarrow Y_i$ tel que
$p_iv_\alpha=uf_\alpha$, où $p_i:Y_i\rightarrow Y$ est le morphisme
canonique. Comme $I$ est grand devant $\Pi$, on peut choisir
$i(\alpha)$ indépendant de $\alpha$, soit $i$. Pour tout couple
d'indices $\alpha$, $\beta\in J$, considérons les composés
$$
\pr_1 v_\alpha, \pr_2
v_\beta:\fprod{T_\alpha}{T_\beta}{X}\rightrightarrows Y_i{\quoi}.
$$
Leurs composés avec $Y_i\rightarrow Y$ sont égaux, donc, comme
$\fprod{T_\alpha}{T_\beta}{X}$ est $\Pi$-préaccessible par
hypothèse, il existe un indice $i'=i(\alpha,\beta)\geq i$ tel que
les composés des flèches envisagées avec $Y_i\rightarrow Y_{i'}$,
soient égales. Comme l'ensemble des couples $\alpha$, $\beta$ est de
cardinal $\leq \Pi^2=\Pi$ ($\Pi$ étant infini), il s'ensuit encore
que l'on peut choisir $i'$ indépendant de $\alpha$, $\beta$. On peut
évidemment supposer $i'=i$. Mais alors, la famille
$(f_\alpha:T_\alpha\rightarrow X)$ étant épimorphique stricte, on
peut trouver un morphisme $u_i:X\rightarrow Y_i$ tel que l'on ait
$u_if_\alpha=v_i$. On a alors $p_iu_i=p$, car pour tout $\alpha$ on
a $(p_iu_i)f_\alpha=p_i(u_if_\alpha)=p_iv_\alpha=uf_\alpha$, et la
famille des $f_\alpha$ est épimorphique. Cela achève la
démonstration de \hyperref[I.lem9.11.2]{9.11.2}.

\begin{enonce*}[remark]{Remarque 9.11.3}\etiquette[9.11.3]{I.rem9.11.3}
On\pageoriginale voit, comme cas particulier de \hyperref[I.prop9.11]{9.11},
que dans la catégorie $E$ tout objet est accessible, dans chacun des
deux cas suivants: 1) $E$ est une $𝒰$-catégorie abélienne à
limites inductives filtrantes exactes et admettant une petite
sous-catégorie génératrice (ici, c'est la \emph{deuxième}
alternative de c), qui s'applique). 2) $E$ est la catégorie des
faisceaux d'ensembles sur un espace topologique $X\in 𝒰$.
Plus généralement, il suffit que $E$ soit la catégorie des faisceaux
d'ensembles sur un $𝒰$-site (II 2.1), ou encore, que $E$
soit un $𝒰$-topos (IV 1.1).
\end{enonce*}

\begin{enonce*}{Définition 9.12}\etiquette[9.12]{I.def9.12}
Soit $E$ une catégorie. On appelle filtration cardinale de $E$ une
filtration croissante $(\Filt^{\Pi}(E))_{\Pi\geq \Pi_\circ}$ de $E$
par des sous-catégories strictement pleines $\Filt^{\Pi}(E)$,
indexée par les cardinaux $\Pi\in 𝒰$ tels que $\Pi\geq
\Pi_\circ$ (où $\Pi_\circ$ est un cardinal infini fixé, dépendant de
la filtration cardinale envisagée), et satisfaisant aux conditions
suivantes:
\begin{enumerate}
\renewcommand{\theenumi}{\alph{enumi}}
\renewcommand{\labelenumi}{\theenumi)}
\item Pour tout $\Pi\geq\Pi_\circ$, $\Filt^{\Pi}(E)$ est
équivalente à une petite catégorie.

\item $E$ satisfait à $L_{\Pi_\circ}$ (\hyperref[I.def9.3]{9.3}), et pour tout
$\Pi\geq \Pi_\circ$, $\Filt^{\Pi}(E)$ est stable dans $E$ pour les
limites inductives filtrantes indexées par des ensembles ordonnés
$I$ grands devant $\Pi_\circ$ tels que $\card I\leq \Pi$.

\item Pour tout $\Pi\geq \Pi_\circ$, et tout $X\in \ob E$, on peut
trouver un isomorphisme
$$
X\simeq \varinjlim_I X_i{\quoi},
$$
où $(X_i)_{i\in I}$ est un système inductif dans $E$ indexé par un
ensemble ordonné $I$\pageoriginale grand devant $\Pi$
(\hyperref[I.def9.1]{9.1}), et où
les $X_i$ sont dans $\Filt^{\Pi}(E)$. De plus, si $X\in \ob
\Filt^{\Pi}(E)$, $\Pi'\geq \Pi$, on peut prendre $I$ tel que $\card
I\leq {\Pi'}^{\Pi}$.
\end{enumerate}
\end{enonce*}

\setcounter{subsection}{12} \setcounter{subsubsection}{0}
\subsubsection{}\etiquette[9.12.1]{I.9.12.1}
On notera qu'il résulte de b) et c) que tout $X\in \ob E$ appartient
à un $\Filt^{\Pi}(E)$, pour $\Pi$ assez grand.

\begin{enonce*}[remark]{Exemple 9.12.2}\etiquette[9.12.2]{I.exem9.12.2}
Prenons $E=(𝒰\text{-}\Ens)$, $\Pi_\circ=\aleph_\circ$,
$\Filt^{\Pi}(E)=$ sous-catégorie pleine de $E$ formée des ensembles
tels que $\card X\leq \Pi$. Plus généralement:
\end{enonce*}

\begin{enonce*}{Proposition 9.13}\etiquette[9.13]{I.prop9.13}
Soit $E$ une $𝒰$-catégorie. On suppose que $E$ est stable
par petites $\varinjlim$ filtrantes, par somme de deux objets et par
conoyaux de doubles flèches, que $E$ est stable par produits fibrés,
et que les morphismes épimorphiques dans $E$ sont épimorphiques
stricts universels. Soit $C$ une petite sous-catégorie pleine de $E$
qui est génératrice par épimorphismes stricts. Soit $\Pi_\circ$ un
cardinal infini $\geq \card F\ell\,  C$. Pour tout cardinal $\Pi\geq
\Pi_\circ$, soit $\Filt^{\Pi}(E)$ la sous-catégorie strictement
pleine de $E$ formée des objets $X$ de $E$ tels qu'il existe une
famille épimorphique stricte $(X_i\rightarrow X)_{i \in I}$ de but
$X$, telle que $\card I\leq \Pi$ et que $X_i\in \ob C$ pour tout
$i\in I$. Alors $(\Filt^{\Pi}(E))_{\Pi\geq \Pi_\circ}$ est une
filtration cardinale de $E$. De plus, pour tout $X\in \ob
\Filt^{\Pi}(E)$, le cardinal de l'ensemble des flèches de $C_{/X}$
(et \textit{a fortiori}, de l'ensemble des objets de $C_{/X}$) est
majoré par $\Pi^{\Pi_\circ}$.
\end{enonce*}

Cette dernière assertion n'est autre que \hyperref[I.prop7.6]{7.6}.

Pour montrer qu'on a une filtration cardinale, il faut prouver les
conditions a), b) et c) de \hyperref[I.def9.12]{9.12}. La condition a)
résulte aussitôt de
7.5.2. Pour b), supposons qu'on ait $X=\varinjlim_I X_i$, avec
$\card I\leq \Pi$ et les\pageoriginale $X_i\in \ob \Filt^{\Pi}(E)$.
Donc la famille des morphismes canoniques $X_i\rightarrow X$ est
épimorphique stricte, et par hypothèse sur les $X_i$, il existe pour
tout $i\in I$ une famille épimorphique stricte $(X_{ij}\rightarrow
X_i)_{j\in J_i}$, avec $\card J_i\leq \Pi$. Passant aux sommes
$\amalg_i X_i$ et $\amalg_{i,j}X_{ij}$, on voit par transitivité des
épimorphismes stricts universels (II 2.5) que la famille des
composés $X_{ij}\rightarrow X_i\rightarrow X$ (indexée par
l'ensemble somme des $J_i$, pour $i\in I$) est épimorphique stricte.
Or on a $\card J\leq \Pi^2=\Pi$, d'où la conclusion. (NB. on a
seulement utilisé le fait que $\varinjlim_I X_i$ existe, et non le
fait que $I$ soit grand devant $\Pi_\circ$ ni même filtrant.)
Prouvons enfin c). Notons que pour tout $X\in \ob E$, la famille des
flèches $X_i\rightarrow X$ de source dans $\ob C$ est strictement
épimorphique, puisque $C$ est génératrice par
épimorphismes stricts,
et évidemment petite, donc il existe un cardinal $\Pi\geq \Pi_\circ$
tel que $X\in \ob \Filt^{\Pi}(E)$. Il reste à prouver que si on a
des cardinaux $\Pi'\geq \Pi\geq \Pi_\circ$, et si $X\in \ob
\Filt^{\Pi'}(E)$, alors on a un isomorphisme
$$
X\simeq \varinjlim_I X_i{\quoi},
$$
avec $I$ grand devant $\Pi$, $\card I\leq {\Pi'}^{\Pi}$, les $X_i$
dans $\ob \Filt^{\Pi}(E)$. Or on a, $C$ étant génératrice par
épimorphismes stricts,
\begin{equation*}
X\simeq \varinjlim_{C_{/X}} T\tag{$\ast$}{\quoi}.
\end{equation*}
Notons aussi %%
qu'en rajoutant successivement à $C$ des sommes et des conoyaux de
doubles flèches de $C$, et de même pour la catégorie ainsi obtenue
etc, on se ramène au cas où $C$ est stable par somme de deux objets
dans $E$ et par conoyaux de doubles flèches dans $E$, et ceci sans
détruire l'hypothèse $\Pi_\circ \geq \card F\ell\, C$. On peut
supposer de plus que, si $E$ contient un\pageoriginale objet initial
fixé $∅_E$, on ait $∅_E\in \ob C$. Ceci implique
que la catégorie $C_{/X}$ est stable par somme de deux objets et
conoyaux de doubles flèches. Elle est alors filtrante, car elle est
non vide, puisque si elle était vide, la relation $(\ast)$
montrerait que $X$ est un objet initial, donc $C_{/X}$ contient la
flèche $∅_E\rightarrow X$, une contradiction. Soit alors
$I$ l'ensemble, ordonné par inclusion, des sous-catégories pleines
$i$ de $C_{/X}$ qui sont filtrantes et telles que $\card \ob i\leq
\Pi$. Pour tout $i\in I$, soit $X_i\in \ob E$ la limite inductive du
foncteur composé $i\rightarrow C_{/X}\rightarrow E$, qui existe par
hypothèse. Alors on a évidemment
$$
\varinjlim_I X_i \simeq \varinjlim_{C/X} T\simeq X{\quoi}.
$$
D'autre part, on a déjà noté que $\card \ob C_{/X}\leq
{\Pi'}^{\Pi_\circ}$. Il s'ensuit que $I$ est grand devant $\Pi$, et
que $\card I\leq
({\Pi'}^{\Pi_\circ})^{\Pi}={\Pi'}^{\Pi\Pi_\circ}={\Pi'}^{\Pi}$ en
vertu du lemme suivant, dont la démonstration est laissé au lecteur
(où on fera $I=C_{/X}$, $c={\Pi'}^{\Pi_\circ}$):

\begin{enonce*}{Lemme 9.13.1}\etiquette[9.13.1]{I.lem9.13.1}
Soient $J$ une catégorie filtrante, $c$ et $\Pi$ deux cardinaux tels
que $c\geq \Sup(\card  F\ell\,  J,\Pi)$, et telle que $\card
\Hom(j,j')\leq \Pi_\circ$ pour tout couple d'objets $j$, $j'$ de
$J$. Soit $I$ l'ensemble des sous-catégories pleines filtrantes $i$
de $J$ telles que $\card \ob i\leq \Pi$. Alors, ordonné par
inclusion, $I$ est grand devant $\Pi$, et on a $\card I\leq
c^{\Pi}$.
\end{enonce*}

Ceci achève la démonstration de \hyperref[I.prop9.13]{9.13}.

\begin{enonce*}[remark]{Remarque 9.13.2}\etiquette[9.13.2]{I.rem9.13.2}
Un\pageoriginale léger effort supplémentaire doit permettre de
remplacer dans \hyperref[I.prop9.13]{9.13} l'hypothèse que $E$ est stable par
petites limites inductives filtrantes par l'hypothèse que $E$
satisfait à $L$ (\hyperref[I.def9.1]{9.1}), si on suppose $E$ stable par
petites sommes, ou que dans $E$ toute \emph{famille} épimorphique
est épimorphique universelle. Il faut alors, dans la démonstration
de c), choisir $\Pi_\circ$ tel que $E$ satisfasse à $L_{\Pi_\circ}$,
et se borner aux sous-catégories pleines $i$ de $C_{/X}$ qui sont
non seulement filtrantes, mais telles que l'ensemble préordonné $\ob
i$ soit grand devant $\Pi_\circ$. Utilisant 8. et l'hypothèse que
$E$ satisfait à $L_{\Pi_\circ}$, on trouve alors que les $X_i$
existent, et on devrait conclure par une variante convenable de
\hyperref[I.lem9.13.1]{9.13.1}, que le rédacteur n'a pas vérifiée.
\end{enonce*}

\begin{enonce*}{Proposition 9.14}\etiquette[9.14]{I.prop9.14}
Soient $f:E\rightarrow F$ un foncteur accessible (\hyperref[I.def9.2]{9.2}) entre deux
catégories munies de filtrations cardinales
$(\Filt^{\Pi}(E))_{\Pi\geq\Pi_\circ}$ et
$(\Filt^{\Pi}(F))_{\Pi\geq\Pi'_\circ}$. Alors il existe un cardinal
$\Pi_1\geq \Sup(\Pi_\circ, \Pi'_\circ)$ tel que pour tout cardinal
$\Pi\geq \Pi_1$, on ait
\begin{equation}
f(\Filt^{\Pi}(E))\subset \Filt^{\Pi^{\Pi_1}}(F){\quoi}.\tag{9.14.1}
\end{equation}
En particulier, si l'on a $\Pi=2^c$ avec $c\geq \Pi_1$, d'où
$\Pi^{\Pi_1}=2^{c\Pi_1}=2^c=\Pi$, on a
\begin{equation}
f(\Filt^{\Pi}(E))\subset \Filt^{\Pi}(F){\quoi}.\tag{9.14.2}
\end{equation}
\end{enonce*}

Signalons tout de suite le

\begin{enonce*}{Corollaire 9.15}\etiquette[9.15]{I.cor9.15}
Soit\pageoriginale $E$ une catégorie, munie de deux filtrations
cardinales\linebreak $(\Filt^{\Pi}(E))_{\Pi\geq\Pi_\circ}$ et
$({\Filt'}^{\Pi}(E))_{\Pi\geq %%
\Pi'_\circ}$. Alors il existe un cardinal $\Pi_1\geq \Sup(\Pi_\circ,
\Pi'_\circ)$ tel que, pour tout cardinal $c\geq\Pi_1$, posant
$\Pi=2^c$, on ait $\Filt^{\Pi}(E)={\Filt'}^{\Pi}(E)$.
\end{enonce*}

\noindent \emph{Preuve de} 9.14. Soit $c$ un cardinal tel que
$c\geq\Sup(\Pi_\circ, \Pi'_\circ)$, et tel que $f$ soit
$c$-admissible. Soit d'autre part $\Pi_1\geq c$ tel que l'on ait
\begin{equation}
f(\Filt^c(E))\subset \Filt^{\Pi_1}(F){\quoi}.\tag{$\ast$}
\end{equation}
Il existe un tel $\Pi_1$, grâce au fait que $\Filt^c(E)$ est
équivalente à une petite catégorie (9.12 a)), donc $f(\Filt^c(E))$
l'est également, de sorte qu'on peut appliquer 9.12.1 aux objets de
cette dernière pour trouver une $\Filt^{\Pi}(F)$ qui les contient %%
tous (compte tenu que les $\Filt^{\Pi}(F)$ sont des sous-catégories
pleines). Soit donc $\Pi\geq \Pi_1$, et $X\in \ob \Filt^{\Pi}(E)$,
prouvons que $f(X)\in \Filt^{\Pi^{\Pi_1}}(F)$. Écrivons en effet
$$
X=\varinjlim_I X_i{\quoi},
$$
avec les $X_i\in \ob \Filt^c(E)$, $I$ grand devant $c$, $\card I\leq
\Pi^c\leq \Pi^{\Pi_1}$ (\hyperref[I.def9.12]{9.12}c)). Comme $f$ est $c$-admissible, $I$
grand devant $c$, on a
$$
f(X)\simeq \varinjlim_I f(X_i){\quoi},
$$
et comme $\card I\leq \Pi^{\Pi_1}$ et $f(X_i)\in \ob
\Filt^{\Pi_1}(F)\subset \ob \Filt^{\Pi^{\Pi_1}}(F)$ par $(\ast)$, on
a $f(X)\in \Filt^{\Pi^{\Pi_1}}(F)$ par \hyperref[I.def9.12]{9.12} b), \cqfd

\setcounter{subsection}{15} \setcounter{subsubsection}{0}
\subsubsection{}\etiquette[9.15.1]{I.sec9.15.1}
La notion de filtration cardinale \hyperref[I.def9.12]{9.12} n'a guère
d'intérêt que lorsque les objets de $E$ sont accessibles. Signalons
qu'il résulte de \hyperref[I.prop9.11]{9.11} que cette condition est satisfaite
lorsque, en plus des hypothèses de\pageoriginale \hyperref[I.prop9.13]{9.13},
on suppose que le foncteur $\Ker$ sur la catégorie des doubles
flèches de $E$ est accessible. Signalons d'autre part:

\begin{enonce*}{Proposition 9.16}\etiquette[9.16]{I.prop9.16}
Soit $E$ une catégorie munie d'une filtration cardinale
$(\Filt^{\Pi}(E))_{\Pi\geq\Pi_\circ}$. Supposons que les éléments de
$\Filt^{\Pi_\circ}(E)$ soient des objets accessibles de $E$; alors
il existe un cardinal $\Pi_1\geq\Pi_\circ$ dans $𝒰$ tel que
les objets de $\Filt^{\Pi_\circ}(E)$ soient $\Pi_1$-accessibles; si
$\Pi_1$ est choisi ainsi, alors pour tout cardinal $\Pi\geq \Pi_1$,
on a (avec les notations de \hyperref[I.sec9.3.1]{9.3.1}):
\begin{equation}
E_{\Pi}\subset \Filt^{\Pi}(E)\subset
E_{(\Pi^{\Pi_\circ})}{\quoi}.\tag{9.16.1}
\end{equation}
\end{enonce*}

L'existence de $\Pi_1$ résulte immédiatement du fait que
$\Filt^{\Pi_\circ}(E)$ est équivalente à une petite
catégorie (\hyperref[I.def9.12]{9.12}a)). Soit alors $\Pi\geq \Pi_1$. Si $X$
  est dans $\Filt^{\Pi}(E)$,
écrivant $X\simeq \varinjlim_I X_i$ avec $\card I\leq
\Pi^{\Pi_\circ}$, $X_i\in \ob \Filt^{\Pi_\circ}(E)$ pour tout $i$
(\hyperref[I.def9.12]{9.12}c)), alors il résulte de \hyperref[I.cor9.9]{9.9} que $X$ est
$\Pi^{\Pi_\circ}$-accessible, d'où la deuxième inclusion (9.16.1).
Supposons que $X$ soit $\Pi$-accessible, et écrivons $X\simeq
\varinjlim_I X_i$, avec $I$ grand devant $\Pi$ et les $X_i\in \ob
\Filt^{\Pi}(E)$ (\hyperref[I.def9.12]{9.12}c)). Par hypothèse sur $X$, l'isomorphisme
donné $X\xrightarrow{\sim} \varinjlim_I X_i$ se factorise par un des
$X_i$, donc $X$ est isomorphe à un facteur direct de cet $X_i$. On
en conclut que $X\in \ob \Filt^{\Pi}(E)$, donc la première inclusion
(9.16.1), grâce au

\begin{enonce*}{Lemme 9.16.2}\etiquette[9.16.2]{I.lem9.16.2}
Tout objet de $E$ qui est un facteur direct d'un objet de
$\Filt^{\Pi}(E)$ est dans $\Filt^{\Pi}(E)$.
\end{enonce*}

En effet, si $X$ est l'image d'un projecteur $p$ dans l'objet $Y$ de
$E$ (\ie d'un endomorphisme $p$ tel que $p^2=p$), et si $I$ est un
ensemble ordonné\pageoriginale  filtrant, $X$ est limite inductive
du système inductif filtrant $(Y_i)_{i\in I}$ défini par $Y_i=Y$
pour tout $i\in I$, $p:Y_i\rightarrow Y_j$ si $i<j$. Prenant $I$
grand devant $\Pi_\circ$ et $\card I\leq \Pi$, on voit donc que si
$Y\in \Filt^{\Pi}(E)$, il en est de même de $X$ en vertu de
\hyperref[I.def9.12]{9.12} b).

\begin{enonce*}{Corollaire 9.17}\etiquette[9.17]{I.cor9.17}
Sous les conditions de \hyperref[I.prop9.16]{9.16}, pour tout cardinal $c\geq
\Pi_1$, posant $\Pi=2^c$, $\Filt^{\Pi}(E)$ est identique à la
sous-catégorie strictement pleine de $E_{\Pi}$ de $E$ formée des
objets $\Pi$-accessibles.
\end{enonce*}

\begin{enonce*}{Corollaire 9.18}\etiquette[9.18]{I.cor9.18}
Soient $E$ une catégorie satisfaisant la condition $L_{\Pi}$
(\hyperref[I.def9.1]{9.1}), où $\Pi$ est un cardinal, et $C$ une
sous-catégorie pleine
de $E$. On désigne par $\Ind(C)_{\Pi}$ la sous-catégorie pleine de
la catégorie $\Ind(C)$ des ind-objets de $C$ (\hyperref[I.sec8.2]{8.2})
formée des
ind-objets de la forme $(X_i)_{i\in I}$, où $I$ est un ensemble
ordonné grand devant $\Pi$. Considérons le foncteur canonique
\begin{equation}
(X_i)_{i\in I}\longmapsto \varinjlim_I
X_i:\Ind(C)_{\Pi}\longrightarrow E{\quoi}.\tag{9.18.1}
\end{equation}
\begin{enumerate}
\renewcommand{\theenumi}{\alph{enumi}}
\renewcommand{\labelenumi}{\theenumi)}
\item Pour que ce foncteur soit pleinement fidèle, il faut et il
suffit que tout objet $X$ de $C$ soit un objet $\Pi$-accessible
(\hyperref[I.def9.3]{9.3}) de $E$.

\item Plaçons-nous sous les conditions de \hyperref[I.cor9.17]{9.17}, en particulier
$\Pi=2^c$, et prenons $C=\Filt^{\Pi}(E)$. Alors le foncteur (9.18.1)
est une équivalence de catégories.
\end{enumerate}
\end{enonce*}

L'assertion a) est une généralisation immédiate de \hyperref[I.prop8.7.5]{8.7.5}
a), et se prouve de la même façon. Alors b) résulte de
\hyperref[I.cor9.17]{9.17} et de la condition \hyperref[I.def9.12]{9.12} c) des
filtrations cardinales, qui implique que le foncteur envisagé est
essentiellement surjectif.

\begin{enonce*}{Corollaire 9.19}\etiquette[9.19]{I.cor9.19}
Sous\pageoriginale les conditions de \hyperref[I.cor9.17]{9.17}, soient
$C=\Filt^{\Pi}(E)$, $F$ une $𝒰$-catégorie, et considérons le
foncteur
\begin{equation}
\Hhom(E,F)_{\Pi}\longrightarrow \Hhom(C,F)\tag{9.19.1}
\end{equation}
induit par le foncteur « restriction à $C$ » $f\mapsto f| C$),
où la source de (9.19.1) est la catégorie des foncteurs
$\Pi$-accessibles de $E$ dans $F$ (\hyperref[I.def9.2]{9.2}). Le foncteur
précédent est
pleinement fidèle. Si $F$ satisfait à la condition $L_{\Pi}$
(\hyperref[I.def9.1]{9.1}), alors le foncteur (9.19.1) est une équivalence de
catégories.
\end{enonce*}

La première assertion se prouve comme \hyperref[I.prop7.8]{7.8}, en utilisant
\hyperref[I.cor9.18]{9.18} b). La deuxième s'obtient en construisant un
foncteur quasi-inverse de (9.19.1), en associant à tout
foncteur $f:C\rightarrow F$ le foncteur $(X_i)_{i\in I}\rightarrow
\varinjlim_i f(X_i)$ de $\Ind(E_{\Pi})$ dans $F$, et en utilisant
\hyperref[I.cor9.18]{9.18} b) pour en déduire un foncteur $\bar{f}:E\rightarrow
F$. Tout revient à %%
montrer que ce%%
dernier est $\Pi$-accessible. Or cela se prouve comme l'assertion
analogue \hyperref[I.prop8.7.3]{8.7.3}.

\begin{enonce*}{Corollaire 9.20}\etiquette[9.20]{I.cor9.20}
Soient $E$ une catégorie admettant une filtration cardinale et telle
que tout objet de $E$ soit accessible (\cf \hyperref[I.prop9.16]{9.16}), $F$
une catégorie. Alors:
\begin{enumerate}
\renewcommand{\theenumi}{\alph{enumi}}
\renewcommand{\labelenumi}{\theenumi)}
\item La catégorie $\Hhom(E,F)_{acc}$ des foncteurs accessibles
de $E$ dans $F$ (\hyperref[I.def9.2]{9.2}) est une
$𝒰$-catégorie. ({\rm NB} on
rappelle (\hyperref[I.sec9.0]{9.0}) que les catégories données $E$, $F$
sont supposées
être des $𝒰$-catégories.)

\item Supposons que $F$ soit stable par petites limites inductives
filtrantes. Pour toute sous-catégorie pleine $C$ de $E$ équivalente
à une petite catégorie, à foncteur d'inclusion $i:C\rightarrow E$,
considérons le foncteur correspondant
$$
i_{!}:\Hhom(C,F)\longrightarrow \Hhom(E,F)
$$
$(5.1)$.\pageoriginale Pour qu'un foncteur $f:E\rightarrow F$ soit
accessible, il faut et il suffit qu'il existe une petite
sous-catégorie $C$ de $E$, telle que $f$ soit dans l'image
essentielle du foncteur précédent $i_{!}$.
\end{enumerate}
\end{enonce*}

\noindent \emph{Démonstration}.
\begin{enumerate}
\renewcommand{\theenumi}{\alph{enumi}}
\renewcommand{\labelenumi}{\theenumi)}
\item Il suffit de prouver que pour tout cardinal $\Pi$ tel que $E$
satisfasse $L_{\Pi}$, la sous-catégorie pleine $\Hhom(E,F)_{\Pi}$ de
$\Hhom(E,F)_{acc}$ est une $𝒰$-catégorie. Il suffit
évidemment de le vérifier pour les cardinaux de la forme $2^c$, avec
$c$ assez grand. Mais alors cela résulte de \hyperref[I.cor9.19]{9.19}, puisque
($C=\Filt^{\Pi}(E)$ étant essentiellement petite) $\Hhom(C,F)$ est
évidemment une $𝒰$-catégorie.

\item Par transitivité de la formation des foncteurs $i_{!}$, on
peut dans l'énoncé se borner aux sous-catégories $C$ de la forme
$\Filt^{\Pi}(E)$, où $\Pi$ est comme dans \hyperref[I.cor9.17]{9.17}. On voit alors
aisément que le foncteur composé $\Hhom(\Filt^{\Pi}(E),F)\rightarrow
\Hhom(E,F)_{\Pi}\rightarrow \Hhom(E,F)$, où la première flèche est
quasi-inverse de (9.19.1), et la deuxième est l'inclusion, n'est
autre que le foncteur $i_{!}$, à isomorphisme près. Donc l'assertion
b) résulte de \hyperref[I.cor9.19]{9.19}.
\end{enumerate}

\begin{enonce*}[remark]{*Exercice 9.20.1}\etiquette[9.20.1]{I.exer9.20.1}
(Le présent exercice utilise les notions de site et de topos,
développés dans les exposés II et IV.) Soient $E$ un
$𝒰$-topos, $𝒱$ un univers tel que $𝒰\in
𝒱$, $C$ une petite sous-catégorie pleine génératrice du
$𝒰$-topos $E$, $\Pi_\circ\in 𝒰$ un cardinal infini
tel que $\Pi_\circ\in \card F\ell\,  C$. Pour tout cardinal $\Pi\geq
\Pi_\circ$, $\Pi\in 𝒰$, soit $\Filt^{\Pi}(E)$ la
sous-catégorie strictement pleine de $E$ formée des objets $X$ tels
qu'il existe une famille épimorphique stricte\pageoriginale
$(X_i\rightarrow X)_{i\in I}$ de but $X$, telle que $\card I\leq
\Pi$ et que $X_i\in \Ob C$ pour tout $i\in I$.
\begin{enumerate}
\renewcommand{\theenumi}{\alph{enumi}}
\renewcommand{\labelenumi}{\theenumi)}
\item Montrer que $(\Filt^{\Pi}(E))_{\Pi\geq \Pi_\circ}$ est une
filtration cardinale de la $𝒰$-catégorie $E$, et qu'on peut
choisir $\Pi_\circ$ tel que pour tout $\Pi\geq \Pi_\circ$, on ait
les inclusions
$$
E_{\Pi}\subset \Filt^{\Pi}(E)\subset E_{\Pi^{\Pi_\circ}}{\quoi},
$$
où pour tout cardinal $c$, $E_c$ désigne la sous-catégorie
strictement pleine de $E$ formée des objets $c$-accessibles.

\item Choisissant un cardinal $\Pi\geq \Pi_\circ$ tel que
$\Pi=\Pi^{\Pi_\circ}$ (par exemple $\Pi$ de la forme $2^c$, avec
$c>\Pi_\circ$), et posant $C=\Filt^{\Pi}(E)$, montrer qu'on a une
équivalence de catégories
$$
\Ind(C)_{\Pi}\xrightarrow{\approx} E{\quoi},
$$
(notation $\Ind(C)_{\Pi}$ de \hyperref[I.cor9.18]{9.18}).

\item Gardons les notations de b), et soient $𝒱$ un univers tel
que $𝒰\subset 𝒱$, $%%
E^{\sim{\phantom{,}}}_{\underline{V}}$ le $𝒱$-topos de
$𝒱$-faisceaux sur le site $E$ (muni de sa topologie
canonique). Désignons par $\Ind(C,𝒱)_{\Pi}$ la catégorie des
$𝒱\text{-}\Ind$-objets de $C$ indexés par des ensembles%%
 d'indices préordonnés \emph{grands devant} $\Pi$. Montrer qu'on a
une équivalence de catégories
$$
\Ind(C,𝒱)_{\Pi}\xrightarrow{\approx}%%
E^{\tilde{\phantom{,}}}_{𝒱}{\quoi}.
$$

\item Soient $c\in 𝒱$ un cardinal, $\Ind(E,𝒱)'_c$ la
sous-catégorie pleine de $\Ind(E,𝒱)$ formée des
$𝒱$-ind-objets de $E$ indexés par un ensemble préordonné qui
est grand devant tout cardinal $<c$. Prenant $c=\card 𝒰$,
montrer\pageoriginale qu'on a une équivalence de catégories
$$
\displaylines{\hfill \Ind(E,𝒱)'_{c}\xrightarrow{\approx}%%
E^{\tilde{\phantom{,}}}_{𝒱}{\quoi}.\hfill{\llap{\textbf{*}}}}
$$
\end{enumerate}
\end{enonce*}

\setcounter{subsection}{20}
\subsection{}\etiquette[9.21]{I.sec9.21}
La présente section \hyperref[I.sec9.21]{9.21} développe des préliminaires
techniques pour
la démonstration du théorème \hyperref[I.them9.22]{9.22} ci-dessous,
qui constitue le résultat principal du présent paragraphe
\hyperref[I.sec9]{9}. Soit
\begin{equation}
p:E\rightarrow B\tag{9.21.1}
\end{equation}
un foncteur fibrant, où $B$ est une petite catégorie, et où les
foncteurs images inverses
$$
f^\ast:E_\beta\longrightarrow E_\alpha
$$
associés aux flèches $f:\alpha\rightarrow \beta$ de $B$ sont
accessibles (\hyperref[I.def9.2]{9.2}). En particulier, les catégories
fibres $E_\alpha$
$(\alpha \in \ob B)$ satisfont à la condition $L$
(\hyperref[I.def9.1]{9.1}), donc, $B$
étant petite, il existe un cardinal $\Pi'_\circ\in 𝒰$ tel
que toutes les catégories $E_\alpha$ satisfont à la condition
$L_{\Pi'_\circ}$. Soit $\Pi_\circ\in 𝒰$ un cardinal infini
$>\Pi'_\circ$, de sorte que l'on a
\begin{equation}
\Pi_{\circ}>\Pi'_{\circ}, E_\alpha \mbox{ satisfait }
L_{\prod'_{\circ}} \mbox{ pour tout } \alpha\in \ob B.\tag{9.21.2}
\end{equation}
D'ailleurs, les hypothèses faites impliquent aussitôt l'existence
d'un cardinal $c\in 𝒰$ satisfaisant aux conditions
suivantes:
\begin{equation}
\left\{
\begin{tabular}{rp{6cm}}
a) & $c$ est infini,\\
b) & $c\geq \card F\ell\, B$,\\
c) & pour toute flèche $f:\alpha\rightarrow \beta$ dans  $B$, le
foncteur  $f^\ast:E_\beta\rightarrow E_\alpha$
         est $c$-accessible.
\end{tabular}\right.
\tag{9.21.3}
\end{equation}
Supposons\pageoriginale de plus qu'on puisse trouver, pour chaque
$a\in \ob B$, une sous-catégorie pleine $C_\alpha$ de $E_\alpha$,
satisfaisant les conditions suivantes:
\begin{equation}
\left\{
\begin{tabular}{rp{6cm}}
a) & Pour tout $X\in \ob C_\alpha$, $X$ est $c$-accessible dans
$E_\alpha$ (9.3).\\
b) & Pour tout $X\in \ob E_\alpha$, on peut trouver un isomorphisme
$X\simeq \varinjlim_I X_i$ dans $E_\alpha$, où les $X_i$ sont dans
$C_\alpha$ et où $I$ est un ensemble ordonné grand devant $c$.
\end{tabular}\right.
\tag{9.21.4}
\end{equation}

Soit $\Pi$ un cardinal tel que l'on ait
\begin{equation}
\Pi\geq \Pi_\circ\quav\mbox{ donc }
\Pi>\Pi'_\circ{\quoi},\tag{9.21.5}
\end{equation}
et pour tout $\alpha\in \ob B$, soit
\begin{equation}
C^{\Pi}_{\alpha}\subset E_{\alpha}\tag{9.21.6}
\end{equation}
formée des objets qui se peuvent représenter sous la forme
$\varinjlim_I X_i$, où $I$ est un ensemble ordonné grand devant
$\Pi'_\circ$, tel que $\card I\leq \Pi$, et où les $X_i$ sont dans
$C_{\alpha}$. Il est clair alors, grâce à \hyperref[I.var7.5.2]{7.5.2}, que
$C^{\Pi}_{\alpha}$ est essentiellement petite, \ie est équivalente à
une petite catégorie.

Soit
\begin{equation}
F=\Hhom_B(B,E)\tag{9.21.7}
\end{equation}
la catégorie des sections de $E$ sur $B$, et soit
\begin{equation}
F^{\Pi}\subset F\tag{9.21.8}
\end{equation}
la sous-catégorie strictement pleine de $F$ formée des sections
$X:\alpha\mapsto X(\alpha)$ telles que pour tout $\alpha\in \ob B$
on ait
$$
X(\alpha)\in C^{\Pi}_{\alpha}{\quoi}.
$$
Il\pageoriginale est clair, les $C^{\Pi}_{\alpha}$ étant
essentiellement petites, qu'il en est de même de la catégorie
$F^{\Pi}$. Nous allons montrer que cette catégorie est génératrice,
et plus précisément:

\begin{enonce*}{Lemme 9.21.9}\etiquette[9.21.9]{I.lem9.21.9}
Sous les conditions et avec les notations précédentes, on a ce qui
suit:
\begin{enumerate}
\renewcommand{\theenumi}{\roman{enumi}}
\renewcommand{\labelenumi}{(\theenumi)}
\item Tout objet de $F$ est accessible (\hyperref[I.def9.3]{9.3}). Si $d$ est un cardinal
$\geq \Pi'_\circ$, et si $\alpha\mapsto X(\alpha)$ est un élément de
$F$ tel que pour tout $\alpha$, $X(\alpha)$ soit $d$-accessible dans
$E_{\alpha}$, alors $X$ est $d$-accessible dans $F$.

\item Supposons $\Pi\leq c$, ou que $\Pi'_\circ$ soit fini (\ie
les $E_{\alpha}$ stables par petites $\varinjlim$ filtrantes). Alors
tout objet $X$ de $F$ est isomorphe à un objet de la forme
$\varinjlim_I X_i$, où les $X_i$ sont dans $F^{\Pi}$ et où $I$ est
grand devant $\Pi$.
\end{enumerate}
\end{enonce*}

Pour prouver (i), notons qu'en vertu de (9.21.4) et de
\hyperref[I.cor9.9]{9.9}, tout
objet de $E_{\alpha}$ est accessible. D'autre part, $F$ satisfait à
la condition $L_{\Pi'_\circ}$, en vertu du


\begin{enonce*}{Lemme 9.21.10}\etiquette[9.21.10]{I.lem9.21.10}
Soit $p:E\rightarrow B$ un foncteur fibrant, $F=\Hhom_B(B,E)$, $I$
une catégorie, $i\mapsto X_i$ un foncteur de $I$ dans $F$. Pour que
$\varinjlim_I X_i$ soit représentable dans $F$, il suffit que pour
tout $\alpha\in \ob B$, $\varinjlim_I X_i(\alpha)$ soit
représentable dans la catégorie fibre $E_{\alpha}$; lorsqu'il en est
ainsi, alors $\varinjlim_I X_i$ « se calcule argument par argument ».
En particulier, si les catégories fibres satisfont à la
condition $L_{\Pi'_\circ}$ ($\Pi'_\circ$ étant un cardinal donné) il
en est de même de $F$.
\end{enonce*}

Nous\pageoriginale laissons le détail de la démonstration (facile)
de \hyperref[I.lem9.21.10]{9.21.10} au lecteur, en nous contentant de remarquer qu'il est
commode d'utiliser le résultat suivant, dont la démonstration est
immédiate:

\begin{enonce*}{Corollaire 9.21.10.1}\etiquette[9.21.10.1]{I.cor9.21.10.1}
Avec les hypothèses et notations de \hyperref[I.lem9.21.10]{9.21.10} pour
$p:E\rightarrow B$, et pour $I$, pour tout $\alpha\in \ob B$, le
foncteur d'inclusion $E_{\alpha}\rightarrow E$ commute aux limites
inductives de type $I$.
\end{enonce*}

Revenant alors aux conditions générales de \hyperref[I.lem9.21.9]{9.21.9}, soit
$X$ un objet de $F$. Comme pour tout $\alpha \in \ob B$, $X(\alpha)$
est accessible dans $E_{\alpha}$, il existe un cardinal $d\in
𝒰$ tel que pour tout $\alpha$, $X(\alpha)$ soit
$d$-accessible dans $E_{\alpha}$. On peut choisir $d\geq
\Pi'_{\circ}$, de sorte que $F$ satisfait à $L_d$ en vertu de
\hyperref[I.lem9.21.10]{9.21.10}. Utilisant encore \hyperref[I.lem9.21.10]{9.21.10} pour le
calcul des limites inductives $\varinjlim_I Y_i$ dans $F$, avec $I$
grand devant $d$, on constate aussitôt que $X$ est $d$-accessible.
Cela prouve (i).

Nous allons prouver maintenant \hyperref[I.lem9.21.9]{9.21.9} (ii) en plusieurs étapes
((9.21.11) à (\hyperref[I.lem9.21.16]{9.21.16})).

Soit donc
$$
X:\alpha\longmapsto X(\alpha)
$$
un objet de $F$. En vertu de (9.21.4 b)), on peut, pour tout
$\alpha\in \ob B$, trouver un ensemble ordonné $I_{\alpha}$ grand
devant $c$, et un isomorphisme $X(\alpha)=\varinjlim_{i\in
I_{\alpha}} X(\alpha)_i$, où $i\mapsto X(\alpha)_i$ est un système
inductif de type $I_{\alpha}$ dans $E_{\alpha}$, les $X(\alpha)_i$
dans $C_{\alpha}$. Considérons l'ensemble ordonné produit
$I=\Pi_{\alpha} I_{\alpha}$; il est clair qu'il est grand devant
$c$, et que les systèmes inductifs\pageoriginale précédents donnent
naissance à des systèmes inductifs dans les $E$,
\begin{equation}
i\longmapsto X(\alpha)_i\in \ob C_{\alpha}, i\in \ob
I{\quoi},\tag{9.21.11}
\end{equation}
indexés par le même ensemble ordonné $I$, et à des isomorphismes
\begin{equation}
X(\alpha)\xrightarrow{\sim}\varinjlim_I X(\alpha)_i\mbox{ dans }
E_{\alpha}{\quoi},\tag{9.21.12}
\end{equation}
pour tout $\alpha\in \ob B$. Nous supposons fixées par la suite des
données (9.21.11) et (9.21.12).

\begin{enonce*}{Lemme 9.21.13}\etiquette[9.21.13]{I.lem9.21.13}
Sous les conditions précédentes, on peut trouver une application
$\varphi:I\rightarrow I$ telle que $\varphi(i)\geq i$ pour tout
$i\in I$, et une application $(i,f)\mapsto \lambda(i,f)$ de $I\times
%%
\mathrm{F}{}l\, B$ dans $%%
\mathrm{F}{}l\, E$, satisfaisant aux conditions suivantes:
\begin{enumerate}
\renewcommand{\theenumi}{\alph{enumi}}
\renewcommand{\labelenumi}{\theenumi)}
\item Pour $i\in I$, $(f:\alpha\rightarrow \beta)\in F\ell B$,
$\lambda(i,f)$ est un $f$-morphisme de $X(\alpha)_i$ dans
$X(\beta)_{\varphi(i)}$, rendant commutatif le diagramme
\[
\xymatrix@C=3cm@R=.8cm{
 X(\alpha)\ar[r]^{X(f)} & X(\beta)\\
           &  X(\beta)_{\varphi(i)}\ar[u]\\
X(\alpha)_{i} \ar[uu]\ar[ru]^{\lambda(i,f)}&
                  \\
\alpha \ar[r]^{f}& \beta\quad , }
\]
où les flèches verticales sont les morphismes canoniques déduits de
(9.21.12).

\item Pour\pageoriginale tout $i\in I$, et $(f:\alpha\rightarrow
  \beta)\in %%
\mathrm{F}{}l\, B$, on a commutativité dans le diagramme
\[
\xymatrix@C=3cm@R=.9cm@L=6pt{
                      & X(\beta)_{\varphi^{2}(i)}\\
  X(\alpha)_{\varphi(i)} \ar[ru]^{\lambda(\varphi(i),f)}  & X(\beta)_{\varphi(i)} \ar[u]\\
X(\alpha)_i  \ar[ru]_{\lambda(i,f)} \ar[u]         \\
\alpha \ar[r]^{f} & \beta\quad . }
\]

\item Pour tout couple de flèches consécutives
$\alpha\xrightarrow{f}\beta\xrightarrow{g}\gamma$ de $B$, on a
commutativité dans le diagramme suivant
\[
\xymatrix@C=2cm@L=6pt@R=.5cm{
X(\alpha) \ar[r]^{X(f)}& X(\beta) \ar[r]^{X(g)}            & X(\gamma)\\
          &                      & X(\gamma)_{\varphi^{2}(i)}\ar[u]\\
          & X(\beta)_{\varphi(i)} \ar[uu]\ar[ru]^{\lambda(\varphi(i),g)} & X(\gamma)_{\varphi(i)}\ar[u]\\
X(\alpha)_i \ar[uuu]\ar[ru]^{\lambda(i,f)} \ar[rru]_{\lambda(i,gf)}&
 \\
\alpha \ar[r]^{f}& \beta \ar[r]^{g}& \gamma\quad , }
\]
où les flèches verticales sont celles déduites de
(9.21.12).
\end{enumerate}
\end{enonce*}

On\pageoriginale notera d'ailleurs que la commutativité des deux
trapèzes supérieurs dans b) est déjà contenu dans a), de sorte que
b) affirme en fait la commutativité du triangle inférieur du
diagramme envisagé.

Prouvons \hyperref[I.lem9.21.13]{9.21.13}. Soient $i\in I$, et $f:\alpha \rightarrow \beta$
une flèche dans $B$. Considérons le composé $X(\alpha_i)\rightarrow
X(\alpha)\xrightarrow{X(f)}X(\beta)\simeq \varinjlim_j X(\beta)_j$,
il définit un morphisme dans $E$:
$$
X(\alpha)_i\rightarrow f^\ast (X(\beta))\simeq f^\ast (\varinjlim_j
X(\beta)_j)\simeq \varinjlim_j f^\ast(X(\beta)_j){\quoi},
$$
où la dernière égalité provient du fait que $f^\ast$ est
$c$-accessible (9.21.3 c)%%
 et que $I$ est grand devant $c$. En vertu
de (9.21.4 a)), comme $X(\alpha)\in \ob C_{\alpha}$, le morphisme
envisagé se factorise par un $f^\ast(X(\alpha)_j)$, où $j$ \textit{a
priori} dépend de $i$ et de $f\in %%
\mathrm{F}{}l\,(B)$. Mais en vertu de (9.21.3 b)), et comme $I$ est
grand devant c), on peut choisir $j$ indépendant de $f$, soit
$j=\varphi(i)$. $I$ étant filtrant, on peut supposer $\varphi(i)\geq
i$. On trouve ainsi, pour $i\in I$ et $f\in %%
\mathrm{F}{}l\, F$, un morphisme
$$
X(\alpha)_i\longrightarrow f^\ast(X(\beta)_{\varphi(i)}){\quoi},
$$
ou ce qui revient au même, un $f$-morphisme
$$
\lambda(i,f):X(\alpha)_i\longrightarrow
X(\beta)_{\varphi(i)}{\quoi},
$$
qui par construction rend commutatif le diagramme
\[\vcenter{
\xymatrix@C=3cm{
X(\alpha)\ar[r]^{X(f)} & X(\beta)\\
          & X(\beta)_{\varphi(i)} \ar[u]\\
X(\alpha)_i \ar[uu] \ar[ru]^{\lambda(i,f)}&\\
\alpha \ar[r]^{f} & \beta \quad . }}\leqno{D(f)}
\]
Considérons\pageoriginale alors, pour $i$, $f$, $g$ donnés, le
triangle inférieur du diagramme envisagé dans 9.21.13
c); il n'est
pas clair qu'il est commutatif, mais les deux composés
$X(\alpha)_i\rightrightarrows X(\beta)_{\varphi^2(i)}$ deviennent
égaux après composition avec $X(\beta)_{\varphi^2(i)}\rightarrow
X(\beta)$, comme il résulte de la commutativité des deux trapèzes
supérieurs, et du trapèze contour global, qui sont les diagrammes
$D(f)$, $D(g)$ et $D(gf)$ respectivement. Donc, comme
$X(\alpha)_{\varphi^2(i)}$ est $c$-accessible (9.21.4 a)) et que
$X(\alpha)\simeq \varinjlim_{j\in I}X(\alpha)_j$, avec $I$ grand
devant $c$, il s'ensuit qu'on peut trouver un élément $j\geq
\varphi^2(i)$ de $I$, tel que les deux flèches envisagées deviennent
égales après composition avec $X(\beta)_{\varphi^2(i)}\rightarrow
X(\beta)_j$. A priori, $j$ dépend de $i$, $f$ et $g$. Mais pour $i$
fixé, l'ensemble des couples possibles $f$, $g$ est de cardinal
$\geq c^2=c$, en vertu de (9.21.3 a) et b)), donc, $I$ étant grand
devant $c$, on peut choisir $j$ indépendant de $f$ et de $g$, soit
$j=\varphi'(i)\geq \varphi^2(i)\geq \varphi(i)$. Soit alors, pour
tout $i\in I$ et $f:\alpha\rightarrow \beta$,
$$
\lambda'(i,f):X(\beta)_i\longrightarrow X(\beta)_{\varphi'(i)}
$$
le $f$-morphisme composé $X(\alpha)_i\rightarrow
X(\beta)_{\varphi(i)}\rightarrow X(\beta)_{\varphi'(i)}$, où la
deuxième flèche est le morphisme de transition. Il est alors
immédiat, par construction, que $(\varphi',\lambda')$ satisfait les
conditions 9.21.13 a) et c), pour $(\varphi,\lambda)$. Procédant de
même pour la condition b), on voit qu'on peut choisir $\varphi'$ de
telle façon que cette condition soit également satisfait pour
$(\varphi',\lambda')$. Cela achève la preuve de (9.21.13).

\setcounter{subsection}{21} \setcounter{subsubsection}{13}
\subsubsection{}\etiquette[9.21.14]{I.sec9.21.14}
Soit\pageoriginale maintenant
$$
J\subset I
$$
une partie de $I$ satisfaisant les conditions suivantes:
\begin{enumerate}
\renewcommand{\theenumi}{\alph{enumi}}
\renewcommand{\labelenumi}{\theenumi)}
\item $J$ est filtrante,

\item pour tout $j\in J$, on a $\varphi(j)\in J$,

\item Pour tout $a\in \ob B$, $X_J(\alpha)=\varinjlim_{i\in
J}X(\alpha)_i$ est représentable dans $E_{\alpha}$.
\end{enumerate}
Les conditions a) et c) sont satisfaites en particulier si $J$ est
grand devant $\Pi'_\circ$ (9.21.2). Il résulte alors de
\hyperref[I.cor9.21.10.1]{9.21.10.1}
que pour tout $\alpha\in \ob B$, $X_J(\alpha)$ est la limite
inductive $\varinjlim_{i\in J}X(\alpha)_i$ dans $E$. Or pour une
flèche $f:\alpha\rightarrow \beta$ de $B$, les flèches
$\lambda(i,f)$, pour $i$ variable dans $J$, définissent grâce à
(\hyperref[I.lem9.21.13]{9.21.13} b)) un morphisme de ind-objets de $(X(\alpha)_i)_{i\in J}$
dans $(X(\beta)_i)_{i\in J}$, d'où un homomorphisme
$$
X_J(f):X_J(\alpha)\rightarrow X_J(\beta)
$$
sur les limites inductives, qui est manifestement un
$f$-homomorphisme. Utilisant (\hyperref[I.lem9.21.13]{9.21.13} (c)), on trouve que l'on a des
relations de transitivité
$$
X_J(gf)=X_J(g)X_J(f){\quoi},
$$
de sorte que l'on a défini une section $X_J\in \ob F$ de $E$ sur
$B$. Enfin (\hyperref[I.lem9.21.13]{9.21.13} a)) nous montre que les homomorphismes
canoniques
$$
X_J(\alpha)\longrightarrow X(\alpha)\quav \alpha\in \ob B{\quoi},
$$
sont fonctoriels en $\alpha$, de sorte qu'on a un homomorphisme
canonique
$$
u_J:X_J\longrightarrow X{\quoi}.
$$
D'autre\pageoriginale part, si
$$
J'\supset J
$$
est une autre partie de $I$ satisfaisant aux conditions a), b), c)
ci-dessus, on trouve un homomorphisme canonique
$$
u_{J',J}:X_J\longrightarrow X_{J'}{\quoi},
$$
de sorte que les $X_J$ forment un système inductif dans $F$,
paramétré par l'ensemble (ordonné par inclusion) des parties $J$ de
$I$ satisfaisant les conditions envisagées. Enfin, les
homomorphismes $u_J$ ci-dessus définissent un homomorphismes de ce
système inductif dans (le système inductif constant défini par) $X$.


\setcounter{subsection}{21} \setcounter{subsubsection}{14}
\subsubsection{}\etiquette[9.21.15]{I.sec9.21.15}
Soit maintenant $K$ un ensemble de parties $J$ de $I$, satisfaisant
aux conditions a), b), c) envisagées dans \hyperref[I.sec9.21.14]{9.21.14}, et
supposons que
$K$ soit filtrant, et de réunion $I$. Alors il est clair que l'on a
$$
\varinjlim_{J\in K}X_J\xrightarrow{\sim}X{\quoi},
$$
en utilisant \hyperref[I.cor9.21.10.1]{9.21.10.1} qui nous ramène à
vérifier qu'on a un
isomorphisme argument par argument.

Prenons par exemple pour $K$ l'ensemble de toutes les parties $J$ de
$I$ qui sont grandes devant $\Pi'_\circ$, stables par $\varphi$, et
telles que $\card J\leq \Pi$. Alors par définition (9.21.6) de
$C^{\Pi}_{\alpha}$, on a, pour tout $\alpha\in \ob B$,
$X_J(\alpha)\in \ob C^{\Pi}_{\alpha}$, donc $X_J\in \ob F^{\Pi}$.
Par suite, \hyperref[I.lem9.21.9]{9.21.9} (ii) sera prouvé si nous
établissons que $K$ est
grand devant $\Pi$ (donc filtrant) et de réunion $I$. Il suffira
évidemment, pour ceci, de prouver que toute partie $S$ de $I$ telle
que $\card S\leq \Pi$ est contenue dans une $J\in K$. Ceci
résultera\pageoriginale en effet de l'hypothèse
préliminaire énoncée
dans \hyperref[I.lem9.21.9]{9.21.9} (ii), et du

\begin{enonce*}{Lemme 9.21.16}\etiquette[9.21.16]{I.lem9.21.16}
Soient $I$ un ensemble ordonné, $\Pi$ un cardinal infini,
$\varphi:I\rightarrow I$ une application de $I$ dans lui-même.
\begin{enumerate}
\renewcommand{\theenumi}{\alph{enumi}}
\renewcommand{\labelenumi}{\theenumi)}
\item Supposons $I$ filtrant. Alors pour toute partie $S$ de $I$ telle
que $\card S\leq \Pi$, il existe une partie filtrante $J\supset S$
de $I$, telle que $\varphi(J)\subset J$ et $\card(J)\leq \Pi$.

\item Soit $\Pi'_{\circ}$ un cardinal $<\Pi$, et supposons $I$ grand
devant $\Pi$. Alors pour toute partie $S$ de $I$ telle que $\card
S\leq \Pi$, il existe une partie $J\supset S$ de $I$ grande devant
$\Pi'_{\circ}$, telle que $\varphi(J)\subset J$ et $\card J\leq
\Pi$.
\end{enumerate}
\end{enonce*}

Prouvons par exemple b) (la démonstration de a) étant analogue et
plus simple). Soit $P$ l'ensemble des parties de $I$ de cardinal
$\leq \Pi$, et soit $T\mapsto i_T:P\rightarrow I$ une application
telle que pour $T\in P$, $i_T$ soit un majorant de $T$ dans $I$;
l'existence de cette application exprime simplement l'hypothèse que
$I$ est grand $\Pi$. Soit $A$ un ensemble bien ordonné tel que
$\card A=\Pi$, et tel que pour tout $a\in A$, l'ensemble des $a'\leq
a$ soit de cardinal $<\Pi$. Il s'ensuit en particulier, comme
$\Pi'_{\circ}<\Pi$, que $A$ \emph{est grand devant} $\Pi'_{\circ}$.
Définissons par récurrence transfinie des applications
$$
S\longmapsto S_a: P\longrightarrow P\quad(a\in A){\quoi},
$$
par les formules suivantes:

$S_{\circ}=S$, $S_{a+1}=S_a\cup \varphi(S_a)\cup\{i_{S_a}\}$,
$S_a=\bigcup_{a<a}S_a$ si a un ordinal limite.

Posons\pageoriginale alors, pour $S\in P$,
$$
J(S)=\bigcup_{a\in A}S_a{\quoi}.
$$
Il est clair alors que $J(S)\supset S$, que $J(S)$ est stable par
$\varphi$, que $\card J(S)\leq \Pi$ (puisque $\card A=\Pi$), enfin
que $J(S)$ est grand devant $\Pi'_{\circ}$ (en utilisant le fait que
$A$ est grand devant $\Pi'_{\circ}$).

Cela démontre \hyperref[I.lem9.21.16]{9.21.16} et achève la
démonstration de \hyperref[I.lem9.21.9]{9.21.9}.

Signalons aussi une variante de \hyperref[I.lem9.21.9]{9.21.9}:

\begin{enonce*}{Lemme 9.21.17}\etiquette[9.21.17]{I.lem9.21.17}
Les notations sont celles de \hyperref[I.lem9.21.9]{9.21.9}. On désigne par $F'$ la
sous-catégorie pleine $\Hhomcart_B(B,E)$ de $F=\Hhom_B(B,E)$ formée
des sections cartésiennes de $E$ sur $F$. Alors:
\begin{enumerate}
\renewcommand{\theenumi}{\roman{enumi}}
\renewcommand{\labelenumi}{(\theenumi)}
\item Tout objet de $F'$ est accessible. Plus précisément, soit
$\alpha\mapsto X(\alpha)$ un élément de $F'$, et soit $d$ un
cardinal tel que $d\geq \Pi'_{\circ}$, $d\geq c$, et tel que pour
tout $\alpha\in \ob B$, $X(\alpha)$ soit un objet $d$-accessible de
$E$. Alors $X$ est un objet $d$-accessible de $F'$.

\item Supposons $\Pi\leq c$ et que les foncteurs
$f^\ast:E_{\beta}\rightarrow E_{\alpha}$ soient
$\Pi'_{\circ}$-accessibles, ou que les catégories $E_{\alpha}$
soient stables par petites limites inductives filtrantes, et que les
foncteurs $f^\ast:E_{\beta}\rightarrow E_{\alpha}$ commutent aux
dites limites inductives. Alors tout objet $X$ de $F'$ est isomorphe
à un objet de la forme $\varinjlim_I X_i$, où pour tout $i\in I$,
$X_i$ est un objet de la sous-catégorie pleine ${F'}^{\Pi}=F'\cap
F^{\Pi_{\circ}}$ de $F'$, et où $I$ est grand devant $\Pi$.
\end{enumerate}
\end{enonce*}

La\pageoriginale démonstration étant toute analogue à celle de
\hyperref[I.lem9.21.9]{9.21.9}, nous nous contentons d'indiquer les points où une
modification de cette dernière est nécessaire. On note d'abord:

\begin{enonce*}{Lemme 9.21.18}\etiquette[9.21.18]{I.lem9.21.18}
L'énoncé \hyperref[I.lem9.21.10]{9.21.10} reste valable lorsqu'on $y$ remplace
$F=\Hhom_B(B,E)$ par $F'=\Hhomcart_B(B,E)$, pourvu que l'on suppose
que les foncteurs images inverse $f^\ast:E_{\beta}\rightarrow
E_{\alpha}$ commutent aux limites inductives de type $I$.
\end{enonce*}

Si donc sous les conditions de \hyperref[I.lem9.21.17]{9.21.17}, $d$ est un cardinal tel que
$d\geq c$ et $d\geq \Pi'_{\circ}$, il résulte de ce qui précède et
de l'hypothèse (9.21.3 c)) que pour tout ensemble
ordonné $I$ grand
devant $d$, les limites inductives de type $I$ sont représentables
dans $F'$ et se calculent argument par argument, d'où la conclusion
\hyperref[I.lem9.21.17]{9.21.17} (i) par un argument immédiat.

Pour prouver (ii), il faut donner un complément à \hyperref[I.lem9.21.13]{9.21.13}:

\begin{enonce*}{Lemme 9.21.19}\etiquette[9.21.19]{I.lem9.21.19}
Sous les conditions de \hyperref[I.lem9.21.13]{9.21.13}, supposons que $X\in \ob F'$ \ie
que $X$ soit une section cartésienne de $E$ sur $B$. Alors on peut
renforcer la conclusion par l'assertion qu'il existe une fonction
$μ$ qui, à tout $i\in I$ et toute flèche $f:\alpha\rightarrow
\beta$ de $B$, associe un morphisme $μ(i,f):f^\ast
X(\beta)_i\rightarrow X(\alpha)_{\varphi(i)}$ dans $E_{\alpha}$, de
telle façon que l'on ait:

\begin{enumerate}
\renewcommand{\theenumi}{\alph{enumi}}
\renewcommand{\labelenumi}{\theenumi)}
\setcounter{enumi}{3}
\item Pour tout $i\in I$ et toute flèche $f:\alpha\rightarrow \beta$
de $B$, les deux diagrammes suivants sont commutatifs:

\begin{displaymath}
\xymatrix@L=5pt@R=1.25cm@C=2.2cm{
                       & f^{\ast}X(\beta)_{\varphi^{2}(i)}\\
X(\alpha)_{\varphi(i)}\ar[ru]^{\lambda(\varphi (i),f)} &\\
                       & f^{\ast}X(\beta)_i \ar[lu]^{\mu(i,f)}\ar[uu]
} \qquad\qquad \xymatrix@L=5pt@R=1.35cm@C=2cm{
X(\alpha)_{\varphi^2(i)} & \\
                         & f^{\ast}X(\alpha)_{\varphi(i)} \ar[lu]_{\mu(\varphi(i),f)}\\
X(\alpha)_i \ar[uu]\ar[ru]_{\lambda(i,j)}             &. }
\end{displaymath}
\end{enumerate}
\end{enonce*}

Pour\pageoriginale le prouver, on procède comme dans 9.21.3, en
considérant le morphisme composé $f^\ast(X(\beta)_i)\rightarrow
f^\ast(X(\beta))\xrightarrow{X(f)^{-1}}X(\alpha)$, où (comme déjà
dans l'écriture de la condition d)) on identifie dans les notations
un $f$-morphisme $R\rightarrow S$ de $E$ avec le morphisme
correspondant $R\rightarrow f^\ast(S)$ de $E_{\alpha}$. Comme
$X(\alpha)=\varinjlim_j X(\alpha)_j$, et $I$ grand devant $c$, on
peut factoriser le morphisme précédent par un des $X(\alpha)_j$, où
$j$ \textit{a priori} dépend de $i$ \emph{et} de $f$, mais peut être
choisi indépendant de $f$, soit $\psi(i)$. Quitte à agrandir la
fonction $\varphi$ de \hyperref[I.lem9.21.13]{9.21.13}, on peut supposer que $\psi=\varphi$.
En procédant comme pour les conditions b) et c) de
\hyperref[I.lem9.21.13]{9.21.13}, on voit
que, quitte à agrandir encore l'application $\varphi$, on peut
supposer que les deux diagrammes de \hyperref[I.lem9.21.19]{9.21.19} d) sont commutatifs.

\setcounter{subsection}{21} \setcounter{subsubsection}{19}
\subsubsection{}\etiquette[9.21.20]{I.sec9.21.20}
L'application $\varphi$ étant choisie comme dans
\hyperref[I.lem9.21.19]{9.21.19}, reprenons
l'argument de 9.21.14, où il faut cependant supposer que $J$
satisfait, en plus des conditions énoncées a) b) c), à la condition:
\begin{enumerate}
\renewcommand{\theenumi}{\alph{enumi}}
\renewcommand{\labelenumi}{\theenumi)}
\setcounter{enumi}{3}
\item Pour tout flèche $f:\alpha\rightarrow \beta$ dans $B$, le
foncteur $f^\ast:E_{\beta}\rightarrow E_{\alpha}$ commute aux
limites inductives de type $J$.
\end{enumerate}

Je\pageoriginale dis que, moyennant cette condition supplémentaire,
la section $X_J$ de $E$ sur $B$ est cartésienne, \ie pour toute
flèche $f:\alpha\rightarrow \beta$ de $B$, le morphisme
\begin{equation}
X_J(\alpha)\longrightarrow f^\ast X_J(\beta)\tag{$\ast$}
\end{equation}
est un isomorphisme. En effet, en vertu de la condition d)
ci-dessus, le but du morphisme envisagé s'identifie à
$\varinjlim_{i\in J}f^\ast(X(\beta)_i)$, et le morphisme s'obtient
par passage aux $\varinjlim_J$ à partir du morphisme de ind-objets
dans $E_{\alpha}$
$$
(X(\alpha)_i)_{i\in J}\longrightarrow (f^\ast)(X(\beta)_i)_{i\in
J}{\quoi},
$$
déduit des morphismes $\lambda(i,f):X(\alpha)_i\rightarrow
f^\ast(X(\beta)_{\varphi(i)})$ pour $i\in J$. Or les conditions
explicitées dans \hyperref[I.lem9.21.19]{9.21.19} nous assurent que
l'homomorphisme précédent
de ind-objets est en fait un isomorphisme, un inverse étant obtenu
par l'homomorphisme déduit du système des $\mu(i,f)$, pour $i\in J$.
Cela prouve notre assertion que $(\ast)$ est un %%
isomorphisme, donc que $X\in \ob F'$.

\subsubsection{}\etiquette[9.21.21]{I.sec9.21.21}
Nous allons supposer maintenant que les foncteurs
$f^\ast:E_{\beta}\rightarrow E_{\alpha}$ sont
$\Pi'_{\circ}$-\emph{accessibles.} Alors les conditions sur $J$
envisagées dans \hyperref[I.sec9.21.20]{9.21.20} sont satisfaites si $J$ est grand devant
$\Pi'_{\circ}$, stable par $\varphi$ et tel que $\card J\leq \Pi$,
et on achève la démonstration de \hyperref[I.lem9.21.17]{9.21.17} comme en
\hyperref[I.sec9.21.15]{9.21.15}.

\begin{enonce*}{Théorème 9.22}\etiquette[9.22]{I.them9.22}
Soit\pageoriginale $p:E\rightarrow B$ un foncteur fibrant, où $B$
est une catégorie essentiellement équivalente à une petite
catégorie, et où pour toute flèche $f:\alpha\rightarrow \beta$ dans
$B$ le foncteur image inverse $f^\ast:E_{\beta}\rightarrow
E_{\alpha}$ est accessible (\hyperref[I.def9.2]{9.2}). Supposons de plus que pour tout
$\alpha\in \ob E$, la catégorie fibre $E_{\alpha}$ admette une
filtration cardinale (\hyperref[I.def9.12]{9.12}) et que tout objet de $E_{\alpha}$ soit
accessible (\hyperref[I.def9.3]{9.3}). Alors chacune des catégories
$F=\Hhom_B(B,E)$ et
$F'=\Hhomcart_B(B,E)$ admet une petite sous-catégorie pleine
génératrice par épimorphismes stricts (\hyperref[I.def7.1]{7.1}), et
chacun de ses objets est accessible.
\end{enonce*}

Il est immédiat que l'énoncé ne change pas essentiellement quand on
remplace $B$ par une sous-catégorie pleine $B_{\circ}$ telle que le
foncteur d'inclusion $B_{\circ}\rightarrow B$ soit une équivalence,
et $E$ par $E_{\circ}=E\times_B B_{\circ}$. Cela nous permet de
supposer que $B$ est une petite catégorie.

Soit, pour tout $\alpha\in \ob B$, $(\Filt^{\Pi}(E))_{\Pi\geq
\Pi_{\alpha}}$ une filtration cardinale de $E_{\alpha}$. Soit
$c_{\circ}\in 𝒰$ un cardinal satisfaisant les conditions
suivantes:
\begin{equation}
\begin{cases}
& c_{\circ}\geq {\displaystyle{\mathop{\Sup}_{\alpha\in ob
B}}}\Pi_{\alpha}{\quav} c_{\circ}\geq \card %%
\mathrm{F}{}l\, B;\\
& \mbox{pour toute flèche $f:\alpha\rightarrow \beta$ de $B$,
$f^\ast:E_{\beta}\rightarrow E_{\alpha}$ est
$c_{\circ}$-accessible;}\\
& \mbox{pour toute $\alpha\in \ob B$, les objets de
$\Filt^{\Pi_{\alpha}}(E_{\alpha})$ sont $c_\circ$-accessibles.}
\end{cases}\tag{9.22.1}
\end{equation}
Posons
$$
c=2^{c_{\circ}}{\quoi},
$$
de sorte que $c>c_{\circ}$, et pour tout $\alpha\in \ob B$, soit
$$
C_{\alpha}=\Filt^c(E_{\alpha}){\quoi}.
$$
Je\pageoriginale dis que les conditions préliminaires à
\hyperref[I.lem9.21.9]{9.21.9} et
\hyperref[I.lem9.21.17]{9.21.17} sont vérifiées, en faisant
$\Pi=\Pi_{\circ}=c$. C'est clair
pour (9.21.2) et (9.21.3). Pour (9.21.4 a)),
cela résulte de \hyperref[I.cor9.17]{9.17},
et (9.21.4 b)) résulte de la condition \hyperref[I.def9.12]{9.12} c)
pour une filtration
cardinale. Notons d'ailleurs que la condition \hyperref[I.def9.12]{9.12} b) des
filtrations cardinales implique que pour tout $\alpha\in \ob B$, on
a $C^{\Pi}_{\alpha}=C_{\alpha}=\Filt^C(E_{\alpha})$, où
$C^{\Pi}_{\alpha}$ est défini dans (9.21.6). Par suite,
\hyperref[I.them9.22]{9.22} résulte
de \hyperref[I.lem9.21.9]{9.21.9} et de \hyperref[I.lem9.21.17]{9.21.17}, et de façon plus
précise, on a prouvé le

\begin{enonce*}{Corollaire 9.23}\etiquette[9.23]{I.cor9.23}
Sous les conditions de \hyperref[I.them9.22]{9.22}, si $c_{\circ}\in 𝒰$ est un
cardinal satisfaisant aux conditions
(9.22.1), et si $c=2^{c_{\circ}}$, alors la sous-catégorie $F^c$ de
$F$ (\resp ${F'}^c$ de $F'$) formée des $X$ tels que $X(\alpha)\in
\ob \Filt^c(E)$ pour tout $\alpha\in \ob B$ est essentiellement
petite et génératrice par épimorphismes stricts; plus précisément,
tout objet $X$ de $F$ (\resp $F'$) est isomorphe à un objet de la
forme $\varinjlim_I X_i$, où les $X_i$ sont dans $F^c$ (\resp dans
${F'}^c$), et où $I$ est un ensemble ordonné grand devant $c$.
\end{enonce*}

Moyennant des hypothèses légèrement plus fortes dans
\hyperref[I.them9.22]{9.22}, on peut
d'ailleurs préciser considérablement \hyperref[I.them9.22]{9.22} et
\hyperref[I.cor9.23]{9.23}:

\begin{enonce*}{Corollaire 9.24}\etiquette[9.24]{I.cor9.24}
Sous les conditions de \hyperref[I.them9.22]{9.22}, supposons que chacune des
catégories
$E_{\alpha}$ $(\alpha\in \ob E)$ est stable par petites limites
inductives filtrantes, et que pour tout $\Pi\geq \Pi_{\alpha}$,
$\Filt^{\Pi}(E_{\alpha})$ est stable par limites inductives
filtrantes indexées par des ensembles ordonnés filtrants $I$ tels
que $\card I\leq \Pi$ (ce qui renforce légèrement la condition
\hyperref[I.def9.12]{9.12} {\rm b)} des\pageoriginale filtrations cardinales). Dans le
cas où c'est $F'$ qu'on considère, supposons de plus que pour toute
flèche $f:\alpha\rightarrow \beta$ de $B$, le foncteur
$f^\ast:E_{\beta}\rightarrow E_{\alpha}$ commute aux petites limites
inductives filtrantes. Soit enfin $c_{\circ}\in 𝒰$ un
cardinal satisfaisant aux conditions (9.22.1), et considérons,
pour tout cardinal $\Pi\geq c=2^{c_{\circ}}$, la sous-catégorie
strictement pleine $F^{\Pi}$ (\resp ${F'}^{\Pi}$) de $F$ (\resp de
$F'$) formée des $X$ tels que l'on ait $X(\alpha)\in
\Filt^{\Pi}(E_{\alpha})$ pour tout $\alpha\in \ob B$. Alors les
sous-catégories envisagées définissent une filtration cardinale
(\hyperref[I.def9.12]{9.12}) de $F$ (\resp de $F'$).
\end{enonce*}

Il faut vérifier les conditions a), b), c) de \hyperref[I.def9.12]{9.12}. Les
conditions
a) et b) résultent des conditions analogues pour les filtrations
cardinales données des $E_{\alpha}$, et de \hyperref[I.lem9.21.10]{9.21.10}
(\resp\hyperref[I.lem9.21.18]{9.21.18},
compte tenu de la deuxième des conditions (9.22.1)). Reste à prouver
c), dont la première partie résulte aussitôt de
\hyperref[I.lem9.21.9]{9.21.9} (\resp
\hyperref[I.lem9.21.17]{9.21.17}), en faisant $\Pi'_{\circ}=0$, $\Pi_{\circ}=c$. Il reste à
prouver que si on a deux cardinaux $\Pi'\geq \Pi\geq c$, alors pour
tout $X$ dans $F^{\Pi'}$ (\resp ${F'}^{\Pi'}$) on a $X\simeq
\varinjlim_{i\in K}X_i$, avec les $X_i$ dans $F^{\Pi}$ (\resp
${F'}^{\Pi}$), $K$ grand devant $\Pi$, et $\card K\leq
{\Pi'}^{\Pi}$. Pour ceci, reprenons les démonstrations de
\hyperref[I.lem9.21.9]{9.21.9}
(ii) et de \hyperref[I.lem9.21.17]{9.21.17} (ii). On peut supposer que chaque $I_{\alpha}$
est de cardinal $\leq \Pi'$ en vertu de la condition \hyperref[I.def9.12]{9.12} c) sur la
filtration cardinale de $E_{\alpha}$, donc on aura
$$
\card I\leq ({\Pi'}^{\Pi})^{\card \ob B}\leq
({\Pi'}^{\Pi})^c={\Pi'}^{\Pi c}={\Pi'}^{\Pi}{\quoi}.
$$
L'ensemble d'indices $K$ utilisé dans 9.21.15 \resp 9.21.21 est
l'ensemble des parties $J$ de $I$ qui sont filtrantes (\ie grandes
devant $\Pi'_{\circ}=0$), stables par $\varphi$ et telles que $\card
J\leq \Pi$. Il est alors immédiat que l'on a
$$
\card K\leq (\card I)^{\Pi}\leq
({\Pi'}^{\Pi})^{\Pi}={\Pi'}^{\Pi^2}={\Pi'}^{\Pi}{\quoi},
$$
ce\pageoriginale qui achève la démonstration de 9.24.

Pour terminer, il convient de donner un énoncé débarrassé des
hypothèses un peu techniques de \hyperref[I.them9.22]{9.22}, en remplaçant
celles-ci par la conjonction, pour les $E_{\alpha}$, des hypothèses
qui interviennent dans \hyperref[I.prop9.11]{9.11} (qui assurent
l'accessibilité des objets des $E$) et dans \hyperref[I.prop9.13]{9.13} (qui
assurent l'existence d'une filtration cardinale dans $E$):

\begin{enonce*}{Corollaire 9.25}\etiquette[9.25]{I.cor9.25}
Soit $p:E\rightarrow B$ un foncteur fibrant, où $B$ est une
catégorie équivalente à une petite catégorie, et
où pour toute
flèche $f:\alpha\rightarrow \beta$ de $B$, le foncteur
$f^\ast:E_{\beta}\rightarrow E_{\alpha}$ est accessible (\hyperref[I.def9.2]{9.2}). On
suppose de plus que, pour tout $\alpha\in \ob B$, la catégorie fibre
$E_{\alpha}$ satisfait aux conditions suivantes:
\begin{enumerate}
\renewcommand{\theenumi}{\alph{enumi}}
\renewcommand{\labelenumi}{\theenumi)}
\item $E_{\alpha}$ admet une petite sous-catégorie pleine,
génératrice par épimorphismes stricts (\hyperref[I.def7.1]{7.1}).

\item $E_{\alpha}$ est stable par produits fibrés, par noyaux de
doubles flèches, par sommes de deux objets et par conoyaux de
doubles flèches, enfin par petites limites inductives filtrantes
\footnote[1]{Cette dernière condition étant probablement inutile,
\cf 9.13.2.}.

\item Le foncteur $\Ker(u,v)$ sur la catégorie des doubles flèches
de $E_{\alpha}$ est accessible (par exemple, commute aux petites
limites inductives filtrantes).

\item Tout épimorphisme strict de $E_{\alpha}$ (\hyperref[I.sec10.2]{10.2}) est un
épimorphisme strict universel.
\end{enumerate}

Sous ces conditions, chacune des catégories $F=\Hhom_B(B,E)$ et
$F'=\Hhomcart_B$ $(B,E)$ admet une petite sous-catégorie génératrice
par épimorphismes stricts, et tous ses objets sont accessibles
(\hyperref[I.def9.3]{9.3}).
\end{enonce*}
De\pageoriginale plus, \hyperref[I.cor9.24]{9.24} nous donne:

\begin{enonce*}{Corollaire 9.26}\etiquette[9.26]{I.cor9.26}
Sous les conditions de $9.25$, et dans le cas où on considère
$F'=\Hhomcart_B(B,E)$, supposons que les foncteurs
$f^\ast:E_{\beta}\rightarrow E_{\alpha}$ commutent aux petites
limites inductives filtrantes. Alors $F$ (\resp $F'$) admet une
filtration cardinale, qu'on peut expliciter par le procédé de
\hyperref[I.cor9.24]{9.24}.
\end{enonce*}

\section{Glossaire}\etiquette[10]{I.sec10}

Pour la commodité du lecteur, nous rassemblons ici les définitions
de quelques termes utilisés dans les numéros précédents. Nous
désignons par $C$ une catégorie.

\subsection{Cartésien, Cocartésien}\etiquette[10.1]{I.sec10.1}

Un diagramme de $C$
\[
\xymatrix@C=3.25cm@R=1.5cm{
A\ar[r] \ar[d]& B \ar[d]\\
C\ar[r] & D }
\]
est dit \emph{cartésien} s'il est commutatif et si le morphisme
canonique de $A$ dans le produit fibré $C\times_D B$ est un
isomorphisme. Il est dit \emph{cocartésien} si le diagramme
correspondant de la catégorie opposée à \sisi{$𝒜$}{$C$} est
cartésien.

\subsection{Épimorphisme, épimorphisme
  strict}\etiquette[10.2]{I.sec10.2}

\etc\pageoriginale \Cf \hyperref[I.sec10.3]{10.3}.

\subsection{Famille épimorphique, épimorphique
  stricte}\etiquette[10.3]{I.sec10.3}

\etc: Une famille
$$
f_i:A_i\longrightarrow B, i\in I{\quoi},
$$
de flèches de même but dans $C$ est dite famille
\emph{épimorphique}, si pour tout objet $C$ de $C$, l'application
induite par les $f_i$:
\begin{equation}
\Hom_{𝒜}(B,C)\longrightarrow \prod_i
\Hom_{𝒜}(A_i,C)\tag{10.3.1}
\end{equation}
est injective.

On dit que la famille envisagée est \emph{épimorphique stricte} si
l'image de l'application (10.3.1) est formée des familles $(g_i)$
telles que pour tout objet $R$ de $C$, tout couple d'indices $i$,
$j\in I$ et tout couple de flèches $u:R\rightarrow A_i$,
$v:R\rightarrow A_j$ avec $f_i u=f_j v$, on a aussi $g_i u=g_iv$.
Lorsque les produits fibrés $A_i\times_B A_j$ sont représentables
pour $i$, $j\in I$, il revient au même de dire que l'image
essentielle de (10.3.1) est formée des $(g_i)$ telles que pour tout
couple d'indices $i$, $j\in I$, on ait $g_i \pr_1=g_i \pr_2$, où
$\pr_1$, $\pr_2$ sont les deux projections de $A_i\times_B A_j$. On
dit que la famille $(f_i)_{i\in I}$ est une \emph{famille
épimorphique effective} si la condition précédente est
vérifiée, \ie
si elle est épimorphique stricte, et si les produits fibrés
$A_i\times_A A_j$ sont représentables.

On dit que la famille $(f_i)_{i\in I}$ est une famille
\emph{épimorphique universelle,} (\resp\break \emph{épimorphique
effective universelle}) si les morphismes $f_i$ sont quarrables
(\hyperref[I.sec10.3]{10.3}), et si pour toute flèche $B'\rightarrow B$, la
famille des
flèches $f'_i:A'_i\rightarrow A_i\times_B B'$, $i\in I$, déduite de
la famille $(f_i)_{i\in I}$ par changement de base $B'\rightarrow
B$, est épimorphique (\resp épimorphique effective).

La\pageoriginale notion de famille \emph{épimorphique stricte
universelle} sera définie dans (II 2.5, 2.6). Notons que lorsque les
morphismes $f_i$ sont quarrables (par exemple si dans $C$ les
produits fibrés sont représentables) cette notion coïncide avec
celle de famille épimorphique effective universelle (utiliser II
2.4). Dans le cas général, on a entre les six variantes envisagées
de la notion d'épimorphicité les implications logiques:
\[
\xymatrix@C=-1cm@R=.5cm{
  &   & \txt<3cm>{épimorphique effective universelle} \ar@{=>}[lld]\ar@{=>}[rrd] &     & &\\
\txt<3cm>{épimorphique effective} \ar@{=>}[rdd]       &   &    &     &  \txt<3cm>{épimorphique stricte universelle} \ar@{=>}[rd] \ar@{=>}[ld]&\\
 &   &    &   \txt{épimorphique\\ stricte}
         \ar@{=>}[lld]&     & \txt{épimorphique \\universelle} \ar@{=>}[llld]\\
         & \txt{épimorphique}&    &     &      &
}
\]
Un morphisme $f:A\rightarrow B$ de $𝒜$ s'appelle un
\emph{épimorphisme (\resp un épimorphisme strict, \resp un
épimorphisme effectif, \resp un épimorphisme universel, \resp un
épimorphisme effectif universel, \resp un épimorphisme strict
universel)} si la famille de morphisme réduit au seul élément $f$
est épimorphique (\resp\ldots).

\subsection{Famille monomorphique, monomorphique
  stricte}\etiquette[10.4]{I.sec10.4}

\etc Une famille de flèches $f_i:A_i\rightarrow B$ de $C$ est dite
\emph{monomorphique (\resp monomorphique stricte,\ldots)} si en tant
que famille de flèches de la catégorie opposée $C^{\circ}$ elle est
épimorphique (\resp épimorphique stricte,\ldots). Une flèche
$f:A\rightarrow B$ de $𝒜$ est appelée un \emph{monomorphisme
(\resp un monomorphisme strict.\ldots)} si la famille réduite à $f$
est monomorphique (\resp\pageoriginale monomorphique
stricte,\ldots), \ie si $f$ en tant que flèche de la catégorie
opposée $C^{\circ}$ est un épimorphisme (\resp un épimorphisme
strict,\ldots).

\subsection{Monomorphisme, monomorphisme
  strict}\etiquette[10.5]{I.sec10.5}

\etc \Cf \hyperref[I.sec10.4]{10.4}.

\subsection{Projecteur, image d'un projecteur, facteur direct d'un
objet}\etiquette[10.6]{I.sec10.6}

Soit $f:X\rightarrow X$ un endomorphisme d'un objet de $C$. On dit
que $f$ est un \emph{projecteur} si on a $f^2=f$. Alors le couple
$(f,\id_X)$ admet un conoyau représentable $X'$ si et seulement si
il admet un noyau représentable $X''$, et lorsqu'il en est ainsi, il
existe un unique isomorphisme $u:X'\rightarrow X''$ dans $C$ tel que
$iup=f$, où $p:X\rightarrow X'$ et $i:X''\rightarrow X$ sont les
morphismes canoniques.\nde{La représentabilité du noyau, ou du conoyau, est
également équivalente à l'existence de morphismes $i$, $p$ et d'un
isomorphisme $u$ tels que $iup=f$ et $pi=u^{-1}$. En particulier, cette
propriété est stable par n'importe quel foncteur.}
On identifie alors généralement $X'$ et
$X''$, et on l'appelle l'\emph{image du projecteur} $f$; on dit que
le projecteur $f$ \emph{admet une image} si $\Ker(f,\id_X)$ est
représentable \ie $\Coker(f,\id_X)$ est représentable. Un sous-objet
(\resp un quotient) $Y$ de $X$ est appelé un \emph{facteur direct}
de $X$ si on peut trouver un projecteur $f$ dans $X$ se factorisant
par $Y$, et admettant $Y$ comme image en tant que sous-objet (\resp
en tant que quotient) de $X$. On dit parfois, par abus de langage,
qu'un objet de $C$ est un \emph{facteur direct de} $X$, s'il est
isomorphe à l'image d'un projecteur dans $X$.

\subsection{Quarrable}\etiquette[10.7]{I.sec10.7}

Une flèche $f:A\rightarrow B$ de $C$ est dite \emph{quarrable} si
pour toute flèche $B'\rightarrow B$ de $C$, le produit fibré
$A\times_B B'$ est représentable dans $C$.

\subsection{Quotient, quotient strict}\etiquette[10.8]{I.sec10.8}

\etc\pageoriginale

Soit $A$ un objet de $C$. Deux épimorphismes $f:A\rightarrow B$,
$f':A\rightarrow B'$ de source $A$ sont dits \emph{équivalents} s'il
existe un isomorphisme $u:B\xrightarrow{\sim}B'$ tel que $uf=f'$. On
obtient ainsi une relation d'équivalence dans l'ensemble des
épimorphismes de source $A$, dont les classes sont appelées les
\emph{quotients, ou objets quotients,} de $A$. Pour tout quotient de
$A$, on suppose généralement choisi un élément de cette classe, soit
$f:A\rightarrow B$, et on parle souvent (par abus de langage) du
quotient $B$ de $A$ (ou du quotient $f:A\rightarrow B$ de $A$). On
dit que $B$ est un \emph{quotient strict} (\resp un \emph{quotient
effectif, \resp un quotient universel,}\ldots) si le morphisme
$f:A\rightarrow B$ est un épimorphisme strict (\resp un épimorphisme
effectif, \resp un épimorphisme universel,\ldots) (\hyperref[I.sec10.3]{10.3}).

\subsection{Relations d'équivalence}\etiquette[10.9]{I.sec10.9}

Soient $F$, $G$ deux préfaisceaux d'ensembles sur $C$. Un diagramme
$p_1$, $p_2:F\rightrightarrows G$ est appelé une \emph{relation
d'équivalence sur} $G$ si pour tout objet $A$ de $C$, l'application
$$
(p_1(A),p_2(A)):F(A)\longrightarrow G(A)\times G(A)
$$
induit une bijection de $F(A)$ sur le graphe d'une relation
d'équivalence sur l'ensemble $G(A)$.

Un diagramme $p_1$, $p_2:B\rightrightarrows C$ dans $C$ est appelé
une relation d'équivalence sur $C$ si le diagramme de préfaisceaux
correspondant est une relation d'équivalence. Lorsque dans $C$ les
produits finis et produits fibrés sont représentables, un foncteur
$\Phi:C\rightarrow C'$ commutant aux produits finis et produits
fibrés transforme les relations d'équivalence sur $C$ en relation
d'équivalence sur $\Phi (C)$.

Soit\pageoriginale $\pi:C\rightarrow D$ un morphisme de $𝒜$
tel que le produit fibré $C\times_D C$ soit représentable. Alors le
diagramme $\pr_1$, $\pr_2:C\times_D C\rightrightarrows C$ est une
relation d'équivalence sur~$C$.

\subsection{Relation d'équivalence effective, effective universelle}\etiquette[10.10]{I.sec10.10}

Une relation d'équivalence $p_1$, $p_2:R\rightrightarrows A$ sur un
objet $A$ de $C$ est dite \emph{relation d'équivalence effective}
s'il existe un morphisme $\pi:A\rightarrow B$ tel que le carré
\[
\xymatrix@C=3.75cm@R=1.5cm{
R\ar[r]^{p_2} \ar[d]_{p_1}& A \ar[d]^{\pi}\\
A\ar[r] & B }
\]
soit cartésien et cocartésien. (Le morphisme $\pi$ est le conoyau du
couple $(p_1,p_2)$, et est donc déterminé à isomorphisme unique
près). Le morphisme $\pi$ est alors un épimorphisme effectif; si
c'est un épimorphisme effectif universel (\hyperref[I.sec10.3]{10.3}) on dit que la
\emph{relation d'équivalence} est \emph{effective universelle.}

\subsection{Sous-objet, sous-objet strict}\etiquette[10.11]{I.sec10.11}

\etc Ce sont les notions duales de celles de quotient, quotient
strict \etc (\hyperref[I.sec10.8]{10.8}).

\begin{thebibliography}{1}
\bibitem[1]{SGA4-01:I.1} \sisi{B. Mitchell, Theory of categories, Academie Press
(1965).}{\textsc{B. Mitchell}, « Theory of categories », Academie
Press (1965).}
\end{thebibliography}


\newpage

\footnotetext[1]{Nous reproduisons ici, avec son accord, des papiers
secrets de N. \textsc{Bourbaki}. Les références de ce texte se
rapportent à son savant ouvrage.}
\section*{II. Appendice: Univers (par
  N. \textsc{Bourbaki}($^*$))}\pageoriginale

\setcounter{section}{0}
\section{Définition et premières propriétés des univers}

\begin{enonce*}{DÉFINITION 1}
Un ensemble $U$ est appelé un univers s'il satisfait aux conditions:
\begin{itemize}
\item[(U.I)] si $x\in U$ et si $y\in x$, alors $y\in U$;

\item[(U.II)] si $x$, $y\in U$, alors $\{x,y\}\in U$;

\item[(U.III)] si $x\in U$, alors $\mathcal{P}(x)\in U$;

\item[(U.IV)] si $(x_{\alpha})_{\alpha\in I}$ est une famille
d'éléments de $U$, et si $I\in U$, alors la réunion
$\bigcup_{\alpha\in I}x_{\alpha}$ appartient à $U$;

\item[(U.?)] si $x$, $y\in U$, alors le couple $(x,y)\in U$.
\end{itemize}
\end{enonce*}
\noindent N.B. Comme il a été je crois décidé pour les prochaines
éditions, on définit le couple à la \sisi{Kuratowaki}{Kuratowski}
par $(x,y)=\{\{x,y\}, \{x\}\}$, la condition (U.?) est inutile car
elle résulte de (U.II).

\begin{enonce*}[remark]{Exemples}
\begin{enumerate}
\renewcommand{\labelenumi}{\theenumi)}
\item L'ensemble \emph{vide} est un univers noté $U_\circ$.

\item Considérons les mots finis non vides formés avec les quatre
symboles $''\{'',''\}'','',''$ et $''∅''$ (\cf Alg. I).
Définissons, par récurrence\pageoriginale sur la longueur $n$ d'un
tel mot, la notion de mot \emph{significatif}:
\begin{enumerate}
\renewcommand{\theenumii}{\alph{enumii}}
\renewcommand{\labelenumii}{(\theenumii)}
\item Pour $n=1$, seul le mot $∅$ est significatif;

\item pour qu'un mot $A$ de longueur $n$ soit significatif, il faut et
suffit qu'il existe $p$ mots significatifs distincts
$A_1,\ldots,A_p$ $(p\geq 1)$ de longueurs $<n$ tels que
$$
A=\{A_1,A_2,\ldots,A_p\}{\quoi}.
$$
\end{enumerate}
\end{enumerate}
\end{enonce*}
\noindent Par exemple $∅$, $\{∅\}$,
$\{\{∅\},\{\{∅\}\},\{∅,\{∅\}\}\}$
sont des mots significatifs. Il est clair, par récurrence sur la
longueur, que tout mot significatif désigne un terme de la Théorie
des Ensembles; par exemple, si $A=\{A_1,A_2,\ldots,A_p\}$, A
désigne l'ensemble dont les éléments sont $A_1$, $A_2,\ldots,$ et
$A_p$; ces termes sont évidemment des \emph{ensembles finis.} Soit
$U_1$ l'ensemble des ensembles ainsi obtenus. On vérifie aisément
que $U_1$ satisfait aux conditions (U.I), (U.II), (U.III) et (U.IV)
de la déf. 1 (mais non à cette idiote de (U.?), ce qui n'est pas
grave si on veut bien décanuler le couple). Donc $U_1$ est un
\emph{univers}. On notera que les éléments de $U_1$ sont finis, et
que $U_1$ est dénombrable.

Dans les énoncés qui suivent, $U$ \emph{désigne un univers.}

\begin{enonce*}{PROPOSITION 1}
Si $X\in U$ et si $y\subset x$, alors $y\in U$.
\end{enonce*}
En effet, on a $\mathcal{P}(x)\in U$ par (U.III), d'où $y\in
\mathcal{P}(x)$ et $y\in U$ par (U.I).

\begin{enonce*}{COROLLAIRE}
Si $x\in U$, tout ensemble quotient $y$ de $x$ est élément de $U$.
\end{enonce*}
En effet $y$ est une partie de $\mathcal{P}(x)$. D'où $y\in U$ par
(U.III) et la prop. 1.

\begin{enonce*}{PROPOSITION 2}
Si\pageoriginale $x\in U$, on a $\{x\}\in U$.
\end{enonce*}
\c Ca résulte de (U.II) appliqué pour $y=x$.

\begin{enonce*}{PROPOSITION 3}
Tout couple, tout triplet, tout quadruplet d'éléments de $U$ est un
élément de $U$.
\end{enonce*}
C'est vrai pour les couples d'après le N.B. ou (U.?). On en déduit
le cas des triplets car $(x,y,z)=((x,y),z)$, puis celui des
quadruplets car $(x,y,z,t)=((x,y,z),t)$.

\begin{enonce*}{PROPOSITION 4}
Si $X$, $Y\in U$, alors $X\times Y\in U$.
\end{enonce*}
En effet, pour $x\in X$, $\{x\}\times Y$ est la réunion de la
famille $\{(x,y)\}_{y\in Y}$; comme on a $\{(x,y)\}\in U$ d'après
(U.I), (U.II) et la prop.3, on a $\{x\}\times Y\in U$ d'après
(U.IV). Enfin $X\times Y$ est la réunion de la famille $(\{x\}\times
Y)_{x\in X}$; c'est donc un élément de $U$ d'après (U.IV) encore.

\begin{enonce*}{COROLLAIRE 1}
Si $X$, $Y$, $Z,\ldots$ sont des éléments de $U$, tous les ensembles
de l'échelle construite sur $X$, $Y$, $Z,\ldots$ sont des éléments
de $U$ {\rm(\cf chap.IV, \S)}.
\end{enonce*}

\c Ca résulte en effet d'applications successives de la prop. 4 et
de (U.III).

\begin{enonce*}{COROLLAIRE 2}
Si $(x_{\alpha})_{\alpha \in I}$ est une famille d'éléments de $U$
et si $I\in U$, l'ensemble somme $\Sigma_{\alpha \in I} x_{\alpha}$
est un élément de $U$.
\end{enonce*}
En effet, cet ensemble somme est une partie du produit
$(\bigcup_{\alpha\in I} x_{\alpha})\times I$ (chap. II), produit
dont les deux facteurs sont éléments de $U$ (par (U.IV) et
l'hypothèse). On applique alors les prop. \hyperref[I.sec4]{4} et 1.

\begin{enonce*}{PROPOSITION 5}
Si\pageoriginale $X$ et $Y$ sont des éléments de $U$, toute
correspondance entre $X$ et $Y$ (en particulier toute application de
$X$ dans $Y$) est un éléments de $U$.
\end{enonce*}
En effet, une telle correspondance $C$ est un triplet $(X,Y,\Gamma)$
où $\Gamma$ est une partie de $X\times Y$ (le graphe de $C$) (chap.
II, \S). On a $\Gamma\in U$ d'après les prop. 4 et 1. D'où $C\in U$
par la prop. 3.

\begin{enonce*}{PROPOSITION 6}
Si $X$, $Y\in U$, tout ensemble $Z$ de correspondances entre $X$ et
$Y$ (en particulier d'applications de $X$ dans $Y$) est élément de
$U$.
\end{enonce*}
En effet, $Z$ est une partie de $\{X\}\times \{Y\}\times
\mathcal{P}(X\times Y)$. Or ce produit est élément de $U$
d'après la
prop.2 et le cor. à la prop.\hyperref[I.sec4]{4}. On a donc $Z\in
U$ par la prop. 1.

\begin{enonce*}{COROLLAIRE}
Si $(x_{\alpha})_{\alpha \in I}$ est une famille d'éléments de $U$,
et si $I\in U$, on a $\Pi_{\alpha\in I}x_{\alpha} \in U$.
\end{enonce*}
En effet, ce produit est un ensemble d'applications de $I$ dans
$\bigcup_{\alpha\in I}x_{\alpha}$, et cette réunion est élément de
$U$ par (U.IV).

\begin{enonce*}{PROPOSITION 7}
Si $X$ est une partie de $U$ dont le cardinal est au plus celui d'un
élément de $U$, alors $X$ est un élément de $U$.
\end{enonce*}
Soit $I$ un élément de $U$ tel que $\card(X)\leq \card(I)$. On a une
surjection $i\mapsto x_i$ d'une partie $I'$ de $I$ sur $X$. Alors
$X$ est la réunion de la famille $(\{x_i\})_{i\in I'}$; or cette
réunion est élément de $U$ par la prop. 1, la prop. 2 et (U.IV).

\begin{enonce*}{COROLLAIRE}
Si\pageoriginale $U$ est non vide toute partie finie de $U$ est un
élément de $U$, et $U$ a des éléments de cardinal fini arbitraire.
\end{enonce*}
Par contre, si $U=∅$, $∅$ est une partie finie
de $∅$ mais non un élément de $∅$.

En effet, si $U$ est non vide, la prop. 2 montre qu'un ensemble
$x_{\circ}$ à un élément appartient à $U$. Par récurrence sur $n$,
posons $x_{n+1}=\mathcal{P}(x_n)$. On %%
a $x_n\in U$ par (U.III), et $\card(x_{n+1})=2^{\card(x_n)}$, de
sorte que $U$ a des éléments de cardinal fini arbitrairement grand.

\begin{enonce*}[remark]{Remarque}
Il résulte de la prop. 2 et du cor. à la prop. 7 que tout univers
non vide $U$ contient l'univers $U_1$ de l'ex. 2; en effet on a
$∅\in U$ d'après la prop. 1. Ainsi $U_1$ est
l'intersection de tous les univers non vides. Plus généralement:
\end{enonce*}

\begin{enonce*}{PROPOSITION 8}
Si $(U_{\lambda})_{\lambda\in L}$ est une famille non vide
d'univers, alors $U=\bigcap_{\lambda\in L}U_{\lambda}$ est un
univers.
\end{enonce*}
Ceci résulte aussitôt de la déf. 1.

\section{Univers et espèces de structures}

Soit $\mathcal{E}$ une espèce de structure; supposons, pour fixer
les idées et alléger l'exposé, que chaque structure d'espèce
$\mathcal{E}$ est définie sur un ensemble de base. Soit $(X,S)$ une
structure d'espèce $\mathcal{E}$ ($X$ étant l'ensemble de base, et
$S$ la structure), et soit $U$ un \emph{univers.} Si $X\in U$, alors
tous les objets constitutifs de la structure $S$ sur $X$ sont
éléments de $U$ (par le cor. 1 de la prop. 4, ${\rm n^{\circ}}$ 1 et
par (U.I)) de sorte que la structure $(X,S)$ est élément de $U$.

Supposons\pageoriginale maintenant que $\mathcal{E}$ soit une espèce
de structure avec \emph{morphismes}. $X$ et $X'$ sont des éléments
de $U$ munis de structures d'espèce $\mathcal{E}$, alors l'ensemble
des morphismes de $X$ dans $X'$ est encore un élément de $U$ (${\rm
n^{\circ}}$ 1, prop. 6).

Considérons alors la \emph{catégorie} $(U\text{-}\mathcal{E})$
définie au \S 1, ${\rm n^{\circ}}$ 2, ex. c). Comme les objets et
les flèches de $(U\text{-}\mathcal{E})$ sont des éléments de $U$,
$(U\text{-}\mathcal{E})$ est un \emph{couple de parties} de $U$,
muni d'un quadruplet d'applications; ce n'est pas, en général, un
élément de $U$. La notation $X\in (U\text{-}\mathcal{E})$ voudra
dire que $X$ est un élément de $U$ muni d'une structure d'espèce
$\mathcal{E}$.

La catégorie $(U\text{-}\mathcal{E})$ est \emph{stable} pour de
nombreuses opérations:

\begin{enumerate}
\renewcommand{\theenumi}{\alph{enumi}}
\renewcommand{\labelenumi}{\theenumi)}
\item Si $X\in(U\text{-}\mathcal{E})$, et si $X'$ est une partie de $X$ qui
admet une \emph{structure induite}, alors $X'\in
(U\text{-}\mathcal{E})$. \c Ca résulte de la prop. 1, ${\rm
n^{\circ}}$ 1.

\item Si $X\in (U\text{-}\mathcal{E})$, et si $X''$ est un ensemble quotient de
$X$ qui admet une \emph{structure quotient}, alors $X''\in
(U\text{-}\mathcal{E})$ (cor. de la prop. 1, ${\rm n^{\circ}}$ 1).

\item Si $X$, $Y\in (U\text{-}\mathcal{E})$, et si $X\times Y$ admet une
\emph{structure produit}, alors $X\times Y\in
(U\text{-}\mathcal{E})$ (prop. 4, ${\rm n^{\circ}}$ 1). Plus
généralement, si $(X_i)_{i\in I}$ est une famille d'éléments de
$(U\text{-}\mathcal{E})$, \emph{si on a} $I\in U$, et si $\Pi_{i\in
I}X_i$ admet une \emph{structure produit}, alors $\Pi_{i\in I}X_i
\in (U\text{-}\mathcal{E})$ (cor. de la prop. 6). Assertions
analogues pour les \emph{structures sommes} (cor. 2 de la prop. 4).

\item Soient $I$ un ensemble préordonné, et $(X_i, f_{ij})_{i,j\in
I}$ un \emph{système projectif} (\resp \emph{inductif}) d'ensembles
munis de structures d'espèce $\mathcal{E}$ et de morphismes; soit
$L$ la limite de ce système, au sens du chap. III. Si $X_i\in U$
pour tout $i$, si $I\in U$, et si $L$ admet une structure limite
projective (\resp inductive), alors on a $L\in
(U\text{-}\mathcal{E})$: en effet $L$ est une\pageoriginale partie
de $\Pi_{i\in I}X_i$ (\resp un ensemble quotient de $\Sigma_{i\in
I}X_i$), et est donc un élément de $U$ d'après le ${\rm n^{\circ}}$
1.
\end{enumerate}

Donnons encore quelques exemples plus particuliers:
\begin{itemize}
\item[* 1)] Soient $X$, $Y\in (U\text{-}\Top)$ deux espaces \emph{localement
compacts}. Alors l'ensemble $\mathscr{C}(X,Y)$ des applications
continues de $X$ dans $Y$, muni de la topologie de la convergence
compacte (Top. Géné. $X$), est un élément de $(U\text{-}\Top)$. En
effet on a $\mathscr{C}(X,Y)\in U$ d'après la prop. 6 du ${\rm
n^{\circ}}$ 1.${}_{\mbox{*}}$

\item[* 2)] Soient $X$, $Y\in (U\text{-}\Ab)$. Alors le groupe
$\Hom_{\ZZ}(X,Y)$ est un élément de $(U\text{-}\Ab)$ par la prop. 6
du ${\rm n^{\circ}}$ 1. Si on suppose de plus qu'on a $\ZZ\in U$, le
produit tensoriel $X\otimes_{\ZZ}Y$ est un élément de
$(U\text{-}\Ab)$; en effet ce produit tensoriel est un quotient de
$\ZZ^{(X\times Y)}$, lequel est une partie de $\ZZ^{X\times Y}$, qui
lui-même est un élément de $U$ par les prop. 4 et 6 du ${\rm
n^{\circ}}$ 1.${}_{\mbox{*}}$

\item[* 3)] Soit $X\in (U\text{-}\Unif)$ un espace uniforme. Alors son
\emph{complété} $\widehat{X}$ est un élément de $(U\text{-}\Unif)$.
En effet $\widehat{X}$ est un ensemble de classes d'équivalences de
filtres de Cauchy sur $X$; or un filtre sur $X$ est un élément de
$\mathcal{P}(\mathcal{P}(X))$, donc une classe d'équivalence de
filtres est un élément de
$\mathcal{P}(\mathcal{P}(\mathcal{P}(X)))$, de sorte que
$\widehat{X}$ est un élément de
$\mathcal{P}(\mathcal{P}(\mathcal{P}(\mathcal{P}(X))))$, donc un
élément de $U$ par (U.III) et (U.I).${}_{\mbox{*}}$
\end{itemize}

N.B. Moralité, Bourbaki devra veiller à bien « canonifier » ses
constructions. Par exemple, dans celles où on adjoint un élément
$\infty$ (compactifié d'Alexandroff, corps projectif), il y aura
intérêt à prendre $\infty =∅$\pageoriginale (car
$∅$ est élément de tout univers non vide) et à former
l'ensemble somme de $\{∅\}$ et de l'ensemble donné. Bien
entendu il faudrait aussi « canonifier » l'ensemble à deux
éléments servant à construire les ensembles sommes; ainsi
$\{∅, \{∅\}\}$ me paraît un bon candidat car il
appartient à tout univers non vide.

\section{Univers et catégories}

Soient $U$ un univers et $C$ une catégorie. Nous écrivons $C\in U$
si on a $\Ob(C)\in U$ et $%%
\mathrm{F}{}l\, (C)\in U$. Cette écriture est justifiée du fait que
$C$ est le sextuplet formé de $\Ob(C)$, de
$\mathrm{F}{}l\,(C)$ et des quatre applications structurales; donc
si $\Ob(C)\in U$ et
$\mathrm{F}{}l\,(C)\in U$, ces quatre applications sont éléments de
$U$ (${\rm n^{\circ}}$ 1, cor. 1 de la prop. 4), donc aussi le
sextuplet $C$ (${\rm n^{\circ}}$ 1, prop. 3, \textit{cum grano
salis}). Ceci est d'ailleurs un cas particulier du ${\rm n^{\circ}}$
2, si l'on considère l'espèce de structure « $\cat$ » de catégorie.
Avec les notations du ${\rm n^{\circ}}$ 2, la relation $C\in U$
s'écrit aussi $C\in (U\text{-}\cat)$.

On notera qu'une catégorie comme $(U\text{-}\Ens)$,
$(U\text{-}\Ord)$, $(U\text{-}\Top)$ ou $(U\text{-}\Gr)$ n'est
\emph{pas}, en général, un élément de $U$.

Soient $C$, $D$ deux catégories et $U$ un univers tels que $C$,
$D\in U$. Si $C'$ est une \emph{sous-catégorie} de $C$, on a $C'\in
U$; la \emph{catégorie produit} $C\times D$, la catégorie somme de
$C$ et $D$, les catégories opposées $C^{\circ}$ et $D^{\circ}$ sont
aussi éléments de $U$ (\cf ${\rm n^{\circ}}$ 2). La \emph{catégorie
de foncteurs} $E=\Hhom(C,D)$ est également un élément de $U$: en
effet $\Ob(E)$ est un ensemble de couples d'applications
$\Ob(C)\rightarrow \Ob(D)$, $F\ell\, (C)\rightarrow F\ell\, (D)$,
donc $\Ob(E)\in U$ par le ${\rm n^{\circ}}$ 1;\pageoriginale quant à
$F\ell\, (E)$, c'est un ensemble de morphisme fonctoriels,
c'est-à-dire d'applications $\Ob(C)\rightarrow F\ell\, (D)$, d'où $
F\ell\, (E)\in U$ par la prop. 6 du ${\rm n^{\circ}}$ 1.

\begin{enonce*}{PROPOSITION 9}
Soient $\mathcal{E}$ une espèce de structure avec morphismes, et
$U$, $U'$ deux univers tels que $U'\subset U$. Alors
$(U'\text{-}\mathcal{E})$ est une sous-catégorie pleine de
$(U\text{-}\mathcal{E})$.
\end{enonce*}

En effet, si $x$ et $y$ sont des objets de
$(U'\text{-}\mathcal{E})$, ce sont par définition des objets de
$(U\text{-}\mathcal{E})$. Quant à $\Hom(x,y)$, c'est le même
ensemble dans les deux catégories (\cf. \S 1, ${\rm n^{\circ}}$ 2,
ex. c)).

N.B. On notera que l'hypothèse que $U$ et $U'$ sont des univers %%
est inutile, de sorte que la prop. \hyperref[I.sec9]{9} remonterait avantageusement au
\S 1, ${\rm n^{\circ}}$ 4. Mais la Tribu a demandé qu'elle soit ici.

\begin{enonce*}[remark]{Remarque}
Il arrive que les axiomes de l'espèce de structure $\mathcal{E}$
impliquent que les \emph{cardinaux} des ensembles munis de
structures d'espèce $\mathcal{E}$ soient \emph{bornés} par un
cardinal fixe, soit $\sheaf{\sheaf{c}}$ (* par exemple l'espèce de
structure de groupe fini, de groupe de type fini, de module de type
fini sur un anneau fixe $A$, ou d'algèbre de type fini sur $A$
${}_{\mbox{*}}$). Supposons alors qu'il existe un élément $z$ de
$U'$ tel que $\sheaf{\sheaf{c}}\leq \card(z)$. Alors, pour tout
ensemble $x$ muni d'une structure d'espèce $\mathcal{E}$, il existe
un élément $x'$ de $U'$ équipotent à $x$, par exemple une partie de
$z$; munissons $x'$ de la structure déduite de celle de $x$ par
transport de structure; on obtient ainsi un élément de
$(U'\text{-}\mathcal{E})$. Il s'ensuit que \emph{le foncteur
d'inclusion de $(U'\text{-}\mathcal{E})$ dans
$(U\text{-}\mathcal{E})$ est alors essentiellement surjectif {\rm
(\S 4, ${\rm n^{\circ}}$ 2, déf. 2)} et est donc une équivalence de
catégorie} (\S 4, ${\rm n^{\circ}}$ 3, th. 1).
\end{enonce*}

\section{L'axiome des univers}\pageoriginale

Les fort agréables résultats de stabilité du ${\rm n^{\circ}}$ 2
n'ont d'intérêt que si on peut les appliquer à autre chose qu'aux
deux petits univers $U_{\circ}$ et $U_1$ décrits au ${\rm
n^{\circ}}$ 1. Nous ajouterons donc aux axiomes de la Théorie des
Ensembles l'axiome suivant:

(A.6) \emph{Pour tout ensemble $x$, il existe un univers $U$ tel que
$x\in U$.}

Cet axiome implique que, si $(x_{\alpha})_{\alpha\in I}$ est une
\emph{famille} d'ensembles, il existe un univers $U$ tel que
$x_{\alpha}\in U$ pour tout $\alpha\in I$: il suffit, en effet,
d'appliquer (A.6) à $x=\bigcup_{\alpha\in I}x_{\alpha}$ et
d'appliquer la prop. 1 du ${\rm n^{\circ}}$ 1.

En particulier, étant donnée une \emph{catégorie} $C$, il existe un
univers $U$ tel que $C\in U$ au sens du ${\rm n^{\circ}}$ 3: on
applique ce qui précède à la famille $(\Ob(C), F\ell\, (C))$. Ceci
s'applique aux catégories de la forme $(V\text{-}\mathcal{E})$ où
$\mathcal{E}$ est une espèce de structure avec morphismes et $V$ un
ensemble, un univers par exemple; en général on a $V\neq U$.

Par exemple, si $V$ est un univers, on a $\Ob(V\text{-}\Ens)=V$ et
$F\ell\, (V\text{-}\Ens)\subset V$ (${\rm n^{\circ}}$ 1, prop. 5).
Les univers $U$ tels que $(V\text{-}\Ens)\in U$ sont donc ceux tels
que $V\in U$. Or la relation $V\in V$ est impossible pour un
univers: en effet, pour toute partie $A$ de $V$, on aurait $A\in V$
(${\rm n^{\circ}}$ 1, prop. 1), d'où $\card(\mathcal{P}(V))\leq
\card(V)$, ce qui est impossible.

\section{Univers et cardinaux fortement inaccessibles}\pageoriginale

Soit $U$ un univers. D'après la condition (U.I) de la déf. 1 tout
élément $x$ de $U$ est une partie de $U$; on a donc $\card(x)\leq
\card(U)$. Comme les cardinaux des parties de $U$ forment un
ensemble bien ordonné, le cardinal
\begin{equation}
\sheaf{\sheaf{c}}(U)={\sup}_{x\in U}\card(x)\tag{1}
\end{equation}
existe. Notons que, pour tout cardinal
$\sheaf{\sheaf{c}}<\sheaf{\sheaf{c}}(U)$, il existe un élément $x$
de $U$ tel que $\card(x)=\sheaf{\sheaf{c}}$; en effet, il existe par
définition $y\in U$ tel que $\sheaf{\sheaf{c}}\leq \card(y)\leq
\sheaf{\sheaf{c}}(U)$, et l'on prend pour $x$ une partie convenable
de $y$. Réciproquement, si $x\in U$, on a
$\card(x)<\sheaf{\sheaf{c}}(U)$; en effet on a $\mathcal{P}(x)\in U$
par (U.III), d'où $2^{\card(x)}\leq \sheaf{\sheaf{c}}(U)$. Le
cardinal $\sheaf{\sheaf{c}}(U)$ a donc les propriétés suivantes:

\begin{enumerate}
\renewcommand{\labelenumi}{\theenumi)}
\item Si $\sheaf{\sheaf{c}}$ est un cardinal $<\sheaf{\sheaf{c}}(U)$, on a
$2^{\sheaf{\sheaf{c}}}<\sheaf{\sheaf{c}}(U)$; en effet, si $x$ est
un élément de $U$ de cardinal $\sheaf{\sheaf{c}}$, on a
$\mathcal{P}(x)\in U$ par application de (U.III), d'où
$2^{\sheaf{\sheaf{c}}}<\sheaf{\sheaf{c}}(U)$.

\item Si $(\sheaf{\sheaf{c}}_{\lambda})_{\lambda\in I}$ est une famille de
cardinaux telle que
$\sheaf{\sheaf{c}}_{\lambda}<\sheaf{\sheaf{c}}(U)$ pour tout
$\lambda\in I$ et que $\card(I)<\sheaf{\sheaf{c}}(U)$, alors le
cardinal somme $\Sigma_{\lambda \in I} \sheaf{\sheaf{c}}_{\lambda}$
est $<\sheaf{\sheaf{c}}(U)$. En effet soit $x_{\lambda}\in U$ tel
que $\card(x_{\lambda})=\sheaf{\sheaf{c}}_{\lambda}$; quitte à
remplacer $I$ par un ensemble d'indices équipotent, on peut supposer
qu'on a $I\in U$; alors l'ensemble somme des $x_{\lambda}$ est
élément de $U$ (cor. 2 de la prop. 4 du ${\rm n^{\circ}}$ 1), ce qui
démontre notre assertion.
\end{enumerate}

Posons la définition suivante:

\begin{enonce*}{DÉFINITION 2}
Un\pageoriginale cardinal $\sheaf{\sheaf{d}}$ est dit fortement
inaccessible si:
\begin{itemize}
\item[(FI.1)] Si $\sheaf{\sheaf{c}}$ est un cardinal tel que
$\sheaf{\sheaf{c}}<\sheaf{\sheaf{d}}$, on a
$2^{\sheaf{\sheaf{c}}}<\sheaf{\sheaf{d}}$.

\item[(FI.2)] Si $(\sheaf{\sheaf{c}}_{\lambda})_{\lambda\in I}$ est une
famille de cardinaux telle que
$\sheaf{\sheaf{c}}_{\lambda}<\sheaf{\sheaf{d}}$ pour tout $\lambda
\in I$ et si $\card(I)<d$, alors $\Sigma_{\lambda\in
I}\sheaf{\sheaf{c}}_{\lambda}<\sheaf{\sheaf{d}}$.
\end{itemize}
\end{enonce*}

\begin{enonce*}[remark]{Exemples}
Le cardinal 0 et le cardinal infini dénombrable sont fortement
inaccessibles. Aucun cardinal fini non nul n'est fortement
inaccessible. Ainsi nous venons de démontrer que l'axiome (A.6) des
univers implique la relation:
\begin{itemize}
\item[${\rm (A'.6)}$] \emph{Tout cardinal est strictement majoré par
un cardinal fortement inaccessible.}
\end{itemize}
Inversement:
\end{enonce*}

\begin{enonce*}{THÉORÈME %%
1} La relation $(A'.6)$ implique l'axiome (A.6) des univers.
\end{enonce*}
En effet soit $A$ un ensemble. Il s'agit de construire un univers
$U$ dont $A$ est élément. Définissons par récurrence une suite
$(A_n)_{n\geq 0}$ d'ensembles au moyen de:
\begin{itemize}
\item[(2)] $A_0=A$, $A_{n+1}=$ réunion des éléments de $A_n=$
ensemble des éléments de $A_n$.
\end{itemize}
\noindent Posons $B=\bigcup^\infty_{n=0} A_n$. Soit, par $(A'.6)$,
$\sheaf{\sheaf{c}}$ un cardinal fortement inaccessible tel que
$\card(B)<\sheaf{\sheaf{c}}$.

Il existe un ensemble \emph{bien ordonné} $I$ tel que
$\card(I)=\sheaf{\sheaf{c}}$. Quitte à remplacer $I$ par son plus
petit segment de cardinal $\sheaf{\sheaf{c}}$, on peut supposer que
\emph{tout segment de $I$ distinct de $I$ a un cardinal
$<\sheaf{\sheaf{c}}$.} Nous\pageoriginale noterons $\varepsilon$ le
\emph{plus petit élément de $I$}. Il résulte de l'hypothèse sur les
segments que $I$ n'a pas de plus grand élément (sinon on
l'enlèverait), donc que tout élément de $I$ admet un successeur;
pour $\alpha \in I$, nous noterons $s(\alpha)$ le \emph{successeur}
de $\alpha$

Ceci étant, définissons, par \emph{récurrence transfinie,} une
famille $(B_{\alpha})_{\alpha\in I}$ d'ensembles au moyen de:
\begin{equation}
\begin{cases}
& B_{\mathcal{E}}=B\\
& B_{s(\alpha)}=B_{\alpha}\cup \mathcal{P}(B_{\alpha})\\
& B_{\alpha}={\displaystyle{\bigcup_{\beta<\alpha}}}B_{\beta}\mbox{
si $\alpha$ n'a pas de prédécesseur.}
\end{cases}\tag{3}
\end{equation}
Posons $U=\bigcup_{\alpha\in I}B_{\alpha}$. Nous allons montrer que
$U$ est l'univers cherché.

On a d'abord $A\in U$, car $A=A_{\circ}$ est une partie de
$B=B_{\mathcal{E}}$, donc un élément de
$\mathcal{P}(B_{\mathcal{E}})\subset B_{s(\mathcal{E})}$.

Notons ensuite que, pour tout $\alpha \in I$, on a:
\begin{equation}
\card(B_{\alpha})<\sheaf{\sheaf{c}}{\quoi}.\tag{4}
\end{equation}
Procédons en effet par récurrence transfinie sur $\alpha$. C'est
vrai pour $\alpha=\mathcal{E}$ par hypothèse. De (4) on déduit
$\card(B_{s(\alpha)})\leq
\card(B_{\alpha})+2^{\card(B_{\alpha})}<\sheaf{\sheaf{c}}+\sheaf{\sheaf{c}}$
(par (FI.1)) = $\sheaf{\sheaf{c}}$ (car $\sheaf{\sheaf{c}}$ est
infini). Enfin, si $\alpha$ n'a pas de prédécesseur, l'ensemble $I'$
des $\beta<\alpha$ a un cardinal $<\sheaf{\sheaf{c}}$ par
construction; d'où (si $\card(B_{\beta})<\sheaf{\sheaf{c}}$ pour
tout $\beta<\alpha$), $\card(B_{\alpha})=\card(\bigcup_{\beta\in I'}
B_{\beta})\leq \Sigma_{\beta\in
I'}\card(B_{\beta})<\sheaf{\sheaf{c}}$ (par (FI.2)). Ceci démontre
(4).

Montrons\pageoriginale maintenant qu'on a
\begin{equation}
\card(x)<\sheaf{\sheaf{c}}\mbox{ pour tout } x\in U{\quoi}.\tag{5}
\end{equation}
En effet, il suffit de montrer que, pour tout $\alpha\in I$, on a
« $\card(x)<\sheaf{\sheaf{c}}$ pour tout $x\in B_{\alpha}$ ».
Procédons encore par récurrence transfinie sur $\alpha$. C'est vrai
pour $\alpha=\mathcal{E}$, car, si $x\in B_{\mathcal{E}}=B$, il
existe $n\geq 0$ tel que $x\in A_n$; d'où $X\subset A_{n+1}$ par
(2), $x\subset B$ et $\card(x)\leq \card(B)<\sheaf{\sheaf{c}}$. Le
cas où $\alpha$ n'a pas de prédécesseur est évident. Enfin, si $x\in
B_{\alpha}$ et si $\alpha=s(\beta)$, on a soit $x\in B_{\beta}$ et
l'assertion « $\card(x)<\sheaf{\sheaf{c}}$ » est vraie par
récurrence, soit $x\in\mathcal{P}(B_{\beta})$, d'où $x\subset
B_{\beta}$ et l'assertion « $\card(x)<\sheaf{\sheaf{c}}$ » est
vraie par (4).

Nous sommes maintenant en mesure de démontrer que $U$ est bien un
\emph{univers}:

\begin{itemize}
\item[(U.I)] Soient $x\in U$ et $y\in x$. Il s'agit de montrer %%
qu'on a $y\in U$. Autrement dit il s'agit de montrer que, pour tout
$\alpha\in I$, on a la relation:
$$
\textrm{« }x\in B_{\alpha}\mbox{ et } y\in x\Longrightarrow y\in
B_{\alpha}\textrm{ »\,}{\quoi}.
$$
On procède encore par récurrence transfinie sur $\alpha$. C'est vrai
pour $\alpha =\mathcal{E}$ car $x\in B_{\mathcal{E}}=B$ implique
$x\in A_n$ pour un certain $n$; donc, si $y\in x$, on a $y\in
A_{n+1}$, d'où $y\in B=B_{\mathcal{E}}$. Le passage à un élément
$\alpha$ sans prédécesseur est évident. Passons enfin de $\alpha$ à
$s(\alpha)$: si $x\in B_{s(\alpha)}$ et si $y\in x$, on a, soit
$x\in B_{\alpha}$, d'où $y\in B_{\alpha}\subset B_{s(\alpha)}$ par
récurrence, soit $x\in \mathcal{P}(B_{\alpha})$%%
, d'où encore $y\in B_{\alpha}\subset B_{s(\alpha)}$.

\item[(U.II)] Soient $x$, $y\in U$. Comme la famille
$(B_{\alpha})_{\alpha\in I}$ est croissante par (3), il existe
$\alpha\in I$ tel que $x$, $y\in B_{\alpha}$. Alors
$\{x,y\}\in\mathcal{P}(B_{\alpha})\subset B_{s(\alpha)}\subset U$.

\item[(U.III)] Soit\pageoriginale $x\in U$. Il existe alors $\alpha\in I$ tel que
$x\in B_{\alpha}$. Il suffit donc de montrer que, pour tout
$\alpha\in I$, on a la relation:
$$
\textrm{« }x\in B_{\alpha}\Longrightarrow \mathcal{P}(x)\in
B_{s(s(\alpha))}\textrm{ »\,}{\quoi}.
$$
On procède encore par récurrence transfinie sur $\alpha$. Si
$\alpha=\mathcal{E}$, on a $x\in A_n$ pour un certain $n$, d'où
$y\subset A_{n+1}\subset B=B_{\mathcal{E}}$ pour toute partie $y$ de
$x$; ainsi on a $y\in \mathcal{P}(B_{\mathcal{E}})\subset
B_{s(\mathcal{E})}$ pour tout $y\in \mathcal{P}(x)$, d'où
$\mathcal{P}(x)\subset B_{s(\mathcal{E})}$ et $\mathcal{P}(x)\in
\mathcal{P}(B_{s(\mathcal{E})})\subset B_{s(s(\mathcal{E}))}$, de
sorte que notre assertion est vraie pour $\alpha= \mathcal{E}$. %%
Le passage à un élément $\alpha$ sans prédécesseur est immédiat, car
$\beta<\alpha$ implique alors $s(s(\beta))<\alpha$. Passons enfin de
$\alpha$ à $s(\alpha)$; soit $x\in B_{s(\alpha)}$; si $x\in
B_{\alpha}$, on a $\mathcal{P}(x)\in B_{s(s(\alpha))}\subset
B_{s(s(s(\alpha)))}$ par récurrence; si $x\in
\mathcal{P}(B_{\alpha})$, on a $x\subset B_{\alpha}$, d'où
$\mathcal{P}(x)\subset \mathcal{P}(B_{\alpha})$ et
$\mathcal{P}(x)\in \mathcal{P}(\mathcal{P}(B_{\alpha}))\in
B_{s(s(s(\alpha)))}$.

\item[(U.IV)] Soit $(x_{\lambda})_{\lambda\in K}$ une famille
d'éléments de $U$ telle que $K\in U$. Il s'agit de montrer que la
réunion $x=\bigcup_{\lambda\in K}x_{\lambda}$ est élément de $U$.
Pour tout $\lambda\in K$, choisissons un $\alpha(\lambda)\in I$ tel
que $x_{\lambda}\in B_{\alpha(\lambda)}$. Montrons que l'ensemble
$\alpha(K)$ des $\alpha(\lambda)$ est \emph{majoré} dans $I$: en
effet, s'il ne l'était pas, on aurait:
$$
I=\bigcup_{\lambda%%
\in K}[\mathcal{E}, \alpha(\lambda)]{\quoi},
$$
ce qui contredirait (FI.2) car $\card(K)<\sheaf{\sheaf{c}}$ et
$\card([\mathcal{E},\alpha(\lambda)])<\sheaf{\sheaf{c}}$ pour tout
$\lambda\in K$. Soit donc $\beta$ un majorant de $\alpha(K)$; on a
$x_{\lambda}\in B_{\beta}$ pour tout $\lambda\in K$ car la famille
$(B_{\alpha})_{\alpha \in I}$ est croissante. On a donc
$x_{\lambda}\subset B_{\beta}$ d'après ce qu'on a vu dans la
démonstration de (U.I), d'où $x=\bigcup_{\lambda\in
K}x_{\lambda}\subset B_{\beta}$. Par (3) il s'ensuit qu'on a $x\in
B_{s(\beta)}$, d'où $x\in U$. \cqfd
\end{itemize}

\begin{enonce*}{DÉFINITION 3}
Soient\pageoriginale $x$, $y$ deux ensembles et $n$ un entier $\geq
0$. On dit que $y$ est un composant d'ordre $n$ de $x$ s'il existe
une suite $(x_j)_{j=0,\ldots,n}$ telle que $x_0=x$, $x_n=y$, et
$x_{j+1}\in x_j$ pour $j=0,\ldots,n-1$.
\end{enonce*}

Ainsi $x$ est le seul composant d'ordre 0 de $x$. Les composants
d'ordre 1 (\resp 2) de $x$ sont les éléments de $x$ (\resp les
éléments des éléments de $x$). On dit que $y$ est un
\emph{composant} de $x$ s'il existe $n\geq 0$ tel que $y$ soit
composant d'ordre $n$ de $x$. La relation « $y$ est un composant
de $x$ » est une \emph{relation de préordre.} D'après le schéma
de sélection-réunion (chap. II, \ldots) la relation « $y$ est un
composant de $x$ » est collectivisante par rapport à $y$, de
sorte que les composants de $x$ forment un \emph{ensemble}.

\begin{enonce*}{DÉFINITION 4}
Soit $\sheaf{\sheaf{c}}$ un cardinal. On dit qu'un ensemble $x$ est
de type $\sheaf{\sheaf{c}}$ (\resp de type strict
$\sheaf{\sheaf{c}}$, de type fini) si tous les composants de $x$ ont
des cardinaux $\leq \sheaf{\sheaf{c}}$ (\resp $<\sheaf{\sheaf{c}}$,
finis).
\end{enonce*}

\begin{enonce*}[remark]{Exemples}
Les éléments de l'univers $U_1$ (${\rm n^{\circ}}$ 1, ex. 2) sont
tous de type fini. Si $\sheaf{\sheaf{c}}$ est un cardinal, et si $x$
est de type $\sheaf{\sheaf{c}}$ (\resp type strict
$\sheaf{\sheaf{c}}$, type fini), alors tout composant de $x$ et
toute partie de $x$ sont de type $\sheaf{\sheaf{c}}$ (\resp type
strict $\sheaf{\sheaf{c}}$, type fini); de même $\mathcal{P}(x)$ est
de type $2^{\sheaf{\sheaf{c}}}$ (\resp de type strict
$2^{\sheaf{\sheaf{c}}}$, de type fini). Si $x$ est de type
$\sheaf{\sheaf{c}}$ et si $\sheaf{\sheaf{c}}\leq
\sheaf{\sheaf{c}}'$, alors $x$ est de type $\sheaf{\sheaf{c}}'$.
\end{enonce*}

\begin{enonce*}{LEMME}
Soient\pageoriginale $\sheaf{\sheaf{c}}$ un cardinal fortement
inaccessible non dénombrable, et $X$ un ensemble de type strict
$\sheaf{\sheaf{c}}$. Alors il existe un cardinal
$\sheaf{\sheaf{d}}<\sheaf{\sheaf{c}}$ tel que $X$ soit de type
$\sheaf{\sheaf{d}}$.
\end{enonce*}

En effet, pour tout entier $n$, notons $\sheaf{\sheaf{d}}_n$ le
cardinal de l'ensemble $X_n$ des composants d'ordre $n$ de $X$. On a
$\sheaf{\sheaf{d}}_0=1$,
$\sheaf{\sheaf{d}}_1=\card(X)<\sheaf{\sheaf{c}}$, et, comme
$X_n=\bigcup_{y\in X_{n-1}} y$, on en déduit, par récurrence sur $n$
et usage de (FI.2), qu'on a $\sheaf{\sheaf{d}}_n<\sheaf{\sheaf{c}}$
pour tout $n$. Alors les $\sheaf{\sheaf{d}}_n$ sont majorés par le
cardinal $\sheaf{\sheaf{d}}=\Sigma_{j\geq 0}\sheaf{\sheaf{d}}_j$,
qui est $<\sheaf{\sheaf{c}}$ par (FI.2) et l'hypothèse de non
dénombrabilité de $\sheaf{\sheaf{c}}$. Alors, si $Y$ est un
composant d'ordre $n$ de $X$, on a $\card(Y)\leq \card(X_{n+1})\leq
\sheaf{\sheaf{d}}$. \cqfd

\begin{enonce*}{PROPOSITION 10}
Soient $U$ un univers et $\sheaf{\sheaf{c}}$ un cardinal fortement
inaccessible. Alors l'ensemble $U'$ des $x\in U$ qui sont de type
strict $\sheaf{\sheaf{c}}$ est un univers.
\end{enonce*}
En effet, si $x$, $y\in U'$ et si $z\in x$, on a évidemment $z\in
U'$ et $\{x,y\}\in U'$. Si $x\in U'$, on a
$\card(x)<\sheaf{\sheaf{c}}$, d'où
$\card(\mathcal{P}(x))<\sheaf{\sheaf{c}}$ par (FI.1); comme un
composant de $\mathcal{P}(x)$ est, ou bien $\mathcal{P}(x)$, ou bien
une partie de $x$, ou bien un composant de $x$, on a
$\mathcal{P}(x)\in U'$. Enfin si $(x_{\alpha})_{\alpha\in I}$ est
une famille d'éléments de $U'$ telle que $I\in U'$, on a
$\card(\bigcup_{\alpha}x_{\alpha})<\sheaf{\sheaf{c}}$ par (FI.2);
comme tout composant d'ordre $>0$ de $\bigcup_{\alpha}x_{\alpha}$
est un composant d'un $x_{\alpha}$, on a bien
$\bigcup_{\alpha}x_{\alpha}\in U'$ \cqfd

\medskip
\noindent \emph{Remarque sur le cardinal} $\sheaf{\sheaf{c}}(U)$.
Soient $U$ un univers et $\sheaf{\sheaf{c}}(U)$ le cardinal défini
par (1). \ie
$$
\sheaf{\sheaf{c}}(U)={\sup}_{x\in U}\card(x){\quoi}.
$$
On\pageoriginale a évidemment $\sheaf{\sheaf{c}}(U)\leq \card(U)$,
car, rappelons-le, $x\in U$ implique $x\subset U$. Mais l'égalité
$\sheaf{\sheaf{c}}(U)=\card(U)$ n'est pas toujours vraie. Par
exemple soit $(x_{\alpha})_{\alpha\in I}$ une famille \emph{non
dénombrable} de symboles avec les axiomes:
$$
\text{«} x_{\alpha}=\{x_{\alpha}\}\mbox{ pour tout } \alpha\in I
\text{»} \quav \text{«} x_{\alpha}\neq x_{\beta}\mbox{ si }
\alpha\neq \beta \text{»}
$$
(autrement dit $x_{\alpha}$ est un ensemble à un seul élément, à
savoir lui-même). Formons, comme dans l'ex. 2 du ${\rm n^{\circ}}$
1, les  « mots significatifs » formés avec les symboles
$∅$, $x_{\alpha}$, $\{,\}$ %%
etc Chacun désigne un ensemble \emph{fini}. Soit $U$ l'ensemble de
ceux-ci. On vérifie aisément que les conditions de la déf. 1 sont
satisfaites, de sorte que $U$ est un \emph{univers}. Comme les mots
significatifs sont des suites finies d'éléments d'un ensemble de
cardinal $\card(I)+4=\card(I)$; on a $\card(U)=\card(I)$ (chap. III;
en fait $\card(U)\geq \card(I)$ suffit, et c'est évident). D'autre
part, $\sheaf{\sheaf{c}}(U)$ est le cardinal dénombrable, d'où
$\sheaf{\sheaf{c}}(U)<\card(U)$.

L'inégalité stricte $\sheaf{\sheaf{c}}(U)<\card(U)$ tient à ce qu'on
a introduit ici des ensembles $x_{\alpha}$ tels que $x_{\alpha}\in
x_{\alpha}$. Si on interdit des horreurs de ce genre, on obtient
$\card(U)=\sheaf{\sheaf{c}}(U)$ pour tout univers $U$, ainsi que
d'autres fort jolis résultats « ne pouvant servir à rien ».
C'est ce qu'on va faire au numéro suivant.

\section{Ensembles et univers artiniens}

\begin{enonce*}{DÉFINITION 5}
On dit qu'un ensemble $x$ est artinien s'il n'existe aucune suite
infinie $(x_n)_{n\geq 0}$ telle que $x_{\circ}=x$ et que $x_{n+1}\in
x_n$ pour tout $n\geq 0$.
\end{enonce*}

\begin{enonce*}[remark]{Exemples}
Les ensembles $∅$, $\{∅\}$,
$\{∅,\{∅\}\}$, plus généralement les éléments de
l'univers $U_1$ du ${\rm n^{\circ}}$ 1, sont artiniens.
\end{enonce*}

En\pageoriginale termes imagés, si $x$ est artinien, et si on prend
un élément $x_1$ de $x$, puis un élément $x_2$ de $x_1$, \etc, le
processus doit s'arrêter, et on arrive à un composant $x_n$ de $x$
qui est \emph{vide}. Autrement dit les ensembles artiniens sont
« construits à partir de $∅$ »; ceci sera précisé plus
tard (\cf (*) dans la démonstration du th. 2).

Les ensembles artiniens jouissent évidemment des propriétés
suivantes:
\begin{itemize}
\item[(AR.I)] Si $x$ est artinien, toute partie de $x$ et tout
composant de $x$ sont artiniens. Pour que $x$ soit artinien, il faut
et il suffit que tout élément de $x$ soit artinien.

\item[(AR.II)] Si $x$ et $y$ sont artiniens, alors $\{x,y\}$ est
artinien.

\item[(AR.III)] Si $x$ est artinien, $\mathcal{P}(x)$ est artinien.

\item[(AR.IV)] Toute réunion d'ensembles artiniens est un ensemble
artinien.
\end{itemize}

Ces propriétés montrent aussitôt qu'on a la

\begin{enonce*}{PROPOSITION 11}
Si $U$ est un univers, l'ensemble des éléments artiniens de $U$ est
un univers, nécessairement artinien.
\end{enonce*}

\begin{enonce*}{COROLLAIRE}
Si $x$ est un ensemble artinien, il existe un univers artinien $U$
tel que $x\in U$.
\end{enonce*}

En effet, par l'axiome (A.6), $x$ est élément d'un univers $V$; on
prend pour $U$ l'ensemble des éléments artiniens de $V$.

La proposition suivante est encore moins utile que le reste du ${\rm
n^{\circ}}$:

\begin{enonce*}{PROPOSITION 12}
Soit\pageoriginale $A$ un ensemble artinien. Alors:
\begin{enumerate}
\renewcommand{\theenumi}{\alph{enumi}}
\renewcommand{\labelenumi}{\theenumi)}
\item Pour tout $x\in A$, on a $x\not\in x$;

\item Si $x$, $y\in A$, on ne peut avoir à la fois $x\in y$ et $y\in
x$;

\item Pour tout $x\in A$, la relation « $y$ est un composant de $z$ » entre composants $y$, $z$ de $x$ est une relation d'ordre;

\item Pour tout élément non vide $x$ de $A$, il existe $y\in x$
tel que $x\cap y=∅$.
\end{enumerate}
\end{enonce*}
En effet la négation de a) (\resp de b)) entraîne l'existence d'une
suite infinie $(x_n)_{n\geq 0}$ contredisant la déf. 5, à savoir
$(x,x,x,\ldots)$ (\resp $(x,y,x,y,x,y,\ldots)$). La négation de c)
veut dire qu'il existe $x\in A$, des composants $y$, $z$ de $x$
distincts, et des suites d'appartenances:
$$
y\in y_1\in\ldots \in y_q\in z\quav z\in z_1\in\ldots\in z_r\in
y{\quoi};
$$
d'où, comme dans b), une suite infinie $(x_n)_{n\geq 0}$
contredisant la déf. 5. Enfin, si d) est fausse, il existe un
élément non vide $x$ de $A$ tel que $x\cap y\neq ∅$ pour
tout $y\in x$; on pose $x_\circ=x$, et on prend pour $x_1$ un
élément de $x$; comme $x\cap x_1\neq ∅$, on prend pour
$x_2$ un élément de $x\cap x_1$ etcétera; plus formellement on
définit par récurrence une suite infinie $(x_n)_{n\geq 1}$
d'éléments de $x$ au moyen de $x_1\in x$, $x_{n+1}\in x_n\cap x$
pour $n\geq 1$; alors la suite $(x_n)_{n\geq 0}$ contredit la déf.
5.

\begin{enonce*}[remark]{Remarque}
Soit $B$ un ensemble. Pour que $B$ soit artinien, il faut et il
suffit que tout ensemble $A$ d'ensembles de composants de $B$
satisfasse à la condition d) de la prop. 12. En effet la nécessité
résulte de (AR.I) et\pageoriginale de la prop. 12. Réciproquement,
si $B$ n'est pas artinien, il existe une suite infinie $(x_n)_{n\geq
0}$ avec $x_{\circ}=B$ et $x_{n+1}\in x_n$ pour tout $n\geq 0$; on
prend alors pour $A$ la partie réduite à l'ensemble $X$ des $x_n$;
ainsi $A$ contient un élément non vide $X$ tel que, pour tout
élément $y$ de $X$, on ait $y\cap X\neq ∅$ (en effet $y$
est de la forme $x_n$, et on a $x_{n+1}\in y\cap X$).
\end{enonce*}

\begin{enonce*}{THÉORÈME 2}
Soit $\sheaf{\sheaf{c}}$ un cardinal infini. Alors:
\begin{itemize}
\item[a)] La relation « $x$ est un ensemble artinien de type $\sheaf{\sheaf{c}}$
(\resp de type strict $\sheaf{\sheaf{c}}$) » %%
est collectivisante par rapport à $x$; l'ensemble
$A_{\sheaf{\sheaf{c}}}$ des ensembles artiniens de type
$\sheaf{\sheaf{c}}$ a pour cardinal $2^{\sheaf{\sheaf{c}}}$.

\item[b)] Si $\sheaf{\sheaf{c}}$ est fortement inaccessible, l'ensemble
$U_{\sheaf{\sheaf{c}}}$ des ensembles artiniens de type strict
$\sheaf{\sheaf{c}}$ est un univers de cardinal $\sheaf{\sheaf{c}}$;
le cardinal $\sheaf{\sheaf{c}}(U_{\sheaf{\sheaf{c}}})=\sup_{x\in
U_{\sheaf{\sheaf{c}}}}\card(x)$ est $\sheaf{\sheaf{c}}$.

\item[c)] Si un univers $U$ admet un élément de cardinal $\sheaf{\sheaf{c}}$,
tout ensemble artinien de type $\sheaf{\sheaf{c}}$ appartient à $U$
(autrement dit $A_{\sheaf{\sheaf{c}}}\subset U$).

\item[${\rm c')}$] Si un univers $U$ est non vide, tout ensemble
artinien de type fini est élément de $U$.
\end{itemize}
\end{enonce*}
Avant de démontrer le th. 2, déduisons en quelques corollaires
illuminants:

\begin{enonce*}{COROLLAIRE 1}
Si un univers $U$ est artinien, alors $\card(U)$ est fortement
inaccessible, et $U$ est l'ensemble des ensembles artiniens de type
strict $\card(U)$.
\end{enonce*}

En\pageoriginale effet posons $\sheaf{\sheaf{c}}(U)=\sup_{x\in
U}\card(x)$. C'est un cardinal fortement inaccessible (début du
${\rm n^{\circ}}$ 5). Supposons le d'abord non dénombrable; alors,
pour tout cardinal infini $\sheaf{\sheaf{c}}<\sheaf{\sheaf{c}}(U)$,
tout ensemble artinien de type $\sheaf{\sheaf{c}}$ est élément de
$U$ par c); donc tout ensemble artinien de type strict
$\sheaf{\sheaf{c}}(U)$ est élément de $U$ par le lemme du ${\rm
n^{\circ}}$ 5. Cette dernière assertion reste valable si
$\sheaf{\sheaf{c}}(U)$ est dénombrable par ${\rm c'}$). Inversement,
d'après (U.I), tout ensemble de $U$ est de type strict
$\sheaf{\sheaf{c}}(U)$. Donc $U$ est l'univers
$U_{\sheaf{\sheaf{c}}(U)}$ de b), d'où
$\sheaf{\sheaf{c}}(U)=\card(U)$ par b).

Il résulte du cor. l qu'on univers artinien est \emph{déterminé} de
façon unique par son cardinal (d'ailleurs fortement inaccessible).
On a donc une « correspondance biunivoque » entre univers
artiniens et cardinaux fortement inaccessibles. En particulier:

\begin{enonce*}{COROLLAIRE 2}
La relation d'inclusion $U\subset U'$ entre univers artiniens est
une relation de bon ordre.
\end{enonce*}

En effet la relation $\sheaf{\sheaf{c}}\leq \sheaf{\sheaf{c}}'$
entre cardinaux est une relation de bon ordre (chap. III) et, avec
les notations du b) du th. 2, les relations $\sheaf{\sheaf{c}}\leq
\sheaf{\sheaf{c}}'$ et $U_{\sheaf{\sheaf{c}}}\subset
U_{\sheaf{\sheaf{c}}'}$ sont équivalentes.

Notons que le th. 2, b) donne une seconde démonstration du th. 1
(${\rm n^{\circ}}$ 5).

Passons à la démonstration du th. 2. Étant donné un ensemble $A$,
nous appellerons \emph{chaîne} de $A$ toute suite finie
$(x_j)_{j=0,\ldots,n}$ telle que $x_0=A$ et que $x_{j+1}\in x_j$
pour $j=0,\ldots,n-1$. Les chaînes de $A$ forment un\pageoriginale
\emph{ensemble}, d'après le schéma de sélection-réunion; notons le
$G(A)$. Étant donnée une chaîne $X=(x_j)_{j=0,\ldots,n}$, les
chaînes de la forme $(x_i)_{i=0,\ldots,q}$ avec $q\leq n$ seront
dites plus petites que $X$; on obtient ainsi, sur $G(A)$, une
structure d'\emph{ensemble ordonné}. Pour que $A$ soit artinien, il
faut et il suffit que $G(A)$ soit un ensemble ordonné
« \emph{noethérien} » (c.à.d. satisfaisant aux conditions
équivalentes du chap. III, \S 6, ${\rm n^{\circ}}$ 5).

Nous allons montrer que:

\noindent (*) \emph{Si $A$ est artinien, il est déterminé de façon
unique par la classe d'isomorphisme de l'ensemble ordonné $G(A)$.}

En effet, étant donnés un ensemble ordonné $G$ et un élément $g\in
G$, nous noterons $S(g)$ l'ensemble des $g'\in G$ tels que $g<g'$ et
que $g\leq h\leq g'$ implique $h=g$ ou $h=g'$ (autrement dit
l'ensemble des « successeurs immédiats » de $g$). Considérons
l'application $\theta$ qui, à toute chaîne $X=(x_0,\ldots,x_n)$ de
$G(A)$ fait correspondre l'ensemble $x_n$; on a alors, pour tout
$X\in G(A)$:
$$
\theta(X)=\{\theta(X')| X'\in S(X)\}{\quoi}.
$$
Comme $A$ est l'image par $\theta$ du plus petit élément $X_0=(A)$
de $G(A)$, il va nous suffire de montrer que $\theta$ est uniquement
déterminée par la classe d'isomorphisme de l'ensemble ordonné
$G(A)$. Or ceci résulte du lemme suivant:

\begin{enonce*}{LEMME 1}
Soient\pageoriginale $G$ un ensemble ordonné noethérien et $\varphi$
une application de $G$ telle que, pour tout $g\in G$, on ait:
\begin{equation}
\varphi(g)=\{\varphi(g')| g'\in S(g)\}{\quoi}.\tag{1}
\end{equation}
Alors $\varphi$ est déterminée de façon unique. De plus, si $U$ est
un univers contenant un élément équipotent à $G$, $\varphi$ prend
ses valeurs dans $U$.
\end{enonce*}

L'hypothèse implique qu'on a $\varphi(g)=∅$ si $g$ est un
élément maximal de $G$.

Soient, en effet, $\varphi$ et $\varphi'$ deux applications telles
que (1) soit vraie; si $\varphi\neq \varphi'$, l'ensemble des $g\in
G$ tels que $\varphi(g)\neq \varphi'(g)$ est non-vide, donc admet un
élément \emph{maximal} $h$ car $G$ est noethérien; on a alors
$\varphi(g)=\varphi'(g)$ pour tout $g>h$, en particulier pour tout
« successeur » $g\in S(h)$; d'où $\varphi(h)=\varphi'(h)$ par
(1), ce qui est une contradiction; on a donc bien
$\varphi=\varphi'$. Soit maintenant $U$ un univers contenant un
élément $x$ équipotent à $G$; montrons que $\varphi$ prend ses
valeurs dans $U$; sinon soit $h$ un élément maximal parmi les $g\in
G$ tels que $\varphi(g)\not\in U$; on a $\varphi(g')\in U$ pour tout
$g'\in S(g)$ de sorte que
$$
\varphi(h)=\{\varphi(g')| g'\in S(h)\}
$$
est une partie de $U$; or, comme son cardinal est inférieur à
$\card(G)$ donc au cardinal d'un élément de $U$, on a $\varphi(h)\in
U$ (${\rm n^{\circ}}$ 1, prop. 7); cette contradiction montre que
$\varphi$ prend ses valeurs dans $U$.

Ceci \pageoriginale %%
étant, démontrons le a) du th. 2. On peut se borner à l'assertion
non-respée, car l'autre en découle aussitôt. Soit
$\sheaf{\sheaf{c}}$ un cardinal \emph{infini}. Si $A$ est un
ensemble de type $\sheaf{\sheaf{c}}$ on a:
\begin{equation}
\card(G(A))\leq \sheaf{\sheaf{c}}{\quoi}%%
\ .\tag{2}
\end{equation}
En effet, si on note $A_n$ l'ensemble des composants d'ordre $n$ de
$A$, on a $\card(A_1)\leq \sheaf{\sheaf{c}}$ et $\card(A_{n+1})\leq
\sheaf{\sheaf{c}}\cdot \card(A_n)$, d'où $\card(A_n)\leq
\sheaf{\sheaf{c}}^n$ par récurrence; or
$\sheaf{\sheaf{c}}^n=\sheaf{\sheaf{c}}$ car $\sheaf{\sheaf{c}}$ est
infini (chap. III); d'où $\card(\bigcup^\infty_{n=0} A_n)\leq
\card(\NN)\cdot \sheaf{\sheaf{c}}=\sheaf{\sheaf{c}}$ (chap. III); or
$G(A)$ est un ensemble de suites \emph{finies} d'éléments de
$\bigcup_n A_n$, de sorte qu'on a bien l'inégalité (2) (chap. III).
Ceci étant, si $E$ est un ensemble de cardinal $\sheaf{\sheaf{c}}$,
la donnée d'une structure d'ordre sur une partie $E'$ de $E$
équivaut à la donnée de la partie de $E\times E$ formée des $(x,y)$
tels que $x\in E'$, $y\in E'$ et $x\leq y$. Donc, en vertu de (2),
les classes d'isomorphisme des ensembles ordonnés $G(A)$ (où $A$ est
de type $\sheaf{\sheaf{c}}$) forment un ensemble
$\mathfrak{F}_{\sheaf{\sheaf{c}}}$, et on a
$\card(\mathfrak{F}_{\sheaf{\sheaf{c}}})\leq
\card(\mathcal{P}(E\times E))=2^{\sheaf{\sheaf{c}}}$. Soit
$\mathfrak{F}'_{\sheaf{\sheaf{c}}}$ la partie de
$\mathfrak{F}_{\sheaf{\sheaf{c}}}$ formée des classes d'ensembles
ordonnés noethériens ayant un plus petit élément; si, à tout $G\in
\mathfrak{F}'_{\sheaf{\sheaf{c}}}$, on fait correspondre la valeur
de la fonction $\varphi$ du lemme 2 au plus petit élément de $G$, on
obtient une application $\theta$ de
$\mathfrak{F}'_{\sheaf{\sheaf{c}}}$ dont l'image contient
\emph{tous} les ensembles artiniens de type $\sheaf{\sheaf{c}}$. Ces
derniers forment donc bien un ensemble $A_{\sheaf{\sheaf{c}}}$, et
on a $\card(A_{\sheaf{\sheaf{c}}})\leq
\card(\mathfrak{F}'_{\sheaf{\sheaf{c}}})\leq
\card(\mathfrak{F}_{\sheaf{\sheaf{c}}})\leq 2^{\sheaf{\sheaf{c}}}$.

Reste %%
\pageoriginale à voir qu'on a
$\card(A_{\sheaf{\sheaf{c}}})=2^{\sheaf{\sheaf{c}}}$. Pour cela il
suffit de voir qu'il existe un ensemble artinien $B$ de type
$\sheaf{\sheaf{c}}$ et de cardinal $\sheaf{\sheaf{c}}$, car les
parties de $B$ seront alors des éléments de $A_{\sheaf{\sheaf{c}}}$.
Cette existence résulte du lemme suivant:

\begin{enonce*}{LEMME 2}
Pour tout cardinal $\sheaf{\sheaf{c}}$, il existe un ensemble
artinien $B$ de type $\sheaf{\sheaf{c}}$ et de cardinal
$\sheaf{\sheaf{c}}$.
\end{enonce*}

On procède par induction transfinie sur $\sheaf{\sheaf{c}}$. Pour
$\sheaf{\sheaf{c}}=0$ on n'a pas le choix, et on prend
$B=∅$. Si $\sheaf{\sheaf{c}}$ à un prédécesseur
$\sheaf{\sheaf{c}}'$, soit $B'$ un ensemble artinien de type
$\sheaf{\sheaf{c}}'$ et de cardinal $\sheaf{\sheaf{c}}'$; on a alors
$\sheaf{\sheaf{c}}\leq 2^{\sheaf{\sheaf{c}}'}$, de sorte qu'il
existe une partie $B$ de $\mathcal{P}(B')$ de cardinal
$\sheaf{\sheaf{c}}$; les éléments de $B$ sont des parties de $B'$ et
ont donc des cardinaux $\leq \sheaf{\sheaf{c}}' \leq
\sheaf{\sheaf{c}}$; les composants d'ordre supérieur de $B$ sont des
composants de $B'$, et ont donc aussi des cardinaux $\leq
\sheaf{\sheaf{c}}' \leq \sheaf{\sheaf{c}}$. Enfin, si
$\sheaf{\sheaf{c}}$ n'a pas de prédécesseur, on choisit, pour tout
cardinal $\sheaf{\sheaf{c}}_{\lambda}<\sheaf{\sheaf{c}}$, un
ensemble artinien $B_{\lambda}$ de type
$\sheaf{\sheaf{c}}_{\lambda}$ et de cardinal
$\sheaf{\sheaf{c}}_{\lambda}$; alors $B=\bigcup _{\lambda}
B_{\lambda}$ répond à la question. Ceci démontre le lemme 2, et
achève la démonstration de la partie a).

Passons à b). Soit $\sheaf{\sheaf{c}}$ un cardinal fortement
inaccessible. On sait déjà, par a), que les ensembles artiniens de
type strict $\sheaf{\sheaf{c}}$ forment un ensemble
$U_{\sheaf{\sheaf{c}}}$. Le fait que $U_{\sheaf{\sheaf{c}}}$ est un
univers résulte aussitôt des propriétés (AR.I) à (AR.IV) des
ensembles artiniens (début du ${\rm n^{\circ}}$), et\pageoriginale
de majorations de cardinaux analogues à celles de la prop. 10 du
${\rm n^{\circ}}$~5. La relation $\sup_{x\in U_{\sheaf{\sheaf{c}}}}
\card(x)=\sheaf{\sheaf{c}}$ résulte%%
 du lemme 2, appliqué aux cardinaux $<\sheaf{\sheaf{c}}$. Enfin,
 pour montrer
que $\card(U_{\sheaf{\sheaf{c}}})=\sheaf{\sheaf{c}}$, supposons
d'abord $\sheaf{\sheaf{c}}$ \emph{non dénombrable}; d'après%%
 le lemme du ${\rm n^{\circ}}$~5,
$U_{\sheaf{\sheaf{c}}}$ est la réunion
$\bigcup_{\sheaf{\sheaf{d}}<\sheaf{\sheaf{c}}}
A_{\sheaf{\sheaf{d}}}$, où $A_{\sheaf{\sheaf{d}}}$ désigne
l'ensemble des ensembles artiniens de type $\sheaf{\sheaf{d}}$; or
on a
$\card(A_{\sheaf{\sheaf{d}}})=2^{\sheaf{\sheaf{d}}}<\sheaf{\sheaf{c}}$
(par a)); d'autre part l'ensemble des cardinaux
$\sheaf{\sheaf{d}}<\sheaf{\sheaf{c}}$ a un cardinal $\leq
\sheaf{\sheaf{c}}$ (chap. III); d'où
$\card(U_{\sheaf{\sheaf{c}}})\leq \sheaf{\sheaf{c}} \cdot
\sheaf{\sheaf{c}}=\sheaf{\sheaf{c}}$, et aussi
$\card(U_{\sheaf{\sheaf{c}}})\geq \sheaf{\sheaf{c}}$ car
$\card(U_{\sheaf{\sheaf{c}}})\geq
\card(A_{\sheaf{\sheaf{d}}})=2^{\sheaf{\sheaf{d}}}$ pour tout
$\sheaf{\sheaf{d}}<\sheaf{\sheaf{c}}$.


Le cas $\sheaf{\sheaf{c}}=0$ étant trivial, reste le cas où
$\sheaf{\sheaf{c}}$ est le \emph{cardinal infini dénombrable.} Dans
ce cas $U_{\sheaf{\sheaf{c}}}$ est l'ensemble des ensembles
artiniens de type fini (\ie finis ainsi que tous leurs composants),
et on utilise un joli résultat de nature combinatoire.

\begin{enonce*}{LEMME 3}
{\rm (D. K\"onig ?)}. Considérons deux suites infinies $(E_n)_{n\geq
1}$, $(f_n)_{n \geq 1}$ d'ensembles finis $E_n$ et d'applications
$f_n:E_{n+1}\rightarrow E_n$. S'il n'existe aucune suite infinie
$(x_n)_{n%%
\geq 1}$ telle que $x_n\in E_n$ et que $f_n(x_{n+1})=x_n$ pour tout
$n$, alors $E_n$ est vide pour $n$ assez grand.
\end{enonce*}

Autrement dit, si toutes les suites $(x_n)$ telles que
$f_n(x_{n+1})=x_n$ sont \emph{finies}, leurs longueurs sont
\emph{bornées}. \c Ca peut s'exprimer en termes de limites
projectives: une limite projective d'ensembles finis non vides est
non vide (\cf Top. Géné., Chap. I, 2e éd. \S 9, ${\rm n^{\circ}}$ 6,
prop. 8, $2^{\circ}$)).

Raisonnons,\pageoriginale en effet, par l'absurde. S'il existe des
$E_n$ non vides pour $n$ arbitrairement grand, aucun $E_n$ n'est
vide (car $E_n=∅$ entraîne $E_{n+1}=∅$ vu
l'existence de $f_n:E_{n+1}\rightarrow E_n$). Appelons
« cohérentes » les suites finies $(x_j)_{1\leq j\leq n}$ telles que
$f_j(x_{j+1})=x_j$ pour $j=1,\ldots,n-1$. Démontrons, par récurrence
sur $n$, l'existence d'une suite cohérente $(a_1,\ldots,a_n)$
$(a_i\in E_i)$ qui, pour tout $q\geq n$, peut être prolongée en une
suite cohérente $(a_1,\ldots,a_n,x_{n+1},\ldots,x_q)$ de longueur
$q$. C'est évident pour $n=0$, car aucun $E_q$ n'est vide. Passons
de $n$ à $n+1$. Si, pour tout $x\in f^{-1}_n(\{a_n\})\subset
E_{n+1}$, toutes des suites cohérentes prolongeant
$(a_1,\ldots,a_n,x)$ avaient des longueurs bornées par un entier
$q(x)$, alors toutes les suites cohérentes prolongeant
$(a_1,\ldots,a_n)$ seraient de longueurs bornées (par $\sup_x q(x)$)
car $E_{n+1}$ est fini; il existe donc $a_{n+1}\in
f^{-1}_n(\{a_n\})$ tel que la suite cohérente
$(a_1,\ldots,a_n,a_{n+1})$ admette des prolongements de longueur
arbitraire. Ceci étant on a une suite infinie $(a_n)_{n%%
\geq 1}$ qui contredit l'hypothèse.

Il résulte du lemme 3 que si $A$ est un ensemble artinien de type
fini, alors l'ensemble ordonné $G(A)$ de ses chaînes est
\emph{fini}: on prend, en effet, pour $E_n$ l'ensemble des chaînes
$(x_n\in x_{n-1}\in\ldots \in x_1\in A)$ à $n+1$ termes (qui est
fini car les composants de $A$ d'ordre $\leq n+1$ sont en nombre
fini et sont tous finis), et pour $f_n$ l'application $(x_{n+1}\in
x_n\in x_{n+1}\ldots)\mapsto (x_n\in x_{n-1}\in\ldots)$. Or
l'ensemble des classes d'isomorphisme\pageoriginale d'ensembles
ordonnés finis est \emph{dénombrable}: en effet, la donnée d'une
structure d'ordre sur une partie finie de $\NN$ équivaut à celle de
son graphe, qui est une partie finie de $\NN\times\NN$; d'autre part
l'ensemble des parties finies d'un ensemble dénombrable est
dénombrable (chap. III). Il résulte de (*) et du lemme 1, que
l'ensemble $\mathfrak{a}$ des ensembles artiniens de type fini est
dénombrable. Il est infini par le lemme 2 (ou, plus simplement, par
usage de la suite $(z_n)_{n\geq 0} $ définie au moyen de
$z_0=∅$, $z_{n+1}=\{z_n\}$). Ceci termine la démonstration
de b).

\begin{enonce*}[remark]{Remarque}
On vient de démontrer que, si $A$ est un ensemble artinien de type
fini, il n'a qu'un nombre fini de composants. Il existe donc un
entier $n$ tel que $A$ soit de type $n$.
\end{enonce*}
Passons à la démonstration de c). Soient $\sheaf{\sheaf{c}}$ un
cardinal infini, $U$ un univers admettant un élément de cardinal
$\sheaf{\sheaf{c}}$, et $A$ un ensemble artinien de type
$\sheaf{\sheaf{c}}$. Alors l'ensemble ordonné $G(A)$ a un cardinal
$\leq \sheaf{\sheaf{c}}$ (formule (2) ci-dessus). La seconde
assertion du lemme 1, appliquée à $G(A)$, montre que l'application
$(x_0,\ldots, x_n)\rightsquigarrow x_n$ de $G(A)$ prend ses valeurs
dans $U$. Autrement dit tous les composants de $A$ sont éléments de
$U$. Ceci démontre c). La démonstration de ${\rm c'}$) est analogue:
si $A$ est un ensemble artinien de type fini, on vient de voir que
$G(A)$ est fini; comme $U$ est non-vide, il contient un élément
équipotent à $G(A)$ par le cor. à la prop. 7 (${\rm n^{\circ}}$ 1).
\cqfd

\section{Remarques métamathématiques vaseuses}\pageoriginale

\begin{enumerate}
\renewcommand{\theenumi}{\alph{enumi}}
\renewcommand{\labelenumi}{\theenumi)}
\item L'axiome « \emph{tout ensemble est artinien} » est
\emph{inoffensif}. En effet, si on a un modèle $M$ de la Théorie des
Ensembles, l'ensemble $M'$ des éléments artiniens de $M$ est aussi
un modèle (\cf prop. 11).

\item L'axiome (A.6) des univers (et l'axiome équivalent $(A'.6)$
des cardinaux fortement inaccessibles) est \emph{indépendant} du
reste de la Théorie des Ensembles. En effet soit $\sheaf{\sheaf{c}}$
le premier cardinal fortement inaccessible non dénombrable.
L'univers $U_{\sheaf{\sheaf{c}}}$ des ensembles artiniens de type
strict $\sheaf{\sheaf{c}}$ (th. 2, b)) est un modèle de la Théorie
des Ensembles sans (A.6): on appelle « ensembles » les éléments
de $U_{\sheaf{\sheaf{c}}}$, la « relation d'appartenance » est
la restriction à $U_{\sheaf{\sheaf{c}}}$ de l'ordinaire, \etc Les
« univers » du modèle sont donc les univers ordinaires qui sont
éléments de $U_{\sheaf{\sheaf{c}}}$. Or on a vu que les seuls
univers qui sont éléments de $U_{\sheaf{\sheaf{c}}}$ sont les deux
peque${\rm\widetilde{n}}$os $U_0=∅$ et $U_1$. Ainsi $U_1$
est un « ensemble » qui n'est élément d'aucun « univers ».
On a donc un modèle de la Théorie des Ensembles où (A.6) est faux.

\item Bourbaki a été trop prudent en se contentant de
« présumer » que l'\emph{axiome de l'infini} (A.5) est
\emph{indépendant} des axiomes et schémas précédents. Il l'est
effectivement, car l'univers dénombrable $U_1$ des ensembles
artiniens de type fini est un modèle où (A.5) est faux, et où les
axiomes et schémas précédents sont vrais.

\item Il serait très intéressant de démontrer que l'axiome (A.6)
des univers est inoffensif. \c Ca paraît difficile et c'est même
indémontrable, dit Paul Cohen.
\end{enumerate}

L'adjectif\pageoriginale « vaseuses » dans le titre veut dire
qu'on ne s'est pas donné la peine, en construisant des modèles, de
canuler le symbole $\tau$ de sorte qu'il n'en fasse pas sortir. La
clef de ça, si on considère un modèle $M$, est de remplacer le
$\tau_x(ℛ(x))$ ordinaire par:
$$
\tau_x(ℛ(x)\mbox{ et } x\in M){\quoi}.
$$
Encore faut-il vérifier que ça transforme bien les quantificateurs
ordinaires en les quantificateurs autrefois dits « typiques »:
$$
(\forall x \in M)\quav (\exists x \in M){\quoi}.
$$

\begin{enonce*}[remark]{Exercices}
\begin{enumerate}
\renewcommand{\labelenumi}{\theenumi)}
\item Soit $n$ un entier $\geq 1$. Montrer que l'ensemble des ensembles
artiniens de type $n$ est infini (les ensembles $z_0=∅$,
$z_1=\{∅\}$, $z_{q+1}=\{z_q\}$ sont de type 1).

\item Appelons \emph{hauteur} d'un ensemble $A$ la borne supérieure
(finie ou infinie) des entiers $n$ tels qu'il existe une suite
$x_n\in x_{n-1}\in \ldots \in x_0=$. Montrer que les ensembles de
hauteur $\leq n$ sont finis et forment un ensemble fini, dont le
cardinal $p_n$ se calcule au moyen de $p_0=1$, $p_{n+1}=2^{p_n}$
(procéder par récurrence sur $n$, en notant que les éléments d'un
ensemble $A$ de hauteur $\leq n$ sont des ensembles de hauteur $\leq
n-1$).

\item Soit $G$ un ensemble ordonné noethérien admettant un plus
petit élément $g_0$ et tel que pour tout $h\in G$, l'ensemble de
$g\leq h$ sont fini et totalement ordonné.
\begin{enumerate}
\renewcommand{\theenumii}{\alph{enumii}}
\renewcommand{\labelenumii}{\theenumii)}
\item Montrer\pageoriginale que tout élément $h\neq g_0$ de $G$
admet un prédécesseur et un seul.

\item Soit $\varphi$ l'application de $G$ définie dans le lemme 1
(\ie $\varphi(g)$ est l'ensemble des $\varphi(g')$ où $g'$ parcourt
l'ensemble des successeurs de $g$). On pose $A=\varphi(g_0)$.
Montrer que $A$ est un ensemble artinien. Pour $g\in G$, soit
$g_0<g_1<\ldots<g_n=g$ la suite des éléments $\leq g$; posons
$f(g)=(\varphi(g_0),\varphi(g_1),\ldots,\varphi(g_n))$; montrer que
$f$ est une application croissante et surjective de $G$ sur
l'ensemble ordonné $G(A)$ du texte. Montrer que, si $f$ est
\emph{injective}, c'est un \emph{isomorphisme} de $G$ sur $G(A)$.

\item Pour $g\in G$, soit $M_g$ l'ensemble des majorants de $g$. On
suppose que, pour tout couple d'éléments distincts $g$, $g'$ ayant
même prédécesseur $p$, les ensembles ordonnés $M_g$ et
$M_{g'}$ sont
non isomorphes. Montrer que l'application $f$ de b) est alors un
\emph{isomorphisme} de $G$ sur $G(A)$. [Si $f(g)=f(g')$ avec $g\neq
g'$ et si $g_0<g_1<\ldots<g_n=g$ et $g'_0<g'_1<\ldots <g'_{n'}=g'$
sont la suite des éléments $\leq g$ et celles des
éléments $\leq
g'$, montrer que $n=n'$, et qu'on peut supposer que $g$ et $g'$ ont
le même prédécesseur $p$; considérer alors l'ensemble
des $p\in G$
tels qu'il existe deux successeurs distincts $g$, $g'$ de $p$ tels
que $f(g)=f(g')$, un élément maximal $q$ de cet ensemble, et deux
successeurs distincts $h$, $h'$ de $q$ tels que $f(h)=f(h')$; noter
que les restrictions de $f$ à $M_h$ et à $M_{h'}$ sont injectives,
donc (par b)) sont des isomorphismes de $M_h$ sur $G(\varphi(h))$ et
de $M_{h'}$ sur $G(\varphi(h'))$; déduire de l'hypothèse de
non-isomorphisme de $M_h$ et $M_{h'}$ qu'on a $\varphi(h)\neq
\varphi(h')$, ce qui contredit $f(h)=f(h')$].
\end{enumerate}
N.B.\pageoriginale L'exercice 3) donne des renseignements très
précis sur la manière dont sont « fabriqués » les ensembles
artiniens. On savait déjà qu'un tel ensemble $A$ est
déterminé par
la classe d'isomorphisme de l'ensemble ordonné $G(A)$ (lemme 1). On
sait maintenant \emph{caractériser} les ensembles ordonnés $G$
isomorphes à des $G(A)$:
\begin{itemize}
\item[$\alpha)$] $G$ est noethérien et admet un plus petit
élément;

\item[$\beta)$] Pour tout $g\in G$, l'ensemble des $h\leq g$ est
totalement ordonné et fini (d'où l'existence et l'unicité du
prédécesseur de $g$);

\item[$\gamma)$] Si $g$, $g'$ sont des éléments distincts ayant
même prédécesseur, l'ensemble $M_g$ des majorants de $g$ et celui
$M_{g'}$ des majorants de $g$ ne sont pas isomorphes.
\end{itemize}

\item Soit $(A_n)_{n\geq 0}$ la suite des \emph{ordinaux} finis,
définie par $A_0=∅$, $A_{n+1}=A_n\cup \{A_n\}$. Montrer
que l'ensemble ordonné $G(A_n)$ a $2^n$ éléments, et que le nombre
de ses éléments de hauteur $q$ (au sens de l'exerc. 2)) est
$\binom{n}{q}$.
\end{enumerate}
\end{enonce*}

\begin{thebibliography}{10}%%
\bibitem[1]{SGA4-01:I.2} \sisi{B. \textsc{Mitchell}, Theory of categories, Academic Press
(1965).}{B. \textsc{Mitchell}, \textit{Theory of categories},
Academic Press, 1965.}
\end{thebibliography}

\ifx\danslelivre\undefined
%\bibliography{../fichiers-configuration-latex/bibliographie}
%\bibliographystyle{../fichiers-configuration-latex/smfalpha}
\end{document}
\fi
